\documentclass{article}
\usepackage[a4paper,margin=1in]{geometry}
\usepackage{xcolor}
\usepackage{hyperref}
\usepackage{titlesec}
\title{ShrijiLang Full Project Code}
\author{Hukum Manikpuri}
\date{\today}
\begin{document}
\maketitle
\tableofcontents
\newpage

\section*{File: ./all_code.txt}
egin{verbatim}
#include <string.h>
#include <stdlib.h>

#include "../include/env.h"
#include "../include/value.h"

/*──────────────────────────────────────────────────────────────
 | SHRIJILANG — ENVIRONMENT STORE (SCOPE STACK)
 |
 | Responsibility:
 |   • Variable storage (Value)
 |   • Scope push / pop
 |   • Get / Set / Update / Exists
 |
 | RULES:
 |   • env NEVER prints
 |   • env NEVER calls shiri
 |   • Only DATA LAYER
 *──────────────────────────────────────────────────────────────*/

/*──────────────────────────────────────────────
 | CREATE NEW ENV
 *──────────────────────────────────────────────*/
Env *new_env(void)
{
    Env *env = (Env *)malloc(sizeof(Env));
    if (!env) return NULL;

    env->depth = 1; /* global scope exists */
    env->frames[0].count = 0;

    return env;
}

/*──────────────────────────────────────────────
 | PUSH NEW SCOPE
 *──────────────────────────────────────────────*/
void env_push_scope(Env *env)
{
    if (!env) return;

    if (env->depth >= ENV_MAX_SCOPES) {
        return; /* silently ignore overflow */
    }

    env->frames[env->depth].count = 0;
    env->depth++;
}

/*──────────────────────────────────────────────
 | POP SCOPE
 *──────────────────────────────────────────────*/
void env_pop_scope(Env *env)
{
    if (!env) return;

    /* never pop global scope */
    if (env->depth <= 1) return;

    int top = env->depth - 1;

    /* free all values in top scope */
    for (int i = 0; i < env->frames[top].count; i++) {
        value_free(&env->frames[top].vars[i].value);
    }

    env->frames[top].count = 0;
    env->depth--;
}

/*──────────────────────────────────────────────
 | CHECK IF VARIABLE EXISTS (search top → bottom)
 *──────────────────────────────────────────────*/
int env_exists(Env *env, const char *name)
{
    if (!env || !name) return 0;

    for (int f = env->depth - 1; f >= 0; f--) {
        for (int i = 0; i < env->frames[f].count; i++) {
            if (strcmp(env->frames[f].vars[i].name, name) == 0) {
                return 1;
            }
        }
    }

    return 0;
}

/*──────────────────────────────────────────────
 | SET VARIABLE
 | - if exists in CURRENT scope: update it
 | - else create new variable in CURRENT scope
 *──────────────────────────────────────────────*/
void env_set(Env *env, const char *name, Value v)
{
    if (!env || !name) return;

    int top = env->depth - 1;
    EnvFrame *frame = &env->frames[top];

    /* update if already exists in current scope */
    for (int i = 0; i < frame->count; i++) {
        if (strcmp(frame->vars[i].name, name) == 0) {
            value_free(&frame->vars[i].value);
            frame->vars[i].value = value_copy(v);
            return;
        }
    }

    /* add new variable in current scope */
    if (frame->count >= ENV_MAX_VARS) {
        return; /* silently ignore overflow */
    }

    strncpy(
        frame->vars[frame->count].name,
        name,
        sizeof(frame->vars[frame->count].name) - 1
    );

    frame->vars[frame->count].name[
        sizeof(frame->vars[frame->count].name) - 1
    ] = '\0';

    frame->vars[frame->count].value = value_copy(v);
    frame->count++;
}

/*──────────────────────────────────────────────
 | UPDATE VARIABLE (search top → bottom)
 | - if exists in ANY scope: update nearest found scope
 | - else create in current scope
 *──────────────────────────────────────────────*/
void env_update(Env *env, const char *name, Value v)
{
    if (!env || !name) return;

    /* search from top scope to global */
    for (int f = env->depth - 1; f >= 0; f--) {
        EnvFrame *frame = &env->frames[f];

        for (int i = 0; i < frame->count; i++) {
            if (strcmp(frame->vars[i].name, name) == 0) {
                value_free(&frame->vars[i].value);
                frame->vars[i].value = value_copy(v);
                return;
            }
        }
    }

    /* if not found, fallback: create in current scope */
    env_set(env, name, v);
}

/*──────────────────────────────────────────────
 | GET VARIABLE VALUE (search top → bottom)
 *──────────────────────────────────────────────*/
Value env_get(Env *env, const char *name)
{
    if (!env || !name) return value_null();

    for (int f = env->depth - 1; f >= 0; f--) {
        for (int i = 0; i < env->frames[f].count; i++) {
            if (strcmp(env->frames[f].vars[i].name, name) == 0) {
                return value_copy(env->frames[f].vars[i].value);
            }
        }
    }

    return value_null();
}
#include "../include/decision_engine.h"

/* SAFE FIXABLE ERRORS */
int is_safe_error(ShrijiErrorCode code)
{
    switch(code)
    {
        case E_PARSE_MISSING_OPERATOR:
        case E_PARSE_DOUBLE_OPERATOR:
            return 1;
        default:
            return 0;
    }
}

/* AMBIGUOUS ERRORS */
int is_ambiguous_error(ShrijiErrorCode code)
{
    switch(code)
    {
        case E_PARSE_INVALID_TOKEN:
        case E_PARSE_02:
            return 1;
        default:
            return 0;
    }
}

/* MAIN DECISION ENGINE */
DecisionType shriji_take_decision(
    int confidence,
    int risk,
    ShrijiErrorCode code
)
{
    (void)confidence;
    (void)risk;

    /* SAFE → always auto fix */
    if (is_safe_error(code))
        return DECISION_AUTO_FIX;

    /* AMBIGUOUS → ask user */
    if (is_ambiguous_error(code))
        return DECISION_EDIT;

    /* REST → reject */
    return DECISION_REJECT;
}
#include "../include/runtime.h"
#include "../include/state.h"

/* GLOBAL POINTER */
ShrijiRuntime *current_runtime = NULL;

void runtime_init(ShrijiRuntime *rt)
{
    if (!rt) return;

    state_init(&rt->state);

    rt->break_flag = 0;
    rt->continue_flag = 0;
    rt->loop_depth = 0;

    rt->return_flag = 0;
    rt->return_value = value_null();

    rt->call_depth = 0;
    rt->error_flag = 0;
}

void runtime_reset(ShrijiRuntime *rt)
{
    if (!rt) return;

    rt->break_flag = 0;
    rt->continue_flag = 0;
    rt->loop_depth = 0;

    rt->return_flag = 0;
    rt->return_value = value_null();

    rt->state.steps_used = 0;

    rt->call_depth = 0;
    rt->error_flag = 0;
}
#include <stdio.h>
#include <string.h>

#include <stdlib.h>

#include "../include/interpreter.h"
#include "../include/parser.h"
#include "../include/ast.h"
#include "../include/env.h"
#include "../include/error.h"
#include "../include/state.h"
#include "../include/event.h"
#include "../include/command.h"
#include "../include/value.h"
#include "../include/ai_router.h"

#include "../include/smriti_personal.h"
#include "../include/smriti_session.h"

#include "../include/runtime.h"

#include "../include/user_config.h"
#include "../include/env.h"

#ifdef SHRIJI_ENABLE_MASTER_TOKENS
#include "lang/token_master.h"
#endif

/*──────────────────────────────────────────────
  helper: unescape string like "a\\nb" -> actual newline
  supports: \n \t \r \\ \"
──────────────────────────────────────────────*/
static char *unescape_string(const char *s)
{
    if (!s) return NULL;

    size_t n = strlen(s);
    char *out = (char *)malloc(n + 1);
    if (!out) return NULL;

    size_t j = 0;
    for (size_t i = 0; i < n; i++) {
        if (s[i] == '\\' && i + 1 < n) {
            char c = s[i + 1];
            if (c == 'n')      { out[j++] = '\n'; i++; continue; }
            else if (c == 't') { out[j++] = '\t'; i++; continue; }
            else if (c == 'r') { out[j++] = '\r'; i++; continue; }
            else if (c == '\\'){ out[j++] = '\\'; i++; continue; }
            else if (c == '"') { out[j++] = '"';  i++; continue; }
        }
        out[j++] = s[i];
    }

    out[j] = '\0';
    return out;
}
/*──────────────────────────────────────────────────────────────
 | SHRIJI TRUTH DIARY — SINGLE SOURCE OF FACT
 *──────────────────────────────────────────────────────────────*/
ShrijiTruthDiary SHRIJI_DIARY = {0};

/*──────────────────────────────────────────────────────────────
 | STATE DECISION HOOK
 *──────────────────────────────────────────────────────────────*/
static int state_allows_execution(ShrijiRuntime *rt)
{
    if (!rt) return 0;

    if (rt->state.safety == STATE_CRITICAL) {
        return 0;
    }
    return 1;
}

/*──────────────────────────────────────────────────────────────
 | SAFE DIVISION
 *──────────────────────────────────────────────────────────────*/
static double safe_div(double a, double b)

{
    if (b == 0) {
        event_fire(EVENT_ERROR, "division by zero");
        shriji_error(
            E_RUNTIME_DIV_ZERO,
            "division",
            "division by zero",
            "denominator must not be zero"
        );
        return 0;
    }
    return a / b;
}

/*──────────────────────────────────────────────────────────────
 | SAFE MODULO
 *──────────────────────────────────────────────────────────────*/
static double safe_mod(double a, double b)
{
    if (b == 0) {
        event_fire(EVENT_ERROR, "modulo by zero");
        shriji_error(
            E_RUNTIME_DIV_ZERO,
            "modulo",
            "modulo by zero",
            "denominator must not be zero"
        );
        return 0;
    }

    long long ai = (long long)a;
    long long bi = (long long)b;

    if (bi == 0)
        return 0;

    return (double)(ai % bi);
}

/*──────────────────────────────────────────────────────────────
 | Value → number helper
 *──────────────────────────────────────────────────────────────*/
static const char *strip_quotes(const char *s)
{
    if (!s) return "";

    while (*s == ' ' || *s == '\t')
        s++;

    size_t len = strlen(s);

    if (len >= 2 && s[0] == '"' && s[len - 1] == '"') {
        static char buf[1024];
        size_t n = len - 2;
        if (n >= sizeof(buf)) n = sizeof(buf) - 1;

        memcpy(buf, s + 1, n);
        buf[n] = '\0';
        return buf;
    }

    return s;
}
/*──────────────────────────────────────────────────────────────
 | MAIN EVALUATOR
 *──────────────────────────────────────────────────────────────*/
Value shriji_execute_line(const char *line, Env *env, ShrijiRuntime *rt)
{
    if (!line || !*line) return value_null();

    while (*line == ' ' || *line == '\t') line++;
    if (*line == '\0') return value_null();

    if (DEV_MODE)
    printf("[KRST] Input: %s\n", line);

    ASTNode *node = parse_program(line);

/*  HARD STOP FIX */
if (!node || error_reported)
{
    if (rt) {
        rt->error_flag = 1;
    }

    error_reported = 0;
    return value_null();
}
    return eval(node, env, rt);
}

Value eval(ASTNode *node, Env *env, ShrijiRuntime *rt)
{
    static int state_initialized = 0;

    if (!state_initialized) {
        state_init(&rt->state);
        state_initialized = 1;
    }

    if (!state_allows_execution(rt) || !node)
        return value_null();

    /* 🧠 GLOBAL STEP COUNTER */
    if (node->type != AST_WHILE &&
        node->type != AST_BREAK &&
        node->type != AST_CONTINUE &&
        node->type != AST_BLOCK &&
        node->type != AST_PROGRAM) {
        rt->state.steps_used++;
    }

    if (rt->state.steps_used > rt->state.max_steps) {
        shriji_error(
            E_PARSE_02,
            "runtime",
            "execution limit exceeded",
            "possible infinite loop detected"
        );
        rt->state.safety = STATE_CRITICAL;
        return value_null();
    }



switch (node->type) {

case AST_BREAK:
    if (rt->loop_depth <= 0) {
        shriji_error(
            E_RUNTIME_01,
            "break",
            "loop ke bahar rukja allowed nahi hai",
            "rukja sirf while loop ke andar use karein"
        );
        return value_null();
    }

    rt->break_flag = 1;
    return value_null();

case AST_CONTINUE:

     if (rt->loop_depth <= 0) {
        shriji_error(
        E_RUNTIME_01,
        "continue",
        "loop ke bahar continue allowed nahi hai",
        "continue sirf while loop ke andar use karein"
    );
    return value_null();
}
    rt->continue_flag = 1;
    return value_null();


case AST_RETURN: {
        Value v = value_null();
        if (node->return_expr)
         v = eval(node->return_expr, env, rt);

       rt->return_flag = 1;
       rt->return_value = v;
        return v;
    }

case AST_UNARY:
{
    ASTNode *cur = node;
    int sign = 1;

    while (cur && cur->type == AST_UNARY) {
        if (cur->op[0] == '-')
            sign *= -1;

        cur = cur->left;
    }

    Value v = eval(cur, env, rt);

    if (v.type != VAL_NUMBER)
        return value_null();

    v.number = v.number * sign;

    return v;
}

/*────────────────────────────────────────
 | INC / DEC (++ / --)
 *────────────────────────────────────────*/

case AST_INC: {

    /* only identifiers allowed */
    if (!node->left || node->left->type != AST_IDENTIFIER) {
        return value_null();
    }

    const char *name = node->left->name;

    Value old = env_get(env, name);

    if (old.type != VAL_NUMBER) {
        return value_null();
    }

    Value newv = old;
    newv.number = old.number + 1;

    env_set(env, name, newv);

    if (node->is_prefix)
        return newv;
    else
        return old;
}

case AST_DEC: {

    /* only identifiers allowed */
    if (!node->left || node->left->type != AST_IDENTIFIER) {
        return value_null();
    }

    const char *name = node->left->name;

    Value old = env_get(env, name);

    if (old.type != VAL_NUMBER) {
        return value_null();
    }

    Value newv = old;
    newv.number = old.number - 1;

    env_set(env, name, newv);

    if (node->is_prefix)
        return newv;
    else
        return old;
}

    /*──────────────────────────────────────────────
      FUNCTIONS
    ──────────────────────────────────────────────*/

    case AST_FUNCTION: {
        /* store function AST node into env */
        env_set(env, node->function_name, value_function(node));
        return value_null();
    }

case AST_RACHNA: {
    env_set(env, node->rachna_name, value_function(node));
    return value_null();
}

case AST_CALL: {

/*────────────────────────────
  STACK GUARD (ENTRY)
────────────────────────────*/
rt->call_depth++;

if (rt->call_depth > 3000) {
    shriji_error(
        E_RUNTIME_01,
        "recursion",
        "maximum recursion depth exceeded",
        "safe limit: 3000"
    );

    rt->call_depth--;
    return value_null();
}

   /*──────────────────────────────────────────────
      BUILT-IN FUNCTIONS (FILE I/O)
      padho("file")                -> returns string
      likho("file", "text")        -> overwrite write -> returns 1/0
      jodo("file", "text")         -> append write   -> returns 1/0

      ✅ polish:
      - supports escape sequences: \n \t \r \\ \"
    ──────────────────────────────────────────────*/

    /* READ */
    if (strcmp(node->function_name, "padho") == 0) {

        if (node->arg_count != 1) {
            shriji_error(
                E_PARSE_02,
                "padho",
                "argument count mismatch",
                "use: padho(\"file.txt\")"
            );
            return value_null();
        }

        Value pathv = eval(node->args[0], env, rt);

        if (pathv.type != VAL_STRING || !pathv.string) {
            value_free(&pathv);
            shriji_error(
                E_PARSE_02,
                "padho",
                "file path must be string",
                "use: padho(\"file.txt\")"
            );
            return value_null();
        }

        const char *path = strip_quotes(pathv.string);

FILE *fp = fopen(path, "rb");
if (!fp) {
    value_free(&pathv);

    shriji_error(
        E_PARSE_02,
        "padho",
        "file not found",
        "check path or create file first"
    );

    return value_string("");
}
        fseek(fp, 0, SEEK_END);
        long sz = ftell(fp);
        fseek(fp, 0, SEEK_SET);

        if (sz < 0) sz = 0;
        if (sz > 1024 * 1024) sz = 1024 * 1024; /* 1MB limit v1 */

        char *buf = (char *)malloc((size_t)sz + 1);
        if (!buf) {
            fclose(fp);
            value_free(&pathv);
            return value_null();
        }

        size_t nread = fread(buf, 1, (size_t)sz, fp);
        buf[nread] = '\0';

        fclose(fp);
        value_free(&pathv);

        Value out = value_string(buf);
        free(buf);
        return out;
    }

    /* WRITE (overwrite) */
    if (strcmp(node->function_name, "likho") == 0) {

        if (node->arg_count != 2) {
            shriji_error(
                E_PARSE_02,
                "likho",
                "argument count mismatch",
                "use: likho(\"file.txt\", \"data\")"
            );
            return value_null();
        }

        Value pathv = eval(node->args[0], env, rt);
        Value datav = eval(node->args[1], env, rt);

        if (pathv.type != VAL_STRING || !pathv.string) {
            value_free(&pathv);
            value_free(&datav);
            shriji_error(
                E_PARSE_02,
                "likho",
                "file path must be string",
                "use: likho(\"file.txt\", \"data\")"
            );
            return value_null();
        }

        const char *path = strip_quotes(pathv.string);

        const char *data = "";
        char tmp[128];
        tmp[0] = '\0';

        char *tmp_data = NULL;

        if (datav.type == VAL_STRING && datav.string) {
            tmp_data = unescape_string(strip_quotes(datav.string));
            data = tmp_data ? tmp_data : "";
        } else if (datav.type == VAL_NUMBER) {
            snprintf(tmp, sizeof(tmp), "%g", datav.number);
            data = tmp;
        }

        FILE *fp = fopen(path, "wb");
        if (!fp) {
            if (tmp_data) free(tmp_data);
            value_free(&pathv);
            value_free(&datav);
            return value_number(0);
        }

        fwrite(data, 1, strlen(data), fp);
        fclose(fp);

        if (tmp_data) free(tmp_data);

        value_free(&pathv);
        value_free(&datav);

        return value_number(1);
    }

    /* APPEND */
    if (strcmp(node->function_name, "jodo") == 0) {

        if (node->arg_count != 2) {
            shriji_error(
                E_PARSE_02,
                "jodo",
                "argument count mismatch",
                "use: jodo(\"file.txt\", \"data\")"
            );
            return value_null();
        }

        Value pathv = eval(node->args[0], env, rt);
        Value datav = eval(node->args[1], env, rt);

        if (pathv.type != VAL_STRING || !pathv.string) {
            value_free(&pathv);
            value_free(&datav);
            shriji_error(
                E_PARSE_02,
                "jodo",
                "file path must be string",
                "use: jodo(\"file.txt\", \"data\")"
            );
            return value_null();
        }

        const char *path = strip_quotes(pathv.string);

        const char *data = "";
        char tmp[128];
        tmp[0] = '\0';

        char *tmp_data = NULL;

        if (datav.type == VAL_STRING && datav.string) {
            tmp_data = unescape_string(strip_quotes(datav.string));
            data = tmp_data ? tmp_data : "";
        } else if (datav.type == VAL_NUMBER) {
            snprintf(tmp, sizeof(tmp), "%g", datav.number);
            data = tmp;
        }

        FILE *fp = fopen(path, "ab");
        if (!fp) {
            if (tmp_data) free(tmp_data);
            value_free(&pathv);
            value_free(&datav);
            return value_number(0);
        }

        fwrite(data, 1, strlen(data), fp);
        fclose(fp);

        if (tmp_data) free(tmp_data);

        value_free(&pathv);
        value_free(&datav);

        return value_number(1);
    }

    /*──────────────────────────────────────────────
      NORMAL FUNCTION CALL
    ──────────────────────────────────────────────*/
Value fnv = env_get(env, node->function_name);

if (fnv.type != VAL_FUNCTION) {
    shriji_error(
        E_PARSE_INVALID_TOKEN,
        "call",
        "function not found",
        "check name or define function"
    );
    return value_null();
}

ASTNode *fn = fnv.function;

    /* validate arg count */
    if (fn->rachna_param_count != node->arg_count) {
        shriji_error(
            E_PARSE_02,
            node->function_name,
            "argument count mismatch",
            "call with correct number of args"
        );
        return value_null();
    }

    /* save outer return state (nested calls safe) */
    int old_return_flag = rt->return_flag;
    Value old_return_value = rt->return_value;

    rt->return_flag = 0;
    rt->return_value = value_null();

    /* new call scope */
    env_push_scope(env);

    /* bind params */
       for (int i = 0; i < fn->rachna_param_count; i++) {
    ASTNode *p = fn->rachna_params[i];
        Value av = eval(node->args[i], env, rt);
        env_set(env, p->name, av);
    }

    /* run body */
    Value result = eval(fn->rachna_body, env, rt);

    if (rt->return_flag)
        result = rt->return_value;

    /* end call scope */
    env_pop_scope(env);

    /* restore outer return state */
    rt->return_flag = old_return_flag;
    rt->return_value = old_return_value;

   rt->call_depth--;
    return result;
}

/*──────────────────────────────────────────────
      CORE VALUES
    ──────────────────────────────────────────────*/
case AST_NUMBER:
    return value_number(node->number_value);

case AST_STRING:
    return value_string(node->string_value);

case AST_BOOL:
    return value_bool(node->bool_value);



    case AST_LIST: {
        int n = node->element_count;

        if (n <= 0 || !node->elements) {
            return value_list(NULL, 0);
        }

        Value *items = (Value *)malloc(sizeof(Value) * n);
        if (!items) return value_null();

        for (int i = 0; i < n; i++) {
            Value ev = eval(node->elements[i], env, rt);
            items[i] = value_copy(ev);
            value_free(&ev);
        }

        return value_list(items, n);
    }


case AST_DICT: {
        int n = node->dict_count;

        if (n <= 0 || !node->dict_keys || !node->dict_values) {
            return value_dict(NULL, NULL, 0);
        }

        Value *keys = (Value *)malloc(sizeof(Value) * n);
        Value *vals = (Value *)malloc(sizeof(Value) * n);

        if (!keys || !vals) {
            if (keys) free(keys);
            if (vals) free(vals);
            return value_null();
        }

        for (int i = 0; i < n; i++) {
            Value kv = eval(node->dict_keys[i], env, rt);
            Value vv = eval(node->dict_values[i], env, rt);

            keys[i] = value_copy(kv);
            vals[i] = value_copy(vv);

            value_free(&kv);
            value_free(&vv);
        }

        return value_dict(keys, vals, n);
    }

case AST_INDEX: {

    Value tv = eval(node->index_target, env, rt);
    Value iv = eval(node->index_expr, env, rt);

    /*──────────────────────────────────────────────
      LIST INDEXING: a[0]
    ──────────────────────────────────────────────*/
    if (tv.type == VAL_LIST) {

        if (!tv.items) {
            value_free(&tv);
            value_free(&iv);
            return value_null();
        }

        if (iv.type != VAL_NUMBER) {
            shriji_error(
                E_RUNTIME_01,
                "list",
                "List index number hona chahiye",
                "example: a[0]"
            );

               value_free(&tv);
               value_free(&iv);
               return value_null();
        }

        long long idx = iv.number;
        value_free(&iv);

        if (idx < 0 || idx >= tv.count) {
            shriji_error(
                E_RUNTIME_01,
                "list",
                "List index out of range",
                "example: a[0]"
            );
           value_free(&tv);
           return value_null();
        }

        Value out = value_copy(tv.items[(int)idx]);
        value_free(&tv);
        return out;
    }

    /*──────────────────────────────────────────────
      DICT INDEXING: d["a"]
    ──────────────────────────────────────────────*/
    if (tv.type == VAL_DICT) {

        if (iv.type != VAL_STRING || !iv.string) {
            shriji_error(
                E_RUNTIME_01,
                "dict",
                "dictionary ki key string honi chahiye",
                "example: d[\"a\"]"
            );
            value_free(&tv);
            value_free(&iv);
            return value_null();
        }

        for (int i = 0; i < tv.dict_count; i++) {

            Value k = tv.dict_keys[i];

            if (k.type == VAL_STRING &&
                strcmp(k.string, iv.string) == 0) {

                Value out = value_copy(tv.dict_values[i]);
                value_free(&tv);
                value_free(&iv);
                return out;
            }
        }

        shriji_error(
            E_RUNTIME_01,
            "dict",
            "dictionary me ye key maujood nahi hai",
            "example: d[\"a\"]"
        );
        value_free(&tv);
        value_free(&iv);
        return value_null();
    }

    /*──────────────────────────────────────────────
      UNSUPPORTED TARGET
    ──────────────────────────────────────────────*/
    value_free(&tv);
    value_free(&iv);
    return value_null();
}


case AST_INDEX_UPDATE: {

    /* target must be identifier so update persists */
    if (!node->index_target || node->index_target->type != AST_IDENTIFIER) {
        return value_null();
    }

    const char *var_name = node->index_target->name;

    /* fetch container directly from env */
    Value tv = env_get(env, var_name);

    Value iv = eval(node->index_expr, env, rt);
    Value vv = eval(node->index_value, env, rt);

    /*──────────────────────────────────────────────
      LIST INDEX UPDATE: a[0] = value
    ──────────────────────────────────────────────*/
    if (tv.type == VAL_LIST) {

        if (iv.type != VAL_NUMBER) goto cleanup;

        long long idx = iv.number;
        if (idx < 0 || idx >= tv.count) goto cleanup;

        value_free(&tv.items[(int)idx]);
        tv.items[(int)idx] = value_copy(vv);

        /* WRITE BACK */
        env_set(env, var_name, tv);

        value_free(&iv);
        return value_copy(vv);
    }

    /*──────────────────────────────────────────────
      DICT INDEX UPDATE: d["a"] = value
    ──────────────────────────────────────────────*/
    if (tv.type == VAL_DICT) {

        if (iv.type != VAL_STRING || !iv.string) goto cleanup;

        /* update if exists */
        for (int i = 0; i < tv.dict_count; i++) {
            Value k = tv.dict_keys[i];
            if (k.type == VAL_STRING && strcmp(k.string, iv.string) == 0) {

                value_free(&tv.dict_values[i]);
                tv.dict_values[i] = value_copy(vv);

                env_set(env, var_name, tv);
                value_free(&iv);
                return value_copy(vv);
            }
        }

        /* append new key */
        tv.dict_keys = realloc(tv.dict_keys,
                               sizeof(Value) * (tv.dict_count + 1));
        tv.dict_values = realloc(tv.dict_values,
                                 sizeof(Value) * (tv.dict_count + 1));

        tv.dict_keys[tv.dict_count]   = value_copy(iv);
        tv.dict_values[tv.dict_count] = value_copy(vv);
        tv.dict_count++;

        env_set(env, var_name, tv);
        value_free(&iv);
        return value_copy(vv);
    }

cleanup:
    value_free(&tv);
    value_free(&iv);
    value_free(&vv);
    return value_null();
}

case AST_IDENTIFIER: {

    const char *name = node->name;

    /* ===== ENV CHECK ===== */
    if (env_exists(env, name)) {
        return env_get(env, name);
    }

    /* ===== SESSION KEYWORD: last ===== */
    if (strcmp(name, "last") == 0) {
        return smriti_session_get_last();
    }

    /* ===== SMRITI FALLBACK ===== */
    const char *mem = smriti_personal_get(name);

    if (mem) {
        Value v;
        v.type = VAL_NUMBER;
        v.number = atof(mem);
        return v;
    }

    /* ===== ERROR ===== */
    shriji_error(
        E_ASSIGN_01,
        name,
        "variable declare nahi hua",
        "use mavi x = ... pehle likhiye"
    );

    return value_null();
}


case AST_ASSIGNMENT: {

    Value value = eval(node->value, env, rt);

    /* ===== FIX: SAFE COPY ===== */
    Value stored = value_copy(value);

    /* ===== ENV SAVE ===== */
    env_update(env, node->name, stored);

    /* ===== SMRITI SAVE ===== */
    if (value.type == VAL_NUMBER) {
        char buf[64];
        snprintf(buf, sizeof(buf), "%g", value.number);

        if (node->name[0] != '\0') {
            /* smriti_personal_set(node->name, buf); */
        }
    }

    event_fire(EVENT_ASSIGNMENT, node->name);
    state_on_success(&rt->state);

    return value;
}



case AST_UPDATE: {

    if (node->name[0] != '\0') {

        if (!env_exists(env, node->name)) {
            shriji_error(
                E_ASSIGN_01,
                node->name,
                "variable not declared",
                "use mavi x = ... first"
            );
            return value_null();
        }

        Value value = eval(node->value, env, rt);

        /*  FIXED: use env_update */
        env_update(env, node->name, value);

        event_fire(EVENT_ASSIGNMENT, node->name);
        state_on_success(&rt->state);

        return value;
    }

    shriji_error(
        E_ASSIGN_01,
        "update",
        "invalid update target",
        "use: x = expr"
    );

    return value_null();
}




case AST_BINARY: {
    Value Lv = eval(node->left, env, rt);
    Value Rv = eval(node->right, env, rt);

    char op = node->op[0];

    /* ───────────────────────────────────────────────
       TYPE CHECK
    ─────────────────────────────────────────────── */
    if (Lv.type != Rv.type) {
        shriji_error(
            E_RUNTIME_TYPE_MISMATCH,
            "binary",
            "type mismatch in expression",
            "same type ke values use karein"
        );
        return value_null();
    }

    /* ───────────────────────────────────────────────
       NUMBER OPERATIONS
    ─────────────────────────────────────────────── */
/*  LOGICAL OPERATORS (GLOBAL — BEFORE TYPE CHECK) */
if (op == '&') {
    return value_bool(
        value_is_truthy(Lv) && value_is_truthy(Rv)
    );
}

if (op == '|') {
    return value_bool(
        value_is_truthy(Lv) || value_is_truthy(Rv)
    );
}


    if (Lv.type == VAL_NUMBER) {

        double L = Lv.number;
        double R = Rv.number;


        if (op == '+') return value_number(L + R);
        if (op == '-') return value_number(L - R);
        if (op == '*') return value_number(L * R);
        if (op == '/') return value_number(safe_div(L, R));
        if (op == '%') return value_number(safe_mod(L, R));

        if (op == '>') return value_bool(L > R);
        if (op == '<') return value_bool(L < R);
        if (op == 'G') return value_bool(L >= R);
        if (op == 'L') return value_bool(L <= R);
        if (op == '=') return value_bool(L == R);
        if (op == '!') return value_bool(L != R);

        shriji_error(
            E_RUNTIME_01,
            "binary",
            "unknown numeric operator",
            NULL
        );
        return value_null();
    }

    /* ───────────────────────────────────────────────
       STRING COMPARISON ONLY
    ─────────────────────────────────────────────── */
    if (Lv.type == VAL_STRING) {

        const char *L = Lv.string ? Lv.string : "";
        const char *R = Rv.string ? Rv.string : "";

        int cmp = strcmp(L, R);

         if (op == '=') return value_bool(cmp == 0);
         if (op == '!') return value_bool(cmp != 0);
         if (op == '>') return value_bool(cmp > 0);
         if (op == '<') return value_bool(cmp < 0);
         if (op == 'G') return value_bool(cmp >= 0);
         if (op == 'L') return value_bool(cmp <= 0);


        shriji_error(
            E_RUNTIME_01,
            "binary",
            "string ke saath sirf comparison allowed hai",
            "use: == != < > <= >="
        );
        return value_null();
    }

    /* ───────────────────────────────────────────────
       UNSUPPORTED TYPE
    ─────────────────────────────────────────────── */
    shriji_error(
        E_PARSE_02,
        "binary",
        "ye operation is type ke liye allowed nahi hai",
        "sirf number ya string comparison supported hai"
    );
    rt->state.safety = STATE_CRITICAL;
                return value_null();
}


       case AST_NOT: {
    Value v = eval(node->left, env, rt);

    int truth = value_is_truthy(v);
    value_free(&v);

    return value_bool(!truth);
    }


case AST_IF: {
    Value condv = eval(node->condition, env, rt);

    Value result = value_null();

    if (value_is_truthy(condv)) {

        env_push_scope(env);

        result = eval(node->left, env, rt);

        env_pop_scope(env);
    }
    else if (node->else_block) {

        env_push_scope(env);

        result = eval(node->else_block, env, rt);

        env_pop_scope(env);
    }

    value_free(&condv);
    return result;
}



case AST_WHILE: {

    rt->loop_depth++;

    int local_guard = 0;

    while (1) {

        Value condv = eval(node->condition, env, rt);

        if (!value_is_truthy(condv)) {
            value_free(&condv);
            break;
        }

        value_free(&condv);

        rt->continue_flag = 0;

        eval(node->body, env, rt);

        if (rt->return_flag) {
            rt->loop_depth--;
            return rt->return_value;
        }

        if (rt->break_flag) {
            rt->break_flag = 0;
            break;
        }

        if (rt->continue_flag) {
            rt->continue_flag = 0;
            continue;
        }

        if (++local_guard > 1000000) {
            shriji_error(
                E_RUNTIME_01,
                "loop",
                "infinite loop detected",
                "use rukja or fix condition"
            );
            rt->state.safety = STATE_CRITICAL;
            break;
        }
    }

    rt->loop_depth--;
    return value_null();
}



case AST_BLOCK: {
    Value last = value_null();


    /*  NEW SCOPE START */
    env_push_scope(env);


    for (int i = 0; i < node->stmt_count; i++) {

        last = eval(node->statements[i], env, rt);

        if (rt->return_flag || rt->break_flag || rt->continue_flag)
            break;
    }

    /*  SCOPE END */
    env_pop_scope(env);


    if (rt->return_flag)
        return rt->return_value;

    return last;
}


case AST_PROGRAM: {
    Value last = value_null();


    for (int i = 0; i < node->stmt_count; i++) {

        last = eval(node->statements[i], env ,rt);

/* 🔥 ONLY break for loop control, not function return */
if (rt->break_flag || rt->continue_flag)
    break;
    }

    if (rt->return_flag)
        return rt->return_value;

    return last;
}

case AST_IMPORT: {

    /* 🔥 ENABLE ANALYSIS MODE */
    shriji_set_error_mode(ERROR_MODE_COLLECT);

    /* import "file.sri" */
    const char *path = node->name;

    if (!path || !*path) {
        shriji_error(
            E_PARSE_02,
            "import",
            "empty import path",
            "use: import \"file.sri\""
        );
        shriji_set_error_mode(ERROR_MODE_IMMEDIATE);
        return value_null();
    }

    /* read file */
    FILE *fp = fopen(path, "rb");
    if (!fp) {
        shriji_error(
            E_PARSE_02,
            "import",
            "file not found",
            "check path or create file first"
        );
        shriji_set_error_mode(ERROR_MODE_IMMEDIATE);
        return value_null();
    }

    fseek(fp, 0, SEEK_END);
    long sz = ftell(fp);
    fseek(fp, 0, SEEK_SET);

    if (sz < 0) sz = 0;
    if (sz > 1024 * 1024) sz = 1024 * 1024;

    char *buf = (char *)malloc((size_t)sz + 1);
    if (!buf) {
        fclose(fp);
        shriji_set_error_mode(ERROR_MODE_IMMEDIATE);
        return value_null();
    }

    size_t nread = fread(buf, 1, (size_t)sz, fp);
    buf[nread] = '\0';
    fclose(fp);

    /* 🔥 FULL FILE PARSE */
    ASTNode *prog = parse_program(buf);
    free(buf);

    if (!prog) {
        shriji_error(
            E_PARSE_02,
            "import",
            "parse failed",
            "check full file syntax"
        );

        shriji_print_all_errors();
        shriji_set_error_mode(ERROR_MODE_IMMEDIATE);
        return value_null();
    }

    /* 🔥 RUN FULL PROGRAM */
    Value result = eval(prog, env, rt);

    /* 🔥 PRINT ALL ERRORS (ONE SHOT) */
    shriji_print_all_errors();

    /* 🔥 BACK TO NORMAL MODE */
    shriji_set_error_mode(ERROR_MODE_IMMEDIATE);

    return result;
}

case AST_COMMAND: {

    CommandType cmd = resolve_command(node->command_name);
    if (cmd == CMD_UNKNOWN) {
        shriji_error(
            E_PARSE_02,
            "command",
            "unknown command",
            "command not recognized"
        );
        return value_null();
    }

if (cmd == CMD_BOLO) {

    /* 🔥 GLOBAL ERROR GUARD */
    if (error_reported) {
        return value_null();
    }

    Value v = value_null();

    if (node->value)
        v = eval(node->value, env , rt);

/*  RUNTIME ERROR GUARD */
if (rt->error_flag) {
    return value_null();
}

if (v.type == VAL_STRING) {
    printf("%s\n", v.string ? v.string : "");
}
else if (v.type == VAL_NUMBER) {
    printf("%g\n", v.number);
}
else if (v.type == VAL_BOOL) {
    printf("%d\n", v.boolean);
}
else if (v.type == VAL_LIST) {
    printf("[");
    for (int i = 0; i < v.count; i++) {
        if (i > 0) printf(", ");

        if (v.items[i].type == VAL_NUMBER)
            printf("%g", v.items[i].number);
        else if (v.items[i].type == VAL_STRING)
            printf("\"%s\"", v.items[i].string);
        else if (v.items[i].type == VAL_BOOL)
            printf("%d", v.items[i].boolean);
        else
            printf("?");
    }
    printf("]\n");
}
else {
    /*  DOUBLE GUARD */
    if (!rt->error_flag && !error_reported) {
        printf("0\n");
    }
}



    smriti_session_set_last(v);

    value_free(&v);
    return value_null();
}

/* AI MODULES → ai_router (STEP-13 REAL OUTPUT) */
if (cmd == CMD_SAKHI || cmd == CMD_NIYU ||
    cmd == CMD_MIRA  || cmd == CMD_KAVYA) {

    Value v = value_null();
    if (node->value)
        v = eval(node->value, env , rt);

    const char *arg = "";
    char tmp[64] = {0};

    if (v.type == VAL_STRING && v.string)
        arg = strip_quotes(v.string);
    else if (v.type == VAL_NUMBER) {
        snprintf(tmp, sizeof(tmp), "%g", v.number);
        arg = tmp;
    }

    ShrijiBridgePacket pkt;
    pkt.text   = arg;
    pkt.lang   = SHRIJI_LANG_AUTO;
    pkt.source = SHRIJI_SRC_REPL;

    L3Response r = ai_router_dispatch(&pkt);
    if (r.text && *r.text)
        printf("%s\n", r.text);

    value_free(&v);
    return value_null();
}

    /* DEFAULT COMMAND */
    return value_number(execute_command(cmd));
}

default:
    return value_null();
}
}

/*──────────────────────────────────────────────────────────────
 | PROGRAM RUNNER
 *──────────────────────────────────────────────────────────────*/
Value run_program(ASTNode *program, Env *env, ShrijiRuntime *rt)
{
    if (!program || program->type != AST_PROGRAM)
        return value_null();

    runtime_reset(rt);
    event_bind_runtime(rt);
    return eval(program, env, rt);
}
#include <string.h>
#include <ctype.h>
#include <stdio.h>
#include <stdio.h>

#include "../include/normalize_engine.h"

#define MAX_WORDS 50
#define MAX_WORD_LEN 64

/* ================================
   BASIC UTILS
================================ */

static void to_lowercase(char *s)
{
    for (int i = 0; s[i]; i++)
        s[i] = tolower(s[i]);
}

static void trim_spaces(char *s)
{
    int i = 0, j = 0;

    while (isspace(s[i])) i++;

    while (s[i])
    {
        if (!(isspace(s[i]) && isspace(s[i+1])))
            s[j++] = s[i];
        i++;
    }

    if (j > 0 && isspace(s[j-1])) j--;
    s[j] = '\0';
}

/* ================================
   TYPO DICTIONARY
================================ */

typedef struct {
    const char *wrong;
    const char *correct;
} TypoRule;

static TypoRule rules[] = {
    {"reeor", "error"},
    {"eror", "error"},
    {"softwere", "software"},
    {"creat", "create"},
    {"keise", "kaise"},
    {"helo", "hello"},
    {"hy", "hi"},
    {"sotwere", "software"},
    {"cteate", "create"},
    {"by", "bye"},

};

static int rule_count = sizeof(rules)/sizeof(TypoRule);

/* ================================
   WORD SPLIT
================================ */

static void split_words(char *input, char words[][MAX_WORD_LEN], int *count)
{
    char *token = strtok(input, " ");
    int i = 0;

    while (token && i < MAX_WORDS)
    {
        strncpy(words[i], token, MAX_WORD_LEN - 1);
        words[i][MAX_WORD_LEN - 1] = '\0';
        i++;
        token = strtok(NULL, " ");
    }

    *count = i;
}

/* ================================
   APPLY TYPO RULES
================================ */

static void apply_typo_rules(char words[][MAX_WORD_LEN], int count)
{
    for (int i = 0; i < count; i++)
    {
        for (int j = 0; j < rule_count; j++)
        {
            if (strcmp(words[i], rules[j].wrong) == 0)
            {
                strncpy(words[i], rules[j].correct, MAX_WORD_LEN - 1);
            }
        }
    }
}

/* ================================
   JOIN WORDS BACK
================================ */

static void join_words(char *output, char words[][MAX_WORD_LEN], int count)
{
    output[0] = '\0';

    for (int i = 0; i < count; i++)
    {
        strncat(output, words[i], MAX_WORD_LEN - 1);
        if (i < count - 1)
            strncat(output, " ", 1);
    }
}

/* ================================
   SPECIAL CASE: EXIT NIRMAN
================================ */

static void fix_joined_commands(char *input)
{
    if (strstr(input, "exitnirman"))
    {
        strcpy(input, "exit nirman");
    }
}

/* ================================
   MAIN FUNCTION
================================ */
void normalize_input(char *input)
{

    if (!input || !*input)
        return;

    char temp[512];
    strncpy(temp, input, sizeof(temp)-1);
    temp[sizeof(temp)-1] = '\0';

    to_lowercase(temp);
    trim_spaces(temp);

    fix_joined_commands(temp);

    char words[MAX_WORDS][MAX_WORD_LEN];
    int count = 0;

    split_words(temp, words, &count);
    apply_typo_rules(words, count);

char result[512];
join_words(result, words, count);

/* safe copy back */
snprintf(input, 512, "%s", result);
}
#include "../include/smriti_personal.h"
#include <string.h>
#include <stdio.h>

#define MAX_KEYS 100

typedef struct {
    char key[32];
    char value[128];
} MemoryItem;

/* ===============================
   INTERNAL MEMORY (RAM)
   =============================== */
static MemoryItem memory[MAX_KEYS];
static int memory_count = 0;

/* ===============================
   FILE PATH
   =============================== */
static const char *SMRITI_FILE = ".shriji/smriti.db";

/* ===============================
   LOAD FROM FILE
   =============================== */
void smriti_personal_load(void)
{
    FILE *fp = fopen(SMRITI_FILE, "r");
    if (!fp) return;

    char line[256];

    while (fgets(line, sizeof(line), fp))
    {
        char *eq = strchr(line, '=');
        if (!eq) continue;

        *eq = '\0';

        char *key = line;
        char *value = eq + 1;

        /* remove newline */
        value[strcspn(value, "\n")] = 0;

        smriti_personal_set(key, value);
    }

    fclose(fp);
}

/* ===============================
   SAVE TO FILE
   =============================== */
void smriti_personal_save(void)
{
    FILE *fp = fopen(SMRITI_FILE, "w");
    if (!fp) return;

    for (int i = 0; i < memory_count; i++)
    {
        fprintf(fp, "%s=%s\n", memory[i].key, memory[i].value);
    }

    fclose(fp);
}

/* ===============================
   SET (WRITE)
   =============================== */
void smriti_personal_set(const char *key, const char *value)
{
    if (!key || !value) return;

    /* overwrite if exists */
    for (int i = 0; i < memory_count; i++)
    {
        if (strcmp(memory[i].key, key) == 0)
        {
            strncpy(memory[i].value, value, sizeof(memory[i].value)-1);
            memory[i].value[sizeof(memory[i].value)-1] = '\0';

            smriti_personal_save();
            return;
        }
    }

    /* new entry */
    if (memory_count < MAX_KEYS)
    {
        strncpy(memory[memory_count].key, key, sizeof(memory[memory_count].key)-1);
        strncpy(memory[memory_count].value, value, sizeof(memory[memory_count].value)-1);

        memory[memory_count].key[sizeof(memory[memory_count].key)-1] = '\0';
        memory[memory_count].value[sizeof(memory[memory_count].value)-1] = '\0';

        memory_count++;

        smriti_personal_save();
    }
}

/* ===============================
   GET (READ)
   =============================== */
const char* smriti_personal_get(const char *key)
{
    if (!key) return NULL;

    for (int i = 0; i < memory_count; i++)
    {
        if (strcmp(memory[i].key, key) == 0)
        {
            return memory[i].value;
        }
    }

    return NULL;
}

/* ===============================
   CLEAR (RAM ONLY)
   =============================== */
void smriti_personal_clear(void)
{
    memory_count = 0;
}
#include <stdio.h>
#include <string.h>
#include "../include/nirman.h"

void nirman_flow_process(char *line)
{
    int intent = nirman_detect_intent(line);

    /* STAGE 0 → START */
    if (intent == 2 && NIRMAN_CTX.flow_stage == 0)
    {
        printf("Aap kya banana chahte ho?\n");
        NIRMAN_CTX.flow_stage = 1;
        return;
    }

    /* STAGE 1 → GOAL CAPTURE */
    if (NIRMAN_CTX.flow_stage == 1)
    {
        if (strstr(line, "software") || strstr(line, "create"))
        {
            printf("Aap kya banana chahte ho?\n");
            return;
        }

        strncpy(NIRMAN_CTX.goal, line, sizeof(NIRMAN_CTX.goal) - 1);
        printf("[FLOW] Target set: %s\n", NIRMAN_CTX.goal);

        NIRMAN_CTX.flow_stage = 2;
        return;
    }

    /* STAGE 2 → ASK TYPE */
    if (NIRMAN_CTX.flow_stage == 2)
    {
        if (strstr(NIRMAN_CTX.goal, "calculator"))
        {
            printf("Aapko kis type ka calculator chahiye?\n");
            NIRMAN_CTX.flow_stage = 3;
            return;
        }
    }

    /* STAGE 3 → HANDLE TYPE */
    if (NIRMAN_CTX.flow_stage == 3)
    {
        if (strstr(NIRMAN_CTX.goal, "calculator"))
        {
            if (strstr(line, "normal"))
            {
                printf("[FLOW] Type selected: normal calculator\n");
                printf("Agla step: features define karte hain\n");

                NIRMAN_CTX.flow_stage = 4;
            }
            else if (strstr(line, "scientific"))
            {
                printf("[FLOW] Type selected: scientific calculator\n");
                printf("Agla step: features define karte hain\n");

                NIRMAN_CTX.flow_stage = 4;
            }
            else
            {
                printf("Type clear nahi hai (normal / scientific)\n");
            }
            return;
        }
    }

    /* STAGE 4 → NEXT STEP (placeholder) */
    if (NIRMAN_CTX.flow_stage == 4)
    {
        printf("[FLOW] Feature planning phase...\n");
        return;
    }
}
#include <stdio.h>
#include <string.h>
#include <stdlib.h>
#include <ctype.h>

#ifndef _WIN32
#include <readline/readline.h>
#include <readline/history.h>
#endif

#include "../include/error.h"
#include "../include/pragya_avastha.h"
#include "../include/fix_rules.h"

/* ===============================
   PREFILL (FOR EDIT MODE)
   =============================== */

static char safe_buffer[256];
static char *prefill_text = NULL;

#ifndef _WIN32
static int prefill_hook(void)
{
    if (prefill_text)
    {
        rl_insert_text(prefill_text);
        rl_point = strlen(prefill_text);
        prefill_text = NULL;
    }
    return 0;
}
#endif

/* ===============================
   USER INPUT (STRICT)
   =============================== */

static char get_user_choice()
{
#ifdef _WIN32
    char buffer[16];

    if (!fgets(buffer, sizeof(buffer), stdin))
        return 'i';

    char c = tolower(buffer[0]);

    if (c == 'e') return 'e';
    if (c == 'i') return 'i';

    return 'i';
#else
    char *line = readline("");

    if (!line || !*line)
    {
        free(line);
        return 'i';
    }

    char c = tolower(line[0]);
    free(line);

    if (c == 'e') return 'e';
    if (c == 'i') return 'i';

    return 'i';
#endif
}

/* ===============================
   POINTER DISPLAY
   =============================== */

static void print_pointer(const char *text, int pos)
{
    if (!text || pos < 0) return;

    printf("%s\n", text);

    for (int i = 0; i < pos; i++)
        printf(" ");

    printf("^\n\n");
}

/* ===============================
   BASIC VALIDATION
   =============================== */

static int contains_invalid_token(const char *input)
{
    for (int i = 0; input[i]; i++)
    {
        char c = input[i];

        if (isalnum(c) || isspace(c) || strchr("+-*/()._{}[]=<>!\"'", c))
            continue;

        return 1;
    }
    return 0;
}

static int has_unmatched_brackets(const char *input)
{
    int count = 0;

    for (int i = 0; input[i]; i++)
    {
        if (input[i] == '(') count++;
        else if (input[i] == ')') count--;
    }

    return count != 0;
}

static int is_valid_input(const char *input)
{
    if (!input || !*input) return 0;
    if (contains_invalid_token(input)) return 0;
    if (has_unmatched_brackets(input)) return 0;
    return 1;
}

/* ===============================
   SIMPLE OPERATOR CLEAN FIX
   =============================== */

static int filter_operators(const char *input, char *output)
{
    int i, j = 0;
    int last_was_op = 0;

    for (i = 0; input[i]; i++)
    {
        char c = input[i];

        if (strchr("+-*/", c))
        {
            if (last_was_op) continue;

            output[j++] = c;
            last_was_op = 1;
        }
        else
        {
            output[j++] = c;
            last_was_op = 0;
        }
    }

    output[j] = '\0';
    return strcmp(input, output) != 0;
}

/* ===============================
   MAIN ERROR INTELLIGENCE
   =============================== */

void shriji_error_intelligence(
    PragyaAvastha *avastha,
    const ShrijiErrorInfo *err,
    void *niyu_result)
{
    (void)niyu_result;

    if (!avastha || !err || !avastha->raw_text)
        return;

    const char *text = avastha->raw_text;

    /* ===== VALIDATION ===== */

    if (!is_valid_input(text))
    {
        avastha->stop_execution = 1;

        int pos = (err && err->col > 0) ? err->col - 1 : 0;
        print_pointer(text, pos);
        return;
    }

    /* ===== RULE BASED FIX ===== */

    FixRule *rule = shriji_get_rule_for_error(err->code);

    if (rule && rule->apply_fix)
    {
        char buffer[256];
        snprintf(buffer, sizeof(buffer), "%s", text);

        if (rule->apply_fix(buffer, sizeof(buffer), err))
        {
            if (rule->safety == FIX_SAFE_DETERMINISTIC)
            {
                printf("\n[Auto Fix Applied]\n");
                printf("Reason: deterministic rule\n");
                printf("Result: %s\n\n", buffer);

                avastha->has_correction = 1;
                avastha->stop_execution = 1;

                snprintf(avastha->corrected_text,
                         sizeof(avastha->corrected_text),
                         "%s", buffer);

                return;
            }

            printf("\nSuggestion: %s\n", buffer);
            printf("\nEdit (E) / Ignore (I): ");

            char choice = get_user_choice();

            if (choice == 'e')
            {
                avastha->stop_execution = 1;

                printf("\nEdit:\n%s\n", text);

#ifndef _WIN32
                snprintf(safe_buffer, sizeof(safe_buffer), "%s", text);
                prefill_text = safe_buffer;

                rl_startup_hook = prefill_hook;
#endif
                return;
            }

            avastha->stop_execution = 1;
            return;
        }
    }

    /* ===== SIMPLE FILTER FIX ===== */

    char filtered[256];

    if (filter_operators(text, filtered))
    {
        printf("\nSuggestion: %s\n", filtered);
        printf("\nEdit (E) / Ignore (I): ");

        char choice = get_user_choice();

        if (choice == 'e')
        {
            avastha->stop_execution = 1;

            printf("\nEdit:\n%s\n", text);

#ifndef _WIN32
            snprintf(safe_buffer, sizeof(safe_buffer), "%s", text);
            prefill_text = safe_buffer;

            rl_startup_hook = prefill_hook;
#endif
            return;
        }

        avastha->stop_execution = 1;
        return;
    }

    /* ===== FALLBACK POINTER ===== */

    int pos = (err && err->col > 0) ? err->col - 1 : 0;

    print_pointer(text, pos);

    avastha->stop_execution = 1;
}
#include <stdio.h>
#include <string.h>

/*
    Build readable example from user input
*/

void shriji_build_example(const char *input, char *out, int size)
{
    int i = 0;
    int j = 0;

    while (input[i] && j < size - 1)
    {
        char c = input[i];

        if (c == '+' || c == '-' || c == '*' || c == '/')
        {
            if (j < size - 3)
            {
                out[j++] = ' ';
                out[j++] = c;
                out[j++] = ' ';
            }
        }
        else
        {
            out[j++] = c;
        }

        i++;
    }

    out[j] = '\0';
}
#include <string.h>
#include <ctype.h>

#include "../include/ai_intent.h"

/* lowercase helper */
static void to_lower(char *s)
{
    for (; *s; s++)
        *s = (char)tolower(*s);
}

/* simple substring match (L3 — lightweight, safe) */
static int has_word(const char *msg, const char *w)
{
    return strstr(msg, w) != NULL;
}

AIIntent ai_detect_intent(const char *message)
{
    if (!message || !*message)
        return AI_INTENT_UNKNOWN;

    char buf[256];
    strncpy(buf, message, sizeof(buf) - 1);
    buf[sizeof(buf) - 1] = '\0';
    to_lower(buf);

    /*──────────────────────────────────────────────
      EXPLAIN / SAMJHAO (LOGIC / TEACHING)
      → Niyu (primary explain engine)
    ──────────────────────────────────────────────*/
    if (has_word(buf, "samjhao") ||
        has_word(buf, "samjha")  ||
        has_word(buf, "samajhao")||
        has_word(buf, "samajh")  ||
        has_word(buf, "explain") ||
        has_word(buf, "what is") ||
        has_word(buf, "kya hota"))
        return AI_INTENT_EXPLAIN;

    /*──────────────────────────────────────────────
      WHY / REASON
    ──────────────────────────────────────────────*/
    if (has_word(buf, "why")  ||
        has_word(buf, "kyun") ||
        has_word(buf, "kyu")  ||
        has_word(buf, "kyon"))
        return AI_INTENT_WHY;

    /*──────────────────────────────────────────────
      CALC / MATH
    ──────────────────────────────────────────────*/
    if (has_word(buf, "calculate") ||
        has_word(buf, "solve")     ||
        has_word(buf, "math")      ||
        has_word(buf, "hisab"))
        return AI_INTENT_CALC;

    /*──────────────────────────────────────────────
      EMOTION / FEELING
    ──────────────────────────────────────────────*/
    if (has_word(buf, "sad")    ||
        has_word(buf, "dukhi")  ||
        has_word(buf, "lonely") ||
        has_word(buf, "akela")  ||
        has_word(buf, "dard"))
        return AI_INTENT_EMOTION;

    return AI_INTENT_UNKNOWN;
}
#include "../../include/lang/grammar.h"
#include "../../include/token.h"

/*
 * SHRIJILANG GRAMMAR CORE
 * Phase-1: token presence validation only
 * No parser logic change yet
 */

void shriji_register_grammar_core(void)
{
    /*
     * IMPORTANT:
     * This function is called once from parse_program()
     * Do NOT put parser logic here.
     * Only grammar / token registration / validation.
     */

#ifdef SHRIJI_ENABLE_MASTER_TOKENS
    /* Phase-1: sanity check hook */
    /* (empty by design for now) */
#endif
}

/* Future phases — intentionally empty */

void shriji_register_grammar_logic(void) {}
void shriji_register_grammar_incdec(void) {}
void shriji_register_grammar_loop(void) {}
void shriji_register_grammar_func(void) {}
#include "lang/grammar.h"

/* Phase-1: empty implementations
 * These will be filled in Phase-2
 */

void shriji_register_grammar_core(void) {
    /* core grammar hooks (placeholder) */
}

void shriji_register_grammar_logic(void) {
    /* && || ! (placeholder) */
}

void shriji_register_grammar_incdec(void) {
    /* ++ -- (placeholder) */
}

void shriji_register_grammar_loop(void) {
    /* jabtak / for (placeholder) */
}

void shriji_register_grammar_func(void) {
    /* kaam / wapas (placeholder) */
}
#include "../../../include/lang/grammar_core.h"

/*
 * Phase-2 Grammar Core
 * Classification only (no parsing / execution)
 */

GrammarKind shriji_grammar_kind(TokenType t)
{
    switch (t) {

        /* ───────── STATEMENTS ───────── */
        case TOKEN_MAVI:
        case TOKEN_KAAM:
        case TOKEN_AGAR:
        case TOKEN_JABTAK:
        case TOKEN_WAPAS:
        case TOKEN_IMPORT:
        case TOKEN_WARNA:
        case TOKEN_RUKJA:
        case TOKEN_CHALU:
            return GRAMMAR_STATEMENT;

        /* ───────── LITERALS ───────── */
        case TOKEN_NUMBER:
        case TOKEN_STRING:
        case TOKEN_TRUE:
        case TOKEN_FALSE:
            return GRAMMAR_LITERAL;

        /* ───────── INC / DEC ───────── */
        case TOKEN_PLUSPLUS:
        case TOKEN_MINUSMINUS:
            return GRAMMAR_INCDEC;

        /* ───────── LOGICAL / ARITHMETIC OPERATORS ───────── */
        case TOKEN_PLUS:
        case TOKEN_MINUS:
        case TOKEN_STAR:
        case TOKEN_SLASH:
        case TOKEN_MOD:

        case TOKEN_EQEQ:
        case TOKEN_NEQ:
        case TOKEN_GT:
        case TOKEN_LT:
        case TOKEN_GTE:
        case TOKEN_LTE:

        case TOKEN_AND:
        case TOKEN_OR:
        case TOKEN_NOT:
            return GRAMMAR_OPERATOR;

        default:
            return GRAMMAR_UNKNOWN;
    }
}

const char *shriji_grammar_kind_name(GrammarKind k)
{
    switch (k) {
        case GRAMMAR_STATEMENT:  return "STATEMENT";
        case GRAMMAR_EXPRESSION: return "EXPRESSION";
        case GRAMMAR_OPERATOR:   return "OPERATOR";
        case GRAMMAR_LITERAL:    return "LITERAL";
        case GRAMMAR_INCDEC:     return "INCDEC";
        case GRAMMAR_RESERVED:   return "RESERVED";
        default:                 return "UNKNOWN";
    }
}
#include <stdio.h>
#include <string.h>

#include "../include/error.h"
#include "../include/runtime.h"
#include "../gyaan/core/gyaan_engine.h"

/*──────────────────────────────────────────────*/
int error_reported = 0;

static ShrijiErrorMode CURRENT_ERROR_MODE = ERROR_MODE_IMMEDIATE;

ShrijiErrorInfo LAST_ERROR = {0};


/* 🔥 ERROR STORAGE */
#define MAX_ERRORS 100

typedef struct {
    int line;
    int col;
    char code[32];
    char context[64];
    char message[256];
    char hint[256];
} StoredError;

static StoredError ERROR_LIST[MAX_ERRORS];
static int ERROR_COUNT = 0;


/*──────────────────────────────────────────────*/
void shriji_set_error_mode(ShrijiErrorMode mode)
{
    CURRENT_ERROR_MODE = mode;
}


/*──────────────────────────────────────────────*/
static const char *error_code_str(ShrijiErrorCode code)
{
    switch (code) {

        case E_ASSIGN_01: return "E_ASSIGN_01";
        case E_ASSIGN_02: return "E_ASSIGN_02";

        case E_PARSE_01: return "E_PARSE_01";
        case E_PARSE_02: return "E_PARSE_02";
        case E_PARSE_EXTRA_OPERATOR: return "E_PARSE_EXTRA_OPERATOR";
        case E_PARSE_MISSING_OPERATOR: return "E_PARSE_MISSING_OPERATOR";
        case E_PARSE_DOUBLE_OPERATOR: return "E_PARSE_DOUBLE_OPERATOR";
        case E_PARSE_OPERATOR_START: return "E_PARSE_OPERATOR_START";
        case E_PARSE_OPERATOR_END: return "E_PARSE_OPERATOR_END";
        case E_PARSE_OPERATOR_CHAIN: return "E_PARSE_OPERATOR_CHAIN";
        case E_PARSE_MISSING_OPERAND: return "E_PARSE_MISSING_OPERAND";

        case E_PARSE_UNMATCHED_PAREN: return "E_PARSE_UNMATCHED_PAREN";
        case E_PARSE_BRACKET_MISSING: return "E_PARSE_BRACKET_MISSING";
        case E_PARSE_BRACKET_EXTRA: return "E_PARSE_BRACKET_EXTRA";
        case E_PARSE_INVALID_TOKEN: return "E_PARSE_INVALID_TOKEN";

        case E_IF_01: return "E_IF_01";

        case E_RUNTIME_01: return "E_RUNTIME_01";
        case E_RUNTIME_DIV_ZERO: return "E_RUNTIME_DIV_ZERO";
        case E_RUNTIME_TYPE_MISMATCH: return "E_RUNTIME_TYPE_MISMATCH";
        case E_RUNTIME_LOOP_LIMIT: return "E_RUNTIME_LOOP_LIMIT";
        case E_RUNTIME_INDEX_ERROR: return "E_RUNTIME_INDEX_ERROR";

        default: return "E_UNKNOWN";
    }
}


/*──────────────────────────────────────────────*/
static const char *error_message_for_code(ShrijiErrorCode code)
{
    switch (code) {

        case E_PARSE_DOUBLE_OPERATOR:
            return "Do operators ek saath aa gaye hain.";

        case E_PARSE_OPERATOR_CHAIN:
            return "Operator sequence galat lag raha hai.";

        case E_PARSE_OPERATOR_START:
            return "Expression operator se start nahi ho sakta.";

        case E_PARSE_OPERATOR_END:
            return "Expression operator par end ho gaya.";

        case E_PARSE_MISSING_OPERAND:
            return "Operator ke liye value missing hai.";

        default:
            return NULL;
    }
}


/*──────────────────────────────────────────────*/
static void store_error(
    ShrijiErrorCode code,
    const char *context,
    const char *message,
    const char *hint,
    int line,
    int col
)
{
    if (ERROR_COUNT >= MAX_ERRORS) return;

    StoredError *e = &ERROR_LIST[ERROR_COUNT++];

    snprintf(e->code, sizeof(e->code), "%s", error_code_str(code));
    snprintf(e->context, sizeof(e->context), "%s", context ? context : "");
    snprintf(e->message, sizeof(e->message), "%s", message ? message : "");
    snprintf(e->hint, sizeof(e->hint), "%s", hint ? hint : "");

    e->line = line;
    e->col = col;
}


/*──────────────────────────────────────────────*/
void shriji_error(
    ShrijiErrorCode code,
    const char *context,
    const char *message,
    const char *hint
)
{
    error_reported = 1;

    if (current_runtime) {
        current_runtime->error_flag = 1;
    }

    const char *auto_msg = error_message_for_code(code);
    if (auto_msg)
        message = auto_msg;

    LAST_ERROR.code = code;
    LAST_ERROR.context = context;
    LAST_ERROR.message = message;
    LAST_ERROR.hint = hint;
    LAST_ERROR.function = NULL;
    LAST_ERROR.has_location = 0;

    if (CURRENT_ERROR_MODE == ERROR_MODE_IMMEDIATE) {

        printf("\n⛔ %s | %s → %s\n",
               error_code_str(code),
               context ? context : "",
               message ? message : "");

        if (hint && hint[0])
            printf("   hint: %s\n", hint);

        printf("\n");

        const ShrijiErrorInfo *last = shriji_last_error();
        if (last) {
            gyaan_print(last);
        }

    } else {
        store_error(code, context, message, hint, 0, 0);
    }
}


/*──────────────────────────────────────────────*/
void shriji_error_at_full(
    Token tok,
    ShrijiErrorCode code,
    const char *context,
    const char *function,
    const char *message,
    const char *hint
)
{
    error_reported = 1;

    if (current_runtime) {
        current_runtime->error_flag = 1;
    }

    const char *auto_msg = error_message_for_code(code);
    if (auto_msg)
        message = auto_msg;

    LAST_ERROR.code = code;
    LAST_ERROR.context = context;
    LAST_ERROR.message = message;
    LAST_ERROR.hint = hint;
    LAST_ERROR.function = function;
    LAST_ERROR.has_location = 1;
    LAST_ERROR.line = tok.line;
    LAST_ERROR.col = tok.col;

    if (CURRENT_ERROR_MODE == ERROR_MODE_IMMEDIATE) {

        printf("\n⛔ %s | %s → %s (line %d, col %d)\n",
               error_code_str(code),
               context ? context : "",
               message ? message : "",
               tok.line,
               tok.col);

        if (hint && hint[0])
            printf("   hint: %s\n", hint);

        printf("\n");

        const ShrijiErrorInfo *last = shriji_last_error();
        if (last) {
            gyaan_print(last);
        }

    } else {
        store_error(code, context, message, hint, tok.line, tok.col);
    }
}


/*──────────────────────────────────────────────*/
void shriji_error_at(
    Token tok,
    ShrijiErrorCode code,
    const char *context,
    const char *message,
    const char *hint
)
{
    shriji_error_at_full(
        tok,
        code,
        context,
        "unknown",
        message,
        hint
    );
}


/*──────────────────────────────────────────────*/
void shriji_print_all_errors(void)
{
    if (ERROR_COUNT == 0) return;

    printf("\nKRST ANALYSIS REPORT\n");
    printf("──────────────────────────────\n");

    for (int i = 0; i < ERROR_COUNT; i++) {
        StoredError *e = &ERROR_LIST[i];

        printf("%d) %s | %s → %s",
               i + 1,
               e->code,
               e->context,
               e->message);

        if (e->line > 0)
            printf(" (line %d, col %d)", e->line, e->col);

        printf("\n");

        if (e->hint[0])
            printf("   hint: %s\n", e->hint);

        printf("\n");
    }

    printf("──────────────────────────────\n\n");

    ERROR_COUNT = 0;

    const ShrijiErrorInfo *last = shriji_last_error();
    if (last) {
        gyaan_print(last);
    }
}


/*──────────────────────────────────────────────*/
const ShrijiErrorInfo *shriji_last_error(void)
{
    if (!error_reported)
        return NULL;

    return &LAST_ERROR;
}
#include "../include/smriti_session.h"
#include <string.h>
#include <stdlib.h>

/* ===============================
   INTERNAL STORAGE
   =============================== */

static char last_input[256] = {0};
static char last_output[256] = {0};

static Value last_value;
static int has_last = 0;

/* ===============================
   LAST INPUT
   =============================== */

void smriti_session_set_last_input(const char *text)
{
    if (!text) return;

    strncpy(last_input, text, sizeof(last_input) - 1);
    last_input[sizeof(last_input) - 1] = '\0';
}

const char* smriti_session_get_last_input(void)
{
    return last_input;
}

/* ===============================
   LAST OUTPUT
   =============================== */

void smriti_session_set_last_output(const char *text)
{
    if (!text) return;

    strncpy(last_output, text, sizeof(last_output) - 1);
    last_output[sizeof(last_output) - 1] = '\0';
}

const char* smriti_session_get_last_output(void)
{
    return last_output;
}

/* ===============================
   LAST VALUE (UPGRADED)
   =============================== */

void smriti_session_set_last(Value v)
{
    if (has_last)
        value_free(&last_value);

    last_value = value_copy(v);
    has_last = 1;
}

Value smriti_session_get_last(void)
{
    if (!has_last)
        return value_null();

    return value_copy(last_value);
}

/* ===============================
   CLEAR SESSION
   =============================== */

void smriti_session_clear(void)
{
    last_input[0] = '\0';
    last_output[0] = '\0';

    if (has_last)
        value_free(&last_value);

    has_last = 0;
}
#include <string.h>

#include "../include/event.h"
#include "../include/state.h"
#include "../include/runtime.h"

/* Internal runtime binding */
static ShrijiRuntime *BOUND_RUNTIME = NULL;

/* Bind runtime once per execution */
void event_bind_runtime(ShrijiRuntime *rt)
{
    BOUND_RUNTIME = rt;
}

/* Fire event */
void event_fire(EventType type, const char *context)
{
    (void)context;

    if (!BOUND_RUNTIME)
        return;

    switch (type)
    {
        case EVENT_EXECUTION_SUCCESS:
        case EVENT_ASSIGNMENT:
            state_on_success(&BOUND_RUNTIME->state);
            break;

        case EVENT_SHUNYA_ACCESS:
        case EVENT_ERROR:
            state_on_error(&BOUND_RUNTIME->state);
            break;

        case EVENT_FATAL_ERROR:
            state_on_error(&BOUND_RUNTIME->state);
            state_on_error(&BOUND_RUNTIME->state);
            break;

        case EVENT_EXECUTION_BLOCKED:
            BOUND_RUNTIME->state.safety = STATE_CRITICAL;
            BOUND_RUNTIME->state.human  = HUMAN_MANDATORY;
            break;

        default:
            break;
    }
}
#include <stdio.h>
#include <string.h>
#include <stdlib.h>

#include "../include/ast.h"

/*──────────────────────────────────────────────────────────────
 |  SHRIJILANG — AST CONSTRUCTOR IMPLEMENTATION (CLEAN & STABLE)
 *──────────────────────────────────────────────────────────────*/

/*──────────────────────────────────────────────────────────────
 | SECTION A — SAFE STRING HELPER
 *──────────────────────────────────────────────────────────────*/
static void safe_copy(char *dest, const char *src, int max)
{
    if (!src) {
        dest[0] = '\0';
        return;
    }
    strncpy(dest, src, max - 1);
    dest[max - 1] = '\0';
}

/*──────────────────────────────────────────────────────────────
 | SECTION B — CORE NODES
 *──────────────────────────────────────────────────────────────*/
ASTNode *new_number_node(double value)
{
    ASTNode *node = calloc(1, sizeof(ASTNode));
    if (!node) return NULL;

    node->type = AST_NUMBER;
    node->number_value = value;
    node->chakra_state = CHAKRA_OK;
    return node;
}

ASTNode *new_bool_node(int value)
{
    ASTNode *node = calloc(1, sizeof(ASTNode));
    if (!node) return NULL;

    node->type = AST_BOOL;
    node->bool_value = value;
    node->chakra_state = CHAKRA_OK;

    return node;
}

ASTNode *new_identifier_node(const char *name)
{
    ASTNode *node = calloc(1, sizeof(ASTNode));
    if (!node) return NULL;

    node->type = AST_IDENTIFIER;
    safe_copy(node->name, name, 128);
    node->chakra_state = CHAKRA_OK;
    return node;
}

ASTNode *new_string_node(const char *str)
{
    ASTNode *node = calloc(1, sizeof(ASTNode));
    if (!node) return NULL;

    node->type = AST_STRING;
    safe_copy(node->string_value, str, 256);
    node->chakra_state = CHAKRA_OK;
    return node;
}

ASTNode *new_binary_node(char op, ASTNode *l, ASTNode *r)
{
    ASTNode *node = calloc(1, sizeof(ASTNode));
    if (!node) return NULL;

    node->type = AST_BINARY;
    node->op[0] = op;
    node->left = l;
    node->right = r;
    node->chakra_state = CHAKRA_OK;
    return node;
}

ASTNode *new_unary_node(Token op, ASTNode *expr)
{
    ASTNode *node = calloc(1, sizeof(ASTNode));
    if (!node) return NULL;

    node->type = AST_UNARY;

    /* store operator */
    node->op[0] = op.start[0];
    node->op[1] = '\0';

    node->left = expr;

    node->chakra_state = CHAKRA_OK;
    return node;
}
/*──────────────────────────────────────────────────────────────
 | SECTION B.1 — ASSIGNMENT / UPDATE
 *──────────────────────────────────────────────────────────────*/
ASTNode *new_assignment_node(const char *name, ASTNode *val)
{
    ASTNode *node = calloc(1, sizeof(ASTNode));
    if (!node) return NULL;

    node->type = AST_ASSIGNMENT;
    safe_copy(node->name, name, 128);
    node->value = val;
    node->chakra_state = CHAKRA_OK;
    return node;
}

ASTNode *new_update_node(const char *name, ASTNode *val)
{
    ASTNode *node = calloc(1, sizeof(ASTNode));
    if (!node) return NULL;

    node->type = AST_UPDATE;
    safe_copy(node->name, name, 128);
    node->value = val;
    node->chakra_state = CHAKRA_OK;
    return node;
}

/*──────────────────────────────────────────────────────────────
 | SECTION C — AI NODES
 *──────────────────────────────────────────────────────────────*/
#define AI_NODE(fn, type_id)                      \
ASTNode *fn(ASTNode *msg) {                       \
    ASTNode *node = calloc(1, sizeof(ASTNode));   \
    if (!node) return NULL;                       \
    node->type = type_id;                         \
    node->value = msg;                            \
    node->chakra_state = CHAKRA_OK;               \
    return node;                                  \
}

AI_NODE(new_sakhi_node, AST_SAKHI)
AI_NODE(new_niyu_node,  AST_NIYU)
AI_NODE(new_shiri_node, AST_SHIRI)
AI_NODE(new_mira_node,  AST_MIRA)
AI_NODE(new_kavya_node, AST_KAVYA)

/*──────────────────────────────────────────────────────────────
 | SECTION D — COMMAND + IMPORT
 *──────────────────────────────────────────────────────────────*/
ASTNode *new_command_node(const char *cmd, ASTNode *arg)
{
    ASTNode *node = calloc(1, sizeof(ASTNode));
    if (!node) return NULL;

    node->type = AST_COMMAND;
    safe_copy(node->command_name, cmd, 128);
    node->value = arg;
    node->chakra_state = CHAKRA_OK;
    return node;
}

ASTNode *new_import_node(const char *module_name)
{
    ASTNode *node = calloc(1, sizeof(ASTNode));
    if (!node) return NULL;

    node->type = AST_IMPORT;
    safe_copy(node->name, module_name, 128);
    node->chakra_state = CHAKRA_OK;
    return node;
}

/*──────────────────────────────────────────────────────────────
 | SECTION D.1 — BLOCK / PROGRAM
 *──────────────────────────────────────────────────────────────*/
ASTNode *new_block_node(ASTNode **stmts, int count)
{
    ASTNode *node = calloc(1, sizeof(ASTNode));
    if (!node) return NULL;

    node->type = AST_BLOCK;
    node->statements = stmts;
    node->stmt_count = count;
    node->chakra_state = CHAKRA_OK;
    return node;
}

ASTNode *new_program_node(ASTNode **stmts, int count)
{
    ASTNode *node = calloc(1, sizeof(ASTNode));
    if (!node) return NULL;

    node->type = AST_PROGRAM;
    node->statements = stmts;
    node->stmt_count = count;
    node->chakra_state = CHAKRA_OK;
    return node;
}

/*──────────────────────────────────────────────────────────────
 | SECTION D.2 — CONTROL FLOW
 *──────────────────────────────────────────────────────────────*/
ASTNode *new_if_node(ASTNode *condition, ASTNode *then_block)
{
    ASTNode *node = calloc(1, sizeof(ASTNode));
    if (!node) return NULL;

    node->type = AST_IF;
    node->condition = condition;
    node->left = then_block;     /* then-block */
    node->chakra_state = CHAKRA_OK;
    return node;
}

ASTNode *new_while_node(ASTNode *condition, ASTNode *body)
{
    ASTNode *node = calloc(1, sizeof(ASTNode));
    if (!node) return NULL;

    node->type = AST_WHILE;
    node->condition = condition;
    node->body = body;
    node->chakra_state = CHAKRA_OK;
    return node;
}

/*──────────────────────────────────────────────────────────────
 | SECTION D.2.3 — LOOP CONTROL (rukja / chalu)
 *──────────────────────────────────────────────────────────────*/
ASTNode *new_break_node(void)
{
    ASTNode *node = calloc(1, sizeof(ASTNode));
    if (!node) return NULL;

    node->type = AST_BREAK;
    node->chakra_state = CHAKRA_OK;
    return node;
}

ASTNode *new_continue_node(void)
{
    ASTNode *node = calloc(1, sizeof(ASTNode));
    if (!node) return NULL;

    node->type = AST_CONTINUE;
    node->chakra_state = CHAKRA_OK;
    return node;
}

/*──────────────────────────────────────────────────────────────
 | SECTION D.3 — FUNCTIONS
 *──────────────────────────────────────────────────────────────*/
ASTNode *new_function_node(
    const char *name,
    ASTNode **params,
    int param_count,
    ASTNode *body
) {
    ASTNode *node = calloc(1, sizeof(ASTNode));
    if (!node) return NULL;

    node->type = AST_FUNCTION;
    safe_copy(node->function_name, name, 128);

    node->params = params;
    node->param_count = param_count;
    node->body = body;

    node->chakra_state = CHAKRA_OK;
    return node;
}

ASTNode *new_rachna_node(
    const char *name,
    ASTNode **params,
    int param_count,
    ASTNode *body
)
{
    ASTNode *node = calloc(1, sizeof(ASTNode));
    if (!node) return NULL;

    node->type = AST_RACHNA;

    /* name */
    safe_copy(node->rachna_name, name, 128);

    /* params */
    node->rachna_params = params;
    node->rachna_param_count = param_count;

    /* body */
    node->rachna_body = body;

    node->chakra_state = CHAKRA_OK;
    return node;
}

ASTNode *new_call_node(
    const char *name,
    ASTNode **args,
    int arg_count
) {
    ASTNode *node = calloc(1, sizeof(ASTNode));
    if (!node) return NULL;

    node->type = AST_CALL;
    safe_copy(node->function_name, name, 128);

    node->args = args;
    node->arg_count = arg_count;

    node->chakra_state = CHAKRA_OK;
    return node;
}

ASTNode *new_return_node(ASTNode *expr)
{
    ASTNode *node = calloc(1, sizeof(ASTNode));
    if (!node) return NULL;

    node->type = AST_RETURN;
    node->return_expr = expr;
    node->chakra_state = CHAKRA_OK;
    return node;
}

/*──────────────────────────────────────────────
 | SECTION E — MEMORY CLEANUP
 *──────────────────────────────────────────────*/
void ast_free(ASTNode *node)
{
    if (!node) return;

    ast_free(node->left);
    ast_free(node->right);
    ast_free(node->value);
    ast_free(node->condition);
    ast_free(node->else_block);

    ast_free(node->body);
    ast_free(node->return_expr);

 /* LIST */
    if (node->elements) {
    for (int i = 0; i < node->element_count; i++) {
       ast_free(node->elements[i]);
    }
       free(node->elements);
    }
    if (node->statements) {
        for (int i = 0; i < node->stmt_count; i++)
            ast_free(node->statements[i]);
        free(node->statements);
    }

    if (node->params) {
        for (int i = 0; i < node->param_count; i++)
            ast_free(node->params[i]);
        free(node->params);
    }

    if (node->args) {
        for (int i = 0; i < node->arg_count; i++)
            ast_free(node->args[i]);
        free(node->args);
    }
/* DICT */
    if (node->dict_keys) {
        for (int i = 0; i < node->dict_count; i++) {
            ast_free(node->dict_keys[i]);
        }
        free(node->dict_keys);
    }

    if (node->dict_values) {
        for (int i = 0; i < node->dict_count; i++) {
            ast_free(node->dict_values[i]);
        }
        free(node->dict_values);
    }
    free(node);
}

/*──────────────────────────────────────────────
 | SECTION F — LOGICAL NOT NODE (nahi)
 *──────────────────────────────────────────────*/
ASTNode *new_not_node(ASTNode *expr)
{
    ASTNode *node = calloc(1, sizeof(ASTNode));
    if (!node) return NULL;

    node->type = AST_NOT;
    node->left = expr;
    node->right = NULL;
    node->chakra_state = CHAKRA_OK;
    return node;
}

/*────────────────────────────────────────
  Inc / Dec (++ / --)
────────────────────────────────────────*/

ASTNode *new_inc_node(ASTNode *expr, int is_prefix)
{
    ASTNode *n = calloc(1, sizeof(ASTNode));
    n->type = AST_INC;

    n->left = expr;          // operand (++x or x++)
    n->is_prefix = is_prefix;

    return n;
}

ASTNode *new_dec_node(ASTNode *expr, int is_prefix)
{
    ASTNode *n = calloc(1, sizeof(ASTNode));
    n->type = AST_DEC;

    n->left = expr;          // operand (--x or x--)
    n->is_prefix = is_prefix;

    return n;
}

/*──────────────────────────────────────────────
  LIST / INDEX NODES
──────────────────────────────────────────────*/

ASTNode *new_list_node(ASTNode **elements, int count)
{
    ASTNode *node = (ASTNode *)malloc(sizeof(ASTNode));
    if (!node) return NULL;

    memset(node, 0, sizeof(ASTNode));
    node->type = AST_LIST;

    node->elements = elements;
    node->element_count = count;
    node->chakra_state = CHAKRA_OK;
    return node;
}

ASTNode *new_index_node(ASTNode *target, ASTNode *index_expr)
{
    ASTNode *node = (ASTNode *)malloc(sizeof(ASTNode));
    if (!node) return NULL;

    memset(node, 0, sizeof(ASTNode));
    node->type = AST_INDEX;

    node->index_target = target;
    node->index_expr   = index_expr;

    return node;
}

ASTNode *new_index_update_node(ASTNode *target, ASTNode *index_expr, ASTNode *value)
{
    ASTNode *node = (ASTNode *)calloc(1, sizeof(ASTNode));
    node->type = AST_INDEX_UPDATE;

    node->index_target = target;
    node->index_expr = index_expr;
    node->index_value = value;

    return node;
}
/*──────────────────────────────────────────────
  DICT / MAP NODES
──────────────────────────────────────────────*/
ASTNode *new_dict_node(ASTNode **keys, ASTNode **values, int count)
{
    ASTNode *node = (ASTNode *)malloc(sizeof(ASTNode));
    if (!node) return NULL;

    memset(node, 0, sizeof(ASTNode));
    node->type = AST_DICT;

    node->dict_keys = keys;
    node->dict_values = values;
    node->dict_count = count;
    node->chakra_state = CHAKRA_OK;
    return node;
}
#include <stdio.h>
#include <string.h>

#include "../include/fix_interactive.h"
#include "../include/error.h"

/* =========================================================
   MEMORY — LAST INPUT TRACKING
========================================================= */

static char last_input[512] = {0};

/* =========================================================
   SAFE LINE INPUT
========================================================= */

static int read_line(char *buffer, int size)
{
    if (!fgets(buffer, size, stdin))
        return 0;

    buffer[strcspn(buffer, "\n")] = '\0';
    return 1;
}

/* =========================================================
   SAME INPUT GUARD
========================================================= */

static int is_same_input(const char *input)
{
    if (strcmp(last_input, input) == 0)
    {
        printf("\n[SYSTEM] Same input detected. Please modify it.\n");
        return 1;
    }

    strncpy(last_input, input, sizeof(last_input) - 1);
    last_input[sizeof(last_input) - 1] = '\0';

    return 0;
}

/* =========================================================
   FIX MENU CORE (SHARED)
========================================================= */

static int show_fix_menu(const char *input, char *output, int size)
{
    while (1)
    {
        printf("\nFix Options:\n\n");
        printf("M) Manual edit\n");
        printf("E) Quick edit\n");
        printf("I) Ignore\n");

        printf("\nSelect: ");

        char choice = getchar();
        getchar(); /* consume newline */

        /* ========================
           M — Manual Edit
        ======================== */
        if (choice == 'M' || choice == 'm')
        {
            printf("Enter new expression: ");

            if (!read_line(output, size))
                return 0;

            if (is_same_input(output))
                continue;

            return 1;
        }

        /* ========================
           E — Quick Edit
        ======================== */
        if (choice == 'E' || choice == 'e')
        {
            printf("Edit:\n%s\n> ", input);

            if (!read_line(output, size))
                return 0;

            if (is_same_input(output))
                continue;

            return 1;
        }

        /* ========================
           I — Ignore
        ======================== */
        if (choice == 'I' || choice == 'i')
        {
            return 0;
        }

        printf("Invalid choice. Try again.\n");
    }
}

/* =========================================================
   PUBLIC WRAPPERS
========================================================= */

int shriji_interactive_double_operator_fix(
    const char *input,
    char *output,
    int size)
{
    return show_fix_menu(input, output, size);
}

int shriji_interactive_bracket_fix(
    const char *input,
    char *output,
    int size)
{
    return show_fix_menu(input, output, size);
}

int shriji_interactive_operator_fix(
    const char *input,
    char *output,
    int size)
{
    return show_fix_menu(input, output, size);
}
/*
──────────────────────────────────────────────
    SHRIJILANG — MAIN ENTRY (FINAL REPL 🌸)
──────────────────────────────────────────────
*/

#include <stdio.h>
#include <stdlib.h>
#include <string.h>

#ifndef _WIN32
#include <readline/readline.h>
#include <readline/history.h>
#endif

#include "../include/env.h"
#include "../include/smriti.h"
#include "../include/smriti_personal.h"
#include "../include/smriti_session.h"
#include "../include/runtime.h"

#include "../krst/krst_core.h"
#include "../include/krst_router.h"
#include "../include/user_config.h"

#include "../include/normalize_engine.h"
#include "../include/nirman.h"
#include "../include/nirman_flow_engine.h"

#include "../include/parser.h"

/* GLOBAL ENV */
Env *GLOBAL_ENV = NULL;

/*──────────────────────────────────────────────*/
static void shriji_init(void)
{
    config_init();
    GLOBAL_ENV = new_env();

    smriti_init();
    smriti_session_clear();
    smriti_personal_load();
}

/*──────────────────────────────────────────────*/
static void run_repl_mode(void)
{
    printf("──────────────────────────────────────────────\n");
    printf("        🌸 ShrijiLang Shell 🌸\n");
    printf("──────────────────────────────────────────────\n");

    printf("Type 'exit' to quit\n\n");

    ShrijiRuntime rt;
    runtime_init(&rt);
    current_runtime = &rt;

    /* MULTI-LINE BUFFER 🌸 */
    char buffer[8192];
    buffer[0] = '\0';

    /* BRACE TRACKER 🌸 */
    int brace_count = 0;

    while (1)
    {
#ifdef _WIN32
        char line_buffer[1024];

        printf("Shiri> ");
        if (!fgets(line_buffer, sizeof(line_buffer), stdin))
            break;

        line_buffer[strcspn(line_buffer, "\n")] = '\0';
        char *line = line_buffer;
#else
        char *line = readline("Shiri> ");
        if (!line) break;

        if (nirman_is_active())
            normalize_input(line);
#endif

        runtime_reset(&rt);

#ifndef _WIN32
        if (*line) add_history(line);
#endif

        /* EXIT */
        if (strcmp(line, "exit") == 0)
        {
            printf("🌸 Radhe Radhe 🌸\n");
            smriti_session_clear();
#ifndef _WIN32
            free(line);
#endif
            break;
        }

        /* DEV MODE SYSTEM 🌸 */

        if (strcmp(line, "dev") == 0)
        {
            printf("Developer mode: %s\n", DEV_MODE ? "ON" : "OFF");
#ifndef _WIN32
            free(line);
#endif
            continue;
        }

        if (strcmp(line, "dev on") == 0)
        {
            DEV_MODE = 1;
            printf("Developer mode enabled\n");
#ifndef _WIN32
            free(line);
#endif
            continue;
        }

        if (strcmp(line, "dev off") == 0)
        {
            DEV_MODE = 0;
            printf("Developer mode disabled\n");
#ifndef _WIN32
            free(line);
#endif
            continue;
        }

        /* EMPTY INPUT */
        if (line[0] == '\0')
        {
#ifndef _WIN32
            free(line);
#endif
            continue;
        }

        /* NIRMAN COMMAND */
        if (strcmp(line, "nirman") == 0)
        {
            nirman_start();
            printf("Nirman mode activated\n");
            printf("Welcome to Shriji World\n");
#ifndef _WIN32
            free(line);
#endif
            continue;
        }

        if (strcmp(line, "exit nirman") == 0 ||
            strcmp(line, "nirman exit") == 0)
        {
            nirman_stop();
            printf("Nirman mode deactivated\n");
#ifndef _WIN32
            free(line);
#endif
            continue;
        }

        /* NIRMAN FLOW */
        if (nirman_is_active())
        {
            nirman_flow_process(line);
#ifndef _WIN32
            free(line);
#endif
            continue;
        }

        /* APPEND TO BUFFER 🌸 */
        strcat(buffer, line);
        strcat(buffer, "\n");

        /* COUNT BRACES 🌸 */
        for (int i = 0; line[i]; i++)
        {
            if (line[i] == '{') brace_count++;
            if (line[i] == '}') brace_count--;
        }

        /* WAIT UNTIL BLOCK COMPLETE 🌸 */
        if (brace_count != 0)
        {
#ifndef _WIN32
            free(line);
#endif
            continue;
        }

        /* PARSE */
        ASTNode *tree = parse_program(buffer);

        if (!tree)
        {
#ifndef _WIN32
            free(line);
#endif
            buffer[0] = '\0';
            brace_count = 0;
            continue;
        }

        /* EXECUTE */
        smriti_session_set_last_input(buffer);

        int fix_attempt = 0;
        const int MAX_FIX_ATTEMPT = 5;

        while (fix_attempt < MAX_FIX_ATTEMPT)
        {
            KRSTRequest req;

            req.input_text = buffer;
            req.has_correction = 0;
            req.corrected_text[0] = '\0';

            krst_route_request(&req);

            if (!req.has_correction)
                break;

            strcpy(buffer, req.corrected_text);
            fix_attempt++;
        }

        /* RESET BUFFER 🌸 */
        buffer[0] = '\0';
        brace_count = 0;

#ifndef _WIN32
        free(line);
#endif
    }
}

/*──────────────────────────────────────────────*/
static int run_file_mode(const char *path)
{
    FILE *fp = fopen(path, "rb");

    if (!fp)
    {
        printf("File not found: %s\n", path);
        return 1;
    }

    fseek(fp, 0, SEEK_END);
    long size = ftell(fp);
    fseek(fp, 0, SEEK_SET);

    char *buffer = malloc(size + 1);
    if (!buffer)
    {
        fclose(fp);
        return 1;
    }

    size_t read_bytes = fread(buffer, 1, size, fp);
    if (read_bytes != (size_t)size)
    {
        printf("File read error\n");
        free(buffer);
        fclose(fp);
        return 1;
    }

    buffer[size] = '\0';
    fclose(fp);

    smriti_session_set_last_input(buffer);

    ShrijiRuntime rt;
    runtime_init(&rt);
    current_runtime = &rt;

    KRSTRequest req;
    req.input_text = buffer;
    req.has_correction = 0;

    krst_route_request(&req);

    if (req.has_correction)
    {
        req.input_text = req.corrected_text;
        krst_route_request(&req);
    }

    free(buffer);
    return 0;
}

/*──────────────────────────────────────────────*/
int main(int argc, char *argv[])
{
    shriji_init();

    if (argc > 1)
        return run_file_mode(argv[1]);

    run_repl_mode();

    return 0;
}
#include <stdio.h>
#include <string.h>
#include <time.h>

#include "../include/command.h"
#include "../include/event.h"

/*──────────────────────────────────────────────────────────────
 | COMMAND ALIAS TABLE
 | 108 public + remaining locked (sample shown)
 *──────────────────────────────────────────────────────────────*/
static CommandAlias COMMAND_ALIASES[] = {

    /* LIST */
    { "ls",        CMD_LIST,   CMD_PUBLIC },
    { "list",      CMD_LIST,   CMD_PUBLIC },
    { "radha",     CMD_LIST,   CMD_PUBLIC },
    { "r",         CMD_LIST,   CMD_PUBLIC },

    /* CD */
    { "cd",        CMD_CD,     CMD_PUBLIC },
    { "chalo",     CMD_CD,     CMD_PUBLIC },

    /* MKDIR */
    { "mkdir",     CMD_MKDIR,  CMD_PUBLIC },
    { "bana",      CMD_MKDIR,  CMD_PUBLIC },

    /* BOLO (PRINT) */
    { "bolo",      CMD_BOLO,   CMD_PUBLIC },

    /* SAMAY (TIME) ✅ */
    { "samay",     CMD_SAMAY,  CMD_PUBLIC },
    { "time",      CMD_SAMAY,  CMD_PUBLIC },
    { "clock",     CMD_SAMAY,  CMD_PUBLIC },

/* AI Modules */
    { "sakhi",     CMD_SAKHI,  CMD_PUBLIC },
    { "niyu",      CMD_NIYU,   CMD_PUBLIC },
    { "mira",      CMD_MIRA,   CMD_PUBLIC },
    { "kavya",     CMD_KAVYA,  CMD_PUBLIC },
    { "shiri",     CMD_SHIRI,  CMD_PUBLIC },

    /* LOCKED / RESERVED */
    { "shyam",     CMD_DELETE, CMD_LOCKED },
    { "keshav",    CMD_COPY,   CMD_LOCKED },

    { NULL, CMD_UNKNOWN, CMD_LOCKED }
};

/*──────────────────────────────────────────────────────────────
 | Resolve user-facing name → internal command
 *──────────────────────────────────────────────────────────────*/
CommandType resolve_command(const char *name)
{
    if (!name)
        return CMD_UNKNOWN;

    for (int i = 0; COMMAND_ALIASES[i].alias; i++) {

        if (strcmp(COMMAND_ALIASES[i].alias, name) == 0) {

            if (COMMAND_ALIASES[i].access == CMD_LOCKED) {
                event_fire(EVENT_ERROR, "command locked");
                return CMD_UNKNOWN;
            }

            return COMMAND_ALIASES[i].cmd;
        }
    }

    return CMD_UNKNOWN;
}

/*──────────────────────────────────────────────────────────────
 | Internal command name (debug / AI explain)
 *──────────────────────────────────────────────────────────────*/
const char *command_name(CommandType cmd)
{
    switch (cmd) {
        case CMD_LIST:   return "CMD_LIST";
        case CMD_CD:     return "CMD_CD";
        case CMD_MKDIR:  return "CMD_MKDIR";
        case CMD_BOLO:   return "CMD_BOLO";
        case CMD_SAMAY:  return "CMD_SAMAY";

        case CMD_SAKHI:  return "CMD_SAKHI";
        case CMD_NIYU:   return "CMD_NIYU";
        case CMD_MIRA:   return "CMD_MIRA";
        case CMD_KAVYA:  return "CMD_KAVYA";
        case CMD_SHIRI:  return "CMD_SHIRI";

        case CMD_DELETE: return "CMD_DELETE";
        case CMD_COPY:   return "CMD_COPY";
        case CMD_MOVE:   return "CMD_MOVE";

        default:         return "CMD_UNKNOWN";
    }
}

/*──────────────────────────────────────────────────────────────
 | Execute internal command
 | NOTE:
 |  • BOLO printing is handled in interpreter (AST_COMMAND)
 |  • This file only supports commands that need no args
 *──────────────────────────────────────────────────────────────*/
int execute_command(CommandType cmd)
{
    switch (cmd) {

        case CMD_LIST:
            event_fire(EVENT_EXECUTION_SUCCESS, "list executed");
            printf("[LIST] executed (stub)\n");
            return 0;

        case CMD_CD:
            event_fire(EVENT_EXECUTION_SUCCESS, "cd executed");
            printf("[CD] executed (stub)\n");
            return 0;

        case CMD_MKDIR:
            event_fire(EVENT_EXECUTION_SUCCESS, "mkdir executed");
            printf("[MKDIR] executed (stub)\n");
            return 0;

        case CMD_BOLO:
            /* printing is done in interpreter */
            event_fire(EVENT_EXECUTION_SUCCESS, "bolo executed");
            return 0;

        case CMD_SAMAY: {
            /* ✅ print system time in 12-hour format (AM/PM) */
            event_fire(EVENT_EXECUTION_SUCCESS, "samay executed");

            time_t now = time(NULL);
            struct tm *t = localtime(&now);

            if (!t) {
                printf("00:00 AM\n");
                return 0;
            }

            int hour = t->tm_hour;
            int minute = t->tm_min;

            const char *ampm = "AM";
            if (hour >= 12) ampm = "PM";

            hour = hour % 12;
            if (hour == 0) hour = 12;

            printf("%02d:%02d %s\n", hour, minute, ampm);
            return 0;
        }

        case CMD_SAKHI:
        case CMD_NIYU:
        case CMD_MIRA:
        case CMD_KAVYA:
        case CMD_SHIRI:
            /* handled inside interpreter because requires args */
            event_fire(EVENT_ERROR, "AI command requires args");
            return -1;
        default:
            event_fire(EVENT_ERROR, "command not implemented");
            return -1;
    }
}
#include "../include/atma_lekha.h"
#include <stdio.h>
#include <time.h>

static const char *ATMA_FILE = ".shriji/atma_lekha.log";

void atma_lekha_init(void)
{
    FILE *fp = fopen(ATMA_FILE, "a");
    if (fp) fclose(fp);
}

void atma_lekha_note(AtmaEventType type,
                     const char *component,
                     const char *detail)
{
    FILE *fp = fopen(ATMA_FILE, "a");
    if (!fp) return;

    time_t now = time(NULL);
    struct tm *t = localtime(&now);

    fprintf(fp,
        "[%04d-%02d-%02d %02d:%02d:%02d] | %d | %s | %s\n",
        t->tm_year + 1900,
        t->tm_mon + 1,
        t->tm_mday,
        t->tm_hour,
        t->tm_min,
        t->tm_sec,
        type,
        component ? component : "unknown",
        detail ? detail : "-"
    );

    fclose(fp);
}
#include <string.h>
#include "../include/typo_engine.h"

/*
    Simple exit typo detection

    Rule:
    - Length must be 4
    - At least 3 characters match with "exit"
*/

int shriji_check_typo(
    const char *input,
    char *suggestion,
    int suggestion_size
)
{
    if (!input || !suggestion)
        return 0;

    const char *target = "exit";

    if (strlen(input) != 4)
        return 0;

    int match = 0;

    for (int i = 0; i < 4; i++) {
        if (input[i] == target[i])
            match++;
    }

    if (match >= 3) {
        strncpy(suggestion, target, suggestion_size - 1);
        suggestion[suggestion_size - 1] = '\0';
        return 1;
    }

    return 0;
}
#include <stdio.h>
#include <stdlib.h>
#include <string.h>

#include "../include/parser.h"
#include "../include/token.h"
#include "../include/ast.h"
#include "../include/error.h"

#include "../include/example_builder.h"

#include "lang/grammar.h"

#ifdef SHRIJI_ENABLE_MASTER_TOKENS
#include "lang/token_master.h"
#endif



/*──────────────────────────────────────────────
 | SHRIJILANG PARSER — V1 STABLE (FULL)
 | Supports:
 |   - mavi (declaration)
 |   - update (x = expr)
 |   - if/else (agar/warna)
 |   - while (jabtak)
 |   - break/continue (rukja/chalu)
 |   - return (wapas)
 |   - functions (kaam)
 |   - function call: add(10, 20)
 |   - newline aware
 |   - list literal: [1, 2, 3]
 |   - indexing: a[0]
 *──────────────────────────────────────────────*/




static ASTNode *parse_statement(void);
static ASTNode *parse_block(void);
static ASTNode *parse_if(void);
static ASTNode *parse_while(void);
static ASTNode *parse_return(void);
static ASTNode *parse_function(void);

static ASTNode *parse_rachna_definition(void);

static ASTNode *parse_assignment(void);

/* expression chain */
static ASTNode *parse_comparison(void);
static ASTNode *parse_expression(void);
static ASTNode *parse_term(void);
static ASTNode *parse_unary(void);
static ASTNode *parse_primary(void);

static ASTNode *parse_value(void);

/* list support */
static ASTNode *parse_list_literal(void);
static ASTNode *parse_dict_literal(void);
static ASTNode *parse_postfix(ASTNode *base);

static ASTNode *parse_import(void);

static Token current;

static Token prev_token;

static const char *current_source = NULL;

static void parser_error(
    Token tok,
    ShrijiErrorCode code,
    const char *context,
    const char *message,
    const char *hint
)
{
    shriji_error_at_full(tok, code, context, __func__, message, hint);
}







/*──────────────────────────────────────────────
 | HELPERS
 *──────────────────────────────────────────────*/
static void advance(void)
{
    prev_token = current;

    do {
        current = scan_token();
    } while (current.type == TOKEN_NEWLINE);
}

static int expect(TokenType t, ShrijiErrorCode c,
                  const char *ctx, const char *msg, const char *hint)
{
    if (current.type == t) {
        advance();
        return 1;
    }

   parser_error(current, c, ctx, msg, hint);
    return 0;
}

static void skip_separators(void)
{
    while (current.type == TOKEN_NEWLINE) {
        advance();
    }
}

static void unescape_into(char *out, int outcap, const char *in, int inlen)
{
    int oi = 0;

    for (int i = 0; i < inlen && oi < outcap - 1; i++) {
        char c = in[i];

        if (c == '\\' && i + 1 < inlen) {
            char n = in[i + 1];
            i++;

            switch (n) {
                case 'n':  out[oi++] = '\n'; break;
                case 't':  out[oi++] = '\t'; break;
                case 'r':  out[oi++] = '\r'; break;
                case '"':  out[oi++] = '"';  break;
                case '\\': out[oi++] = '\\'; break;
                default:
                    out[oi++] = n;
                    break;
            }
            continue;
        }

        out[oi++] = c;
    }

    out[oi] = 0;
}






/*──────────────────────────────────────────────
 | PARSER SYNCHRONIZE
 *──────────────────────────────────────────────*/







static void parser_sync(void)
{
    while (current.type != TOKEN_NEWLINE &&
           current.type != TOKEN_EOF)
    {
        advance();
    }
}








/*──────────────────────────────────────────────
 | PROGRAM
 *──────────────────────────────────────────────*/

ASTNode *parse_program(const char *source)
{
   current_source = source;

#ifdef SHRIJI_ENABLE_MASTER_TOKENS
    shriji_register_grammar_core();   // 👈 यही FINAL जगह है
#endif

    init_tokenizer(source);
    advance();

    ASTNode **stmts = NULL;
    int count = 0;

    while (current.type != TOKEN_EOF) {
   error_reported = 0;
        skip_separators();
        if (current.type == TOKEN_EOF)
            break;






ASTNode *s = parse_statement();

if (!s) {
    // 🔥 ERROR RECOVERY
    while (current.type != TOKEN_NEWLINE &&
           current.type != TOKEN_EOF) {
        advance();
    }

    skip_separators();
    continue;
}
        ASTNode **tmp = realloc(stmts, sizeof(ASTNode*) * (count + 1));
        if (!tmp) {
            free(stmts);
            return NULL;
        }
        stmts = tmp;

        stmts[count++] = s;

        skip_separators();
    }

    return new_program_node(stmts, count);
}










/*──────────────────────────────────────────────
  | STATEMENT
  *──────────────────────────────────────────────*/
static ASTNode *parse_statement(void)
{
    ASTNode *node = NULL;

    if (current.type == TOKEN_LEFT_BRACE) return parse_block();

     if (current.type == TOKEN_RACHNA)
    return parse_rachna_definition();

    if (current.type == TOKEN_AGAR) {
        ASTNode *i = parse_if();
        if (!i) printf("DEBUG: IF NULL\n");
        return i;
    }

    if (current.type == TOKEN_JABTAK) {
        ASTNode *w = parse_while();
        if (!w) printf("DEBUG: WHILE NULL\n");
        return w;
    }

    if (current.type == TOKEN_WAPAS)      return parse_return();
    if (current.type == TOKEN_KAAM)       return parse_function();
    if (current.type == TOKEN_RACHNA)    return parse_rachna_definition();
    if (current.type == TOKEN_MAVI)       return parse_assignment();
    if (current.type == TOKEN_IMPORT)     return parse_import();
    /* samay */
    if (current.type == TOKEN_SAMAY) {
        advance();
        skip_separators();
        return new_command_node("samay", NULL);
    }

    /* rukja */
    if (current.type == TOKEN_RUKJA) {
        advance();
        skip_separators();
        return new_break_node();
    }

    /* chalu */
    if (current.type == TOKEN_CHALU) {
        advance();
        skip_separators();
        return new_continue_node();
    }

/* bolo <expr> */
if (current.type == TOKEN_BOLO) {
    advance();
    skip_separators();

    node = parse_expression();

    if (!node) {
        /* expression layer already handled error */
        return NULL;
    }

    skip_separators();
    return new_command_node("bolo", node);
}

/* AI commands */
if (current.type == TOKEN_SAKHI ||
    current.type == TOKEN_NIYU  ||
    current.type == TOKEN_MIRA  ||
    current.type == TOKEN_KAVYA ||
    current.type == TOKEN_SHIRI) {

    char cmd[32];
    int len = current.length;

    if (len >= (int)sizeof(cmd))
        len = (int)sizeof(cmd) - 1;

    strncpy(cmd, current.start, len);
    cmd[len] = '\0';   /* ✅ VERY IMPORTANT */

    advance();
    skip_separators();

    ASTNode *arg = parse_value();
    if (!arg) return NULL;

    skip_separators();
    return new_command_node(cmd, arg);
}









if (current.type == TOKEN_IDENTIFIER) {

    Token name_token = current;

    advance();
    skip_separators();

    /* CASE 1: UPDATE (x = ...) */
    if (current.type == TOKEN_EQUAL) {

        advance();
        skip_separators();

        ASTNode *val = parse_expression();
        if (!val) return NULL;

        char name[128];
        strncpy(name, name_token.start, name_token.length);
        name[name_token.length] = '\0';

        skip_separators();
        return new_update_node(name, val);
    }

    /* CASE 2: FUNCTION CALL (test(...)) */
    if (current.type == TOKEN_LEFT_PAREN) {

        /* rollback */
        current = name_token;

        ASTNode *call = parse_expression();
        if (!call) return NULL;

        skip_separators();
        return call;
    }

    /* CASE 3: VARIABLE */
    current = name_token;

    ASTNode *expr = parse_expression();
    if (!expr) return NULL;

    skip_separators();
    return expr;
}







/* raw expression */
 if (current.type == TOKEN_NUMBER ||
    current.type == TOKEN_STRING ||
    current.type == TOKEN_TRUE   ||
    current.type == TOKEN_FALSE  ||
    current.type == TOKEN_LEFT_PAREN ||
    current.type == TOKEN_NAHI   ||
    current.type == TOKEN_LEFT_BRACKET ||
    current.type == TOKEN_PLUS ||
    current.type == TOKEN_MINUS ||
    current.type == TOKEN_STAR ||
    current.type == TOKEN_SLASH)
{

    node = parse_expression();

    if (!node) {
        if (!error_reported) {
            parser_error(
                current,
                E_PARSE_INVALID_TOKEN,
                "expr",
                "expression samajh nahi aaya",
                "example: 5 + 2"
            );
        }
        return NULL;
    }

    /*  FIXED CHECK */
    if (current.type != TOKEN_NEWLINE &&
        current.type != TOKEN_EOF &&
        current.type != TOKEN_RIGHT_BRACE &&
        current.type != TOKEN_RIGHT_PAREN)
    {
        char example[64];
        shriji_build_example(current_source, example, sizeof(example));

        shriji_error_at(
            current,
            E_PARSE_OPERATOR_CHAIN,
            "expr",
            "Expression ke baad extra tokens bach gaye hain",
            example
        );

        return NULL;
    }

    return node;
}







    /* '=' can never start a statement */
    if (current.type == TOKEN_EQUAL) {

char example[64];
shriji_build_example(current_source, example, sizeof(example));

shriji_error_at(
            current,
            E_PARSE_INVALID_TOKEN,
            "statement",
            "Assignment yahin se start nahi ho sakta. Variable aur '=' ek hi line me likho.",
            example
        );

        return NULL;
    }









/* FINAL FALLBACK — ONLY IF NO ERROR ALREADY REPORTED */
if (!error_reported) {
    parser_error(
        current,
        E_PARSE_INVALID_TOKEN,
        "statement",
        "statement samajh nahi aaya",
        "use: bolo / mavi / agar"
    );
}

return NULL;
}






/*──────────────────────────────────────────────
 | BLOCK { ... }
 *──────────────────────────────────────────────*/
static ASTNode *parse_block(void)
{
    if (!expect(TOKEN_LEFT_BRACE, E_PARSE_BRACKET_MISSING, "{", "missing {", "{ ... }"))
        return NULL;

    ASTNode **stmts = NULL;
    int count = 0;

    while (current.type != TOKEN_RIGHT_BRACE &&
           current.type != TOKEN_EOF)
    {

//  FORCE SKIP INVALID TOKENS (CRITICAL FIX)
if (current.type == TOKEN_ERROR) {
    advance();
    continue;
}

        skip_separators();

        if (current.type == TOKEN_RIGHT_BRACE)
            break;

        ASTNode *s = parse_statement();

        if (!s) {
            printf("DEBUG: statement NULL at block\n");

            while (current.type != TOKEN_NEWLINE &&
                   current.type != TOKEN_EOF &&
                   current.type != TOKEN_RIGHT_BRACE)
            {
                advance();
            }

            skip_separators();
            continue;
        }

        ASTNode **tmp = realloc(stmts, sizeof(ASTNode*) * (count + 1));
        if (!tmp) {
            free(stmts);
            return NULL;
        }

        stmts = tmp;
        stmts[count++] = s;

        skip_separators();
    }

    if (!expect(TOKEN_RIGHT_BRACE, E_PARSE_BRACKET_MISSING, "}", "missing }", "{ ... }")) {
        free(stmts);
        return NULL;
    }

    skip_separators();

    return new_block_node(stmts, count);
}









/*──────────────────────────────────────────────
 | IF / ELSE
 *──────────────────────────────────────────────*/
static ASTNode *parse_if(void)
{
    advance(); /* consume agar */

    ASTNode *cond = parse_expression();
    if (!cond) return NULL;

    skip_separators();

    ASTNode *then_block = parse_block();
    if (!then_block) return NULL;

    skip_separators();

    ASTNode *node = new_if_node(cond, then_block);

    if (current.type == TOKEN_WARNA) {
        advance(); /* consume warna */
        skip_separators();

        node->else_block = parse_block();
        if (!node->else_block) return NULL;

        skip_separators();
    }

    return node;
}








/*──────────────────────────────────────────────
 | WHILE
 *──────────────────────────────────────────────*/
static ASTNode *parse_while(void)
{
    advance(); /* consume jabtak */

    ASTNode *cond = parse_expression();
    if (!cond) return NULL;

    skip_separators();

    ASTNode *body = parse_block();
    if (!body) return NULL;

    return new_while_node(cond, body);
}









/*──────────────────────────────────────────────
 | IMPORT "file.sri"
 *──────────────────────────────────────────────*/
static ASTNode *parse_import(void)
{
    advance(); /* consume import */
    skip_separators();

    if (current.type != TOKEN_STRING) {
char example[64];
shriji_build_example(current_source, example, sizeof(example));

shriji_error_at(
            current,
            E_IMPORT_PATH_INVALID,
            "import",
            "expected string path",
            example
        );

        return NULL;
    }

    char mod[256];
    int len = current.length;

    const char *p = current.start;
    if (len >= 2 && p[0] == '"' && p[len - 1] == '"') {
        p++;
        len -= 2;
    }

    if (len >= (int)sizeof(mod)) len = (int)sizeof(mod) - 1;
    strncpy(mod, p, len);
    mod[len] = 0;

    advance(); /* consume string */

    skip_separators();
    return new_import_node(mod);
}










/*──────────────────────────────────────────────
 | RETURN / WAPAS
 *──────────────────────────────────────────────*/
static ASTNode *parse_return(void)
{
    advance(); /* consume wapas */

    if (current.type == TOKEN_NEWLINE ||
        current.type == TOKEN_RIGHT_BRACE ||
        current.type == TOKEN_EOF) {

        skip_separators();
        return new_return_node(NULL);
    }

       ASTNode *val = parse_value();
    if (!val) return NULL;

    skip_separators();
    return new_return_node(val);
}










/*──────────────────────────────────────────────
 | FUNCTION / KAAM
 *──────────────────────────────────────────────*/
static ASTNode *parse_function(void)
{
    advance(); /* consume kaam */
    skip_separators();

    if (current.type != TOKEN_IDENTIFIER) {
        shriji_error_at(current, E_FUNCTION_NAME_INVALID, "kaam", "expected function name",
                        "kaam add(a, b) { ... }");
        return NULL;
    }

    char fname[128];
    strncpy(fname, current.start, current.length);
    fname[current.length] = 0;
    advance(); /* consume function name */

    skip_separators();

    if (!expect(TOKEN_LEFT_PAREN, E_PARSE_02, "(", "missing (",
                "kaam add(a, b) { ... }"))
        return NULL;

    ASTNode **params = NULL;
    int param_count = 0;

    skip_separators();

    if (current.type != TOKEN_RIGHT_PAREN) {
        while (1) {

            if (current.type != TOKEN_IDENTIFIER) {
                shriji_error_at(current, E_FUNCTION_PARAM_INVALID, "params",
                                "expected parameter name",
                                "kaam add(a, b) { ... }");
                free(params);
                return NULL;
            }

            char p[128];
            strncpy(p, current.start, current.length);
            p[current.length] = 0;
            advance();

            ASTNode **tmp = realloc(params, sizeof(ASTNode*) * (param_count + 1));
            if (!tmp) {
                free(params);
                return NULL;
            }
            params = tmp;

            params[param_count++] = new_identifier_node(p);

            skip_separators();

            if (current.type == TOKEN_COMMA) {
                advance();
                skip_separators();
            }

            if (current.type == TOKEN_RIGHT_PAREN)
                break;
        }
    }

    if (!expect(TOKEN_RIGHT_PAREN, E_PARSE_BRACKET_MISSING, ")", "missing )",
                "kaam add(a, b) { ... }")) {
        free(params);
        return NULL;
    }

    skip_separators();

    ASTNode *body = parse_block();
    if (!body) return NULL;

    skip_separators();
    return new_function_node(fname, params, param_count, body);
}

static ASTNode *parse_rachna_definition(void)
{
    advance(); // consume 'rachna'
    skip_separators();

    /* name */
    if (current.type != TOKEN_IDENTIFIER) {
        parser_error(current, E_PARSE_INVALID_TOKEN,
                     "rachna",
                     "name expected",
                     "example: rachna add() { }");
        return NULL;
    }

    char name[128];
    strncpy(name, current.start, current.length);
    name[current.length] = '\0';

    advance();
    skip_separators();

    /* '(' */
    if (!expect(TOKEN_LEFT_PAREN,
                E_PARSE_INVALID_TOKEN,
                "rachna",
                "'(' expected",
                "example: rachna add() { }"))
        return NULL;

    skip_separators();

    /* params */
    ASTNode **params = NULL;
    int param_count = 0;

    if (current.type != TOKEN_RIGHT_PAREN) {
        while (1) {

            if (current.type != TOKEN_IDENTIFIER) {
                parser_error(current, E_PARSE_INVALID_TOKEN,
                             "rachna",
                             "parameter name expected",
                             "example: rachna add(a,b) { }");
                return NULL;
            }

            char p[128];
            strncpy(p, current.start, current.length);
            p[current.length] = '\0';

            ASTNode **tmp = realloc(params, sizeof(ASTNode*) * (param_count + 1));
            if (!tmp) return NULL;
            params = tmp;

            params[param_count++] = new_identifier_node(p);

            advance();
            skip_separators();

            if (current.type == TOKEN_COMMA) {
                advance();
                skip_separators();
                continue;
            }

            break;
        }
    }

    if (!expect(TOKEN_RIGHT_PAREN,
                E_PARSE_INVALID_TOKEN,
                "rachna",
                "')' expected",
                "example: rachna add(a,b) { }"))
        return NULL;

    skip_separators();

    /* body */
    ASTNode *body = parse_block();
    if (!body) return NULL;

    skip_separators();

    return new_rachna_node(name, params, param_count, body);
}

















/*──────────────────────────────────────────────
 | IDENTIFIER / UPDATE
 *──────────────────────────────────────────────*/

static ASTNode *parse_value(void)
{
    return parse_comparison();
}












/*──────────────────────────────────────────────
  ASSIGNMENT (MAVI)
──────────────────────────────────────────────*/
static ASTNode *parse_assignment(void)
{
    advance(); /* consume 'mavi' */
    skip_separators();

    if (current.type != TOKEN_IDENTIFIER) {
        shriji_error_at(
            current,
            E_ASSIGN_01,
            "mavi",
            "Declaration samajh li, par variable ka naam clear nahi hua.",
            "mavi x  ya  mavi x = 10"
        );
        return NULL;
    }

    char name[128];
    strncpy(name, current.start, current.length);
    name[current.length] = 0;
    advance();

    skip_separators();

    if (current.type == TOKEN_EQUAL) {
        advance();
        skip_separators();

        if (current.type == TOKEN_NEWLINE ||
            current.type == TOKEN_EOF ||
            current.type == TOKEN_RIGHT_BRACE) {

            shriji_error_at(
                current,
                E_ASSIGN_02,
                "assignment",
                "'=' ke baad value missing hai",
                "example: mavi x = 10"
            );
            return NULL;
        }

        ASTNode *val = parse_value();
        if (!val) return NULL;

        skip_separators();
        return new_assignment_node(name, val);
    }

    return new_assignment_node(name, NULL);
}






static ASTNode *parse_logical_and(void)
{
    ASTNode *node = parse_comparison();

    while (current.type == TOKEN_AND)
    {
        char op = '&';
        advance();

        ASTNode *right = parse_comparison();

        node = new_binary_node(op, node, right);
    }

    return node;
}


static ASTNode *parse_logical_or(void)
{
    ASTNode *node = parse_logical_and();

    while (current.type == TOKEN_OR)
    {
        char op = '|';
        advance();

        ASTNode *right = parse_logical_and();
        node = new_binary_node(op, node, right);
    }

    return node;
}

/*──────────────────────────────────────────────
 | EXPRESSIONS
 | precedence:
 |   primary -> unary -> term(* / %) -> expression(+ -) -> comparison
 *──────────────────────────────────────────────*/

static ASTNode *parse_comparison(void)
{
    ASTNode *n = parse_term();
    if (!n) return NULL;

    while (current.type == TOKEN_GT  ||
           current.type == TOKEN_LT  ||
           current.type == TOKEN_GTE ||
           current.type == TOKEN_LTE ||
           current.type == TOKEN_EQEQ ||
           current.type == TOKEN_NEQ)
    {
        char op =
            (current.type == TOKEN_GT)    ? '>' :
            (current.type == TOKEN_LT)    ? '<' :
            (current.type == TOKEN_GTE)   ? 'G' :
            (current.type == TOKEN_LTE)   ? 'L' :
            (current.type == TOKEN_EQEQ)  ? '=' :
            (current.type == TOKEN_NEQ)   ? '!' :
            '?';

        advance();

        ASTNode *right = parse_term();
        if (!right) return NULL;

        n = new_binary_node(op, n, right);
    }

    return n;
}


static ASTNode *parse_expression(void)
{
   ASTNode *n = parse_logical_or();
    if (!n) return NULL;




/*──────────────────────────────────────────────
 | OPERATOR START DETECTION
 *──────────────────────────────────────────────*/
if (prev_token.type == TOKEN_NEWLINE ||
    prev_token.type == TOKEN_EOF)
{
    if (current.type == TOKEN_PLUS ||
        current.type == TOKEN_MINUS ||
        current.type == TOKEN_STAR ||
        current.type == TOKEN_SLASH)
    {

        char example[64];
        shriji_build_example(current_source, example, sizeof(example));

        shriji_error_at(
            current,
            E_PARSE_OPERATOR_START,
            "expr",
            "Expression operator se start nahi ho sakta",
            example
        );

        parser_sync();
        return NULL;
    }
}



while (current.type == TOKEN_PLUS  ||
       current.type == TOKEN_MINUS)
{
    char op;

    if (current.type == TOKEN_PLUS) op = '+';
    else op = '-';

    advance();







/* operator classification */

if (current.type == TOKEN_PLUS  ||
    current.type == TOKEN_MINUS ||
    current.type == TOKEN_STAR  ||
    current.type == TOKEN_SLASH ||
    current.type == TOKEN_AND   ||
    current.type == TOKEN_OR)
{
    ShrijiErrorCode code = E_PARSE_DOUBLE_OPERATOR;

    if (prev_token.type == current.type)
    {
        code = E_PARSE_DOUBLE_OPERATOR;
    }
    else
    {
        code = E_PARSE_OPERATOR_CHAIN;
    }






char example[64];
shriji_build_example(current_source, example, sizeof(example));

shriji_error_at(
    current,
    code,
    "expr",
    "Operator chain detect hui hai.",
    example
);

return NULL;

}












/* operator end detection */

if (current.type == TOKEN_NEWLINE ||
    current.type == TOKEN_EOF ||
    current.type == TOKEN_RIGHT_BRACE)
{
char example[64];
shriji_build_example(current_source, example, sizeof(example));

shriji_error_at(
    current,
    E_PARSE_OPERATOR_END,
    "expr",
    "Operator ke baad value missing hai",
    example
);

    return NULL;
}

        ASTNode *right = parse_term();
        if (!right) return NULL;

        n = new_binary_node(op, n, right);
    }

    return n;
}


static ASTNode *parse_term(void)
{
    ASTNode *n = parse_unary();
    if (!n) return NULL;

    while (current.type == TOKEN_STAR ||
           current.type == TOKEN_SLASH ||
           current.type == TOKEN_MOD)
    {
        char op =
            (current.type == TOKEN_STAR)  ? '*' :
            (current.type == TOKEN_SLASH) ? '/' : '%';

        advance();










/* double operator detection */

if (current.type == TOKEN_STAR ||
    current.type == TOKEN_SLASH ||
    current.type == TOKEN_MOD)
{
    char example[64];
    shriji_build_example(current_source, example, sizeof(example));

    shriji_error_at(
        current,
        E_PARSE_DOUBLE_OPERATOR,
        "expr",
        "Do operators ek saath aa gaye hain.",
        example
    );

    return NULL;
}









/* operator end detection for * / % */

if (current.type == TOKEN_NEWLINE ||
    current.type == TOKEN_EOF ||
    current.type == TOKEN_RIGHT_BRACE)
{

char example[64];
shriji_build_example(current_source, example, sizeof(example));

shriji_error_at(
    current,
    E_PARSE_OPERATOR_END,
    "expr",
    "Expression operator par end ho gaya",
    example
);

skip_separators();
return NULL;

}

        ASTNode *right = parse_unary();
        if (!right) return NULL;

        n = new_binary_node(op, n, right);
    }

    return n;
}









/*──────────────────────────────────────────────
 | UNARY
 |   - prefix ++x / --x
 |   - unary + / -
 |   - nahi (logical not)
 *──────────────────────────────────────────────*/
static ASTNode *parse_unary(void)
{
    /* prefix ++ */
    if (current.type == TOKEN_PLUSPLUS) {
        advance();

        ASTNode *expr = parse_unary();
        if (!expr) return NULL;

        Token fake = {0};
        fake.start = "+";
        fake.length = 1;

        return new_unary_node(fake, expr);
    }

    /* prefix -- */
    if (current.type == TOKEN_MINUSMINUS) {
        advance();

        ASTNode *expr = parse_unary();
        if (!expr) return NULL;

        Token fake = {0};
        fake.start = "-";
        fake.length = 1;

        return new_unary_node(fake, expr);
    }

    /* unary + / - */
    if (current.type == TOKEN_PLUS || current.type == TOKEN_MINUS) {

        Token op = current;
        advance();

        /* prevent invalid operator chain */
        if (current.type == TOKEN_PLUS ||
            current.type == TOKEN_MINUS ||
            current.type == TOKEN_STAR ||
            current.type == TOKEN_SLASH)
        {
            shriji_error_at(
                current,
                E_PARSE_OPERATOR_CHAIN,
                "expr",
                "operator chain galat hai",
                "example: +5 ya -3"
            );
            return NULL;
        }

        ASTNode *expr = parse_primary();
        if (!expr) return NULL;

        return new_unary_node(op, expr);
    }

    /* nahi (logical not) */
    if (current.type == TOKEN_NAHI) {
        advance();

        ASTNode *expr = parse_unary();
        if (!expr) return NULL;

        return new_not_node(expr);
    }

    /* default primary */
    ASTNode *node = parse_primary();
    if (!node) return NULL;

    /* detect missing operator: 6(5+9) */
    if (current.type == TOKEN_LEFT_PAREN)
    {
        char example[64];
        shriji_build_example(current_source, example, sizeof(example));

        shriji_error_at(
            current,
            E_PARSE_MISSING_OPERATOR,
            "expr",
            "Do values ke beech operator missing hai",
            example
        );

        return NULL;
    }

    return node;
}







/*──────────────────────────────────────────────
 | LIST LITERAL: [a, b, c]
 *──────────────────────────────────────────────*/
static ASTNode *parse_list_literal(void)
{
    if (!expect(TOKEN_LEFT_BRACKET, E_LIST_SYNTAX_ERROR, "[", "missing [", "[1,2,3]"))
        return NULL;

    skip_separators();

    ASTNode **elements = NULL;
    int count = 0;

    /* empty list: [] */
    if (current.type != TOKEN_RIGHT_BRACKET) {

        while (1) {

            ASTNode *elem = parse_expression();
            if (!elem) {
                free(elements);
                return NULL;
            }

            ASTNode **tmp = realloc(elements, sizeof(ASTNode*) * (count + 1));
            if (!tmp) {
                free(elements);
                return NULL;
            }
            elements = tmp;
            elements[count++] = elem;

            skip_separators();

            if (current.type == TOKEN_RIGHT_BRACKET)
                break;

            if (!expect(TOKEN_COMMA, E_LIST_SYNTAX_ERROR, ",", "expected ','", "[1, 2, 3]")) {
                free(elements);
                return NULL;
            }

            skip_separators();
        }
    }

    if (!expect(TOKEN_RIGHT_BRACKET, E_PARSE_BRACKET_MISSING, "]", "missing ]", "[1,2,3]")) {
        free(elements);
        return NULL;
    }

    skip_separators();
    return new_list_node(elements, count);
}












/*──────────────────────────────────────────────
 | POSTFIX:
 |   - indexing: a[0]
 |   - postfix inc/dec: x++ , x--
 *──────────────────────────────────────────────*/
static ASTNode *parse_postfix(ASTNode *base)
{
    ASTNode *node = base;

    while (1) {

        /* ───── indexing: a[expr] ───── */
        if (current.type == TOKEN_LEFT_BRACKET) {

            advance(); /* consume '[' */
            skip_separators();

            ASTNode *idx = parse_expression();
            if (!idx) return NULL;

            skip_separators();

            if (!expect(TOKEN_RIGHT_BRACKET, E_PARSE_BRACKET_MISSING,
                        "]", "missing ]", "a[0]"))
                return NULL;

            skip_separators();

            node = new_index_node(node, idx);
            continue;
        }

        break;
    }

    return node;
}










/*──────────────────────────────────────────────
 | PRIMARY
 *──────────────────────────────────────────────*/
static ASTNode *parse_primary(void)
{
if (current.type == TOKEN_EOF ||
        current.type == TOKEN_NEWLINE)
    {
        char example[64];
        shriji_build_example(current_source, example, sizeof(example));

        shriji_error_at(
            current,
            E_PARSE_OPERATOR_END,
            "expr",
            "Expression incomplete hai",
            example
        );

        return NULL;
    }


    if (current.type == TOKEN_TRUE) {
         advance();
           return new_bool_node(1);
        }

   if (current.type == TOKEN_FALSE) {
        advance();
           return new_bool_node(0);
        }








if (current.type == TOKEN_NUMBER) {

    char temp[64];
    int len = current.length;

    if (len >= 63)
        len = 63;

    strncpy(temp, current.start, len);
    temp[len] = '\0';

    double v = atof(temp);

    advance();

    return new_number_node(v);
}









if (current.type == TOKEN_STRING) {
    char s[256];
    int len = current.length;

    const char *p = current.start;
    if (len >= 2 && p[0] == '"' && p[len - 1] == '"') {
        p++;
        len -= 2;
    }

    unescape_into(s, sizeof(s), p, len);

    advance();
    return new_string_node(s);
}


    /* list literal: [ ... ] */
    if (current.type == TOKEN_LEFT_BRACKET) {
        ASTNode *lst = parse_list_literal();
        if (!lst) return NULL;
        return parse_postfix(lst);
    }





/* dict literal: { ... } */
if (current.type == TOKEN_LEFT_BRACE) {
    ASTNode *d = parse_dict_literal();
    if (!d) return NULL;
    ASTNode *node = parse_postfix(d);
if (!node) return NULL;
return node;
}

    if (current.type == TOKEN_IDENTIFIER) {

        char name[128];
        strncpy(name, current.start, current.length);
        name[current.length] = 0;
        advance();






/* function call: name(...) */
if (current.type == TOKEN_LEFT_PAREN) {

    advance(); /* consume '(' */
    skip_separators();

    ASTNode **args = NULL;
    int arg_count = 0;

    if (current.type != TOKEN_RIGHT_PAREN) {
        while (1) {

            ASTNode *arg = parse_expression();
            if (!arg) {
                free(args);
                return NULL;
            }

            ASTNode **tmp = realloc(args, sizeof(ASTNode*) * (arg_count + 1));
            if (!tmp) {
                free(args);
                return NULL;
            }
            args = tmp;
            args[arg_count++] = arg;

            skip_separators();

            if (current.type == TOKEN_RIGHT_PAREN)
                break;

            if (current.type != TOKEN_COMMA) {
                shriji_error_at(
                    current,
                    E_FUNCTION_PARAM_INVALID,
                    "call",
                    "expected ',' between arguments",
                    "example: add(5, 6)"
                );
                free(args);
                return NULL;
            }

            advance();
            skip_separators();
        }
    }

    if (!expect(TOKEN_RIGHT_PAREN, E_PARSE_01, ")", "missing ')'",
                "example: add(10, 20)")) {
        free(args);
        return NULL;
    }

    skip_separators();

    ASTNode *node = new_call_node(name, args, arg_count);
    node = parse_postfix(node);
    if (!node) return NULL;

    return node;
}


        ASTNode *node = new_identifier_node(name);
node = parse_postfix(node);
if (!node) return NULL;
return node;
    }

    if (current.type == TOKEN_LEFT_PAREN) {
        advance();
        ASTNode *n = parse_comparison();
        if (!n) return NULL;

        if (!expect(TOKEN_RIGHT_PAREN, E_PARSE_UNMATCHED_PAREN, ")", "missing )", "(expr)"))
            return NULL;

        ASTNode *node = parse_postfix(n);
if (!node) return NULL;
return node;
    }

    shriji_error_at(current, E_PARSE_INVALID_TOKEN, "expr", "bad token", "check syntax");
    return NULL;
}










/*──────────────────────────────────────────────
 | DICT / MAP LITERAL: { "a": 1, "b": 2 }
 *──────────────────────────────────────────────*/
static ASTNode *parse_dict_literal(void)
{
    if (!expect(TOKEN_LEFT_BRACE, E_PARSE_BRACKET_MISSING, "{", "missing {", "{\"a\":1}"))
        return NULL;

    skip_separators();

    ASTNode **keys = NULL;
    ASTNode **vals = NULL;
    int count = 0;

    /* empty dict: {} */
    if (current.type != TOKEN_RIGHT_BRACE) {

        while (1) {

            /* key must be string for v1 */
            if (current.type != TOKEN_STRING) {
                shriji_error_at(
                    current,
                    E_DICT_KEY_INVALID,
                    "dict",
                    "dict key must be string",
                    "example: {\"a\": 1, \"b\": 2}"
                );
                free(keys);
                free(vals);
                return NULL;
            }

            /* parse key as string node */
            ASTNode *k = parse_primary();
            if (!k) {
                free(keys);
                free(vals);
                return NULL;
            }

            skip_separators();

            if (!expect(TOKEN_COLON, E_DICT_SYNTAX_ERROR, ":", "missing ':'",
                        "example: {\"a\": 1}")) {
                free(keys);
                free(vals);
                return NULL;
            }

            skip_separators();

            /* parse value expression */
            ASTNode *v = parse_expression();
            if (!v) {
                free(keys);
                free(vals);
                return NULL;
            }








/* push pair */
ASTNode **tmpk = realloc(keys, sizeof(ASTNode*) * (count + 1));
if (!tmpk) {
    free(keys);
    free(vals);
    return NULL;
}
keys = tmpk;

ASTNode **tmpv = realloc(vals, sizeof(ASTNode*) * (count + 1));
if (!tmpv) {
    free(vals);
    return NULL;
}
vals = tmpv;

keys[count] = k;
vals[count] = v;
count++;

skip_separators();






/* end? */
if (current.type == TOKEN_RIGHT_BRACE)
    break;

if (!expect(TOKEN_COMMA, E_DICT_SYNTAX_ERROR, ",", "expected ','",
            "example: {\"a\": 1, \"b\": 2}")) {
    free(keys);
    free(vals);
    return NULL;
}

        }
    }

    if (!expect(TOKEN_RIGHT_BRACE, E_PARSE_BRACKET_MISSING, "}", "missing '}'",
                "example: {\"a\": 1, \"b\": 2}")) {
        free(keys);
        free(vals);
        return NULL;
    }

    skip_separators();
    return new_dict_node(keys, vals, count);
}

/* SHREE_RADHIKA_RANI */

#include <stdio.h>
#include <string.h>
#include <ctype.h>

#include "../include/parser.h"
#include "../include/error.h"
#include "../include/fix_engine.h"
#include "../include/fix_rules.h"

#define MAX_WORKING_BUFFER 512


/*──────────────────────────────────────────────
   PARSE ONLY (STRICT MODE)
──────────────────────────────────────────────*/

ASTNode *parse_once(const char *input)
{
    error_reported = 0;
    shriji_set_error_mode(ERROR_MODE_IMMEDIATE);

    return parse_program(input);
}


/*──────────────────────────────────────────────
   MAIN EXECUTION (SAFE VERSION)
──────────────────────────────────────────────*/

ASTNode *language_execute_with_fix(
    const char *input,
    int *was_fixed,
    int *penalty_out
)
{
    if (!input)
        return NULL;

    if (was_fixed)
        *was_fixed = 0;

    if (penalty_out)
        *penalty_out = 0;

    error_reported = 0;
    shriji_set_error_mode(ERROR_MODE_IMMEDIATE);

    /* FIRST PARSE */
    ASTNode *root = parse_program(input);

    /* ERROR CASE */
    if (!root || error_reported)
    {
        char fixed[MAX_WORKING_BUFFER];
        memset(fixed, 0, sizeof(fixed));

        FixType type = fix_apply(input, fixed);

        /* SAFE FIX */
        if (type == FIX_SAFE)
        {
            if (was_fixed)
                *was_fixed = 1;

            error_reported = 0;

            ASTNode *new_root = parse_program(fixed);

            if (!new_root || error_reported)
                return NULL;

            return new_root;
        }

        /* no safe fix */
        return NULL;
    }

    return root;
}


/* ================================
   HELPER: invalid token
================================ */

static int has_invalid_char(const char *input)
{
    for (int i = 0; input[i]; i++)
    {
        char c = input[i];

        if (isdigit(c) || strchr("+-*/(). ", c))
            continue;

        return 1;
    }
    return 0;
}


/* ================================
   HELPER: operator chain length (FIXED)
================================ */

static int operator_chain_len(const char *input, int pos)
{
    int len = 1;

    while (input[pos + len] != '\0' &&
           strchr("+-*/", input[pos + len]))
    {
        len++;
    }

    return len;
}


/* ================================
   MAIN FIX ENGINE
================================ */

FixType fix_apply(const char *input, char *output)
{
    if (!input || !*input)
        return FIX_NONE;

    if (has_invalid_char(input))
        return FIX_REJECT;

    int i, j = 0;
    int changed = 0;

    for (i = 0; input[i]; i++)
    {
        /* implicit multiplication */
        if (isdigit(input[i]) && input[i+1] == '(')
        {
            output[j++] = input[i];
            output[j++] = '*';
            changed = 1;
            continue;
        }

        if (input[i] == ')' && isdigit(input[i+1]))
        {
            output[j++] = input[i];
            output[j++] = '*';
            changed = 1;
            continue;
        }

        /* operator chain */
        if (strchr("+-*/", input[i]))
        {
            int len = operator_chain_len(input, i);

            if (len == 2 && input[i] == input[i+1])
            {
                output[j++] = input[i];
                i++;
                changed = 1;
                continue;
            }

            if (len > 2 || (len == 2 && input[i] != input[i+1]))
            {
                return FIX_DOUBTFUL;
            }
        }

        output[j++] = input[i];
    }

    output[j] = '\0';

    if (changed)
        return FIX_SAFE;

    return FIX_NONE;
}
#include "../include/nirman.h"
#include <string.h>

/* ===============================
   GLOBAL CONTEXT
   =============================== */
NirmanContext NIRMAN_CTX = {0};

/* ===============================
   LIFECYCLE
   =============================== */

void nirman_start(void)
{
    NIRMAN_CTX.active = 1;

    /* reset flow */
    NIRMAN_CTX.flow_stage = 0;
    NIRMAN_CTX.goal[0] = '\0';
}

void nirman_stop(void)
{
    NIRMAN_CTX.active = 0;

    /* reset flow */
    NIRMAN_CTX.flow_stage = 0;
    NIRMAN_CTX.goal[0] = '\0';
}

int nirman_is_active(void)
{
    return NIRMAN_CTX.active;
}

/* ===============================
   INTERNAL HELPERS
   =============================== */

static void to_lower(char *s)
{
    for (int i = 0; s[i]; i++)
    {
        if (s[i] >= 'A' && s[i] <= 'Z')
            s[i] += 32;
    }
}

/* ===============================
   INTENT ENGINE
   =============================== */

int nirman_detect_intent(const char *input)
{
    if (!input)
        return 0;

    char temp[256];
    strncpy(temp, input, sizeof(temp) - 1);
    temp[255] = '\0';

    to_lower(temp);

    /* 1 = TALK */
    if (strstr(temp, "kaise") ||
        strstr(temp, "hello") ||
        strstr(temp, "hi") ||
        strstr(temp, "hy") ||
        strstr(temp, "radhe"))
    {
        return 1;
    }

    /* 2 = BUILD */
    if (strstr(temp, "bana") ||
        strstr(temp, "create") ||
        strstr(temp, "app") ||
        strstr(temp, "software"))
    {
        return 2;
    }

    /* 3 = ERROR / HELP */
    if (strstr(temp, "error") ||
        strstr(temp, "problem") ||
        strstr(temp, "issue"))
    {
        return 3;
    }

    return 0;
}
#include "../include/fix_rules.h"
#include "../include/fix_interactive.h"

#include <string.h>
#include <stdio.h>
#include <ctype.h>

/*──────────────────────────────────────────────
    FIX FUNCTIONS
──────────────────────────────────────────────*/

/* FIX: double operator (6++6 → 6+6) */
static int fix_double_operator(
    char *buffer,
    size_t buffer_size,
    const ShrijiErrorInfo *err)
{
    (void)buffer_size;
    (void)err;

    int changed = 0;

    for (int i = 0; buffer[i]; i++)
    {
        if ((buffer[i] == '+' || buffer[i] == '-' ||
             buffer[i] == '*' || buffer[i] == '/')
            && buffer[i] == buffer[i+1])
        {
            memmove(&buffer[i], &buffer[i+1], strlen(&buffer[i+1]) + 1);
            changed = 1;
            i--;
        }
    }

    return changed;
}

/* FIX: missing operator (6(7+7) → 6*(7+7)) */
static int fix_missing_operator(
    char *buffer,
    size_t buffer_size,
    const ShrijiErrorInfo *err)
{
    (void)err;

    for (int i = 0; buffer[i]; i++)
    {
        if (isdigit(buffer[i]) && buffer[i+1] == '(')
        {
            if (strlen(buffer) + 2 >= buffer_size)
                return 0;

            memmove(&buffer[i+2], &buffer[i+1], strlen(&buffer[i+1]) + 1);
            buffer[i+1] = '*';

            return 1;
        }
    }

    return 0;
}

/* FIX: missing value after '=' */
static int fix_missing_value_after_equal(
    char *buffer,
    size_t buffer_size,
    const ShrijiErrorInfo *err)
{
    (void)err;

    size_t len = strlen(buffer);

    if (len + 3 >= buffer_size)
        return 0;

    while (len > 0 && isspace((unsigned char)buffer[len - 1]))
        len--;

    if (len > 0 && buffer[len - 1] == '=')
    {
        strcat(buffer, " 0");
        return 1;
    }

    return 0;
}

/*──────────────────────────────────────────────
    FIX RULE TABLE
──────────────────────────────────────────────*/

static FixRule fix_rules[] = {

    /* deterministic: always correct */
    {
        .error_code = E_PARSE_MISSING_OPERATOR,
        .category = FIXCAT_MISSING_OPERATOR,
        .safety = FIX_SAFE_DETERMINISTIC,
        .confidence_penalty = 5,
        .max_attempt = 1,
        .reparse_required = 1,
        .apply_fix = fix_missing_operator
    },

    /* deterministic: assignment must have value */
    {
        .error_code = E_ASSIGN_02,
        .category = FIXCAT_MISSING_VALUE,
        .safety = FIX_SAFE_DETERMINISTIC,
        .confidence_penalty = 6,
        .max_attempt = 1,
        .reparse_required = 1,
        .apply_fix = fix_missing_value_after_equal
    },

    /* structural: operator cleanup (not always guaranteed) */
    {
        .error_code = E_PARSE_OPERATOR_CHAIN,
        .category = FIXCAT_EXTRA_OPERATOR,
        .safety = FIX_SAFE_STRUCTURAL,
        .confidence_penalty = 8,
        .max_attempt = 1,
        .reparse_required = 1,
        .apply_fix = fix_double_operator
    }

};

/*──────────────────────────────────────────────
    RULE LOOKUP
──────────────────────────────────────────────*/

static int fix_rule_count =
    sizeof(fix_rules) / sizeof(FixRule);

FixRule *shriji_get_rule_for_error(ShrijiErrorCode code)
{
    for (int i = 0; i < fix_rule_count; i++)
    {
        if (fix_rules[i].error_code == code)
            return &fix_rules[i];
    }

    return NULL;
}
#include "../include/smriti.h"
#include <string.h>

/* ------------------------------
   INTERNAL STATE (SILENT)
   ------------------------------ */

/* continuation (existing) */
static int    has_cont = 0;
static double cont_val = 0.0;
static char   cont_op  = 0;

/* STEP-11 context */
static unsigned long turn_counter = 0;
static char last_input_buf[512];
static int  last_intent_val = 0;

/* ------------------------------
   LIFECYCLE
   ------------------------------ */
void smriti_init(void)
{
    has_cont = 0;
    cont_val = 0.0;
    cont_op  = 0;

    turn_counter = 0;
    last_input_buf[0] = '\0';
    last_intent_val = 0;
}

/* ------------------------------
   CONTINUATION (AS-IS)
   ------------------------------ */
int smriti_has_continuation(void)
{
    return has_cont;
}

void smriti_set_continuation(double value, char op)
{
    has_cont = 1;
    cont_val = value;
    cont_op  = op;
}

void smriti_clear_continuation(void)
{
    has_cont = 0;
    cont_val = 0.0;
    cont_op  = 0;
}

double smriti_get_continuation_value(void)
{
    return cont_val;
}

char smriti_get_continuation_op(void)
{
    return cont_op;
}

/* ------------------------------
   STEP-11 CONTEXT (PASSIVE)
   ------------------------------ */
void smriti_next_turn(void)
{
    turn_counter++;
}

unsigned long smriti_get_turn(void)
{
    return turn_counter;
}

void smriti_set_last_input(const char *text)
{
    if (!text)
    {
        last_input_buf[0] = '\0';
        return;
    }
    /* safe snapshot */
    strncpy(last_input_buf, text, sizeof(last_input_buf) - 1);
    last_input_buf[sizeof(last_input_buf) - 1] = '\0';
}

const char* smriti_get_last_input(void)
{
    return last_input_buf;
}

void smriti_set_last_intent(int intent)
{
    last_intent_val = intent;
}

int smriti_get_last_intent(void)
{
    return last_intent_val;
}
#include <string.h>
#include <stdlib.h>
#include <ctype.h>
#include "../include/translator.h"

// --------------------------------------------------------
// INTERNAL STRUCT FOR KEYWORD MAPPING
// --------------------------------------------------------
typedef struct {
    const char *world;
    const char *shriji;
} TranslateEntry;

// --------------------------------------------------------
// WORLD → SHRIJILANG KEYWORD MAP (SMART ROUTING V3)
// --------------------------------------------------------
static TranslateEntry keyword_map[] = {

    // PRINT / OUTPUT
    {"print", "aaradhya"},
    {"echo", "aaradhya"},
    {"say", "sakhi"},
    {"speak", "sakhi"},
    {"tell", "sakhi"},

    // MIRA — Mathematical Solver
    {"solve", "mira"},
    {"calculate", "mira"},
    {"compute", "mira"},

    // SHIRI — Soft Guidance
    {"guide", "shiri"},
    {"help", "shiri"},
    {"direction", "shiri"},

    // NIYU — Logic / Why AI
    {"why", "niyu"},
    {"explain", "niyu"},
    {"reason", "niyu"},

    // KAVYA — Creative Engine
    {"create", "kavya"},
    {"write", "kavya"},
    {"story", "kavya"},

    // DEVOTIONAL → DIRECT KAVYA
    {"radhe", "kavya"},
    {"hare", "kavya"},
    {"krishna", "kavya"},
    {"shyam", "kavya"},

    // VARIABLE SYSTEM
    {"var", "mavi"},
    {"let", "mavi"},
    {"int", "mavi"},
    {"float", "mavi"}
};

// --------------------------------------------------------
// GET FIRST WORD
// --------------------------------------------------------
char* get_first_word(const char *source) {
    while (*source == ' ') source++;

    const char *start = source;
    while (*source && !isspace(*source))
        source++;

    int len = source - start;
    char *word = malloc(len + 1);
    strncpy(word, start, len);
    word[len] = '\0';

    return word;
}

// --------------------------------------------------------
// MAIN KEYWORD TRANSLATION
// --------------------------------------------------------
const char* translate_keyword(const char *word) {

for (size_t i = 0; i < sizeof(keyword_map)/sizeof(keyword_map[0]); i++) {

        if (strcmp(word, keyword_map[i].world) == 0)
            return keyword_map[i].shriji;
    }

    return NULL;
}

// --------------------------------------------------------
// SMART AUTOCORRECT FOR NATURAL INPUT
// --------------------------------------------------------
const char* mira_autocorrect(const char *word) {

    // Logical → Mira
    if (strcmp(word, "prit") == 0) return "mira";

    // Devotional → Kavya
    if (strcmp(word, "radhe") == 0) return "kavya";
    if (strcmp(word, "hare") == 0) return "kavya";
    if (strcmp(word, "krishna") == 0) return "kavya";

    // Emotional pain → Sakhi
    if (strcmp(word, "mera") == 0) return "sakhi";
    if (strcmp(word, "mann") == 0) return "sakhi";
    if (strcmp(word, "toot") == 0) return "sakhi";
    if (strcmp(word, "dukhi") == 0) return "sakhi";
    if (strcmp(word, "hurt") == 0) return "sakhi";
    if (strcmp(word, "sad") == 0) return "sakhi";

    // Creative autocorrect
    if (strcmp(word, "creat") == 0) return "kavya";

    // Shiri autocorrect
    if (strcmp(word, "gide") == 0) return "shiri";

    return NULL;
}
#include <stdio.h>
#include <string.h>

#include "../include/ai_router.h"
#include "../include/ai_intent.h"

#include "../pragya/contracts/l3_request.h"
#include "../pragya/contracts/l3_response.h"
#include "../pragya/contracts/l3_intent.h"

/*
   Legacy AI Router

   NOTE:
   ShrijiLang now uses KRST + Pragya Router
   for all intelligence routing.

   This module is preserved only for
   future conversational AI layer.
*/

L3Response ai_router_dispatch(const ShrijiBridgePacket *pkt)
{
    L3Response resp;

    resp.text = NULL;
    resp.success = 0;

    if (!pkt || !pkt->text)
        return resp;

    /* For now just acknowledge input */

    resp.text = "AI router currently inactive.";
    resp.success = 1;

    return resp;
}
#include <stdio.h>
#include "../include/state.h"

/*──────────────────────────────────────────────────────────────
 | SHRIJILANG — EXECUTION STATE ENGINE
 |
 | Purpose:
 |   - Runtime safety monitoring
 |   - Confidence / risk state management
 |   - Execution limits (anti-freeze)
 |   - Human confirmation escalation policy
 |
 | NOTE:
 |   This layer must remain deterministic.
 |   No external dependencies, no printing from core logic.
 *──────────────────────────────────────────────────────────────*/

/*──────────────────────────────────────────────
 | INTERNAL CONSTANTS (Tuning)
 *──────────────────────────────────────────────*/

/* Global execution budget:
 * Prevents accidental freeze when user writes an infinite loop.
 * Keep large enough so normal programs can run.
 */
#define SHRIJI_DEFAULT_MAX_STEPS   100000000LL  /* 10 crore */

/* Error escalation thresholds */
#define SHRIJI_WARN_THRESHOLD      3
#define SHRIJI_CRITICAL_THRESHOLD  5

/*──────────────────────────────────────────────
 | INIT
 *──────────────────────────────────────────────*/
void state_init(ExecutionState *state)
{
    if (!state) return;

    /* runtime safety status */
    state->safety = STATE_SAFE;

    /* confidence model */
    state->confidence = CONFIDENCE_HIGH;

    /* pressure model */
    state->error_pressure = 0;

    /* human control mode */
    state->human = HUMAN_AUTO_OK;

    /* execution budget */
    state->max_steps  = SHRIJI_DEFAULT_MAX_STEPS;
    state->steps_used = 0;
}

/*──────────────────────────────────────────────
 | RESET
 *──────────────────────────────────────────────*/
void state_reset(ExecutionState *state)
{
    state_init(state);
}

/*──────────────────────────────────────────────
 | ERROR HOOK
 *──────────────────────────────────────────────*/
void state_on_error(ExecutionState *state)
{
    if (!state) return;

    /* increase error pressure */
    if (state->error_pressure < 1000000)
        state->error_pressure++;

    /* stage-1: warning / risk */
    if (state->error_pressure >= SHRIJI_WARN_THRESHOLD) {
        state->confidence = CONFIDENCE_LOW;
        state->safety = STATE_RISK;
        state->human = HUMAN_CONFIRM_REQUIRED;
    }

    /* stage-2: critical stop */
    if (state->error_pressure >= SHRIJI_CRITICAL_THRESHOLD) {
        state->safety = STATE_CRITICAL;
        state->human = HUMAN_MANDATORY;
    }
}

/*──────────────────────────────────────────────
 | SUCCESS HOOK
 *──────────────────────────────────────────────*/
void state_on_success(ExecutionState *state)
{
    if (!state) return;

    /* decay pressure gradually */
    if (state->error_pressure > 0)
        state->error_pressure--;

    /* fully recover when stable again */
    if (state->error_pressure == 0) {
        state->confidence = CONFIDENCE_HIGH;
        state->safety = STATE_SAFE;
        state->human = HUMAN_AUTO_OK;
    }
}
/*=============================================================
  SHRIJILANG — TOKENIZER (PRODUCTION STABLE VERSION)
=============================================================*/

#include <stdio.h>
#include <string.h>
#include <ctype.h>

#include "../include/token.h"

/*──────────────────────────────────────────────
  INTERNAL STATE
──────────────────────────────────────────────*/

static const char *source = NULL;
static int current = 0;
static int line = 1;
static int col  = 1;

/*──────────────────────────────────────────────
  HELPERS
──────────────────────────────────────────────*/

static char peek(void)
{
    return source[current];
}

static char peek_next(void)
{
    if (source[current] == '\0') return '\0';
    return source[current + 1];
}

static char advance(void)
{
    char c = source[current++];

    if (c == '\n') {
        line++;
        col = 1;
    } else {
        col++;
    }

    return c;
}

static int match(char expected)
{
    if (peek() != expected) return 0;
    advance();
    return 1;
}

static Token make_token(TokenType type, const char *start, int length)
{
    Token tok;
    tok.type   = type;
    tok.start  = start;
    tok.length = length;
    tok.line   = line;
    tok.col    = col - length;
    return tok;
}

/*──────────────────────────────────────────────
  SKIP WHITESPACE + COMMENTS
──────────────────────────────────────────────*/

static void skip_whitespace(void)
{
    for (;;) {
        char c = peek();

        switch (c) {

            case ' ':
            case '\r':
            case '\t':
                advance();
                break;

            case '\n':
                advance();
                return;

            case '#':
                while (peek() != '\n' && peek() != '\0') {
                    advance();
                }
                if (peek() == '\n') advance();
                break;

            default:
                return;
        }
    }
}

/*──────────────────────────────────────────────
  NUMBER
──────────────────────────────────────────────*/

static Token number(void)
{
    const char *start = &source[current - 1];

    while (isdigit(peek()))
        advance();

    if (peek() == '.' && isdigit(peek_next())) {
        advance();
        while (isdigit(peek()))
            advance();
    }

    int length = &source[current] - start;
    return make_token(TOKEN_NUMBER, start, length);
}

/*──────────────────────────────────────────────
  IDENTIFIER / KEYWORD
──────────────────────────────────────────────*/

static TokenType check_keyword(const char *start, int length)
{
    if (length == 4 && strncmp(start, "mavi", 4) == 0) return TOKEN_MAVI;
    if (length == 4 && strncmp(start, "agar", 4) == 0) return TOKEN_AGAR;
    if (length == 5 && strncmp(start, "warna", 5) == 0) return TOKEN_WARNA;
    if (length == 4 && strncmp(start, "kaam", 4) == 0) return TOKEN_KAAM;
    if (length == 6 && strncmp(start, "jabtak", 6) == 0) return TOKEN_JABTAK;
    if (length == 5 && strncmp(start, "rukja", 5) == 0) return TOKEN_RUKJA;
    if (length == 5 && strncmp(start, "chalu", 5) == 0) return TOKEN_CHALU;
    if (length == 4 && strncmp(start, "nahi", 4) == 0) return TOKEN_NAHI;
    if (length == 4 && strncmp(start, "sahi", 4) == 0) return TOKEN_TRUE;
    if (length == 5 && strncmp(start, "galat", 5) == 0) return TOKEN_FALSE;


    if (length == 4 && strncmp(start, "bolo", 4) == 0) return TOKEN_BOLO;
    if (length == 6 && strncmp(start, "rachna", 6) == 0) return TOKEN_RACHNA;
    if (length == 5 && strncmp(start, "wapas", 5) == 0) return TOKEN_WAPAS;


    if (length == 5 && strncmp(start, "sakhi", 5) == 0) return TOKEN_SAKHI;
    if (length == 4 && strncmp(start, "niyu", 4) == 0) return TOKEN_NIYU;
    if (length == 4 && strncmp(start, "mira", 4) == 0) return TOKEN_MIRA;
    if (length == 5 && strncmp(start, "kavya", 5) == 0) return TOKEN_KAVYA;
    if (length == 5 && strncmp(start, "shiri", 5) == 0) return TOKEN_SHIRI;

    if (length == 6 && strncmp(start, "import", 6) == 0) return TOKEN_IMPORT;

    return TOKEN_IDENTIFIER;
}

static Token identifier(void)
{
    const char *start = &source[current - 1];

    while (isalnum(peek()) || peek() == '_')
        advance();

    int length = &source[current] - start;
    TokenType type = check_keyword(start, length);

    return make_token(type, start, length);
}

/*──────────────────────────────────────────────
  STRING
──────────────────────────────────────────────*/

static Token string(void)
{
    const char *start = &source[current - 1];

    while (peek() != '"' && peek() != '\0') {
        advance();
    }

    if (peek() == '\0') {
        return make_token(TOKEN_ERROR, start, 0);
    }

    advance();

    int length = &source[current] - start;
    return make_token(TOKEN_STRING, start, length);
}

/*──────────────────────────────────────────────
  INIT
──────────────────────────────────────────────*/

void init_tokenizer(const char *src)
{
    source  = src;
    current = 0;
    line    = 1;
    col     = 1;
}

/*──────────────────────────────────────────────
  MAIN SCANNER
──────────────────────────────────────────────*/

Token scan_token(void)
{
    skip_whitespace();

    const char *start = &source[current];

    if (peek() == '\0')
        return make_token(TOKEN_EOF, start, 0);

    char c = advance();

    if (isdigit(c))
        return number();

    if (isalpha(c) || c == '_')
        return identifier();

    switch (c) {

        case '+':
         if (match('+')) return make_token(TOKEN_PLUSPLUS, start, 2);
                   return make_token(TOKEN_PLUS, start, 1);

        case '-':
         if (match('-')) return make_token(TOKEN_MINUSMINUS, start, 2);
                   return make_token(TOKEN_MINUS, start, 1);

        case '*': return make_token(TOKEN_STAR, start, 1);
        case '/': return make_token(TOKEN_SLASH, start, 1);
        case '&':
    if (match('&')) return make_token(TOKEN_AND, start, 2);
             break;

        case '|':
    if (match('|')) return make_token(TOKEN_OR, start, 2);
             break;

        case '%': return make_token(TOKEN_MOD, start, 1);

        case '(': return make_token(TOKEN_LEFT_PAREN, start, 1);
        case ')': return make_token(TOKEN_RIGHT_PAREN, start, 1);
        case '{': return make_token(TOKEN_LEFT_BRACE, start, 1);
        case '}': return make_token(TOKEN_RIGHT_BRACE, start, 1);
        case '[': return make_token(TOKEN_LEFT_BRACKET, start, 1);
        case ']': return make_token(TOKEN_RIGHT_BRACKET, start, 1);

        case ',': return make_token(TOKEN_COMMA, start, 1);
        case ':': return make_token(TOKEN_COLON, start, 1);

        case '=':
            return make_token(match('=') ? TOKEN_EQEQ : TOKEN_EQUAL, start, 1);

        case '!':
            return make_token(match('=') ? TOKEN_NEQ : TOKEN_NOT, start, 1);

        case '>':
            return make_token(match('=') ? TOKEN_GTE : TOKEN_GT, start, 1);

        case '<':
            return make_token(match('=') ? TOKEN_LTE : TOKEN_LT, start, 1);

        case '\n':
            return make_token(TOKEN_NEWLINE, start, 1);

        case '"':
            return string();
    }

    /* Unknown / Unicode characters skip */
    return scan_token();
}
#include "../include/user_config.h"

/* ===== GLOBAL CONFIG ===== */

/* 🌸 DEV MODE */
int DEV_MODE = 0;  // 0 = USER | 1 = DEV

/* 🌸 LANGUAGE MODE */
int CURRENT_MODE = MODE_INDIA;

/* 🌸 PERSONALITY */
int CURRENT_PERSONALITY = P_SAKHI;

/* ===== INIT ===== */
void config_init()
{
    /* future config load */
}

/* ===== SETTERS ===== */
void config_set_personality(int p)
{
    CURRENT_PERSONALITY = p;
}

void config_set_mode(int mode)
{
    CURRENT_MODE = mode;
}
#ifndef SHRIJI_FIX_RULES_H
#define SHRIJI_FIX_RULES_H

#include "../include/error.h"

/*
    SHRIJILANG — FIX CATEGORY SYSTEM
    ---------------------------------
    Ye layer error code ko fix category me convert karegi.

    IMPORTANT:
    Fix engine directly error code check nahi karega.
    Pehle category milegi.
    Fir category ke basis pe rule apply hoga.
*/

/* ─────────────────────────────────────
   FIX CATEGORY ENUM
───────────────────────────────────── */

typedef enum {

    FIX_NONE = 0,              // No fix possible
    FIX_MISSING_VALUE,         // '=' ke baad value missing
    FIX_BAD_TOKEN,             // Unexpected / invalid token
    FIX_STRUCTURAL_SAFE        // Safe structural auto-correction

} FixCategory;


/* ─────────────────────────────────────
   CLASSIFICATION FUNCTION
───────────────────────────────────── */

/*
    Error code → Fix category mapping
*/
FixCategory shriji_classify_error(ShrijiErrorCode code);


#endif
/*=============================================================
  SHRIJI GYAAN ENGINE — FINAL COMPLETE RULE SET
=============================================================*/

#include <stddef.h>
#include "../../include/error.h"

/*──────────────────────────────────────────────
  STRUCT
──────────────────────────────────────────────*/
typedef struct {
    int code;

    const char *module;
    const char *file;
    const char *function;

    const char *rule;
    const char *explain;
    const char *hint;

} GyaanRule;

/*=============================================================
  🧠 RULE DATABASE START
=============================================================*/
static GyaanRule GYAAN_RULES[] = {

/*──────────── PARSER CORE ────────────*/

{E_PARSE_OPERATOR_START,"parser","parser.c","parse_expression",
"expression cannot start with operator",
"operator se expression start nahi hota",
"example: 10 + 2"},

{E_PARSE_OPERATOR_END,"parser","parser.c","parse_expression",
"operator must have value after it",
"operator ke baad value missing hai",
"example: 10 + 2"},

{E_PARSE_DOUBLE_OPERATOR,"parser","parser.c","parse_expression",
"two operators cannot come together",
"do operator ek saath nahi aate",
"example: 10 + 2"},

{E_PARSE_OPERATOR_CHAIN,"parser","parser.c","parse_expression",
"operator chain is invalid",
"operator chain galat hai",
"example: 10 + 2"},

{E_PARSE_INVALID_TOKEN,"parser","parser.c","parse_statement",
"statement must start with valid token",
"line ki shuruaat galat token se hui hai",
"use: bolo / mavi / agar"},

{E_PARSE_BRACKET_MISSING,"parser","parser.c","parse_primary",
"closing bracket missing",
"bracket close nahi hua",
"example: (10 + 2)"},

/*──────────── PARSER ADVANCED ────────────*/

{E_PARSE_EXTRA_OPERATOR,"parser","parser.c","parse_expression",
"extra operator found",
"extra operator detect hua",
"example: 10 + 2"},

{E_PARSE_MISSING_OPERATOR,"parser","parser.c","parse_expression",
"operator missing between values",
"values ke beech operator missing hai",
"example: 10 + 2"},

{E_PARSE_UNMATCHED_PAREN,"parser","parser.c","parse_primary",
"parentheses mismatch",
"bracket match nahi hua",
"example: (10 + 2)"},

{E_PARSE_BRACKET_EXTRA,"parser","parser.c","parse_primary",
"extra closing bracket",
"extra bracket mila",
"bracket remove karein"},

{E_PARSE_MISSING_OPERAND,"parser","parser.c","parse_expression",
"operand missing",
"value missing hai operator ke liye",
"example: 10 + 2"},

/*──────────── ASSIGNMENT ────────────*/

{E_ASSIGN_01,"runtime","interpreter.c","AST_IDENTIFIER",
"variable must be declared before use",
"variable declare nahi hua",
"use: mavi x = 10"},

{E_ASSIGN_02,"parser","parser.c","parse_assignment",
"assignment must have value",
"assignment me value missing hai",
"example: mavi x = 10"},

/*──────────── FUNCTION ────────────*/

{E_PARSE_01,"parser","parser.c","general",
"syntax error",
"syntax samajh nahi aaya",
"statement check karein"},

{E_PARSE_02,"system","multiple","general",
"invalid usage",
"usage galat hai",
"syntax check karein"},

/*──────────── CONTROL FLOW ────────────*/

{E_IF_01,"parser","parser.c","parse_if",
"invalid if condition",
"if condition galat hai",
"example: agar x > 0 { ... }"},

{E_RUNTIME_LOOP_LIMIT,"runtime","interpreter.c","AST_WHILE",
"loop exceeded safe limit",
"loop limit cross ho gaya",
"condition fix karein"},

/*──────────── DATA STRUCTURES ────────────*/

{E_RUNTIME_INDEX_ERROR,"runtime","interpreter.c","AST_INDEX",
"index out of bounds",
"index range ke bahar hai",
"example: a[0]"},

/*──────────── RUNTIME ────────────*/

{E_RUNTIME_DIV_ZERO,"runtime","interpreter.c","safe_div",
"division by zero not allowed",
"0 se divide nahi kar sakte",
"example: 10 / 2"},

{E_RUNTIME_TYPE_MISMATCH,"runtime","interpreter.c","AST_BINARY",
"type mismatch in expression",
"different type values use hue hain",
"same type use karein"},

{E_RUNTIME_01,"runtime","interpreter.c","general",
"runtime error occurred",
"runtime me issue aaya",
"logic check karein"}

};

/*=============================================================
  🧠 RULE DATABASE END
=============================================================*/

/*──────────────────────────────────────────────
  COUNT
──────────────────────────────────────────────*/
static int GYAAN_COUNT =
sizeof(GYAAN_RULES) / sizeof(GYAAN_RULES[0]);

/*──────────────────────────────────────────────
  LOOKUP
──────────────────────────────────────────────*/
const GyaanRule* gyaan_lookup(int code)
{
    for (int i = 0; i < GYAAN_COUNT; i++) {
        if (GYAAN_RULES[i].code == code) {
            return &GYAAN_RULES[i];
        }
    }
    return NULL;
}
#ifndef GYAAN_ENGINE_H
#define GYAAN_ENGINE_H

#include "../../include/error.h"

typedef struct {

    ShrijiErrorCode error;

    const char *module;
    const char *file;
    const char *function;

    const char *rule;
    const char *explain;
    const char *hint;

} GyaanEntry;

const GyaanEntry* gyaan_lookup(ShrijiErrorCode code);

void gyaan_print(const ShrijiErrorInfo *err);

#endif
/*=============================================================
  SHRIJI GYAAN ENGINE — CORE (HINGLISH + CLEAR UX VERSION)
=============================================================*/

#include <stdio.h>
#include <string.h>
#include "../../include/error.h"

typedef struct {
    int code;

    const char *module;
    const char *file;
    const char *function;

    const char *rule;
    const char *explain;
    const char *hint;

} GyaanRule;

extern const GyaanRule* gyaan_lookup(int code);

static const char* safe(const char *s)
{
    return (s && *s) ? s : "N/A";
}

void gyaan_print(const ShrijiErrorInfo *err)
{
    if (!err) return;

    const GyaanRule *r = gyaan_lookup(err->code);

    /* SOFT FUNCTION MATCH */
    if (r && err->function && r->function) {
        if (strcmp(err->function, r->function) != 0) {
            /* allow fallback */
        }
    }

    /* FALLBACK (NO RULE FOUND) */
    if (!r) {
        printf("\nSHRIJI INSIGHT\n");
        printf("──────────────────────────────\n");

        printf("System ko is error ka exact explanation nahi mila\n");

        printf("Samjho   : input ka structure (syntax) sahi format me nahi hai\n");

        printf("Hint     : statement ko proper format me likho\n");
        printf("Example  : bolo 5 + 2\n");

        if (err->has_location) {
            printf("\nLocation : line %d, col %d\n", err->line, err->col);
        }

        printf("──────────────────────────────\n");
        return;
    }

    /* NORMAL RULE PRINT */
    printf("\nSHRIJI INSIGHT\n");
    printf("──────────────────────────────\n");

    printf("Module   : %s\n", safe(r->module));
    printf("File     : %s\n", safe(r->file));
    printf("Function : %s\n\n", safe(r->function));

    printf("Rule     : %s\n", safe(r->rule));
    printf("Samjho   : %s\n", safe(r->explain));
    printf("Hint     : %s\n", safe(r->hint));

    switch (err->code)
    {
        case E_PARSE_MISSING_OPERATOR:
            printf("\nNote     : yahan values ke beech operator missing tha\n");
            break;

        case E_PARSE_DOUBLE_OPERATOR:
            printf("\nNote     : do operator ek saath aa gaye the\n");
            break;

        case E_PARSE_OPERATOR_START:
            printf("\nNote     : expression operator se start nahi hota\n");
            break;

        case E_PARSE_OPERATOR_END:
            printf("\nNote     : expression operator par khatam ho gaya\n");
            break;

        default:
            break;
    }

    if (err->has_location) {
        printf("\nLocation : line %d, col %d\n", err->line, err->col);
    }

    printf("──────────────────────────────\n");
}
#include <stdio.h>
#include <string.h>

#include "../include/krst_router.h"
#include "../krst/krst_core.h"
#include "../include/pragya_avastha.h"

#include "../include/error.h"
#include "../include/error_intelligence.h"

/* ===============================
   INPUT NORMALIZATION (SAFE)
   =============================== */
static void normalize_input(const char *input, char *output)
{
    if (!input || !output) return;

    int i = 0;
    for (; input[i]; i++) {
        if (input[i] >= 'A' && input[i] <= 'Z')
            output[i] = input[i] + 32;
        else
            output[i] = input[i];
    }
    output[i] = '\0';
}

/* ===============================
   COMMAND DETECTION
   =============================== */
static char detect_command(const char *input)
{
    if (!input) return 'G';

    if (input[0] == 'g' && input[1] == ' ')
        return 'G';

    if (input[0] == 'e' && input[1] == ' ')
        return 'E';

    return 'G';
}

/* ===============================
   CONTEXT REWRITE (C2 COMPLETE)
   =============================== */
static int rewrite_context_input(KRSTRequest *req)
{
    if (!req || !req->input_text) return 0;

    int num = 0;
    static char buffer[128];

    const char *input = req->input_text;

    /* ===== ADD ===== */
    if ((strstr(input, "add") || strstr(input, "jodo")) &&
        sscanf(input, "usme %d", &num) == 1) {

        snprintf(buffer, sizeof(buffer), "bolo last + %d", num);
        req->input_text = buffer;
        return 1;
    }

    /* ===== MINUS ===== */
    if ((strstr(input, "minus") || strstr(input, "ghatao") || strstr(input, "subtract")) &&
        (sscanf(input, "usme %d", &num) == 1 ||
         sscanf(input, "usme se %d", &num) == 1)) {

        snprintf(buffer, sizeof(buffer), "bolo last - %d", num);
        req->input_text = buffer;
        return 1;
    }

    /* ===== MULTIPLY ===== */
    if ((strstr(input, "multiply") || strstr(input, "guna") || strstr(input, "into")) &&
        sscanf(input, "usme %d", &num) == 1) {

        snprintf(buffer, sizeof(buffer), "bolo last * %d", num);
        req->input_text = buffer;
        return 1;
    }

    /* ===== DIVIDE ===== */
    if ((strstr(input, "divide") || strstr(input, "bhaag") || strstr(input, "divided")) &&
        (sscanf(input, "usme %d", &num) == 1 ||
         sscanf(input, "usme se %d", &num) == 1)) {

        snprintf(buffer, sizeof(buffer), "bolo last / %d", num);
        req->input_text = buffer;
        return 1;
    }

    return 0;
}

/* ===============================
   MAIN ROUTER
   =============================== */
void krst_route_request(KRSTRequest *req)
{
    if (!req || !req->input_text)
        return;

    /* ===== SAFE NORMALIZATION ===== */
    static char normalized[256];

    normalize_input(req->input_text, normalized);
    req->input_text = normalized;

    /* ===== COMMAND DETECTION ===== */
    char cmd = detect_command(req->input_text);

    /* ===== COMMAND ROUTING ===== */

     if (cmd == 'E') {
    req->input_text += 2;

    static char explain_buf[256];
   snprintf(explain_buf, sizeof(explain_buf), "bolo %.240s", req->input_text);
    req->input_text = explain_buf;

    printf("[EXPLAIN MODE]\n");
}

    else if (cmd == 'G') {
        if (req->input_text[0] == 'g' && req->input_text[1] == ' ')
            req->input_text += 2;
    }

    /* ===== CONTEXT REWRITE ===== */
    rewrite_context_input(req);

    /* ===== CREATE AVASTHA ===== */
    PragyaAvastha avastha;
    memset(&avastha, 0, sizeof(PragyaAvastha));

    avastha.raw_text = req->input_text;

      avastha.raw_text = req->input_text;

/*  NEW: pass input mode */
    avastha.is_program_mode = req->is_program_mode;

/* ===== KRST CORE ===== */
krst_process(&avastha);

/* ===== ERROR INTELLIGENCE FALLBACK ===== */
if (!avastha.has_correction && error_reported)
{
    shriji_error_intelligence(&avastha, &LAST_ERROR, NULL);
}

/* ===== RETURN CORRECTION ===== */
if (avastha.has_correction)
{
    req->has_correction = 1;

    strncpy(req->corrected_text,
            avastha.corrected_text,
            sizeof(req->corrected_text) - 1);

    req->corrected_text[
        sizeof(req->corrected_text) - 1] = '\0';
}
}
/* RADHE_RADHE_SHREEJI_RADHE_RANI */

#include <stdio.h>
#include <string.h>
#include <ctype.h>

#include "../include/interpreter.h"
#include "../include/error.h"
#include "../include/fix_engine.h"
#include "../include/nirman.h"
#include "../include/parser.h"
#include "../include/value.h"
#include "../include/env.h"
#include "../include/runtime.h"
#include "../include/pragya_avastha.h"
#include "../include/decision_engine.h"

#include "krst_core.h"
#include "krst_types.h"

#include "../pragya/router/l3_router.h"
#include "../include/error_intelligence.h"

#include "../include/user_config.h"

extern Env *GLOBAL_ENV;

/* ================= SESSION ================= */

static int session_confidence = 100;
static int session_risk = 4;
static int success_streak = 0;

/* ================= LABEL ================= */

static const char* teach_label(int level)
{
    switch(level)
    {
        case KRST_TEACH_SILENT:  return "SILENT";
        case KRST_TEACH_HINT:    return "HINT";
        case KRST_TEACH_EXPLAIN: return "EXPLAIN";
        case KRST_TEACH_DEEP:    return "DEEP";
        case KRST_TEACH_TRAIN:   return "TRAIN";
        default: return "SILENT";
    }
}

static const char* tone_label(int tone)
{
    switch(tone)
    {
        case KRST_TONE_NEUTRAL: return "NEUTRAL";
        case KRST_TONE_STRICT:  return "STRICT";
        case KRST_TONE_ALERT:   return "ALERT";
        default: return "NEUTRAL";
    }
}

/* ================= THRESHOLD ================= */

static void apply_thresholds(KRSTDecision *d)
{
    int c = d->confidence_score;

    if (c <= 20)      d->teaching_level = KRST_TEACH_TRAIN;
    else if (c <= 40) d->teaching_level = KRST_TEACH_DEEP;
    else if (c <= 60) d->teaching_level = KRST_TEACH_EXPLAIN;
    else if (c <= 80) d->teaching_level = KRST_TEACH_HINT;
    else              d->teaching_level = KRST_TEACH_SILENT;

    if (session_risk >= 76)
    {
        d->tone = KRST_TONE_ALERT;
        d->escalate = 1;
    }
    else if (session_risk >= 30)
    {
        d->tone = KRST_TONE_STRICT;
        d->escalate = 0;
    }
    else
    {
        d->tone = KRST_TONE_NEUTRAL;
        d->escalate = 0;
    }
}

/* ================= MAIN ================= */

int krst_process(PragyaAvastha *avastha)
{
    if (!avastha || !avastha->raw_text)
        return 1;

const char *input = avastha->raw_text;

/*  NIRMAN START */
if (strcmp(input, "nirman") == 0) {
    nirman_start();
    printf("Nirman mode activated\n");
    printf("Welcome to Shriji World\n");
    return 0;
}

/*  NIRMAN STOP */
if (strcmp(input, "exit nirman") == 0) {
    nirman_stop();
    printf("Nirman mode deactivated\n");
    return 0;
}

/*  NIRMAN MODE ROUTING */
if (nirman_is_active()) {

    int intent = nirman_detect_intent(input);

    printf("SHRIJI (intent=%d)\n", intent);

    return 0;
}


    if (DEV_MODE)
        printf("[KRST] Input: %s\n", input);

    error_reported = 0;
    avastha->stop_execution = 0;

    avastha->confidence = session_confidence;
    avastha->risk = session_risk;

    /* ================= PARSE ================= */

    int was_fixed = 0;
    int fix_penalty = 0;

    ASTNode *root = NULL;

    if (avastha->is_program_mode)
        root = parse_program(input);
    else
        root = language_execute_with_fix(input, &was_fixed, &fix_penalty);

    /* ================= ERROR ================= */

    if (!root || error_reported)
    {
        avastha->ast = NULL;
        avastha->stop_execution = 1;
        success_streak = 0;

        const ShrijiErrorInfo *err_ptr = shriji_last_error();

        if (!err_ptr)
            return 0;

        ShrijiErrorInfo err_safe = *err_ptr;

        DecisionType d_type = shriji_take_decision(
            session_confidence,
            session_risk,
            err_safe.code
        );

        char fixed[256];
        FixType ftype = fix_apply(input, fixed);

        /* AUTO FIX ONLY */
        if (ftype == FIX_SAFE && d_type == DECISION_AUTO_FIX)
        {
            printf("\n[Shriji] auto fix → %s\n\n", fixed);

            avastha->has_correction = 1;

            snprintf(avastha->corrected_text,
                     sizeof(avastha->corrected_text),
                     "%s", fixed);

            return 0;
        }

        /*  NO EI CALL HERE */
        return 0;
    }

    /* ================= SUCCESS ================= */
   if (!root || error_reported || avastha->stop_execution)
{
    return 0;
}
    avastha->ast = root;

    ShrijiRuntime runtime;
    runtime_init(&runtime);

    Value result = run_program(root, GLOBAL_ENV, &runtime);

    if (!error_reported && !runtime.error_flag)
{
        success_streak++;

        int recovery = 5 + (success_streak * 2);
        if (recovery > 25) recovery = 25;

        session_confidence += recovery;
        session_risk -= recovery;

        if (was_fixed)
        {
            session_confidence -= fix_penalty;
            session_risk += fix_penalty;
        }

        if (result.type == VAL_NUMBER)
            printf("%g\n", result.number);
        else if (result.type == VAL_STRING && result.string)
            printf("%s\n", result.string);
    }

    value_free(&result);

    if (session_confidence < 0) session_confidence = 0;
    if (session_confidence > 100) session_confidence = 100;

    if (session_risk < 4) session_risk = 4;
    if (session_risk > 100) session_risk = 100;

    KRSTDecision decision = {0};
    decision.confidence_score = session_confidence;

    apply_thresholds(&decision);

    if (DEV_MODE)
    {
        printf("[KRST] Confidence: %d | Risk: %d | Teach: %s | Tone: %s | Esc: %d\n",
               session_confidence,
               session_risk,
               teach_label(decision.teaching_level),
               tone_label(decision.tone),
               decision.escalate);
    }

    avastha->confidence = session_confidence;
    avastha->risk = session_risk;
    avastha->teach_level = decision.teaching_level;
    avastha->tone = decision.tone;

    pragya_route(avastha);

    return 0;
}
#ifndef KRST_TYPES_H
#define KRST_TYPES_H

#include "../include/env.h"

/* ==============================
   KRST ACTION ENUM
============================== */

typedef enum {
    KRST_ACTION_ALLOW = 0,
    KRST_ACTION_BLOCK,
    KRST_ACTION_TEACH,
    KRST_ACTION_ESCALATE
} KRSTAction;

/* ==============================
   KRST TEACHING LEVEL ENUM
============================== */

typedef enum {
    KRST_TEACH_SILENT = 0,
    KRST_TEACH_HINT,
    KRST_TEACH_EXPLAIN,
    KRST_TEACH_DEEP,
    KRST_TEACH_TRAIN
} KRSTTeachingLevel;

/* ==============================
   KRST TONE ENUM
============================== */

typedef enum {
    KRST_TONE_NEUTRAL = 0,
    KRST_TONE_SUPPORTIVE,
    KRST_TONE_GUIDING,
    KRST_TONE_STRICT,
    KRST_TONE_ALERT
} KRSTTone;

/* ==============================
   KRST REQUEST STRUCT
============================== */
typedef struct
{
    const char *input_text;

     int is_program_mode;
    /* correction support */

    char corrected_text[512];
    int has_correction;

} KRSTRequest;
/* ==============================
   KRST DECISION STRUCT
============================== */

typedef struct {
    KRSTAction action;

    int confidence_score;   /* 0–100 */
    int risk_score;         /* 0–100 */

    int allow_execution;
    int reason_code;

    KRSTTeachingLevel teaching_level;
    KRSTTone tone;

    int escalate;   /* 1 = escalation required */

} KRSTDecision;

#endif
#ifndef KRST_CORE_H
#define KRST_CORE_H

#include "../include/pragya_avastha.h"

/*
    KRST Core Interface

    Role:
    - Main processing engine
    - Takes PragyaAvastha as input
    - Updates it (error, correction, teaching, etc.)
*/

/* Main processing function */
int krst_process(PragyaAvastha *avastha);

#endif
#include <string.h>
#include "niyu_engine.h"

int niyu_analyze(const char *input)
{
    if (!input)
        return 0;

    if (strlen(input) == 0)
        return 0;

    return 0;
}
# ============================================================
# SHRIJILANG — UNIFIED BUILD SYSTEM (KRST + PRAGYA)
# ============================================================

CC = gcc

CFLAGS = -Wall -Wextra -std=c11 -O2 \
	-Iinclude \
	-Isrc \
	-Ienv \
	-Ipragya \
	-Ikrst

# ============================================================
# CORE SOURCE FILES
# ============================================================

SRC = \
	src/token.c \
	src/ast.c \
	src/parser.c \
	gyaan/core/gyaan_engine.c \
	gyaan/data/gyaan_rules.c \
	src/interpreter.c \
	src/runtime.c \
	src/error.c \
	src/state.c \
	src/event.c \
	src/command.c \
	src/translator.c \
	src/example_builder.c \
	src/ai_router.c \
	src/ai_intent.c \
	src/atma_lekha.c \
	src/smriti.c \
        src/smriti_session.c \
	src/smriti_personal.c \
	src/main.c \
	env/env.c \
	src/lang/grammar/grammar_registry.c \
	src/lang/grammar/grammar_core.c \
	pragya/router/l3_router.c \
	pragya/niyu/niyu_core.c \
	pragya/sakhi/sakhi_core.c \
	pragya/mira/mira_core.c \
	krst/krst_core.c \
	krst/krst_router.c \
	krst/niyu_engine.c \
	src/fix_engine.c \
	src/fix_rules.c \
	src/fix_interactive.c \
	src/typo_engine.c \
	src/error_intelligence.c \
	src/decision_engine.c \
src/user_config.c \
src/nirman.c \
src/normalize_engine.c \
src/nirman_flow_engine.c \

# ============================================================
# OUTPUT
# ============================================================

OUT = shrijilang

# ============================================================
# BUILD RULES
# ============================================================

$(OUT): $(SRC)
	$(CC) $(CFLAGS) $(SRC) -o $(OUT) -lreadline

clean:
	rm -f $(OUT)
	rm -f *.o
	rm -f src/*.o
	rm -f env/*.o
	rm -f src/lang/grammar/*.o
	rm -f pragya/router/*.o
	rm -f pragya/niyu/*.o
	rm -f pragya/sakhi/*.o
	rm -f pragya/mira/*.o
	rm -f krst/*.o

run: $(OUT)
	./$(OUT)

rebuild: clean $(OUT)
#include "sakhi_core.h"
#include <stdio.h>

L3Response sakhi_speak(const L3Request *request,
                       const NiyuResult *result)
{
    (void)result;

    L3Response res;

    if (!request || !request->raw_input)
    {
        res.text = "Input samajh nahi aaya.";
        res.success = 0;
        return res;
    }

    printf("[PRAGYA] Main sun rahi hoon. Thoda clear likhoge to aur achha ho jayega 🙂\n");

    res.text = request->raw_input;
    res.success = 1;

    return res;
}
#ifndef SAKHI_CORE_H
#define SAKHI_CORE_H

#include "sakhi_contract.h"

/*
 * SAKHI CORE
 *
 * Implementation file will:
 * - handle tone / language / clarity
 * - never inspect raw logic
 *
 * This header intentionally exposes no helpers.
 */

#endif /* SAKHI_CORE_H */
#ifndef SAKHI_CONTRACT_H
#define SAKHI_CONTRACT_H

#include "../contracts/l3_request.h"
#include "../contracts/l3_response.h"
#include "../niyu/niyu_contract.h"

/*
 * SAKHI CONTRACT
 *
 * Sakhi is the voice of the system.
 * - Receives decided intent + thinking result (opaque)
 * - Produces final human-readable response
 *
 * Sakhi never:
 * - thinks or calculates
 * - decides intent
 * - teaches (that is Mira)
 */

/* Forward declaration (opaque thinking result) */
/*
 * sakhi_speak
 *
 * Final response generator.
 * - Called only by Router
 * - Uses intent already decided
 * - Formats language and tone
 */
L3Response sakhi_speak(const L3Request *request,
                       const NiyuResult *result);

#endif /* SAKHI_CONTRACT_H */
#ifndef L3_INTENT_H
#define L3_INTENT_H

/* 
 * L3 Intent:
 * Router decides intent.
 * Niyu never decides intent.
 * Sakhi/Mira never guess intent.
 */

typedef enum {
    L3_INTENT_UNKNOWN = 0,

    /* Explanation / reasoning */
    L3_INTENT_EXPLAIN,
    L3_INTENT_WHY,

    /* Pure calculation / logic */
    L3_INTENT_CALC,

    /* Teaching / learning */
    L3_INTENT_TEACH,

    /* Emotional / tone-based */
    L3_INTENT_EMOTION,

    /* Explicit command */
    L3_INTENT_COMMAND

} L3Intent;

#endif /* L3_INTENT_H */
#ifndef L3_RULES_H
#define L3_RULES_H

/*
 * L3 NON-NEGOTIABLE LAWS
 *
 * 1. Router decides intent.
 * 2. Niyu THINKS, never SPEAKS.
 * 3. Sakhi SPEAKS, never THINKS.
 * 4. Mira TEACHES only when intent == L3_INTENT_TEACH.
 *    No external permission, no memory switch.
 * 5. No direct Niyu <-> Sakhi calls.
 * 6. All communication flows through Router.
 * 7. Contracts are stable forever.
 * 8. Engine may change, behavior must not.
 */

#define L3_RULE_ROUTER_IS_AUTHORITY 1
#define L3_RULE_NIYU_SILENT         1
#define L3_RULE_SAKHI_NO_LOGIC      1
#define L3_RULE_MIRA_EXPLICIT_ONLY  1

#endif /* L3_RULES_H */
#ifndef L3_REQUEST_H
#define L3_REQUEST_H

#include "l3_intent.h"

/*
 * L3Request:
 * Immutable input structure.
 * Router decides intent.
 * No teaching_mode flag here.
 * Behavior intent se hi control hoga.
 */

typedef struct {
    const char *raw_input;     /* original user text */
    L3Intent    intent;        /* decided by router */
} L3Request;

#endif /* L3_REQUEST_H */
#ifndef L3_RESPONSE_H
#define L3_RESPONSE_H

/*
 * L3Response:
 * Only Sakhi or Mira can create final response.
 * Niyu can never write directly here.
 */

typedef struct {
    const char *text;      /* final output */
    int         success;   /* logical success/failure */
} L3Response;

#endif /* L3_RESPONSE_H */
// JAI SHREE KRISHNA
#include "mira_contract.h"
#include "../niyu/niyu_contract.h"

#include "../../include/error.h"
#include "../../include/pragya_avastha.h"

#include <stdio.h>
#include <string.h>


/* =========================================
   POINTER PRINTER
========================================= */

static void print_pointer(const char *text, int pos)
{
    if (!text || pos < 0)
        return;

    printf("%s\n", text);

    for (int i = 0; i < pos; i++)
        printf(" ");

    printf("^\n\n");
}


/* =========================================
   MIRA TEACHING ENGINE
========================================= */

void mira_teach_avastha(PragyaAvastha *a, NiyuResult *res)
{
    if (!a || !res)
        return;

    const char *text = a->raw_text;


    /* valid input → silent */

    if (res->error_type == NIYU_ERR_VALID)
        return;


    switch (a->error_code)
    {

    /* =========================================
       DOUBLE OPERATOR
    ========================================= */

    case E_PARSE_DOUBLE_OPERATOR:
    case E_PARSE_OPERATOR_CHAIN:

        printf("[SHRI_JI] Lagta hai yahaan extra operator aa gaya hai \n\n");

        print_pointer(text, res->error_pos);

        printf("Do operators ek saath allowed nahi hote.\n\n");

        printf("Example:\n");
        printf("6 + 5\n");

        break;


    /* =========================================
       OPERATOR AT START
    ========================================= */

    case E_PARSE_OPERATOR_START:

        printf("[SHRI_JI] Expression operator se start nahi hota \n\n");

        print_pointer(text, res->error_pos);

        printf("Pehle ek value honi chahiye.\n\n");

        printf("Example:\n");
        printf("6 + 5\n");

        break;


    /* =========================================
       OPERATOR AT END
    ========================================= */

    case E_PARSE_OPERATOR_END:

        printf("[SHRI_JI] Expression incomplete lag raha hai \n\n");

        print_pointer(text, res->error_pos);

        printf("Operator ke baad ek value expected hoti hai.\n\n");

        printf("Example:\n");
        printf("6 + 5\n");

        break;


    /* =========================================
       MISSING OPERATOR
    ========================================= */

    case E_PARSE_MISSING_OPERATOR:

        printf("[SHRI_JI] Lagta hai operator missing hai \n\n");

        print_pointer(text, res->error_pos);

        printf("Do values ke beech operator hona chahiye.\n\n");

        printf("Example:\n");
        printf("5 + 6\n");

        break;


    /* =========================================
       ASSIGNMENT ERROR
    ========================================= */

    case E_ASSIGN_02:

        printf("[SHRI_JI] Assignment incomplete lag raha hai.\n\n");

        print_pointer(text, res->error_pos);

        printf("'=' ke baad value expected hai.\n\n");

        printf("Example:\n");
        printf("mavi x = 10\n");

        break;


    /* =========================================
       GENERIC SYNTAX ERROR
    ========================================= */

    case E_PARSE_02:

        printf("[SHRI_JI] Syntax error detect hua hai.\n\n");

        print_pointer(text, res->error_pos);

        printf("Yahan invalid token ya unsupported symbol mila hai.\n\n");

        printf("Example valid syntax:\n");
        printf("6 + 5\n");

        break;


    /* =========================================
       DEFAULT
    ========================================= */

default:
    return;
    }
}
#ifndef MIRA_CONTRACT_H
#define MIRA_CONTRACT_H

#include "../contracts/l3_request.h"
#include "../contracts/l3_response.h"
#include "../niyu/niyu_contract.h"
#include "../../include/pragya_avastha.h"

/*
 * MIRA CONTRACT
 *
 * Mira is the teaching layer.
 * - Activated only when teaching_mode == 1
 * - Provides structured explanation / lesson
 *
 * Mira never:
 * - handles normal conversation
 * - decides intent
 * - performs deep logic (that is Niyu)
 */

/*
 * Legacy interface (old L3 system)
 */
L3Response mira_teach(const L3Request *request);

/*
 * New ShrijiLang teaching interface
 * Used by Pragya Router
 */

void mira_teach_avastha(
    PragyaAvastha *avastha,
    NiyuResult *result
);

#endif /* MIRA_CONTRACT_H */
#ifndef MIRA_CORE_H
#define MIRA_CORE_H

#include "mira_contract.h"

/*
 * MIRA CORE
 *
 * Implementation file will:
 * - manage lesson flow
 * - provide explanations
 *
 * No other module may call Mira directly.
 */

#endif /* MIRA_CORE_H */
L3 (KR̥ṢṬ) — FOUNDATION DOCUMENT
===============================

This document is the final, non-negotiable foundation of L3.
It defines roles, laws, flow, and future guarantees.
Implementation may evolve. Behavior must not.

----------------------------------------------------------------
1. PURPOSE
----------------------------------------------------------------
L3 is the intelligence orchestration layer of ShrijiLang.

It separates:
- Thinking (logic, calculation)
- Speaking (human-facing output)
- Teaching (explicit learning mode)

L3 is designed to work today with lightweight logic
and tomorrow with machine-backed, high-compute intelligence
without breaking language behavior.

----------------------------------------------------------------
2. CORE ROLES (FROZEN)
----------------------------------------------------------------

Router
------
- Central authority of L3
- Decides intent
- Controls full request lifecycle
- Chooses which module participates

Router never:
- Thinks deeply
- Formats language
- Teaches directly

Niyu (THINK)
------------
- Silent thinking engine
- Performs logic, reasoning, calculation
- Produces internal result only

Niyu never:
- Speaks to user
- Decides intent
- Formats language

Sakhi (SPEAK)
-------------
- Final voice of the system
- Produces human-readable output
- Applies tone and clarity

Sakhi never:
- Thinks or calculates
- Decides intent
- Teaches

Mira (TEACH)
------------
- Teaching and explanation layer
- Activated only when explicitly allowed
- Provides structured learning responses

Mira never:
- Handles normal conversation
- Decides intent
- Performs deep logic

----------------------------------------------------------------
3. FLOW (CANONICAL)
----------------------------------------------------------------

User Input
   |
   v
Router
   |
   |-- decide intent
   |-- check teaching_mode
   |
   |--> Niyu (if logic / calculation needed)
   |        |
   |        v
   |   NiyuResult (internal)
   |
   |--> Sakhi (normal response)
   |        OR
   |--> Mira  (teaching response)
   |
   v
Final L3Response to User

No module may bypass the Router.

----------------------------------------------------------------
4. CONTRACTS (IMMUTABLE)
----------------------------------------------------------------
- l3_request.h  : defines input shape
- l3_intent.h   : defines intent space
- l3_response.h : defines output shape
- l3_rules.h    : defines laws

Contracts are stable forever.
Engines may change. Contracts must not.

----------------------------------------------------------------
5. NON-NEGOTIABLE LAWS
----------------------------------------------------------------
1. Router is the sole authority.
2. Niyu thinks, never speaks.
3. Sakhi speaks, never thinks.
4. Mira teaches only when explicitly enabled.
5. No direct Niyu <-> Sakhi communication.
6. All flows pass through Router.
7. Behavior consistency is mandatory.
8. Future machine intelligence must fit this structure.

----------------------------------------------------------------
6. FUTURE GUARANTEE
----------------------------------------------------------------
L3 is intentionally designed as a replaceable intelligence layer.

Today:
- Local logic
- Lightweight reasoning

Tomorrow:
- Dedicated machines
- Heavy computation
- Real-world modeling

The language behavior, contracts, and flow
must remain unchanged across this evolution.

----------------------------------------------------------------
7. FINAL NOTE
----------------------------------------------------------------
This document is the NĒEV (foundation) of L3.

Any change that violates this document
is considered a breaking change and is not allowed.

Implementation follows this document.
This document never follows implementation.
#include "l3_router.h"

#include "../niyu/niyu_contract.h"
#include "../sakhi/sakhi_contract.h"
#include "../mira/mira_contract.h"

#include "../../include/pragya_avastha.h"

#include <stdio.h>

/*
=========================================================
   PRAGYA ROUTER (L3)

   ROLE:
   - KRST se aayi hui avastha ko process karna
   - Niyu → logic analysis
   - Mira → teaching layer
   - Sakhi → user response

   NOTE:
   - stop_execution flag ka respect MUST hai
=========================================================
*/

void pragya_route(PragyaAvastha *avastha)
{
    /* ========================= */
    /* BASIC SAFETY CHECK        */
    /* ========================= */

    if (!avastha)
        return;

    /* ========================= */
    /* 🔥 STOP EXECUTION GUARD   */
    /* ========================= */

    if (avastha->stop_execution)
        return;

    /* ========================= */
    /* SILENT MODE               */
    /* ========================= */

    if (avastha->teach_level == 0)
        return;

    /* ========================= */
    /* NIYU ANALYSIS             */
    /* ========================= */

    NiyuResult *result = niyu_think_avastha(avastha);

    if (!result)
        return;

    /* ========================= */
    /* MIRA TEACHING             */
    /* ========================= */

    if (avastha->teach_level > 1)
    {
        mira_teach_avastha(avastha, result);
    }

    /* ========================= */
    /* SAKHI RESPONSE            */
    /* ========================= */

    if (avastha->error_code != 0)
    {
        L3Request request;

        request.raw_input = avastha->raw_text;
        request.intent = 0;

        sakhi_speak(&request, result);
    }

    /* ========================= */
    /* CLEANUP                   */
    /* ========================= */

    niyu_free_result(result);
}
#ifndef PRAGYA_ROUTER_H
#define PRAGYA_ROUTER_H

#include "../../include/pragya_avastha.h"

void pragya_route(PragyaAvastha *avastha);

#endif
#ifndef NIYU_CONTRACT_H
#define NIYU_CONTRACT_H

#include "../../include/pragya_avastha.h"

/* =============================== */
/* NIYU ERROR TYPES                */
/* =============================== */

typedef enum {

    NIYU_ERR_VALID = 0,
    NIYU_ERR_SYNTAX,
    NIYU_ERR_MISSING_OPERAND,
    NIYU_ERR_EXTRA_OPERATOR,
    NIYU_ERR_INVALID_IDENTIFIER,
    NIYU_ERR_UNDEFINED_VARIABLE,
    NIYU_ERR_TYPE_MISMATCH,
    NIYU_ERR_DIV_ZERO,
    NIYU_ERR_UNEXPECTED_TOKEN,
    NIYU_ERR_EMPTY_INPUT,
    NIYU_ERR_INVALID_COMMAND,
    NIYU_ERR_RUNTIME

} NiyuErrorType;


/* =============================== */
/* NIYU INTENT TYPES               */
/* =============================== */

typedef enum {

    NIYU_INTENT_CALCULATION = 0,
    NIYU_INTENT_DECLARATION,
    NIYU_INTENT_PRINT,
    NIYU_INTENT_UNKNOWN

} NiyuIntent;


/* =============================== */
/* NIYU RESULT STRUCT              */
/* =============================== */

typedef struct {

    NiyuIntent intent;
    NiyuErrorType error_type;

    int severity;

    int error_pos;

} NiyuResult;


/* =============================== */
/* NIYU CORE API                   */
/* =============================== */

NiyuResult* niyu_think_avastha(PragyaAvastha *avastha);

void niyu_free_result(NiyuResult *res);

#endif
#include "niyu_contract.h"
#include "../../include/pragya_avastha.h"

#include <stdlib.h>
#include <string.h>
#include <ctype.h>

/* ============================= */
/* Helper functions              */
/* ============================= */

static int is_operator(char c)
{
    return (c=='+' || c=='-' || c=='*' || c=='/');
}

static int is_digit_char(char c)
{
    return isdigit((unsigned char)c);
}

static int starts_with(const char *text, const char *word)
{
    return strncmp(text, word, strlen(word)) == 0;
}

/* ============================= */
/* Skip spaces                   */
/* ============================= */

static int next_non_space(const char *t, int i)
{
    while (t[i] && isspace((unsigned char)t[i]))
        i++;

    return i;
}


/* ============================= */
/* NIYU THINK ENGINE             */
/* ============================= */

NiyuResult* niyu_think_avastha(PragyaAvastha *a)
{
    if (!a || !a->raw_text)
        return NULL;

    const char *t = a->raw_text;

    NiyuResult *res = malloc(sizeof(NiyuResult));

    if (!res)
        return NULL;

    res->intent = NIYU_INTENT_UNKNOWN;
    res->error_type = NIYU_ERR_VALID;
    res->severity = 0;
    res->error_pos = -1;

    int len = strlen(t);

    if (len == 0)
    {
        res->error_type = NIYU_ERR_EMPTY_INPUT;
        res->severity = 1;
        return res;
    }

    /* ============================= */
    /* INTENT DETECTION              */
    /* ============================= */

    if (starts_with(t, "mavi"))
    {
        res->intent = NIYU_INTENT_DECLARATION;
        return res;
    }

    if (starts_with(t, "bolo"))
    {
        res->intent = NIYU_INTENT_PRINT;
        return res;
    }

    res->intent = NIYU_INTENT_CALCULATION;

    /* ============================= */
    /* ERROR PATTERN DETECTION       */
    /* ============================= */

    for (int i = 0; i < len; i++)
    {
        int j = next_non_space(t, i+1);

        if (!t[j])
            break;

        char c1 = t[i];
        char c2 = t[j];

        /* operator operator */

        if (is_operator(c1) && is_operator(c2))
        {
            res->error_type = NIYU_ERR_EXTRA_OPERATOR;
            res->severity = 2;
            res->error_pos = j;
            return res;
        }

        /* digit digit (missing operator) */

        if (is_digit_char(c1) && is_digit_char(c2))
        {
            if (j > i+1)
            {
                res->error_type = NIYU_ERR_SYNTAX;
                res->severity = 2;
                res->error_pos = j;
                return res;
            }
        }
    }

    /* ============================= */
    /* Missing operand               */
    /* ============================= */

    int last = len-1;

    while (last >= 0 && isspace((unsigned char)t[last]))
        last--;

    if (last >= 0 && is_operator(t[last]))
    {
        res->error_type = NIYU_ERR_MISSING_OPERAND;
        res->severity = 2;
        res->error_pos = last;
        return res;
    }

    /* ============================= */
    /* Undefined variable heuristic  */
    /* ============================= */

    if (isalpha((unsigned char)t[0]))
    {
        res->error_type = NIYU_ERR_UNDEFINED_VARIABLE;
        res->severity = 3;
        res->error_pos = 0;
        return res;
    }

    return res;
}


/* ============================= */
/* FREE RESULT                   */
/* ============================= */

void niyu_free_result(NiyuResult *res)
{
    if (res)
        free(res);
}
#ifndef NIYU_CORE_H
#define NIYU_CORE_H

#include "niyu_contract.h"

/*
 * NIYU CORE
 *
 * Implementation file will:
 * - define NiyuResult structure
 * - implement internal logic
 *
 * This header exposes NOTHING extra.
 * Future machine-backed engine can replace niyu_core.c
 * without touching any other file.
 */

#endif /* NIYU_CORE_H */
#ifndef USER_CONFIG_H
#define USER_CONFIG_H

// Personality mapping
#define P_SAKHI 1
#define P_NIYU  2
#define P_MIRA  3
#define P_KAVYA 4
#define P_SHIRI 5
#define P_ALL   6   // Unified AI

// Mode mapping
#define MODE_INDIA  1   // Hindi-English mix
#define MODE_GLOBAL 2   // Full English

// Developer mode flag (ON by default)
extern int DEV_MODE;

// Current personality (default = Sakhi)
extern int CURRENT_PERSONALITY;

// Current mode (default = INDIA)
extern int CURRENT_MODE;

void config_init();
void config_set_personality(int p);
void config_set_mode(int mode);

#endif
#ifndef NIYU_ENGINE_H
#define NIYU_ENGINE_H

int niyu_analyze(const char *input);

#endif
#ifndef EVENT_H
#define EVENT_H

#include "runtime.h"
/*──────────────────────────────────────────────────────────────
 |  SHRIJILANG — EVENT SYSTEM
 |
 |  Purpose:
 |   • Represent what just happened in the system
 |   • Bridge Interpreter → State → AI
 |
 |  Introduced in STEP-6
 |  Extended in STEP-7.3 (COMMAND events)
 *────────────────────────────────────────────────────────────*/

/*──────────────────────────────────────────────────────────────
 |  EVENT TYPES
 *────────────────────────────────────────────────────────────*/
typedef enum {

    EVENT_NONE = 0,

    /* Execution lifecycle */
    EVENT_EXECUTION_START,
    EVENT_EXECUTION_BLOCKED,
    EVENT_EXECUTION_SUCCESS,

    /* Variable & logic */
    EVENT_ASSIGNMENT,
    EVENT_SHUNYA_ACCESS,

    /* Command system (STEP-7.3) */
    EVENT_COMMAND,          /* ls / cd / mkdir / aliases */

    /* Errors */
    EVENT_ERROR,
    EVENT_FATAL_ERROR

} EventType;


/*──────────────────────────────────────────────────────────────
 |  EVENT STRUCTURE
 *────────────────────────────────────────────────────────────*/
typedef struct {

    EventType type;

    /* Optional short context (safe capped) */
    char context[64];

} Event;


/*──────────────────────────────────────────────────────────────
 |  EVENT API
 *────────────────────────────────────────────────────────────*/
void event_fire(EventType type, const char *context);
void event_bind_runtime(ShrijiRuntime *rt);

#endif /* EVENT_H */
/* RADHE_RADHE_SHREEJI */

/* =========================================================
   SHRIJILANG — EXECUTION ENGINE (CORE CONTRACT)
   ========================================================= */

#ifndef SHRIJI_ENGINE_H
#define SHRIJI_ENGINE_H

#include "ast.h"
#include "value.h"
#include "env.h"
#include "runtime.h"
#include "error.h"

/* =========================================================
   ENGINE CONFIG
   ========================================================= */

/* Maximum number of fix attempts allowed */
#define ENGINE_MAX_FIX_ATTEMPTS 5

/* =========================================================
   ENGINE STATUS
   ========================================================= */

typedef enum
{
    ENGINE_OK = 0,
    ENGINE_PARSE_ERROR,
    ENGINE_RUNTIME_ERROR,
    ENGINE_FIX_FAILED
} EngineStatus;

/* =========================================================
   ENGINE RESULT STRUCTURE
   ========================================================= */

typedef struct
{
    /* ---------- Language Layer ---------- */

    ASTNode *ast;              /* Final AST (after fix if any) */

    /* ---------- Runtime Layer ---------- */

    Value result;              /* Final evaluated result */

    /* ---------- Fix Engine Meta ---------- */

    int was_fixed;             /* 1 if any fix applied */
    int fix_attempts;          /* number of attempts used */
    int confidence_penalty;    /* total penalty from fixes */

    /* ---------- Status ---------- */

    EngineStatus status;       /* final execution status */

} EngineResult;

/* =========================================================
   ENGINE ENTRY POINT
   ========================================================= */

/*
   Main execution pipeline:

   Input
     → Parse
     → Fix (if needed)
     → Re-parse
     → Runtime
     → Result

   NOTE:
   - This function is PURE (no printing ideally in future)
   - KRST should control how it is used
*/
EngineResult shriji_engine_execute(
    const char *input,
    Env *env
);

/* =========================================================
   LOW-LEVEL CONTROL (INTERNAL USE)
   ========================================================= */

/* Single parse attempt (no fix) */
ASTNode *engine_parse_once(const char *input);

/* Apply fix rules */
int engine_apply_fix(
    char *buffer,
    size_t size,
    const ShrijiErrorInfo *err
);

/* =========================================================
   SAFETY UTILITIES
   ========================================================= */

/* Reset result safely */
void engine_result_init(EngineResult *res);

/* Free result safely */
void engine_result_free(EngineResult *res);

/* =========================================================
   LEGACY FIX ENGINE (KRST COMPAT)
   ========================================================= */

ASTNode *language_execute_with_fix(
    const char *input,
    int *was_fixed,
    int *penalty_out
);
typedef enum {
    FIX_NONE = 0,
    FIX_SAFE,
    FIX_DOUBTFUL,
    FIX_REJECT
} FixType;

FixType fix_apply(const char *input, char *output);

#endif /* SHRIJI_ENGINE_H */
#ifndef SHRIJI_TYPO_ENGINE_H
#define SHRIJI_TYPO_ENGINE_H

/*
    Phase-1 Typo Suggestion Engine
    Only checks for "exit" typo
*/

int shriji_check_typo(
    const char *input,
    char *suggestion,
    int suggestion_size
);

#endif
#ifndef SMRITI_PERSONAL_H
#define SMRITI_PERSONAL_H

/* ===============================
   PERSONAL SMRITI API
   =============================== */

/* load from file (startup) */
void smriti_personal_load(void);

/* save one key */
void smriti_personal_set(const char *key, const char *value);

/* get value */
const char* smriti_personal_get(const char *key);

/* clear memory (RAM only) */
void smriti_personal_clear(void);

/* force save to file */
void smriti_personal_save(void);

#endif
#ifndef AST_H
#define AST_H
#include "../include/token.h"
#ifdef SHRIJI_ENABLE_MASTER_TOKENS
#include "lang/token_master.h"
#endif

/*──────────────────────────────────────────────────────────────
 |  SHRIJILANG — ABSTRACT SYNTAX TREE (AST)
 |
 |  Used by:
 |   • Parser
 |   • Interpreter
 |   • AI Engines (Sakhi / Niyu / Mira / Kavya / Shiri)
 |   • Universal Command Engine (sri_*)
 |   • Future Grammar (block, if, loop, function)
 *────────────────────────────────────────────────────────────*/

/*──────────────────────────────────────────────────────────────
 |  SECTION A — CHAKRA DIAGNOSTIC ENGINE
 *────────────────────────────────────────────────────────────*/
typedef enum {
    CHAKRA_OK = 0,
    CHAKRA_MOOLADHARA,
    CHAKRA_SWADHISTHAN,
    CHAKRA_MANIPUR,
    CHAKRA_ANAHATA,
    CHAKRA_VISHUDDHI,
    CHAKRA_AJNA,
    CHAKRA_SAHASRARA
} ChakraState;


/*──────────────────────────────────────────────────────────────
 |  SECTION B — AST NODE TYPES
 *────────────────────────────────────────────────────────────*/
typedef enum {

    /* Core literals */
    AST_NUMBER,
    AST_IDENTIFIER,
    AST_STRING,
    AST_BOOL,        // true / false
    /* Expressions */
    AST_BINARY,
    AST_UNARY,
    AST_ASSIGNMENT,
    AST_UPDATE,        // x = expression
    AST_LIST,          // [a, b, c]
    AST_INDEX,         // a[0]
    AST_DICT,          // {"a":1, "b":2}
    AST_INDEX_UPDATE,
/* Control / Grammar */
    AST_IF,
    AST_BLOCK,
    AST_PROGRAM,     // script / file root
    AST_IMPORT,      // module import
    AST_LOOP,        // future
    AST_WHILE,       // jabtak loop
    AST_BREAK,       // rukja (break)
    AST_CONTINUE,    // chalu (continue)
    AST_RETURN,      // wapas (return)
    AST_NOT,    //logical NOT (Nahi)
    AST_INC,         // ++
    AST_DEC,         // --

    /* AI Nodes */
    AST_SAKHI,
    AST_NIYU,
    AST_SHIRI,
    AST_MIRA,
    AST_KAVYA,

/* Functions */
    AST_FUNCTION,      // function definition
    AST_CALL,          // function call
    AST_RACHNA,
    /* Universal command (sri_*) */
    AST_COMMAND

} ASTNodeType;


/*─────────────────────────────────────────────────────────────
 | SECTION C — AST NODE STRUCTURE
 *─────────────────────────────────────────────────────────────*/
typedef struct ASTNode {

    ASTNodeType type;

    /*────────────────────────
      Primitive values
    ────────────────────────*/
    double  number_value;
    int  bool_value;     // 1 = true, 0 = false
    char string_value[256];
    char name[128];

    /*────────────────────────
      Binary expression
    ────────────────────────*/
    struct ASTNode *left;
    struct ASTNode *right;
    char op[4];                 // + - * / > < >= <= == !=

    int is_prefix;   /* for ++x / x++ / --x / x-- */

    /*────────────────────────
      Assignment / Command
    ────────────────────────*/
    struct ASTNode *value;
    char command_name[128];

    /*────────────────────────
      IF / ELSE
    ────────────────────────*/
    struct ASTNode *condition;   // agar condition
    struct ASTNode *else_block;  // warna block

    /*────────────────────────
      Block
    ────────────────────────*/
    struct ASTNode **statements;
    int stmt_count;

/*────────────────────────
      Function
    ────────────────────────*/
    char function_name[128];

    struct ASTNode **params;
    int param_count;

    struct ASTNode *body;

/*────────────────────────
  Rachna
────────────────────────*/
char rachna_name[128];
struct ASTNode **rachna_params;
int rachna_param_count;
struct ASTNode *rachna_body;

/*────────────────────────
      List / Index
    ────────────────────────*/
    struct ASTNode **elements;
    int element_count;

    struct ASTNode *index_target;   // a[...]
    struct ASTNode *index_expr;     // ... inside []

    struct ASTNode *index_value;    // a[expr] = value
/*────────────────────────
      Dict / Map
    ────────────────────────*/
    struct ASTNode **dict_keys;
    struct ASTNode **dict_values;
    int dict_count;

    /*────────────────────────
      Call
    ────────────────────────*/
    struct ASTNode **args;
    int arg_count;

    /*────────────────────────
      Return
    ────────────────────────*/
    struct ASTNode *return_expr;

    /*────────────────────────
      Diagnostics
    ────────────────────────*/
    ChakraState chakra_state;

} ASTNode;


/*──────────────────────────────────────────────────────────────
 |  SECTION D — CONSTRUCTORS
 *────────────────────────────────────────────────────────────*/

/* Core */
ASTNode *new_number_node(double value);
ASTNode *new_bool_node(int value);
ASTNode *new_identifier_node(const char *name);
ASTNode *new_string_node(const char *str);

/* Expressions */
ASTNode *new_binary_node(char op, ASTNode *left, ASTNode *right);
ASTNode *new_unary_node(Token op, ASTNode *expr);
ASTNode *new_assignment_node(const char *name, ASTNode *value);
ASTNode *new_update_node(const char *name, ASTNode *value);

ASTNode *new_list_node(ASTNode **elements, int count);
ASTNode *new_index_node(ASTNode *target, ASTNode *index_expr);

ASTNode *new_dict_node(ASTNode **keys, ASTNode **values, int count);

ASTNode *new_index_update_node(ASTNode *target, ASTNode *index_expr, ASTNode *value);

/* AI */
ASTNode *new_sakhi_node(ASTNode *msg);
ASTNode *new_niyu_node(ASTNode *msg);
ASTNode *new_shiri_node(ASTNode *msg);
ASTNode *new_mira_node(ASTNode *msg);
ASTNode *new_kavya_node(ASTNode *msg);

/* Universal command */
ASTNode *new_command_node(const char *cmd_name, ASTNode *arg);
ASTNode *new_import_node(const char *module_name);

/* Block */
ASTNode *new_block_node(ASTNode **stmts, int count);
ASTNode *new_program_node(ASTNode **stmts, int count);

/* Grammar */
ASTNode *new_if_node(ASTNode *condition, ASTNode *then_block);
ASTNode *new_while_node(ASTNode *condition, ASTNode *body);
ASTNode *new_break_node(void);
ASTNode *new_continue_node(void);
ASTNode *new_return_node(ASTNode *expr);

/* Functions */
ASTNode *new_function_node(
    const char *name,
    ASTNode **params,
    int param_count,
    ASTNode *body
);

ASTNode *new_rachna_node(
    const char *name,
    ASTNode **params,
    int param_count,
    ASTNode *body
);
ASTNode *new_call_node(
    const char *name,
    ASTNode **args,
    int arg_count
);

ASTNode *new_return_node(ASTNode *expr);

/* Cleanup */
void ast_free(ASTNode *node);

/* Logical */
ASTNode *new_not_node(ASTNode *expr);

/* Inc / Dec */
ASTNode *new_inc_node(ASTNode *expr, int is_prefix);
ASTNode *new_dec_node(ASTNode *expr, int is_prefix);

#endif /* AST_H */
#ifndef AI_INTENT_H
#define AI_INTENT_H

/*──────────────────────────────────────────────
  AI INTENT — L3 (KR̥ṢṬ)
  Purpose:
  - User ke words ko canonical intent me normalize karna
  - Routing decision ko simplify karna
  - Interpreter / execution ko bilkul affect NAHI karta
──────────────────────────────────────────────*/

typedef enum {
    AI_INTENT_UNKNOWN = 0,
    AI_INTENT_EXPLAIN,   /* samjhao, explain, batao */
    AI_INTENT_WHY,       /* why, kyun, kyu */
    AI_INTENT_CALC,      /* calculate, solve, math */
    AI_INTENT_EMOTION    /* sad, dukhi, lonely */
} AIIntent;

/* raw text -> detected intent */
AIIntent ai_detect_intent(const char *message);

#endif /* AI_INTENT_H */
#ifndef DECISION_ENGINE_H
#define DECISION_ENGINE_H

#include "error.h"

typedef enum {
    DECISION_AUTO_FIX,
    DECISION_EDIT,
    DECISION_REJECT
} DecisionType;

/* classification */
int is_safe_error(ShrijiErrorCode code);
int is_ambiguous_error(ShrijiErrorCode code);

/* main decision */
DecisionType shriji_take_decision(
    int confidence,
    int risk,
    ShrijiErrorCode code
);

#endif
#ifndef SHRIJI_LANG_CONFIG_H
#define SHRIJI_LANG_CONFIG_H

#define SHRIJI_ENABLE_MASTER_TOKENS

#endif
#ifndef SHRIJI_TOKEN_MASTER_H
#define SHRIJI_TOKEN_MASTER_H

/*
  SHRIJILANG — MASTER TOKEN REGISTRY
  ---------------------------------
  This file is a DESIGN + FUTURE registry.

  RULES:
  1. token.h is the runtime truth
  2. Only tokens present in token.h are ACTIVE
  3. Others are FUTURE / PARKED (no runtime meaning yet)
*/

/* =====================================================
   CORE / ACTIVE TOKENS (MIRROR OF token.h)
   ===================================================== */

typedef enum {

    /* ----- Special ----- */
    MT_TOKEN_EOF = 0,
    MT_TOKEN_ERROR,
    MT_TOKEN_NEWLINE,

    /* ----- Literals ----- */
    MT_TOKEN_NUMBER,
    MT_TOKEN_IDENTIFIER,
    MT_TOKEN_STRING,
    MT_TOKEN_TRUE,
    MT_TOKEN_FALSE,

    /* ----- Core Keywords ----- */
    MT_TOKEN_MAVI,
    MT_TOKEN_AGAR,
    MT_TOKEN_WARNA,
    MT_TOKEN_KAAM,
    MT_TOKEN_JABTAK,
    MT_TOKEN_RUKJA,
    MT_TOKEN_CHALU,
    MT_TOKEN_NAHI,
    MT_TOKEN_BOLO,
    MT_TOKEN_WAPAS,
    MT_TOKEN_IMPORT,

    /* ----- AI / System Commands ----- */
    MT_TOKEN_SAKHI,
    MT_TOKEN_NIYU,
    MT_TOKEN_MIRA,
    MT_TOKEN_KAVYA,
    MT_TOKEN_SHIRI,
    MT_TOKEN_SAMAY,

    /* ----- Operators ----- */
    MT_TOKEN_PLUS,
    MT_TOKEN_MINUS,
    MT_TOKEN_STAR,
    MT_TOKEN_SLASH,
    MT_TOKEN_MOD,

    MT_TOKEN_EQUAL,
    MT_TOKEN_EQEQ,
    MT_TOKEN_NEQ,
    MT_TOKEN_GT,
    MT_TOKEN_LT,
    MT_TOKEN_GTE,
    MT_TOKEN_LTE,

    /* ----- Increment / Decrement ----- */
    MT_TOKEN_PLUSPLUS,
    MT_TOKEN_MINUSMINUS,

    /* ----- Logical ----- */
    MT_TOKEN_AND,
    MT_TOKEN_OR,
    MT_TOKEN_NOT,

    /* ----- Punctuation ----- */
    MT_TOKEN_LEFT_PAREN,
    MT_TOKEN_RIGHT_PAREN,
    MT_TOKEN_LEFT_BRACE,
    MT_TOKEN_RIGHT_BRACE,
    MT_TOKEN_LEFT_BRACKET,
    MT_TOKEN_RIGHT_BRACKET,
    MT_TOKEN_COMMA,
    MT_TOKEN_COLON,

/* =====================================================
   FUTURE / PARKED TOKENS (NOT IN token.h YET)
   ===================================================== */

    /* ----- Declarations ----- */
    MT_TOKEN_CONST,
    MT_TOKEN_LET,

    /* ----- Types ----- */
    MT_TOKEN_TYPE_INT,
    MT_TOKEN_TYPE_FLOAT,
    MT_TOKEN_TYPE_STRING,
    MT_TOKEN_TYPE_BOOL,
    MT_TOKEN_TYPE_LIST,
    MT_TOKEN_TYPE_DICT,
    MT_TOKEN_TYPE_ANY,

    /* ----- Control Flow (future) ----- */
    MT_TOKEN_FOR,
    MT_TOKEN_DO,
    MT_TOKEN_SWITCH,
    MT_TOKEN_CASE,
    MT_TOKEN_DEFAULT,

    /* ----- Assignment Variants ----- */
    MT_TOKEN_PLUSEQ,
    MT_TOKEN_MINUSEQ,
    MT_TOKEN_STAREQ,
    MT_TOKEN_SLASHEQ,
    MT_TOKEN_MODEQ,

    /* ----- Strict / Advanced Compare ----- */
    MT_TOKEN_STRICT_EQ,
    MT_TOKEN_STRICT_NEQ,

    /* ----- Bitwise ----- */
    MT_TOKEN_BIT_AND,
    MT_TOKEN_BIT_OR,
    MT_TOKEN_BIT_XOR,
    MT_TOKEN_BIT_NOT,
    MT_TOKEN_SHIFT_LEFT,
    MT_TOKEN_SHIFT_RIGHT,

    /* ----- Ternary ----- */
    MT_TOKEN_QUESTION,

    /* ----- Structure (future sugar) ----- */
    MT_TOKEN_DOT,
    MT_TOKEN_ARROW,

    /* ----- Data Structures ----- */
    MT_TOKEN_LIST,
    MT_TOKEN_DICT,
    MT_TOKEN_SET,
    MT_TOKEN_TUPLE,

    /* ----- Sentinel ----- */
    MT_TOKEN__COUNT

} ShrijiMasterToken;

#endif /* SHRIJI_TOKEN_MASTER_H */
#ifndef SHRIJI_LANG_GRAMMAR_CORE_H
#define SHRIJI_LANG_GRAMMAR_CORE_H

#include "../token.h"

/*
 * Phase-1 Grammar Core
 * Only classification — no parsing logic
 */

typedef enum {
    GRAMMAR_STATEMENT,
    GRAMMAR_EXPRESSION,
    GRAMMAR_OPERATOR,
    GRAMMAR_LITERAL,

    GRAMMAR_INCDEC,     /* ++ -- */

    GRAMMAR_RESERVED,
    GRAMMAR_UNKNOWN
} GrammarKind;

/* Query API */
GrammarKind shriji_grammar_kind(TokenType t);
const char *shriji_grammar_kind_name(GrammarKind k);

#endif /* SHRIJI_LANG_GRAMMAR_CORE_H */
#ifndef SHRIJI_LANG_GRAMMAR_H
#define SHRIJI_LANG_GRAMMAR_H

/* Grammar extension registry
 * Phase-1: only hooks, no grammar rules
 */

/* Called once during parser init */
void shriji_register_grammar_core(void);

/* Optional future grammars */
void shriji_register_grammar_logic(void);
void shriji_register_grammar_incdec(void);
void shriji_register_grammar_loop(void);
void shriji_register_grammar_func(void);

#endif /* SHRIJI_LANG_GRAMMAR_H */
#ifndef VALUE_H
#define VALUE_H

#include <string.h>
#include <stdlib.h>

#include "ast.h"

/*──────────────────────────────────────────────
  VALUE TYPES
──────────────────────────────────────────────*/
typedef enum {
    VAL_NULL = 0,
    VAL_NUMBER,
    VAL_STRING,
    VAL_BOOL,
    VAL_LIST,
    VAL_FUNCTION,
    VAL_DICT
} ValueType;

typedef struct Value {
    ValueType type;

    double number;
    char *string;
    int boolean;
    /* ✅ list value */
    struct Value *items;
    int count;

    /* ✅ function value */
    ASTNode *function;

    /* ✅ dict value */
    struct Value *dict_keys;
    struct Value *dict_values;
    int dict_count;

} Value;

/*──────────────────────────────────────────────
  CONSTRUCTORS
──────────────────────────────────────────────*/
static inline Value value_null(void)
{
    Value v;
    v.type = VAL_NULL;

    v.number = 0;
    v.string = NULL;

    v.items = NULL;
    v.count = 0;

    v.function = NULL;

    v.dict_keys = NULL;
    v.dict_values = NULL;
    v.dict_count = 0;

    return v;
}

static inline Value value_number(double n)
{
    Value v = value_null();
    v.type = VAL_NUMBER;
    v.number = n;
    return v;
}

static inline Value value_string(const char *s)
{
    Value v = value_null();
    v.type = VAL_STRING;

    if (!s) {
        v.string = NULL;
        return v;
    }

    v.string = (char *)malloc(strlen(s) + 1);
    if (!v.string) return value_null();

    strcpy(v.string, s);
    return v;
}

/* BOOL SUPPORT */
static inline Value value_bool(int b)
{
    Value v = value_null();
    v.type = VAL_BOOL;
    v.boolean = b ? 1 : 0;
    return v;
}

static inline int value_is_truthy(Value v)
{
    switch (v.type) {
        case VAL_BOOL:   return v.boolean;
        case VAL_NUMBER: return v.number != 0;
        case VAL_STRING: return v.string && v.string[0] != '\0';
        case VAL_NULL:   return 0;
        default:         return 1;  // list, dict, function → truthy
    }
}


/* ✅ list constructor */
static inline Value value_list(Value *items, int count)
{
    Value v = value_null();
    v.type = VAL_LIST;
    v.items = items;
    v.count = count;
    return v;
}

/* ✅ function constructor */
static inline Value value_function(ASTNode *fn)
{
    Value v = value_null();
    v.type = VAL_FUNCTION;
    v.function = fn;
    return v;
}

/* ✅ dict constructor */
static inline Value value_dict(Value *keys, Value *values, int count)
{
    Value v = value_null();
    v.type = VAL_DICT;

    v.dict_keys = keys;
    v.dict_values = values;
    v.dict_count = count;

    return v;
}

/*──────────────────────────────────────────────
  COPY
──────────────────────────────────────────────*/
static inline Value value_copy(Value src)
{
    if (src.type == VAL_STRING) {
        return value_string(src.string ? src.string : "");
    }

    if (src.type == VAL_NUMBER) {
        return value_number(src.number);
    }

    if (src.type == VAL_BOOL) {
         return value_bool(src.boolean);
    }


    if (src.type == VAL_LIST) {

        if (src.count <= 0 || !src.items) {
            return value_list(NULL, 0);
        }

        Value *items = (Value *)malloc(sizeof(Value) * src.count);
        if (!items) return value_null();

        for (int i = 0; i < src.count; i++) {
            items[i] = value_copy(src.items[i]);
        }

        return value_list(items, src.count);
    }

    if (src.type == VAL_FUNCTION) {
        /* function node pointer copy (AST is owned elsewhere) */
        return value_function(src.function);
    }

    if (src.type == VAL_DICT) {

        if (src.dict_count <= 0 || !src.dict_keys || !src.dict_values) {
            return value_dict(NULL, NULL, 0);
        }

        Value *keys = (Value *)malloc(sizeof(Value) * src.dict_count);
        Value *values = (Value *)malloc(sizeof(Value) * src.dict_count);

        if (!keys || !values) {
            if (keys) free(keys);
            if (values) free(values);
            return value_null();
        }

        for (int i = 0; i < src.dict_count; i++) {
            keys[i] = value_copy(src.dict_keys[i]);
            values[i] = value_copy(src.dict_values[i]);
        }

        return value_dict(keys, values, src.dict_count);
    }

    return value_null();
}

/*──────────────────────────────────────────────
  FREE
──────────────────────────────────────────────*/
static inline void value_free(Value *v)
{
    if (!v) return;

    if (v->type == VAL_STRING && v->string) {
        free(v->string);
        v->string = NULL;
    }

    if (v->type == VAL_LIST && v->items) {
        for (int i = 0; i < v->count; i++) {
            value_free(&v->items[i]);
        }
        free(v->items);
        v->items = NULL;
        v->count = 0;
    }

    if (v->type == VAL_DICT) {
        if (v->dict_keys) {
            for (int i = 0; i < v->dict_count; i++) {
                value_free(&v->dict_keys[i]);
            }
            free(v->dict_keys);
            v->dict_keys = NULL;
        }

        if (v->dict_values) {
            for (int i = 0; i < v->dict_count; i++) {
                value_free(&v->dict_values[i]);
            }
            free(v->dict_values);
            v->dict_values = NULL;
        }

        v->dict_count = 0;
    }

    /* function is not freed here */
    v->function = NULL;

    /* reset */
    v->type = VAL_NULL;
    v->number = 0.0;
    v->string = NULL;

    v->items = NULL;
    v->count = 0;

    v->dict_keys = NULL;
    v->dict_values = NULL;
    v->dict_count = 0;
}


#endif
#ifndef LOG_MODE_H
#define LOG_MODE_H

#include <stdio.h>
#include "user_config.h"

/* ---------------------------------------------------------
 * LOGGING MACROS
 * DEV_MODE=1  => verbose prints
 * DEV_MODE=0  => silent
 * --------------------------------------------------------- */

#define NIYU_LOG(...)  do { if (DEV_MODE) printf(__VA_ARGS__); } while (0)
#define SAKHI_LOG(...) do { if (DEV_MODE) printf(__VA_ARGS__); } while (0)
#define MIRA_LOG(...)  do { if (DEV_MODE) printf(__VA_ARGS__); } while (0)
#define KAVYA_LOG(...) do { if (DEV_MODE) printf(__VA_ARGS__); } while (0)

#endif
#ifndef SHRIJILANG_AI_ROUTER_H
#define SHRIJILANG_AI_ROUTER_H

#include "../pragya/contracts/l3_response.h"

/*──────────────────────────────────────────────
  SHRIJILANG AI ROUTER — PUBLIC INTERFACE (L3)
──────────────────────────────────────────────*/

typedef enum {
    SHRIJI_LANG_AUTO = 0,
    SHRIJI_LANG_HINDI_MIX,
    SHRIJI_LANG_ENGLISH
} ShrijiLangHint;

typedef enum {
    SHRIJI_SRC_FILE = 0,
    SHRIJI_SRC_REPL
} ShrijiSource;

typedef struct {
    const char      *text;
    ShrijiLangHint   lang;
    ShrijiSource     source;
} ShrijiBridgePacket;

/* NEW: L3 dispatcher */
L3Response ai_router_dispatch(const ShrijiBridgePacket *packet);

/* old compatibility */
void ai_router_process(const char *message);

#endif /* SHRIJILANG_AI_ROUTER_H */
#ifndef SHRIJILANG_ATMA_LEKHA_H
#define SHRIJILANG_ATMA_LEKHA_H

/*──────────────────────────────────────────────────────────────
  ATMA-LEKHA — SYSTEM SELF AUDIT (L3-C3)

  Purpose:
  - Record system-side failures & misrouting
  - NO user data
  - Silent by default
  - No behavior change

  This is NOT Smriti.
──────────────────────────────────────────────────────────────*/

typedef enum {
    ATMA_ROUTER_FALLBACK = 1,
    ATMA_INTENT_UNKNOWN,
    ATMA_COMMAND_UNKNOWN,
    ATMA_INTERNAL_MISROUTE
} AtmaEventType;

/* initialize system self-audit */
void atma_lekha_init(void);

/* record an event (silent) */
void atma_lekha_note(AtmaEventType type,
                     const char *component,
                     const char *detail);

#endif
#ifndef ERROR_H
#define ERROR_H

#include "../include/token.h"

typedef enum {

    E_UNKNOWN = 0,

    E_ASSIGN_01,
    E_ASSIGN_02,

    E_PARSE_01,
    E_PARSE_02,

    E_PARSE_EXTRA_OPERATOR,
    E_PARSE_MISSING_OPERATOR,
    E_PARSE_DOUBLE_OPERATOR,
    E_PARSE_OPERATOR_START,
    E_PARSE_OPERATOR_END,
    E_PARSE_OPERATOR_CHAIN,
    E_PARSE_MISSING_OPERAND,

    E_PARSE_BRACKET_MISSING,
    E_PARSE_BRACKET_EXTRA,
    E_PARSE_UNMATCHED_PAREN,
    E_PARSE_UNMATCHED_BRACKET,
    E_PARSE_UNMATCHED_BRACE,

    E_PARSE_INVALID_TOKEN,
    E_PARSE_EMPTY_EXPRESSION,

    E_FUNCTION_NAME_INVALID,
    E_FUNCTION_PARAM_INVALID,

    E_IMPORT_PATH_INVALID,

    E_LIST_SYNTAX_ERROR,
    E_DICT_KEY_INVALID,
    E_DICT_SYNTAX_ERROR,

    E_IF_01,

    E_RUNTIME_01,
    E_RUNTIME_DIV_ZERO,
    E_RUNTIME_TYPE_MISMATCH,
    E_RUNTIME_LOOP_LIMIT,
    E_RUNTIME_INDEX_ERROR

} ShrijiErrorCode;


/*──────────────────────────────────────────────*/

typedef struct {

    ShrijiErrorCode code;

    const char *context;
    const char *message;
    const char *hint;

    const char *function;   // NEW (optional)

    int has_location;
    int line;
    int col;

} ShrijiErrorInfo;


/*──────────────────────────────────────────────*/

/* OLD STYLE (5 ARGUMENT) */
void shriji_error_at(
    Token tok,
    ShrijiErrorCode code,
    const char *context,
    const char *message,
    const char *hint
);

/* NEW STYLE (6 ARGUMENT) */
void shriji_error_at_full(
    Token tok,
    ShrijiErrorCode code,
    const char *context,
    const char *function,
    const char *message,
    const char *hint
);

const ShrijiErrorInfo *shriji_last_error(void);

extern int error_reported;
extern ShrijiErrorInfo LAST_ERROR;

/*──────────────────────────────────────────────*/

typedef enum {
    ERROR_MODE_IMMEDIATE,
    ERROR_MODE_SILENT,
    ERROR_MODE_COLLECT
} ShrijiErrorMode;

void shriji_set_error_mode(ShrijiErrorMode mode);

void shriji_print_all_errors(void);

void shriji_error(
    ShrijiErrorCode code,
    const char *context,
    const char *message,
    const char *hint
);

#endif
// SHRIJI_RADHR RADHR

#ifndef SHRIJI_TOKEN_LIST_H
#define SHRIJI_TOKEN_LIST_H

/*
──────────────────────────────────────────────
 ShrijiLang Token System V2
 Dynamic Token Storage (Compiler Grade)
──────────────────────────────────────────────
*/

#include "token.h"

/*──────────────────────────────────────────────
  TOKEN LIST STRUCTURE
──────────────────────────────────────────────*/

typedef struct
{
    Token *data;      // token array
    int count;        // number of tokens used
    int capacity;     // allocated memory

} TokenList;


/*──────────────────────────────────────────────
  TOKEN LIST API
──────────────────────────────────────────────*/

/* create new list */
TokenList token_list_create(void);

/* add token */
void token_list_push(TokenList *list, Token tok);

/* get token */
Token token_list_get(TokenList *list, int index);

/* number of tokens */
int token_list_size(TokenList *list);

/* free memory */
void token_list_free(TokenList *list);


#endif
#ifndef TRANSLATOR_H
#define TRANSLATOR_H

// Translate English/world command → ShrijiLang keyword
const char* translate_keyword(const char *word);

// Autocorrect system for mistyped user commands
const char* mira_autocorrect(const char *word);

// Extract first word from input (helper for parser)
char* get_first_word(const char *source);

#endif
#ifndef NIRMAN_FLOW_ENGINE_H
#define NIRMAN_FLOW_ENGINE_H

void nirman_flow_process(char *line);

#endif
#ifndef NIRMAN_H
#define NIRMAN_H

/* ===============================
   CONTEXT STRUCT
   =============================== */
typedef struct {
    int active;           /* nirman mode ON/OFF */
    int flow_stage;       /* state machine stage */
    char goal[256];       /* user goal */
} NirmanContext;

/* ===============================
   GLOBAL ACCESS
   =============================== */
extern NirmanContext NIRMAN_CTX;

/* ===============================
   LIFECYCLE
   =============================== */
void nirman_start(void);
void nirman_stop(void);
int nirman_is_active(void);

/* ===============================
   INTENT ENGINE
   =============================== */
int nirman_detect_intent(const char *input);

#endif
#ifndef PRAGYA_AVASTHA_H
#define PRAGYA_AVASTHA_H

#include "ast.h"
#include "error.h"
#include "value.h"

/*
   PragyaAvastha

   Shriji system ki vartamaan sthiti.
   Yeh packet KRST se Pragya tak jata hai
   aur fir Niyu analysis ke liye use karti hai.
*/

typedef struct {

/* User input */
const char *raw_text;

/*  NEW: input mode */
int is_program_mode;

/* NEW: correction system */
int has_correction;
char corrected_text[256];
int stop_execution;

    /* Parsed structure */
    ASTNode *ast;

    /* Error state */
    ShrijiErrorCode error_code;

    int error_pos;
    int error_len;

    /* KRST metrics */
    int confidence;
    int risk;

    int teach_level;
    int tone;

    /* Runtime result */
    Value runtime_result;

} PragyaAvastha;

#endif
#ifndef SHRIJI_EXAMPLE_BUILDER_H
#define SHRIJI_EXAMPLE_BUILDER_H

void shriji_build_example(
    const char *input,
    char *out,
    int size
);

#endif
#ifndef NORMALIZE_ENGINE_H
#define NORMALIZE_ENGINE_H

void normalize_input(char *input);

#endif
#ifndef STATE_H
#define STATE_H

/*──────────────────────────────────────────────────────────────
 |  SHRIJILANG — EXECUTION STATE SYSTEM
 |
 |  Purpose:
 |   • Track current execution condition
 |   • Provide memory of risk, confidence, and errors
 |
 |  Used by:
 |   • Interpreter
 |   • Decision Engine (future)
 *────────────────────────────────────────────────────────────*/

/*──────────────────────────────────────────────────────────────
 |  SAFETY LEVEL — how risky the current situation is
 *────────────────────────────────────────────────────────────*/
typedef enum {
    STATE_SAFE = 0,
    STATE_ALERT,
    STATE_RISK,
    STATE_CRITICAL
} SafetyLevel;


/*──────────────────────────────────────────────────────────────
 |  CONFIDENCE LEVEL — how sure the system is
 *────────────────────────────────────────────────────────────*/
typedef enum {
    CONFIDENCE_HIGH = 0,
    CONFIDENCE_MEDIUM,
    CONFIDENCE_LOW
} ConfidenceLevel;


/*──────────────────────────────────────────────────────────────
 |  HUMAN DEPENDENCY — when human input is required
 *────────────────────────────────────────────────────────────*/
typedef enum {
    HUMAN_AUTO_OK = 0,
    HUMAN_CONFIRM_REQUIRED,
    HUMAN_MANDATORY
} HumanDependency;


/*──────────────────────────────────────────────────────────────
 |  STATE STRUCTURE — complete execution snapshot
 *────────────────────────────────────────────────────────────*/
typedef struct {
    SafetyLevel      safety;
    ConfidenceLevel  confidence;
    int              error_pressure;
    HumanDependency  human;
long long max_steps;
long long steps_used;

} ExecutionState;


/*──────────────────────────────────────────────────────────────
 |  STATE API (implemented in state.c)
 *────────────────────────────────────────────────────────────*/
void state_init(ExecutionState *state);
void state_reset(ExecutionState *state);
void state_on_error(ExecutionState *state);
void state_on_success(ExecutionState *state);

#endif /* STATE_H */
#ifndef SHRIJI_FIX_RULES_H
#define SHRIJI_FIX_RULES_H

#include <stddef.h>
#include "error.h"

/* ─────────────────────────────────────────────
   FIX CATEGORY
───────────────────────────────────────────── */

typedef enum {
    FIXCAT_NONE = 0,
    FIXCAT_MISSING_VALUE,
    FIXCAT_BAD_TOKEN,
    FIXCAT_MISSING_OPERATOR,
    FIXCAT_MISSING_IDENTIFIER,
    FIXCAT_EXTRA_OPERATOR
} FixCategory;

/* ─────────────────────────────────────────────
   FIX SAFETY LEVEL
───────────────────────────────────────────── */

typedef enum {
    FIX_SAFE_DETERMINISTIC = 0,   /* 100% sure → auto apply */
    FIX_SAFE_STRUCTURAL,          /* safe but not guaranteed */
    FIX_RISKY_CONTEXTUAL          /* user decision required */
} FixSafetyLevel;

/* ─────────────────────────────────────────────
   FIX FUNCTION SIGNATURE
───────────────────────────────────────────── */

typedef int (*FixApplyFunction)(
    char *buffer,
    size_t buffer_size,
    const ShrijiErrorInfo *err
);

/* ─────────────────────────────────────────────
   FIX RULE STRUCTURE
───────────────────────────────────────────── */

typedef struct {

    ShrijiErrorCode error_code;

    FixCategory category;
    FixSafetyLevel safety;

    int confidence_penalty;
    int max_attempt;
    int reparse_required;

    FixApplyFunction apply_fix;

} FixRule;

/* ─────────────────────────────────────────────
   RULE LOOKUP
───────────────────────────────────────────── */

FixRule *shriji_get_rule_for_error(ShrijiErrorCode code);

#endif
#ifndef SAKHI_ENGINE_H
#define SAKHI_ENGINE_H

void sakhi_speak(const char *msg);

#endif
#ifndef ENV_H
#define ENV_H

#include "value.h"

/*──────────────────────────────────────────────
 | SHRIJILANG — ENVIRONMENT (VARIABLE STORE)
 |
 | Env = Scope stack
 |  - Global scope always exists
 |  - Blocks { } push/pop scopes
 |
 | Rules:
 |  - env_set()     : create/update in CURRENT scope only
 |  - env_update()  : update in nearest existing scope (top → bottom)
 |  - env_get()     : read from nearest scope (top → bottom)
 |  - env_exists()  : check existence (top → bottom)
 *──────────────────────────────────────────────*/

#define ENV_MAX_VARS     1024
#define ENV_MAX_SCOPES   64

/* one variable */
typedef struct {
    char name[64];
    Value value;
} EnvVar;

/* one scope frame */
typedef struct {
    EnvVar vars[ENV_MAX_VARS];
    int count;
} EnvFrame;

/* environment = stack of frames */
typedef struct {
    EnvFrame frames[ENV_MAX_SCOPES];
    int depth; /* >= 1 always (global scope) */
} Env;

/* create */
Env *new_env(void);

/* scope control */
void env_push_scope(Env *env);
void env_pop_scope(Env *env);

/* variables */
void env_set(Env *env, const char *name, Value v);       /* current scope only */
void env_update(Env *env, const char *name, Value v);    /* nearest scope update */
Value env_get(Env *env, const char *name);               /* nearest scope read */
int env_exists(Env *env, const char *name);

#endif /* ENV_H */
#ifndef SHRIJI_ERROR_INTELLIGENCE_H
#define SHRIJI_ERROR_INTELLIGENCE_H

#include "error.h"
#include "../pragya/niyu/niyu_contract.h"
#include "../include/pragya_avastha.h"

/* error clarity levels */

#define ERR_CONFUSED 0
#define ERR_PARTIAL  1
#define ERR_CLEAR    2


/* understanding engine */

int shriji_understand_error(
    ShrijiErrorCode code,
    const NiyuResult *niyu
);


/* pointer extractor */

void shriji_extract_pointer(
    const ShrijiErrorInfo *err,
    int *pos,
    int *len
);


/* main intelligence entry */

void shriji_error_intelligence(
    PragyaAvastha *avastha,
    const ShrijiErrorInfo *err,
    const NiyuResult *niyu
);

int shriji_get_example(const char *input, int pos, char *out);
#endif
#ifndef TOKEN_H
#define TOKEN_H

/*──────────────────────────────────────────────
  SHRIJILANG — TOKEN SYSTEM (LOCKED)
  Now supports: line/col for professional errors
──────────────────────────────────────────────*/

typedef enum {

    TOKEN_EOF = 0,

    /* Literals */
    TOKEN_NUMBER,
    TOKEN_IDENTIFIER,
    TOKEN_STRING,
    TOKEN_TRUE,
    TOKEN_FALSE,

    /* Keywords */
    TOKEN_MAVI,
    TOKEN_AGAR,
    TOKEN_WARNA,
    TOKEN_KAAM,
    TOKEN_JABTAK,
    TOKEN_RUKJA,
    TOKEN_CHALU,
    TOKEN_NAHI,
    TOKEN_BOLO,

    TOKEN_SAKHI,
    TOKEN_NIYU,
    TOKEN_MIRA,
    TOKEN_KAVYA,
    TOKEN_SHIRI,

    TOKEN_SAMAY,
    TOKEN_WAPAS,

    TOKEN_IMPORT,
    TOKEN_RACHNA,

    /* Operators */
    TOKEN_PLUS,
    TOKEN_MINUS,
    TOKEN_STAR,
    TOKEN_SLASH,
    TOKEN_MOD,

    TOKEN_EQUAL,
    TOKEN_EQEQ,
    TOKEN_NEQ,
    TOKEN_GT,
    TOKEN_LT,
    TOKEN_GTE,
    TOKEN_LTE,

/* Increment / Decrement */
TOKEN_PLUSPLUS,     // ++
TOKEN_MINUSMINUS,   // --

/* Logical Operators */
TOKEN_AND,          // &&
TOKEN_OR,           // ||
TOKEN_NOT,          // !

    /* Punctuation */
    TOKEN_LEFT_PAREN,
    TOKEN_RIGHT_PAREN,
    TOKEN_LEFT_BRACKET,
    TOKEN_RIGHT_BRACKET,
    TOKEN_LEFT_BRACE,
    TOKEN_RIGHT_BRACE,
    TOKEN_COMMA,
    TOKEN_COLON,
    TOKEN_NEWLINE,

    /* Special */
    TOKEN_ERROR

} TokenType;

typedef struct {
    TokenType type;
    const char *start;
    int length;
    int line;
    int col;
} Token;

/* Scanner API */
void init_tokenizer(const char *source);
Token scan_token(void);

//* #ifdef SHRIJI_ENABLE_MASTER_TOKENS
//* #include "lang/token_master.h"
//* #endif

#endif // TOKEN_H
#ifndef SHRIJI_FIX_INTERACTIVE_H
#define SHRIJI_FIX_INTERACTIVE_H

#include <stddef.h>
#include "error.h"


/* =========================================
   INTERACTIVE OPERATOR FIX
   Example: 5 5 → 5 + 5
   ========================================= */

int shriji_interactive_operator_fix(
    const char *input,
    char *output,
    int size
);


/* =========================================
   INTERACTIVE BRACKET FIX
   Example: (5+6 → (5+6)
            (5+6 → 5+6
   ========================================= */

int shriji_interactive_bracket_fix(
    const char *input,
    char *output,
    int size
);

int shriji_interactive_double_operator_fix(
    const char *input,
    char *output,
    int size
);

int fix_missing_value_after_equal_interactive(
    char *buffer,
    size_t buffer_size,
    const ShrijiErrorInfo *err
);
#endif
#ifndef MIRA_ENGINE_H
#define MIRA_ENGINE_H

void mira_teach(const char *msg);

#endif
#ifndef INTERPRETER_H
#define INTERPRETER_H

#include "env.h"
#include "ast.h"
#include "value.h"
#include "runtime.h"   // 🔥 NEW — runtime struct access

//──────────────────────────────────────────────────────────────
// SHRIJI TRUTH DIARY (SINGLE SOURCE OF TRUTH)
//──────────────────────────────────────────────────────────────
typedef struct {
    char last_var[64];
    long long last_value;
    int has_assignment;
    int had_runtime_error;
} ShrijiTruthDiary;

// defined in src/interpreter.c
extern ShrijiTruthDiary SHRIJI_DIARY;


//──────────────────────────────────────────────────────────────
// MAIN EVALUATOR
//──────────────────────────────────────────────────────────────
//
// 🔥 IMPORTANT CHANGE:
// eval now receives runtime pointer
//
Value eval(ASTNode *node, Env *env, ShrijiRuntime *rt);


//──────────────────────────────────────────────────────────────
// PROGRAM ENTRY
//──────────────────────────────────────────────────────────────
//
// 🔥 IMPORTANT CHANGE:
// run_program now receives runtime pointer
//
Value run_program(ASTNode *program, Env *env, ShrijiRuntime *rt);


#endif /* INTERPRETER_H */
#ifndef KRST_ROUTER_H
#define KRST_ROUTER_H

#include "../krst/krst_types.h"

void krst_route_request(KRSTRequest *req);

#endif
#ifndef PARSER_H
#define PARSER_H

//──────────────────────────────────────────────────────────────
//  SHRIJILANG PARSER — HEADER (UNIVERSE-CLASS VERSION)
//
//  This module converts raw source code strings into AST Nodes.
//  
//  Core features supported in V1:
//    ✔ mavi assignment       → mavi x = 10 + 20
//    ✔ arithmetic expressions
//    ✔ sakhi "msg"
//    ✔ niyu "msg"
//    ✔ shiri "msg"
//    ✔ mira "msg"
//    ✔ kavya "msg"
//    ✔ sri_<command> "msg"
//    ✔ GAURI Super-List Engine (V1.5)
//
//  Upcoming in V2:
//    ✔ multi-statement blocks { }
//    ✔ loops
//    ✔ conditions
//──────────────────────────────────────────────────────────────

#include "ast.h"

//──────────────────────────────────────────────────────────────
//  PARSE ENTRY POINT
//──────────────────────────────────────────────────────────────
//
//  parse_program()
//      Converts a full input string into a root AST node.
//      Returns NULL if syntax error.
//
ASTNode *parse_program(const char *source);


//──────────────────────────────────────────────────────────────
//  GAURI DATA-STRUCTURE PARSER DECLARATIONS (NEW)
//──────────────────────────────────────────────────────────────
//
//  Gauri commands allow creation + mutation of "living lists"
//  Example:
//      gauri_list x = "[1,2,3]"
//      gauri_push x 50
//      gauri_pop x
//      gauri_get x 2
//      gauri_set x 2 99
//      gauri_sum x
//      gauri_sort x
//      gauri_reverse x
//
ASTNode* parse_gauri_list();
ASTNode* parse_gauri_push();
ASTNode* parse_gauri_pop();
ASTNode* parse_gauri_get();
ASTNode* parse_gauri_set();
ASTNode* parse_gauri_sum();
ASTNode* parse_gauri_sort();
ASTNode* parse_gauri_reverse();

#endif
#ifndef RUNTIME_H
#define RUNTIME_H

#include "state.h"
#include "value.h"

/*──────────────────────────────────────────────
 | SHRIJI RUNTIME — ISOLATED EXECUTION BRAIN
 *──────────────────────────────────────────────*/

typedef struct {

    ExecutionState state;

    int break_flag;
    int continue_flag;
    int loop_depth;

    int return_flag;
    Value return_value;

    int call_depth;

    int error_flag;   // NEW

} ShrijiRuntime;

/* GLOBAL POINTER */
extern ShrijiRuntime *current_runtime;

/* Init once */
void runtime_init(ShrijiRuntime *rt);

/* Reset before each run */
void runtime_reset(ShrijiRuntime *rt);

#endif
#ifndef SMRITI_H
#define SMRITI_H

/* lifecycle */
void smriti_init(void);

/* math continuation (existing behavior) */
int  smriti_has_continuation(void);
void smriti_set_continuation(double value, char op);
void smriti_clear_continuation(void);

double smriti_get_continuation_value(void);
char   smriti_get_continuation_op(void);

/* ===== STEP-11: context (NO logic, NO decision) ===== */

/* turn tracking */
void smriti_next_turn(void);
unsigned long smriti_get_turn(void);

/* last input snapshot */
void smriti_set_last_input(const char *text);
const char* smriti_get_last_input(void);

/* last intent snapshot (L3 intent enum value passed as int) */
void smriti_set_last_intent(int intent);
int  smriti_get_last_intent(void);

#endif
#ifndef SMRITI_SESSION_H
#define SMRITI_SESSION_H

/* ===============================
   SESSION MEMORY (SHORT TERM)
   =============================== */

#include "value.h"

/* last input */
void smriti_session_set_last_input(const char *text);
const char* smriti_session_get_last_input(void);

/* last output */
void smriti_session_set_last_output(const char *text);
const char* smriti_session_get_last_output(void);

/* last evaluated value (UPGRADED) */
void smriti_session_set_last(Value v);
Value smriti_session_get_last(void);

/* clear session */
void smriti_session_clear(void);

#endif
#ifndef COMMAND_H
#define COMMAND_H

/*──────────────────────────────────────────────────────────────
 | INTERNAL COMMAND IDENTIFIERS (TRUTH)
 *──────────────────────────────────────────────────────────────*/
typedef enum {

    CMD_UNKNOWN = 0,

    /* Filesystem */
    CMD_LIST,        /* ls */
    CMD_CD,          /* cd */
    CMD_MKDIR,       /* mkdir */

    /* Output / I/O */
    CMD_BOLO,        /* bolo (print) */
    CMD_SAMAY,

/* AI Modules */
    CMD_SAKHI,
    CMD_NIYU,
    CMD_MIRA,
    CMD_KAVYA,
    CMD_SHIRI,

    /* Future expansion */
    CMD_DELETE,
    CMD_COPY,
    CMD_MOVE



} CommandType;


/*──────────────────────────────────────────────────────────────
 | COMMAND ACCESS POLICY
 *──────────────────────────────────────────────────────────────*/
typedef enum {
    CMD_PUBLIC = 1,     /* visible to everyone */
    CMD_LOCKED = 0      /* reserved / internal */
} CommandAccess;


/*──────────────────────────────────────────────────────────────
 | COMMAND ALIAS (1008 NAMES LIVE HERE)
 *──────────────────────────────────────────────────────────────*/
typedef struct {
    const char     *alias;     /* user-facing name */
    CommandType     cmd;       /* internal truth */
    CommandAccess   access;    /* public / locked */
} CommandAlias;


/*──────────────────────────────────────────────────────────────
 | API
 *──────────────────────────────────────────────────────────────*/
CommandType resolve_command(const char *name);
const char *command_name(CommandType cmd);
int execute_command(CommandType cmd);

#endif /* COMMAND_H */
ShrijiLang v1.0 (Core Stable)

──────────────────────────────────────────────
Introduction
──────────────────────────────────────────────
ShrijiLang ek custom programming language hai
jo Hinglish based syntax par kaam karti hai.

ShrijiLang is a custom programming language
built with Hinglish-style syntax for intuitive learning.

Ye version "Core Stable" hai.
This version is the "Core Stable" release.

──────────────────────────────────────────────
Features
──────────────────────────────────────────────
✔ rachna (functions)
✔ multi-argument support
✔ nested function calls
✔ recursion support
✔ KRST intelligence system
✔ error insight system

✔ Function support (rachna)
✔ Multiple arguments
✔ Nested execution
✔ Recursion support
✔ KRST intelligent execution layer
✔ Error insight system

──────────────────────────────────────────────
How to Run (Linux)
──────────────────────────────────────────────
1. Terminal open karo / Open terminal
2. Navigate to bin folder:

   cd release/linux/bin

3. Run:

   ./shrijilang

──────────────────────────────────────────────
Example
──────────────────────────────────────────────
rachna add(a, b) {
    wapas a + b
}

bolo add(5, 3)

Output:
8

──────────────────────────────────────────────
Limitations (Current Version)
──────────────────────────────────────────────
- Named arguments supported nahi hain
  (example: add(a=5, b=10))

- Default parameters supported nahi hain
  (example: rachna add(a, b=5))

- Named arguments are not supported yet
- Default parameters are not supported yet

──────────────────────────────────────────────
Future Updates
──────────────────────────────────────────────
v1.1 → Named arguments
v1.2 → Default parameters
v2.0 → Advanced execution engine

──────────────────────────────────────────────
Author
──────────────────────────────────────────────
MR._MR Manikpuri

Mission:
Thee Angel 😇

Vision:
To build a new generation intelligent programming ecosystem.
nd{verbatim}
\section*{File: ./env/env.c}
egin{verbatim}
#include <string.h>
#include <stdlib.h>

#include "../include/env.h"
#include "../include/value.h"

/*──────────────────────────────────────────────────────────────
 | SHRIJILANG — ENVIRONMENT STORE (SCOPE STACK)
 |
 | Responsibility:
 |   • Variable storage (Value)
 |   • Scope push / pop
 |   • Get / Set / Update / Exists
 |
 | RULES:
 |   • env NEVER prints
 |   • env NEVER calls shiri
 |   • Only DATA LAYER
 *──────────────────────────────────────────────────────────────*/

/*──────────────────────────────────────────────
 | CREATE NEW ENV
 *──────────────────────────────────────────────*/
Env *new_env(void)
{
    Env *env = (Env *)malloc(sizeof(Env));
    if (!env) return NULL;

    env->depth = 1; /* global scope exists */
    env->frames[0].count = 0;

    return env;
}

/*──────────────────────────────────────────────
 | PUSH NEW SCOPE
 *──────────────────────────────────────────────*/
void env_push_scope(Env *env)
{
    if (!env) return;

    if (env->depth >= ENV_MAX_SCOPES) {
        return; /* silently ignore overflow */
    }

    env->frames[env->depth].count = 0;
    env->depth++;
}

/*──────────────────────────────────────────────
 | POP SCOPE
 *──────────────────────────────────────────────*/
void env_pop_scope(Env *env)
{
    if (!env) return;

    /* never pop global scope */
    if (env->depth <= 1) return;

    int top = env->depth - 1;

    /* free all values in top scope */
    for (int i = 0; i < env->frames[top].count; i++) {
        value_free(&env->frames[top].vars[i].value);
    }

    env->frames[top].count = 0;
    env->depth--;
}

/*──────────────────────────────────────────────
 | CHECK IF VARIABLE EXISTS (search top → bottom)
 *──────────────────────────────────────────────*/
int env_exists(Env *env, const char *name)
{
    if (!env || !name) return 0;

    for (int f = env->depth - 1; f >= 0; f--) {
        for (int i = 0; i < env->frames[f].count; i++) {
            if (strcmp(env->frames[f].vars[i].name, name) == 0) {
                return 1;
            }
        }
    }

    return 0;
}

/*──────────────────────────────────────────────
 | SET VARIABLE
 | - if exists in CURRENT scope: update it
 | - else create new variable in CURRENT scope
 *──────────────────────────────────────────────*/
void env_set(Env *env, const char *name, Value v)
{
    if (!env || !name) return;

    int top = env->depth - 1;
    EnvFrame *frame = &env->frames[top];

    /* update if already exists in current scope */
    for (int i = 0; i < frame->count; i++) {
        if (strcmp(frame->vars[i].name, name) == 0) {
            value_free(&frame->vars[i].value);
            frame->vars[i].value = value_copy(v);
            return;
        }
    }

    /* add new variable in current scope */
    if (frame->count >= ENV_MAX_VARS) {
        return; /* silently ignore overflow */
    }

    strncpy(
        frame->vars[frame->count].name,
        name,
        sizeof(frame->vars[frame->count].name) - 1
    );

    frame->vars[frame->count].name[
        sizeof(frame->vars[frame->count].name) - 1
    ] = '\0';

    frame->vars[frame->count].value = value_copy(v);
    frame->count++;
}

/*──────────────────────────────────────────────
 | UPDATE VARIABLE (search top → bottom)
 | - if exists in ANY scope: update nearest found scope
 | - else create in current scope
 *──────────────────────────────────────────────*/
void env_update(Env *env, const char *name, Value v)
{
    if (!env || !name) return;

    /* search from top scope to global */
    for (int f = env->depth - 1; f >= 0; f--) {
        EnvFrame *frame = &env->frames[f];

        for (int i = 0; i < frame->count; i++) {
            if (strcmp(frame->vars[i].name, name) == 0) {
                value_free(&frame->vars[i].value);
                frame->vars[i].value = value_copy(v);
                return;
            }
        }
    }

    /* if not found, fallback: create in current scope */
    env_set(env, name, v);
}

/*──────────────────────────────────────────────
 | GET VARIABLE VALUE (search top → bottom)
 *──────────────────────────────────────────────*/
Value env_get(Env *env, const char *name)
{
    if (!env || !name) return value_null();

    for (int f = env->depth - 1; f >= 0; f--) {
        for (int i = 0; i < env->frames[f].count; i++) {
            if (strcmp(env->frames[f].vars[i].name, name) == 0) {
                return value_copy(env->frames[f].vars[i].value);
            }
        }
    }

    return value_null();
}
nd{verbatim}
\section*{File: ./src/decision_engine.c}
egin{verbatim}
#include "../include/decision_engine.h"

/* SAFE FIXABLE ERRORS */
int is_safe_error(ShrijiErrorCode code)
{
    switch(code)
    {
        case E_PARSE_MISSING_OPERATOR:
        case E_PARSE_DOUBLE_OPERATOR:
            return 1;
        default:
            return 0;
    }
}

/* AMBIGUOUS ERRORS */
int is_ambiguous_error(ShrijiErrorCode code)
{
    switch(code)
    {
        case E_PARSE_INVALID_TOKEN:
        case E_PARSE_02:
            return 1;
        default:
            return 0;
    }
}

/* MAIN DECISION ENGINE */
DecisionType shriji_take_decision(
    int confidence,
    int risk,
    ShrijiErrorCode code
)
{
    (void)confidence;
    (void)risk;

    /* SAFE → always auto fix */
    if (is_safe_error(code))
        return DECISION_AUTO_FIX;

    /* AMBIGUOUS → ask user */
    if (is_ambiguous_error(code))
        return DECISION_EDIT;

    /* REST → reject */
    return DECISION_REJECT;
}
nd{verbatim}
\section*{File: ./src/runtime.c}
egin{verbatim}
#include "../include/runtime.h"
#include "../include/state.h"

/* GLOBAL POINTER */
ShrijiRuntime *current_runtime = NULL;

void runtime_init(ShrijiRuntime *rt)
{
    if (!rt) return;

    state_init(&rt->state);

    rt->break_flag = 0;
    rt->continue_flag = 0;
    rt->loop_depth = 0;

    rt->return_flag = 0;
    rt->return_value = value_null();

    rt->call_depth = 0;
    rt->error_flag = 0;
}

void runtime_reset(ShrijiRuntime *rt)
{
    if (!rt) return;

    rt->break_flag = 0;
    rt->continue_flag = 0;
    rt->loop_depth = 0;

    rt->return_flag = 0;
    rt->return_value = value_null();

    rt->state.steps_used = 0;

    rt->call_depth = 0;
    rt->error_flag = 0;
}
nd{verbatim}
\section*{File: ./src/interpreter.c}
egin{verbatim}
#include <stdio.h>
#include <string.h>

#include <stdlib.h>

#include "../include/interpreter.h"
#include "../include/parser.h"
#include "../include/ast.h"
#include "../include/env.h"
#include "../include/error.h"
#include "../include/state.h"
#include "../include/event.h"
#include "../include/command.h"
#include "../include/value.h"
#include "../include/ai_router.h"

#include "../include/smriti_personal.h"
#include "../include/smriti_session.h"

#include "../include/runtime.h"

#include "../include/user_config.h"
#include "../include/env.h"

#ifdef SHRIJI_ENABLE_MASTER_TOKENS
#include "lang/token_master.h"
#endif

/*──────────────────────────────────────────────
  helper: unescape string like "a\\nb" -> actual newline
  supports: \n \t \r \\ \"
──────────────────────────────────────────────*/
static char *unescape_string(const char *s)
{
    if (!s) return NULL;

    size_t n = strlen(s);
    char *out = (char *)malloc(n + 1);
    if (!out) return NULL;

    size_t j = 0;
    for (size_t i = 0; i < n; i++) {
        if (s[i] == '\\' && i + 1 < n) {
            char c = s[i + 1];
            if (c == 'n')      { out[j++] = '\n'; i++; continue; }
            else if (c == 't') { out[j++] = '\t'; i++; continue; }
            else if (c == 'r') { out[j++] = '\r'; i++; continue; }
            else if (c == '\\'){ out[j++] = '\\'; i++; continue; }
            else if (c == '"') { out[j++] = '"';  i++; continue; }
        }
        out[j++] = s[i];
    }

    out[j] = '\0';
    return out;
}
/*──────────────────────────────────────────────────────────────
 | SHRIJI TRUTH DIARY — SINGLE SOURCE OF FACT
 *──────────────────────────────────────────────────────────────*/
ShrijiTruthDiary SHRIJI_DIARY = {0};

/*──────────────────────────────────────────────────────────────
 | STATE DECISION HOOK
 *──────────────────────────────────────────────────────────────*/
static int state_allows_execution(ShrijiRuntime *rt)
{
    if (!rt) return 0;

    if (rt->state.safety == STATE_CRITICAL) {
        return 0;
    }
    return 1;
}

/*──────────────────────────────────────────────────────────────
 | SAFE DIVISION
 *──────────────────────────────────────────────────────────────*/
static double safe_div(double a, double b)

{
    if (b == 0) {
        event_fire(EVENT_ERROR, "division by zero");
        shriji_error(
            E_RUNTIME_DIV_ZERO,
            "division",
            "division by zero",
            "denominator must not be zero"
        );
        return 0;
    }
    return a / b;
}

/*──────────────────────────────────────────────────────────────
 | SAFE MODULO
 *──────────────────────────────────────────────────────────────*/
static double safe_mod(double a, double b)
{
    if (b == 0) {
        event_fire(EVENT_ERROR, "modulo by zero");
        shriji_error(
            E_RUNTIME_DIV_ZERO,
            "modulo",
            "modulo by zero",
            "denominator must not be zero"
        );
        return 0;
    }

    long long ai = (long long)a;
    long long bi = (long long)b;

    if (bi == 0)
        return 0;

    return (double)(ai % bi);
}

/*──────────────────────────────────────────────────────────────
 | Value → number helper
 *──────────────────────────────────────────────────────────────*/
static const char *strip_quotes(const char *s)
{
    if (!s) return "";

    while (*s == ' ' || *s == '\t')
        s++;

    size_t len = strlen(s);

    if (len >= 2 && s[0] == '"' && s[len - 1] == '"') {
        static char buf[1024];
        size_t n = len - 2;
        if (n >= sizeof(buf)) n = sizeof(buf) - 1;

        memcpy(buf, s + 1, n);
        buf[n] = '\0';
        return buf;
    }

    return s;
}
/*──────────────────────────────────────────────────────────────
 | MAIN EVALUATOR
 *──────────────────────────────────────────────────────────────*/
Value shriji_execute_line(const char *line, Env *env, ShrijiRuntime *rt)
{
    if (!line || !*line) return value_null();

    while (*line == ' ' || *line == '\t') line++;
    if (*line == '\0') return value_null();

    if (DEV_MODE)
    printf("[KRST] Input: %s\n", line);

    ASTNode *node = parse_program(line);

/*  HARD STOP FIX */
if (!node || error_reported)
{
    if (rt) {
        rt->error_flag = 1;
    }

    error_reported = 0;
    return value_null();
}
    return eval(node, env, rt);
}

Value eval(ASTNode *node, Env *env, ShrijiRuntime *rt)
{
    static int state_initialized = 0;

    if (!state_initialized) {
        state_init(&rt->state);
        state_initialized = 1;
    }

    if (!state_allows_execution(rt) || !node)
        return value_null();

    /* 🧠 GLOBAL STEP COUNTER */
    if (node->type != AST_WHILE &&
        node->type != AST_BREAK &&
        node->type != AST_CONTINUE &&
        node->type != AST_BLOCK &&
        node->type != AST_PROGRAM) {
        rt->state.steps_used++;
    }

    if (rt->state.steps_used > rt->state.max_steps) {
        shriji_error(
            E_PARSE_02,
            "runtime",
            "execution limit exceeded",
            "possible infinite loop detected"
        );
        rt->state.safety = STATE_CRITICAL;
        return value_null();
    }



switch (node->type) {

case AST_BREAK:
    if (rt->loop_depth <= 0) {
        shriji_error(
            E_RUNTIME_01,
            "break",
            "loop ke bahar rukja allowed nahi hai",
            "rukja sirf while loop ke andar use karein"
        );
        return value_null();
    }

    rt->break_flag = 1;
    return value_null();

case AST_CONTINUE:

     if (rt->loop_depth <= 0) {
        shriji_error(
        E_RUNTIME_01,
        "continue",
        "loop ke bahar continue allowed nahi hai",
        "continue sirf while loop ke andar use karein"
    );
    return value_null();
}
    rt->continue_flag = 1;
    return value_null();


case AST_RETURN: {
        Value v = value_null();
        if (node->return_expr)
         v = eval(node->return_expr, env, rt);

       rt->return_flag = 1;
       rt->return_value = v;
        return v;
    }

case AST_UNARY:
{
    ASTNode *cur = node;
    int sign = 1;

    while (cur && cur->type == AST_UNARY) {
        if (cur->op[0] == '-')
            sign *= -1;

        cur = cur->left;
    }

    Value v = eval(cur, env, rt);

    if (v.type != VAL_NUMBER)
        return value_null();

    v.number = v.number * sign;

    return v;
}

/*────────────────────────────────────────
 | INC / DEC (++ / --)
 *────────────────────────────────────────*/

case AST_INC: {

    /* only identifiers allowed */
    if (!node->left || node->left->type != AST_IDENTIFIER) {
        return value_null();
    }

    const char *name = node->left->name;

    Value old = env_get(env, name);

    if (old.type != VAL_NUMBER) {
        return value_null();
    }

    Value newv = old;
    newv.number = old.number + 1;

    env_set(env, name, newv);

    if (node->is_prefix)
        return newv;
    else
        return old;
}

case AST_DEC: {

    /* only identifiers allowed */
    if (!node->left || node->left->type != AST_IDENTIFIER) {
        return value_null();
    }

    const char *name = node->left->name;

    Value old = env_get(env, name);

    if (old.type != VAL_NUMBER) {
        return value_null();
    }

    Value newv = old;
    newv.number = old.number - 1;

    env_set(env, name, newv);

    if (node->is_prefix)
        return newv;
    else
        return old;
}

    /*──────────────────────────────────────────────
      FUNCTIONS
    ──────────────────────────────────────────────*/

    case AST_FUNCTION: {
        /* store function AST node into env */
        env_set(env, node->function_name, value_function(node));
        return value_null();
    }

case AST_RACHNA: {
    env_set(env, node->rachna_name, value_function(node));
    return value_null();
}

case AST_CALL: {

/*────────────────────────────
  STACK GUARD (ENTRY)
────────────────────────────*/
rt->call_depth++;

if (rt->call_depth > 3000) {
    shriji_error(
        E_RUNTIME_01,
        "recursion",
        "maximum recursion depth exceeded",
        "safe limit: 3000"
    );

    rt->call_depth--;
    return value_null();
}

   /*──────────────────────────────────────────────
      BUILT-IN FUNCTIONS (FILE I/O)
      padho("file")                -> returns string
      likho("file", "text")        -> overwrite write -> returns 1/0
      jodo("file", "text")         -> append write   -> returns 1/0

      ✅ polish:
      - supports escape sequences: \n \t \r \\ \"
    ──────────────────────────────────────────────*/

    /* READ */
    if (strcmp(node->function_name, "padho") == 0) {

        if (node->arg_count != 1) {
            shriji_error(
                E_PARSE_02,
                "padho",
                "argument count mismatch",
                "use: padho(\"file.txt\")"
            );
            return value_null();
        }

        Value pathv = eval(node->args[0], env, rt);

        if (pathv.type != VAL_STRING || !pathv.string) {
            value_free(&pathv);
            shriji_error(
                E_PARSE_02,
                "padho",
                "file path must be string",
                "use: padho(\"file.txt\")"
            );
            return value_null();
        }

        const char *path = strip_quotes(pathv.string);

FILE *fp = fopen(path, "rb");
if (!fp) {
    value_free(&pathv);

    shriji_error(
        E_PARSE_02,
        "padho",
        "file not found",
        "check path or create file first"
    );

    return value_string("");
}
        fseek(fp, 0, SEEK_END);
        long sz = ftell(fp);
        fseek(fp, 0, SEEK_SET);

        if (sz < 0) sz = 0;
        if (sz > 1024 * 1024) sz = 1024 * 1024; /* 1MB limit v1 */

        char *buf = (char *)malloc((size_t)sz + 1);
        if (!buf) {
            fclose(fp);
            value_free(&pathv);
            return value_null();
        }

        size_t nread = fread(buf, 1, (size_t)sz, fp);
        buf[nread] = '\0';

        fclose(fp);
        value_free(&pathv);

        Value out = value_string(buf);
        free(buf);
        return out;
    }

    /* WRITE (overwrite) */
    if (strcmp(node->function_name, "likho") == 0) {

        if (node->arg_count != 2) {
            shriji_error(
                E_PARSE_02,
                "likho",
                "argument count mismatch",
                "use: likho(\"file.txt\", \"data\")"
            );
            return value_null();
        }

        Value pathv = eval(node->args[0], env, rt);
        Value datav = eval(node->args[1], env, rt);

        if (pathv.type != VAL_STRING || !pathv.string) {
            value_free(&pathv);
            value_free(&datav);
            shriji_error(
                E_PARSE_02,
                "likho",
                "file path must be string",
                "use: likho(\"file.txt\", \"data\")"
            );
            return value_null();
        }

        const char *path = strip_quotes(pathv.string);

        const char *data = "";
        char tmp[128];
        tmp[0] = '\0';

        char *tmp_data = NULL;

        if (datav.type == VAL_STRING && datav.string) {
            tmp_data = unescape_string(strip_quotes(datav.string));
            data = tmp_data ? tmp_data : "";
        } else if (datav.type == VAL_NUMBER) {
            snprintf(tmp, sizeof(tmp), "%g", datav.number);
            data = tmp;
        }

        FILE *fp = fopen(path, "wb");
        if (!fp) {
            if (tmp_data) free(tmp_data);
            value_free(&pathv);
            value_free(&datav);
            return value_number(0);
        }

        fwrite(data, 1, strlen(data), fp);
        fclose(fp);

        if (tmp_data) free(tmp_data);

        value_free(&pathv);
        value_free(&datav);

        return value_number(1);
    }

    /* APPEND */
    if (strcmp(node->function_name, "jodo") == 0) {

        if (node->arg_count != 2) {
            shriji_error(
                E_PARSE_02,
                "jodo",
                "argument count mismatch",
                "use: jodo(\"file.txt\", \"data\")"
            );
            return value_null();
        }

        Value pathv = eval(node->args[0], env, rt);
        Value datav = eval(node->args[1], env, rt);

        if (pathv.type != VAL_STRING || !pathv.string) {
            value_free(&pathv);
            value_free(&datav);
            shriji_error(
                E_PARSE_02,
                "jodo",
                "file path must be string",
                "use: jodo(\"file.txt\", \"data\")"
            );
            return value_null();
        }

        const char *path = strip_quotes(pathv.string);

        const char *data = "";
        char tmp[128];
        tmp[0] = '\0';

        char *tmp_data = NULL;

        if (datav.type == VAL_STRING && datav.string) {
            tmp_data = unescape_string(strip_quotes(datav.string));
            data = tmp_data ? tmp_data : "";
        } else if (datav.type == VAL_NUMBER) {
            snprintf(tmp, sizeof(tmp), "%g", datav.number);
            data = tmp;
        }

        FILE *fp = fopen(path, "ab");
        if (!fp) {
            if (tmp_data) free(tmp_data);
            value_free(&pathv);
            value_free(&datav);
            return value_number(0);
        }

        fwrite(data, 1, strlen(data), fp);
        fclose(fp);

        if (tmp_data) free(tmp_data);

        value_free(&pathv);
        value_free(&datav);

        return value_number(1);
    }

    /*──────────────────────────────────────────────
      NORMAL FUNCTION CALL
    ──────────────────────────────────────────────*/
Value fnv = env_get(env, node->function_name);

if (fnv.type != VAL_FUNCTION) {
    shriji_error(
        E_PARSE_INVALID_TOKEN,
        "call",
        "function not found",
        "check name or define function"
    );
    return value_null();
}

ASTNode *fn = fnv.function;

    /* validate arg count */
    if (fn->rachna_param_count != node->arg_count) {
        shriji_error(
            E_PARSE_02,
            node->function_name,
            "argument count mismatch",
            "call with correct number of args"
        );
        return value_null();
    }

    /* save outer return state (nested calls safe) */
    int old_return_flag = rt->return_flag;
    Value old_return_value = rt->return_value;

    rt->return_flag = 0;
    rt->return_value = value_null();

    /* new call scope */
    env_push_scope(env);

    /* bind params */
       for (int i = 0; i < fn->rachna_param_count; i++) {
    ASTNode *p = fn->rachna_params[i];
        Value av = eval(node->args[i], env, rt);
        env_set(env, p->name, av);
    }

    /* run body */
    Value result = eval(fn->rachna_body, env, rt);

    if (rt->return_flag)
        result = rt->return_value;

    /* end call scope */
    env_pop_scope(env);

    /* restore outer return state */
    rt->return_flag = old_return_flag;
    rt->return_value = old_return_value;

   rt->call_depth--;
    return result;
}

/*──────────────────────────────────────────────
      CORE VALUES
    ──────────────────────────────────────────────*/
case AST_NUMBER:
    return value_number(node->number_value);

case AST_STRING:
    return value_string(node->string_value);

case AST_BOOL:
    return value_bool(node->bool_value);



    case AST_LIST: {
        int n = node->element_count;

        if (n <= 0 || !node->elements) {
            return value_list(NULL, 0);
        }

        Value *items = (Value *)malloc(sizeof(Value) * n);
        if (!items) return value_null();

        for (int i = 0; i < n; i++) {
            Value ev = eval(node->elements[i], env, rt);
            items[i] = value_copy(ev);
            value_free(&ev);
        }

        return value_list(items, n);
    }


case AST_DICT: {
        int n = node->dict_count;

        if (n <= 0 || !node->dict_keys || !node->dict_values) {
            return value_dict(NULL, NULL, 0);
        }

        Value *keys = (Value *)malloc(sizeof(Value) * n);
        Value *vals = (Value *)malloc(sizeof(Value) * n);

        if (!keys || !vals) {
            if (keys) free(keys);
            if (vals) free(vals);
            return value_null();
        }

        for (int i = 0; i < n; i++) {
            Value kv = eval(node->dict_keys[i], env, rt);
            Value vv = eval(node->dict_values[i], env, rt);

            keys[i] = value_copy(kv);
            vals[i] = value_copy(vv);

            value_free(&kv);
            value_free(&vv);
        }

        return value_dict(keys, vals, n);
    }

case AST_INDEX: {

    Value tv = eval(node->index_target, env, rt);
    Value iv = eval(node->index_expr, env, rt);

    /*──────────────────────────────────────────────
      LIST INDEXING: a[0]
    ──────────────────────────────────────────────*/
    if (tv.type == VAL_LIST) {

        if (!tv.items) {
            value_free(&tv);
            value_free(&iv);
            return value_null();
        }

        if (iv.type != VAL_NUMBER) {
            shriji_error(
                E_RUNTIME_01,
                "list",
                "List index number hona chahiye",
                "example: a[0]"
            );

               value_free(&tv);
               value_free(&iv);
               return value_null();
        }

        long long idx = iv.number;
        value_free(&iv);

        if (idx < 0 || idx >= tv.count) {
            shriji_error(
                E_RUNTIME_01,
                "list",
                "List index out of range",
                "example: a[0]"
            );
           value_free(&tv);
           return value_null();
        }

        Value out = value_copy(tv.items[(int)idx]);
        value_free(&tv);
        return out;
    }

    /*──────────────────────────────────────────────
      DICT INDEXING: d["a"]
    ──────────────────────────────────────────────*/
    if (tv.type == VAL_DICT) {

        if (iv.type != VAL_STRING || !iv.string) {
            shriji_error(
                E_RUNTIME_01,
                "dict",
                "dictionary ki key string honi chahiye",
                "example: d[\"a\"]"
            );
            value_free(&tv);
            value_free(&iv);
            return value_null();
        }

        for (int i = 0; i < tv.dict_count; i++) {

            Value k = tv.dict_keys[i];

            if (k.type == VAL_STRING &&
                strcmp(k.string, iv.string) == 0) {

                Value out = value_copy(tv.dict_values[i]);
                value_free(&tv);
                value_free(&iv);
                return out;
            }
        }

        shriji_error(
            E_RUNTIME_01,
            "dict",
            "dictionary me ye key maujood nahi hai",
            "example: d[\"a\"]"
        );
        value_free(&tv);
        value_free(&iv);
        return value_null();
    }

    /*──────────────────────────────────────────────
      UNSUPPORTED TARGET
    ──────────────────────────────────────────────*/
    value_free(&tv);
    value_free(&iv);
    return value_null();
}


case AST_INDEX_UPDATE: {

    /* target must be identifier so update persists */
    if (!node->index_target || node->index_target->type != AST_IDENTIFIER) {
        return value_null();
    }

    const char *var_name = node->index_target->name;

    /* fetch container directly from env */
    Value tv = env_get(env, var_name);

    Value iv = eval(node->index_expr, env, rt);
    Value vv = eval(node->index_value, env, rt);

    /*──────────────────────────────────────────────
      LIST INDEX UPDATE: a[0] = value
    ──────────────────────────────────────────────*/
    if (tv.type == VAL_LIST) {

        if (iv.type != VAL_NUMBER) goto cleanup;

        long long idx = iv.number;
        if (idx < 0 || idx >= tv.count) goto cleanup;

        value_free(&tv.items[(int)idx]);
        tv.items[(int)idx] = value_copy(vv);

        /* WRITE BACK */
        env_set(env, var_name, tv);

        value_free(&iv);
        return value_copy(vv);
    }

    /*──────────────────────────────────────────────
      DICT INDEX UPDATE: d["a"] = value
    ──────────────────────────────────────────────*/
    if (tv.type == VAL_DICT) {

        if (iv.type != VAL_STRING || !iv.string) goto cleanup;

        /* update if exists */
        for (int i = 0; i < tv.dict_count; i++) {
            Value k = tv.dict_keys[i];
            if (k.type == VAL_STRING && strcmp(k.string, iv.string) == 0) {

                value_free(&tv.dict_values[i]);
                tv.dict_values[i] = value_copy(vv);

                env_set(env, var_name, tv);
                value_free(&iv);
                return value_copy(vv);
            }
        }

        /* append new key */
        tv.dict_keys = realloc(tv.dict_keys,
                               sizeof(Value) * (tv.dict_count + 1));
        tv.dict_values = realloc(tv.dict_values,
                                 sizeof(Value) * (tv.dict_count + 1));

        tv.dict_keys[tv.dict_count]   = value_copy(iv);
        tv.dict_values[tv.dict_count] = value_copy(vv);
        tv.dict_count++;

        env_set(env, var_name, tv);
        value_free(&iv);
        return value_copy(vv);
    }

cleanup:
    value_free(&tv);
    value_free(&iv);
    value_free(&vv);
    return value_null();
}

case AST_IDENTIFIER: {

    const char *name = node->name;

    /* ===== ENV CHECK ===== */
    if (env_exists(env, name)) {
        return env_get(env, name);
    }

    /* ===== SESSION KEYWORD: last ===== */
    if (strcmp(name, "last") == 0) {
        return smriti_session_get_last();
    }

    /* ===== SMRITI FALLBACK ===== */
    const char *mem = smriti_personal_get(name);

    if (mem) {
        Value v;
        v.type = VAL_NUMBER;
        v.number = atof(mem);
        return v;
    }

    /* ===== ERROR ===== */
    shriji_error(
        E_ASSIGN_01,
        name,
        "variable declare nahi hua",
        "use mavi x = ... pehle likhiye"
    );

    return value_null();
}


case AST_ASSIGNMENT: {

    Value value = eval(node->value, env, rt);

    /* ===== FIX: SAFE COPY ===== */
    Value stored = value_copy(value);

    /* ===== ENV SAVE ===== */
    env_update(env, node->name, stored);

    /* ===== SMRITI SAVE ===== */
    if (value.type == VAL_NUMBER) {
        char buf[64];
        snprintf(buf, sizeof(buf), "%g", value.number);

        if (node->name[0] != '\0') {
            /* smriti_personal_set(node->name, buf); */
        }
    }

    event_fire(EVENT_ASSIGNMENT, node->name);
    state_on_success(&rt->state);

    return value;
}



case AST_UPDATE: {

    if (node->name[0] != '\0') {

        if (!env_exists(env, node->name)) {
            shriji_error(
                E_ASSIGN_01,
                node->name,
                "variable not declared",
                "use mavi x = ... first"
            );
            return value_null();
        }

        Value value = eval(node->value, env, rt);

        /*  FIXED: use env_update */
        env_update(env, node->name, value);

        event_fire(EVENT_ASSIGNMENT, node->name);
        state_on_success(&rt->state);

        return value;
    }

    shriji_error(
        E_ASSIGN_01,
        "update",
        "invalid update target",
        "use: x = expr"
    );

    return value_null();
}




case AST_BINARY: {
    Value Lv = eval(node->left, env, rt);
    Value Rv = eval(node->right, env, rt);

    char op = node->op[0];

    /* ───────────────────────────────────────────────
       TYPE CHECK
    ─────────────────────────────────────────────── */
    if (Lv.type != Rv.type) {
        shriji_error(
            E_RUNTIME_TYPE_MISMATCH,
            "binary",
            "type mismatch in expression",
            "same type ke values use karein"
        );
        return value_null();
    }

    /* ───────────────────────────────────────────────
       NUMBER OPERATIONS
    ─────────────────────────────────────────────── */
/*  LOGICAL OPERATORS (GLOBAL — BEFORE TYPE CHECK) */
if (op == '&') {
    return value_bool(
        value_is_truthy(Lv) && value_is_truthy(Rv)
    );
}

if (op == '|') {
    return value_bool(
        value_is_truthy(Lv) || value_is_truthy(Rv)
    );
}


    if (Lv.type == VAL_NUMBER) {

        double L = Lv.number;
        double R = Rv.number;


        if (op == '+') return value_number(L + R);
        if (op == '-') return value_number(L - R);
        if (op == '*') return value_number(L * R);
        if (op == '/') return value_number(safe_div(L, R));
        if (op == '%') return value_number(safe_mod(L, R));

        if (op == '>') return value_bool(L > R);
        if (op == '<') return value_bool(L < R);
        if (op == 'G') return value_bool(L >= R);
        if (op == 'L') return value_bool(L <= R);
        if (op == '=') return value_bool(L == R);
        if (op == '!') return value_bool(L != R);

        shriji_error(
            E_RUNTIME_01,
            "binary",
            "unknown numeric operator",
            NULL
        );
        return value_null();
    }

    /* ───────────────────────────────────────────────
       STRING COMPARISON ONLY
    ─────────────────────────────────────────────── */
    if (Lv.type == VAL_STRING) {

        const char *L = Lv.string ? Lv.string : "";
        const char *R = Rv.string ? Rv.string : "";

        int cmp = strcmp(L, R);

         if (op == '=') return value_bool(cmp == 0);
         if (op == '!') return value_bool(cmp != 0);
         if (op == '>') return value_bool(cmp > 0);
         if (op == '<') return value_bool(cmp < 0);
         if (op == 'G') return value_bool(cmp >= 0);
         if (op == 'L') return value_bool(cmp <= 0);


        shriji_error(
            E_RUNTIME_01,
            "binary",
            "string ke saath sirf comparison allowed hai",
            "use: == != < > <= >="
        );
        return value_null();
    }

    /* ───────────────────────────────────────────────
       UNSUPPORTED TYPE
    ─────────────────────────────────────────────── */
    shriji_error(
        E_PARSE_02,
        "binary",
        "ye operation is type ke liye allowed nahi hai",
        "sirf number ya string comparison supported hai"
    );
    rt->state.safety = STATE_CRITICAL;
                return value_null();
}


       case AST_NOT: {
    Value v = eval(node->left, env, rt);

    int truth = value_is_truthy(v);
    value_free(&v);

    return value_bool(!truth);
    }


case AST_IF: {
    Value condv = eval(node->condition, env, rt);

    Value result = value_null();

    if (value_is_truthy(condv)) {

        env_push_scope(env);

        result = eval(node->left, env, rt);

        env_pop_scope(env);
    }
    else if (node->else_block) {

        env_push_scope(env);

        result = eval(node->else_block, env, rt);

        env_pop_scope(env);
    }

    value_free(&condv);
    return result;
}



case AST_WHILE: {

    rt->loop_depth++;

    int local_guard = 0;

    while (1) {

        Value condv = eval(node->condition, env, rt);

        if (!value_is_truthy(condv)) {
            value_free(&condv);
            break;
        }

        value_free(&condv);

        rt->continue_flag = 0;

        eval(node->body, env, rt);

        if (rt->return_flag) {
            rt->loop_depth--;
            return rt->return_value;
        }

        if (rt->break_flag) {
            rt->break_flag = 0;
            break;
        }

        if (rt->continue_flag) {
            rt->continue_flag = 0;
            continue;
        }

        if (++local_guard > 1000000) {
            shriji_error(
                E_RUNTIME_01,
                "loop",
                "infinite loop detected",
                "use rukja or fix condition"
            );
            rt->state.safety = STATE_CRITICAL;
            break;
        }
    }

    rt->loop_depth--;
    return value_null();
}



case AST_BLOCK: {
    Value last = value_null();


    /*  NEW SCOPE START */
    env_push_scope(env);


    for (int i = 0; i < node->stmt_count; i++) {

        last = eval(node->statements[i], env, rt);

        if (rt->return_flag || rt->break_flag || rt->continue_flag)
            break;
    }

    /*  SCOPE END */
    env_pop_scope(env);


    if (rt->return_flag)
        return rt->return_value;

    return last;
}


case AST_PROGRAM: {
    Value last = value_null();


    for (int i = 0; i < node->stmt_count; i++) {

        last = eval(node->statements[i], env ,rt);

/* 🔥 ONLY break for loop control, not function return */
if (rt->break_flag || rt->continue_flag)
    break;
    }

    if (rt->return_flag)
        return rt->return_value;

    return last;
}

case AST_IMPORT: {

    /* 🔥 ENABLE ANALYSIS MODE */
    shriji_set_error_mode(ERROR_MODE_COLLECT);

    /* import "file.sri" */
    const char *path = node->name;

    if (!path || !*path) {
        shriji_error(
            E_PARSE_02,
            "import",
            "empty import path",
            "use: import \"file.sri\""
        );
        shriji_set_error_mode(ERROR_MODE_IMMEDIATE);
        return value_null();
    }

    /* read file */
    FILE *fp = fopen(path, "rb");
    if (!fp) {
        shriji_error(
            E_PARSE_02,
            "import",
            "file not found",
            "check path or create file first"
        );
        shriji_set_error_mode(ERROR_MODE_IMMEDIATE);
        return value_null();
    }

    fseek(fp, 0, SEEK_END);
    long sz = ftell(fp);
    fseek(fp, 0, SEEK_SET);

    if (sz < 0) sz = 0;
    if (sz > 1024 * 1024) sz = 1024 * 1024;

    char *buf = (char *)malloc((size_t)sz + 1);
    if (!buf) {
        fclose(fp);
        shriji_set_error_mode(ERROR_MODE_IMMEDIATE);
        return value_null();
    }

    size_t nread = fread(buf, 1, (size_t)sz, fp);
    buf[nread] = '\0';
    fclose(fp);

    /* 🔥 FULL FILE PARSE */
    ASTNode *prog = parse_program(buf);
    free(buf);

    if (!prog) {
        shriji_error(
            E_PARSE_02,
            "import",
            "parse failed",
            "check full file syntax"
        );

        shriji_print_all_errors();
        shriji_set_error_mode(ERROR_MODE_IMMEDIATE);
        return value_null();
    }

    /* 🔥 RUN FULL PROGRAM */
    Value result = eval(prog, env, rt);

    /* 🔥 PRINT ALL ERRORS (ONE SHOT) */
    shriji_print_all_errors();

    /* 🔥 BACK TO NORMAL MODE */
    shriji_set_error_mode(ERROR_MODE_IMMEDIATE);

    return result;
}

case AST_COMMAND: {

    CommandType cmd = resolve_command(node->command_name);
    if (cmd == CMD_UNKNOWN) {
        shriji_error(
            E_PARSE_02,
            "command",
            "unknown command",
            "command not recognized"
        );
        return value_null();
    }

if (cmd == CMD_BOLO) {

    /* 🔥 GLOBAL ERROR GUARD */
    if (error_reported) {
        return value_null();
    }

    Value v = value_null();

    if (node->value)
        v = eval(node->value, env , rt);

/*  RUNTIME ERROR GUARD */
if (rt->error_flag) {
    return value_null();
}

if (v.type == VAL_STRING) {
    printf("%s\n", v.string ? v.string : "");
}
else if (v.type == VAL_NUMBER) {
    printf("%g\n", v.number);
}
else if (v.type == VAL_BOOL) {
    printf("%d\n", v.boolean);
}
else if (v.type == VAL_LIST) {
    printf("[");
    for (int i = 0; i < v.count; i++) {
        if (i > 0) printf(", ");

        if (v.items[i].type == VAL_NUMBER)
            printf("%g", v.items[i].number);
        else if (v.items[i].type == VAL_STRING)
            printf("\"%s\"", v.items[i].string);
        else if (v.items[i].type == VAL_BOOL)
            printf("%d", v.items[i].boolean);
        else
            printf("?");
    }
    printf("]\n");
}
else {
    /*  DOUBLE GUARD */
    if (!rt->error_flag && !error_reported) {
        printf("0\n");
    }
}



    smriti_session_set_last(v);

    value_free(&v);
    return value_null();
}

/* AI MODULES → ai_router (STEP-13 REAL OUTPUT) */
if (cmd == CMD_SAKHI || cmd == CMD_NIYU ||
    cmd == CMD_MIRA  || cmd == CMD_KAVYA) {

    Value v = value_null();
    if (node->value)
        v = eval(node->value, env , rt);

    const char *arg = "";
    char tmp[64] = {0};

    if (v.type == VAL_STRING && v.string)
        arg = strip_quotes(v.string);
    else if (v.type == VAL_NUMBER) {
        snprintf(tmp, sizeof(tmp), "%g", v.number);
        arg = tmp;
    }

    ShrijiBridgePacket pkt;
    pkt.text   = arg;
    pkt.lang   = SHRIJI_LANG_AUTO;
    pkt.source = SHRIJI_SRC_REPL;

    L3Response r = ai_router_dispatch(&pkt);
    if (r.text && *r.text)
        printf("%s\n", r.text);

    value_free(&v);
    return value_null();
}

    /* DEFAULT COMMAND */
    return value_number(execute_command(cmd));
}

default:
    return value_null();
}
}

/*──────────────────────────────────────────────────────────────
 | PROGRAM RUNNER
 *──────────────────────────────────────────────────────────────*/
Value run_program(ASTNode *program, Env *env, ShrijiRuntime *rt)
{
    if (!program || program->type != AST_PROGRAM)
        return value_null();

    runtime_reset(rt);
    event_bind_runtime(rt);
    return eval(program, env, rt);
}
nd{verbatim}
\section*{File: ./src/normalize_engine.c}
egin{verbatim}
#include <string.h>
#include <ctype.h>
#include <stdio.h>
#include <stdio.h>

#include "../include/normalize_engine.h"

#define MAX_WORDS 50
#define MAX_WORD_LEN 64

/* ================================
   BASIC UTILS
================================ */

static void to_lowercase(char *s)
{
    for (int i = 0; s[i]; i++)
        s[i] = tolower(s[i]);
}

static void trim_spaces(char *s)
{
    int i = 0, j = 0;

    while (isspace(s[i])) i++;

    while (s[i])
    {
        if (!(isspace(s[i]) && isspace(s[i+1])))
            s[j++] = s[i];
        i++;
    }

    if (j > 0 && isspace(s[j-1])) j--;
    s[j] = '\0';
}

/* ================================
   TYPO DICTIONARY
================================ */

typedef struct {
    const char *wrong;
    const char *correct;
} TypoRule;

static TypoRule rules[] = {
    {"reeor", "error"},
    {"eror", "error"},
    {"softwere", "software"},
    {"creat", "create"},
    {"keise", "kaise"},
    {"helo", "hello"},
    {"hy", "hi"},
    {"sotwere", "software"},
    {"cteate", "create"},
    {"by", "bye"},

};

static int rule_count = sizeof(rules)/sizeof(TypoRule);

/* ================================
   WORD SPLIT
================================ */

static void split_words(char *input, char words[][MAX_WORD_LEN], int *count)
{
    char *token = strtok(input, " ");
    int i = 0;

    while (token && i < MAX_WORDS)
    {
        strncpy(words[i], token, MAX_WORD_LEN - 1);
        words[i][MAX_WORD_LEN - 1] = '\0';
        i++;
        token = strtok(NULL, " ");
    }

    *count = i;
}

/* ================================
   APPLY TYPO RULES
================================ */

static void apply_typo_rules(char words[][MAX_WORD_LEN], int count)
{
    for (int i = 0; i < count; i++)
    {
        for (int j = 0; j < rule_count; j++)
        {
            if (strcmp(words[i], rules[j].wrong) == 0)
            {
                strncpy(words[i], rules[j].correct, MAX_WORD_LEN - 1);
            }
        }
    }
}

/* ================================
   JOIN WORDS BACK
================================ */

static void join_words(char *output, char words[][MAX_WORD_LEN], int count)
{
    output[0] = '\0';

    for (int i = 0; i < count; i++)
    {
        strncat(output, words[i], MAX_WORD_LEN - 1);
        if (i < count - 1)
            strncat(output, " ", 1);
    }
}

/* ================================
   SPECIAL CASE: EXIT NIRMAN
================================ */

static void fix_joined_commands(char *input)
{
    if (strstr(input, "exitnirman"))
    {
        strcpy(input, "exit nirman");
    }
}

/* ================================
   MAIN FUNCTION
================================ */
void normalize_input(char *input)
{

    if (!input || !*input)
        return;

    char temp[512];
    strncpy(temp, input, sizeof(temp)-1);
    temp[sizeof(temp)-1] = '\0';

    to_lowercase(temp);
    trim_spaces(temp);

    fix_joined_commands(temp);

    char words[MAX_WORDS][MAX_WORD_LEN];
    int count = 0;

    split_words(temp, words, &count);
    apply_typo_rules(words, count);

char result[512];
join_words(result, words, count);

/* safe copy back */
snprintf(input, 512, "%s", result);
}
nd{verbatim}
\section*{File: ./src/smriti_personal.c}
egin{verbatim}
#include "../include/smriti_personal.h"
#include <string.h>
#include <stdio.h>

#define MAX_KEYS 100

typedef struct {
    char key[32];
    char value[128];
} MemoryItem;

/* ===============================
   INTERNAL MEMORY (RAM)
   =============================== */
static MemoryItem memory[MAX_KEYS];
static int memory_count = 0;

/* ===============================
   FILE PATH
   =============================== */
static const char *SMRITI_FILE = ".shriji/smriti.db";

/* ===============================
   LOAD FROM FILE
   =============================== */
void smriti_personal_load(void)
{
    FILE *fp = fopen(SMRITI_FILE, "r");
    if (!fp) return;

    char line[256];

    while (fgets(line, sizeof(line), fp))
    {
        char *eq = strchr(line, '=');
        if (!eq) continue;

        *eq = '\0';

        char *key = line;
        char *value = eq + 1;

        /* remove newline */
        value[strcspn(value, "\n")] = 0;

        smriti_personal_set(key, value);
    }

    fclose(fp);
}

/* ===============================
   SAVE TO FILE
   =============================== */
void smriti_personal_save(void)
{
    FILE *fp = fopen(SMRITI_FILE, "w");
    if (!fp) return;

    for (int i = 0; i < memory_count; i++)
    {
        fprintf(fp, "%s=%s\n", memory[i].key, memory[i].value);
    }

    fclose(fp);
}

/* ===============================
   SET (WRITE)
   =============================== */
void smriti_personal_set(const char *key, const char *value)
{
    if (!key || !value) return;

    /* overwrite if exists */
    for (int i = 0; i < memory_count; i++)
    {
        if (strcmp(memory[i].key, key) == 0)
        {
            strncpy(memory[i].value, value, sizeof(memory[i].value)-1);
            memory[i].value[sizeof(memory[i].value)-1] = '\0';

            smriti_personal_save();
            return;
        }
    }

    /* new entry */
    if (memory_count < MAX_KEYS)
    {
        strncpy(memory[memory_count].key, key, sizeof(memory[memory_count].key)-1);
        strncpy(memory[memory_count].value, value, sizeof(memory[memory_count].value)-1);

        memory[memory_count].key[sizeof(memory[memory_count].key)-1] = '\0';
        memory[memory_count].value[sizeof(memory[memory_count].value)-1] = '\0';

        memory_count++;

        smriti_personal_save();
    }
}

/* ===============================
   GET (READ)
   =============================== */
const char* smriti_personal_get(const char *key)
{
    if (!key) return NULL;

    for (int i = 0; i < memory_count; i++)
    {
        if (strcmp(memory[i].key, key) == 0)
        {
            return memory[i].value;
        }
    }

    return NULL;
}

/* ===============================
   CLEAR (RAM ONLY)
   =============================== */
void smriti_personal_clear(void)
{
    memory_count = 0;
}
nd{verbatim}
\section*{File: ./src/nirman_flow_engine.c}
egin{verbatim}
#include <stdio.h>
#include <string.h>
#include "../include/nirman.h"

void nirman_flow_process(char *line)
{
    int intent = nirman_detect_intent(line);

    /* STAGE 0 → START */
    if (intent == 2 && NIRMAN_CTX.flow_stage == 0)
    {
        printf("Aap kya banana chahte ho?\n");
        NIRMAN_CTX.flow_stage = 1;
        return;
    }

    /* STAGE 1 → GOAL CAPTURE */
    if (NIRMAN_CTX.flow_stage == 1)
    {
        if (strstr(line, "software") || strstr(line, "create"))
        {
            printf("Aap kya banana chahte ho?\n");
            return;
        }

        strncpy(NIRMAN_CTX.goal, line, sizeof(NIRMAN_CTX.goal) - 1);
        printf("[FLOW] Target set: %s\n", NIRMAN_CTX.goal);

        NIRMAN_CTX.flow_stage = 2;
        return;
    }

    /* STAGE 2 → ASK TYPE */
    if (NIRMAN_CTX.flow_stage == 2)
    {
        if (strstr(NIRMAN_CTX.goal, "calculator"))
        {
            printf("Aapko kis type ka calculator chahiye?\n");
            NIRMAN_CTX.flow_stage = 3;
            return;
        }
    }

    /* STAGE 3 → HANDLE TYPE */
    if (NIRMAN_CTX.flow_stage == 3)
    {
        if (strstr(NIRMAN_CTX.goal, "calculator"))
        {
            if (strstr(line, "normal"))
            {
                printf("[FLOW] Type selected: normal calculator\n");
                printf("Agla step: features define karte hain\n");

                NIRMAN_CTX.flow_stage = 4;
            }
            else if (strstr(line, "scientific"))
            {
                printf("[FLOW] Type selected: scientific calculator\n");
                printf("Agla step: features define karte hain\n");

                NIRMAN_CTX.flow_stage = 4;
            }
            else
            {
                printf("Type clear nahi hai (normal / scientific)\n");
            }
            return;
        }
    }

    /* STAGE 4 → NEXT STEP (placeholder) */
    if (NIRMAN_CTX.flow_stage == 4)
    {
        printf("[FLOW] Feature planning phase...\n");
        return;
    }
}
nd{verbatim}
\section*{File: ./src/error_intelligence.c}
egin{verbatim}
#include <stdio.h>
#include <string.h>
#include <stdlib.h>
#include <ctype.h>

#ifndef _WIN32
#include <readline/readline.h>
#include <readline/history.h>
#endif

#include "../include/error.h"
#include "../include/pragya_avastha.h"
#include "../include/fix_rules.h"

/* ===============================
   PREFILL (FOR EDIT MODE)
   =============================== */

static char safe_buffer[256];
static char *prefill_text = NULL;

#ifndef _WIN32
static int prefill_hook(void)
{
    if (prefill_text)
    {
        rl_insert_text(prefill_text);
        rl_point = strlen(prefill_text);
        prefill_text = NULL;
    }
    return 0;
}
#endif

/* ===============================
   USER INPUT (STRICT)
   =============================== */

static char get_user_choice()
{
#ifdef _WIN32
    char buffer[16];

    if (!fgets(buffer, sizeof(buffer), stdin))
        return 'i';

    char c = tolower(buffer[0]);

    if (c == 'e') return 'e';
    if (c == 'i') return 'i';

    return 'i';
#else
    char *line = readline("");

    if (!line || !*line)
    {
        free(line);
        return 'i';
    }

    char c = tolower(line[0]);
    free(line);

    if (c == 'e') return 'e';
    if (c == 'i') return 'i';

    return 'i';
#endif
}

/* ===============================
   POINTER DISPLAY
   =============================== */

static void print_pointer(const char *text, int pos)
{
    if (!text || pos < 0) return;

    printf("%s\n", text);

    for (int i = 0; i < pos; i++)
        printf(" ");

    printf("^\n\n");
}

/* ===============================
   BASIC VALIDATION
   =============================== */

static int contains_invalid_token(const char *input)
{
    for (int i = 0; input[i]; i++)
    {
        char c = input[i];

        if (isalnum(c) || isspace(c) || strchr("+-*/()._{}[]=<>!\"'", c))
            continue;

        return 1;
    }
    return 0;
}

static int has_unmatched_brackets(const char *input)
{
    int count = 0;

    for (int i = 0; input[i]; i++)
    {
        if (input[i] == '(') count++;
        else if (input[i] == ')') count--;
    }

    return count != 0;
}

static int is_valid_input(const char *input)
{
    if (!input || !*input) return 0;
    if (contains_invalid_token(input)) return 0;
    if (has_unmatched_brackets(input)) return 0;
    return 1;
}

/* ===============================
   SIMPLE OPERATOR CLEAN FIX
   =============================== */

static int filter_operators(const char *input, char *output)
{
    int i, j = 0;
    int last_was_op = 0;

    for (i = 0; input[i]; i++)
    {
        char c = input[i];

        if (strchr("+-*/", c))
        {
            if (last_was_op) continue;

            output[j++] = c;
            last_was_op = 1;
        }
        else
        {
            output[j++] = c;
            last_was_op = 0;
        }
    }

    output[j] = '\0';
    return strcmp(input, output) != 0;
}

/* ===============================
   MAIN ERROR INTELLIGENCE
   =============================== */

void shriji_error_intelligence(
    PragyaAvastha *avastha,
    const ShrijiErrorInfo *err,
    void *niyu_result)
{
    (void)niyu_result;

    if (!avastha || !err || !avastha->raw_text)
        return;

    const char *text = avastha->raw_text;

    /* ===== VALIDATION ===== */

    if (!is_valid_input(text))
    {
        avastha->stop_execution = 1;

        int pos = (err && err->col > 0) ? err->col - 1 : 0;
        print_pointer(text, pos);
        return;
    }

    /* ===== RULE BASED FIX ===== */

    FixRule *rule = shriji_get_rule_for_error(err->code);

    if (rule && rule->apply_fix)
    {
        char buffer[256];
        snprintf(buffer, sizeof(buffer), "%s", text);

        if (rule->apply_fix(buffer, sizeof(buffer), err))
        {
            if (rule->safety == FIX_SAFE_DETERMINISTIC)
            {
                printf("\n[Auto Fix Applied]\n");
                printf("Reason: deterministic rule\n");
                printf("Result: %s\n\n", buffer);

                avastha->has_correction = 1;
                avastha->stop_execution = 1;

                snprintf(avastha->corrected_text,
                         sizeof(avastha->corrected_text),
                         "%s", buffer);

                return;
            }

            printf("\nSuggestion: %s\n", buffer);
            printf("\nEdit (E) / Ignore (I): ");

            char choice = get_user_choice();

            if (choice == 'e')
            {
                avastha->stop_execution = 1;

                printf("\nEdit:\n%s\n", text);

#ifndef _WIN32
                snprintf(safe_buffer, sizeof(safe_buffer), "%s", text);
                prefill_text = safe_buffer;

                rl_startup_hook = prefill_hook;
#endif
                return;
            }

            avastha->stop_execution = 1;
            return;
        }
    }

    /* ===== SIMPLE FILTER FIX ===== */

    char filtered[256];

    if (filter_operators(text, filtered))
    {
        printf("\nSuggestion: %s\n", filtered);
        printf("\nEdit (E) / Ignore (I): ");

        char choice = get_user_choice();

        if (choice == 'e')
        {
            avastha->stop_execution = 1;

            printf("\nEdit:\n%s\n", text);

#ifndef _WIN32
            snprintf(safe_buffer, sizeof(safe_buffer), "%s", text);
            prefill_text = safe_buffer;

            rl_startup_hook = prefill_hook;
#endif
            return;
        }

        avastha->stop_execution = 1;
        return;
    }

    /* ===== FALLBACK POINTER ===== */

    int pos = (err && err->col > 0) ? err->col - 1 : 0;

    print_pointer(text, pos);

    avastha->stop_execution = 1;
}
nd{verbatim}
\section*{File: ./src/example_builder.c}
egin{verbatim}
#include <stdio.h>
#include <string.h>

/*
    Build readable example from user input
*/

void shriji_build_example(const char *input, char *out, int size)
{
    int i = 0;
    int j = 0;

    while (input[i] && j < size - 1)
    {
        char c = input[i];

        if (c == '+' || c == '-' || c == '*' || c == '/')
        {
            if (j < size - 3)
            {
                out[j++] = ' ';
                out[j++] = c;
                out[j++] = ' ';
            }
        }
        else
        {
            out[j++] = c;
        }

        i++;
    }

    out[j] = '\0';
}
nd{verbatim}
\section*{File: ./src/ai_intent.c}
egin{verbatim}
#include <string.h>
#include <ctype.h>

#include "../include/ai_intent.h"

/* lowercase helper */
static void to_lower(char *s)
{
    for (; *s; s++)
        *s = (char)tolower(*s);
}

/* simple substring match (L3 — lightweight, safe) */
static int has_word(const char *msg, const char *w)
{
    return strstr(msg, w) != NULL;
}

AIIntent ai_detect_intent(const char *message)
{
    if (!message || !*message)
        return AI_INTENT_UNKNOWN;

    char buf[256];
    strncpy(buf, message, sizeof(buf) - 1);
    buf[sizeof(buf) - 1] = '\0';
    to_lower(buf);

    /*──────────────────────────────────────────────
      EXPLAIN / SAMJHAO (LOGIC / TEACHING)
      → Niyu (primary explain engine)
    ──────────────────────────────────────────────*/
    if (has_word(buf, "samjhao") ||
        has_word(buf, "samjha")  ||
        has_word(buf, "samajhao")||
        has_word(buf, "samajh")  ||
        has_word(buf, "explain") ||
        has_word(buf, "what is") ||
        has_word(buf, "kya hota"))
        return AI_INTENT_EXPLAIN;

    /*──────────────────────────────────────────────
      WHY / REASON
    ──────────────────────────────────────────────*/
    if (has_word(buf, "why")  ||
        has_word(buf, "kyun") ||
        has_word(buf, "kyu")  ||
        has_word(buf, "kyon"))
        return AI_INTENT_WHY;

    /*──────────────────────────────────────────────
      CALC / MATH
    ──────────────────────────────────────────────*/
    if (has_word(buf, "calculate") ||
        has_word(buf, "solve")     ||
        has_word(buf, "math")      ||
        has_word(buf, "hisab"))
        return AI_INTENT_CALC;

    /*──────────────────────────────────────────────
      EMOTION / FEELING
    ──────────────────────────────────────────────*/
    if (has_word(buf, "sad")    ||
        has_word(buf, "dukhi")  ||
        has_word(buf, "lonely") ||
        has_word(buf, "akela")  ||
        has_word(buf, "dard"))
        return AI_INTENT_EMOTION;

    return AI_INTENT_UNKNOWN;
}
nd{verbatim}
\section*{File: ./src/lang/grammar.c}
egin{verbatim}
#include "../../include/lang/grammar.h"
#include "../../include/token.h"

/*
 * SHRIJILANG GRAMMAR CORE
 * Phase-1: token presence validation only
 * No parser logic change yet
 */

void shriji_register_grammar_core(void)
{
    /*
     * IMPORTANT:
     * This function is called once from parse_program()
     * Do NOT put parser logic here.
     * Only grammar / token registration / validation.
     */

#ifdef SHRIJI_ENABLE_MASTER_TOKENS
    /* Phase-1: sanity check hook */
    /* (empty by design for now) */
#endif
}

/* Future phases — intentionally empty */

void shriji_register_grammar_logic(void) {}
void shriji_register_grammar_incdec(void) {}
void shriji_register_grammar_loop(void) {}
void shriji_register_grammar_func(void) {}
nd{verbatim}
\section*{File: ./src/lang/grammar/grammar_registry.c}
egin{verbatim}
#include "lang/grammar.h"

/* Phase-1: empty implementations
 * These will be filled in Phase-2
 */

void shriji_register_grammar_core(void) {
    /* core grammar hooks (placeholder) */
}

void shriji_register_grammar_logic(void) {
    /* && || ! (placeholder) */
}

void shriji_register_grammar_incdec(void) {
    /* ++ -- (placeholder) */
}

void shriji_register_grammar_loop(void) {
    /* jabtak / for (placeholder) */
}

void shriji_register_grammar_func(void) {
    /* kaam / wapas (placeholder) */
}
nd{verbatim}
\section*{File: ./src/lang/grammar/grammar_core.c}
egin{verbatim}
#include "../../../include/lang/grammar_core.h"

/*
 * Phase-2 Grammar Core
 * Classification only (no parsing / execution)
 */

GrammarKind shriji_grammar_kind(TokenType t)
{
    switch (t) {

        /* ───────── STATEMENTS ───────── */
        case TOKEN_MAVI:
        case TOKEN_KAAM:
        case TOKEN_AGAR:
        case TOKEN_JABTAK:
        case TOKEN_WAPAS:
        case TOKEN_IMPORT:
        case TOKEN_WARNA:
        case TOKEN_RUKJA:
        case TOKEN_CHALU:
            return GRAMMAR_STATEMENT;

        /* ───────── LITERALS ───────── */
        case TOKEN_NUMBER:
        case TOKEN_STRING:
        case TOKEN_TRUE:
        case TOKEN_FALSE:
            return GRAMMAR_LITERAL;

        /* ───────── INC / DEC ───────── */
        case TOKEN_PLUSPLUS:
        case TOKEN_MINUSMINUS:
            return GRAMMAR_INCDEC;

        /* ───────── LOGICAL / ARITHMETIC OPERATORS ───────── */
        case TOKEN_PLUS:
        case TOKEN_MINUS:
        case TOKEN_STAR:
        case TOKEN_SLASH:
        case TOKEN_MOD:

        case TOKEN_EQEQ:
        case TOKEN_NEQ:
        case TOKEN_GT:
        case TOKEN_LT:
        case TOKEN_GTE:
        case TOKEN_LTE:

        case TOKEN_AND:
        case TOKEN_OR:
        case TOKEN_NOT:
            return GRAMMAR_OPERATOR;

        default:
            return GRAMMAR_UNKNOWN;
    }
}

const char *shriji_grammar_kind_name(GrammarKind k)
{
    switch (k) {
        case GRAMMAR_STATEMENT:  return "STATEMENT";
        case GRAMMAR_EXPRESSION: return "EXPRESSION";
        case GRAMMAR_OPERATOR:   return "OPERATOR";
        case GRAMMAR_LITERAL:    return "LITERAL";
        case GRAMMAR_INCDEC:     return "INCDEC";
        case GRAMMAR_RESERVED:   return "RESERVED";
        default:                 return "UNKNOWN";
    }
}
nd{verbatim}
\section*{File: ./src/error.c}
egin{verbatim}
#include <stdio.h>
#include <string.h>

#include "../include/error.h"
#include "../include/runtime.h"
#include "../gyaan/core/gyaan_engine.h"

/*──────────────────────────────────────────────*/
int error_reported = 0;

static ShrijiErrorMode CURRENT_ERROR_MODE = ERROR_MODE_IMMEDIATE;

ShrijiErrorInfo LAST_ERROR = {0};


/* 🔥 ERROR STORAGE */
#define MAX_ERRORS 100

typedef struct {
    int line;
    int col;
    char code[32];
    char context[64];
    char message[256];
    char hint[256];
} StoredError;

static StoredError ERROR_LIST[MAX_ERRORS];
static int ERROR_COUNT = 0;


/*──────────────────────────────────────────────*/
void shriji_set_error_mode(ShrijiErrorMode mode)
{
    CURRENT_ERROR_MODE = mode;
}


/*──────────────────────────────────────────────*/
static const char *error_code_str(ShrijiErrorCode code)
{
    switch (code) {

        case E_ASSIGN_01: return "E_ASSIGN_01";
        case E_ASSIGN_02: return "E_ASSIGN_02";

        case E_PARSE_01: return "E_PARSE_01";
        case E_PARSE_02: return "E_PARSE_02";
        case E_PARSE_EXTRA_OPERATOR: return "E_PARSE_EXTRA_OPERATOR";
        case E_PARSE_MISSING_OPERATOR: return "E_PARSE_MISSING_OPERATOR";
        case E_PARSE_DOUBLE_OPERATOR: return "E_PARSE_DOUBLE_OPERATOR";
        case E_PARSE_OPERATOR_START: return "E_PARSE_OPERATOR_START";
        case E_PARSE_OPERATOR_END: return "E_PARSE_OPERATOR_END";
        case E_PARSE_OPERATOR_CHAIN: return "E_PARSE_OPERATOR_CHAIN";
        case E_PARSE_MISSING_OPERAND: return "E_PARSE_MISSING_OPERAND";

        case E_PARSE_UNMATCHED_PAREN: return "E_PARSE_UNMATCHED_PAREN";
        case E_PARSE_BRACKET_MISSING: return "E_PARSE_BRACKET_MISSING";
        case E_PARSE_BRACKET_EXTRA: return "E_PARSE_BRACKET_EXTRA";
        case E_PARSE_INVALID_TOKEN: return "E_PARSE_INVALID_TOKEN";

        case E_IF_01: return "E_IF_01";

        case E_RUNTIME_01: return "E_RUNTIME_01";
        case E_RUNTIME_DIV_ZERO: return "E_RUNTIME_DIV_ZERO";
        case E_RUNTIME_TYPE_MISMATCH: return "E_RUNTIME_TYPE_MISMATCH";
        case E_RUNTIME_LOOP_LIMIT: return "E_RUNTIME_LOOP_LIMIT";
        case E_RUNTIME_INDEX_ERROR: return "E_RUNTIME_INDEX_ERROR";

        default: return "E_UNKNOWN";
    }
}


/*──────────────────────────────────────────────*/
static const char *error_message_for_code(ShrijiErrorCode code)
{
    switch (code) {

        case E_PARSE_DOUBLE_OPERATOR:
            return "Do operators ek saath aa gaye hain.";

        case E_PARSE_OPERATOR_CHAIN:
            return "Operator sequence galat lag raha hai.";

        case E_PARSE_OPERATOR_START:
            return "Expression operator se start nahi ho sakta.";

        case E_PARSE_OPERATOR_END:
            return "Expression operator par end ho gaya.";

        case E_PARSE_MISSING_OPERAND:
            return "Operator ke liye value missing hai.";

        default:
            return NULL;
    }
}


/*──────────────────────────────────────────────*/
static void store_error(
    ShrijiErrorCode code,
    const char *context,
    const char *message,
    const char *hint,
    int line,
    int col
)
{
    if (ERROR_COUNT >= MAX_ERRORS) return;

    StoredError *e = &ERROR_LIST[ERROR_COUNT++];

    snprintf(e->code, sizeof(e->code), "%s", error_code_str(code));
    snprintf(e->context, sizeof(e->context), "%s", context ? context : "");
    snprintf(e->message, sizeof(e->message), "%s", message ? message : "");
    snprintf(e->hint, sizeof(e->hint), "%s", hint ? hint : "");

    e->line = line;
    e->col = col;
}


/*──────────────────────────────────────────────*/
void shriji_error(
    ShrijiErrorCode code,
    const char *context,
    const char *message,
    const char *hint
)
{
    error_reported = 1;

    if (current_runtime) {
        current_runtime->error_flag = 1;
    }

    const char *auto_msg = error_message_for_code(code);
    if (auto_msg)
        message = auto_msg;

    LAST_ERROR.code = code;
    LAST_ERROR.context = context;
    LAST_ERROR.message = message;
    LAST_ERROR.hint = hint;
    LAST_ERROR.function = NULL;
    LAST_ERROR.has_location = 0;

    if (CURRENT_ERROR_MODE == ERROR_MODE_IMMEDIATE) {

        printf("\n⛔ %s | %s → %s\n",
               error_code_str(code),
               context ? context : "",
               message ? message : "");

        if (hint && hint[0])
            printf("   hint: %s\n", hint);

        printf("\n");

        const ShrijiErrorInfo *last = shriji_last_error();
        if (last) {
            gyaan_print(last);
        }

    } else {
        store_error(code, context, message, hint, 0, 0);
    }
}


/*──────────────────────────────────────────────*/
void shriji_error_at_full(
    Token tok,
    ShrijiErrorCode code,
    const char *context,
    const char *function,
    const char *message,
    const char *hint
)
{
    error_reported = 1;

    if (current_runtime) {
        current_runtime->error_flag = 1;
    }

    const char *auto_msg = error_message_for_code(code);
    if (auto_msg)
        message = auto_msg;

    LAST_ERROR.code = code;
    LAST_ERROR.context = context;
    LAST_ERROR.message = message;
    LAST_ERROR.hint = hint;
    LAST_ERROR.function = function;
    LAST_ERROR.has_location = 1;
    LAST_ERROR.line = tok.line;
    LAST_ERROR.col = tok.col;

    if (CURRENT_ERROR_MODE == ERROR_MODE_IMMEDIATE) {

        printf("\n⛔ %s | %s → %s (line %d, col %d)\n",
               error_code_str(code),
               context ? context : "",
               message ? message : "",
               tok.line,
               tok.col);

        if (hint && hint[0])
            printf("   hint: %s\n", hint);

        printf("\n");

        const ShrijiErrorInfo *last = shriji_last_error();
        if (last) {
            gyaan_print(last);
        }

    } else {
        store_error(code, context, message, hint, tok.line, tok.col);
    }
}


/*──────────────────────────────────────────────*/
void shriji_error_at(
    Token tok,
    ShrijiErrorCode code,
    const char *context,
    const char *message,
    const char *hint
)
{
    shriji_error_at_full(
        tok,
        code,
        context,
        "unknown",
        message,
        hint
    );
}


/*──────────────────────────────────────────────*/
void shriji_print_all_errors(void)
{
    if (ERROR_COUNT == 0) return;

    printf("\nKRST ANALYSIS REPORT\n");
    printf("──────────────────────────────\n");

    for (int i = 0; i < ERROR_COUNT; i++) {
        StoredError *e = &ERROR_LIST[i];

        printf("%d) %s | %s → %s",
               i + 1,
               e->code,
               e->context,
               e->message);

        if (e->line > 0)
            printf(" (line %d, col %d)", e->line, e->col);

        printf("\n");

        if (e->hint[0])
            printf("   hint: %s\n", e->hint);

        printf("\n");
    }

    printf("──────────────────────────────\n\n");

    ERROR_COUNT = 0;

    const ShrijiErrorInfo *last = shriji_last_error();
    if (last) {
        gyaan_print(last);
    }
}


/*──────────────────────────────────────────────*/
const ShrijiErrorInfo *shriji_last_error(void)
{
    if (!error_reported)
        return NULL;

    return &LAST_ERROR;
}
nd{verbatim}
\section*{File: ./src/smriti_session.c}
egin{verbatim}
#include "../include/smriti_session.h"
#include <string.h>
#include <stdlib.h>

/* ===============================
   INTERNAL STORAGE
   =============================== */

static char last_input[256] = {0};
static char last_output[256] = {0};

static Value last_value;
static int has_last = 0;

/* ===============================
   LAST INPUT
   =============================== */

void smriti_session_set_last_input(const char *text)
{
    if (!text) return;

    strncpy(last_input, text, sizeof(last_input) - 1);
    last_input[sizeof(last_input) - 1] = '\0';
}

const char* smriti_session_get_last_input(void)
{
    return last_input;
}

/* ===============================
   LAST OUTPUT
   =============================== */

void smriti_session_set_last_output(const char *text)
{
    if (!text) return;

    strncpy(last_output, text, sizeof(last_output) - 1);
    last_output[sizeof(last_output) - 1] = '\0';
}

const char* smriti_session_get_last_output(void)
{
    return last_output;
}

/* ===============================
   LAST VALUE (UPGRADED)
   =============================== */

void smriti_session_set_last(Value v)
{
    if (has_last)
        value_free(&last_value);

    last_value = value_copy(v);
    has_last = 1;
}

Value smriti_session_get_last(void)
{
    if (!has_last)
        return value_null();

    return value_copy(last_value);
}

/* ===============================
   CLEAR SESSION
   =============================== */

void smriti_session_clear(void)
{
    last_input[0] = '\0';
    last_output[0] = '\0';

    if (has_last)
        value_free(&last_value);

    has_last = 0;
}
nd{verbatim}
\section*{File: ./src/event.c}
egin{verbatim}
#include <string.h>

#include "../include/event.h"
#include "../include/state.h"
#include "../include/runtime.h"

/* Internal runtime binding */
static ShrijiRuntime *BOUND_RUNTIME = NULL;

/* Bind runtime once per execution */
void event_bind_runtime(ShrijiRuntime *rt)
{
    BOUND_RUNTIME = rt;
}

/* Fire event */
void event_fire(EventType type, const char *context)
{
    (void)context;

    if (!BOUND_RUNTIME)
        return;

    switch (type)
    {
        case EVENT_EXECUTION_SUCCESS:
        case EVENT_ASSIGNMENT:
            state_on_success(&BOUND_RUNTIME->state);
            break;

        case EVENT_SHUNYA_ACCESS:
        case EVENT_ERROR:
            state_on_error(&BOUND_RUNTIME->state);
            break;

        case EVENT_FATAL_ERROR:
            state_on_error(&BOUND_RUNTIME->state);
            state_on_error(&BOUND_RUNTIME->state);
            break;

        case EVENT_EXECUTION_BLOCKED:
            BOUND_RUNTIME->state.safety = STATE_CRITICAL;
            BOUND_RUNTIME->state.human  = HUMAN_MANDATORY;
            break;

        default:
            break;
    }
}
nd{verbatim}
\section*{File: ./src/ast.c}
egin{verbatim}
#include <stdio.h>
#include <string.h>
#include <stdlib.h>

#include "../include/ast.h"

/*──────────────────────────────────────────────────────────────
 |  SHRIJILANG — AST CONSTRUCTOR IMPLEMENTATION (CLEAN & STABLE)
 *──────────────────────────────────────────────────────────────*/

/*──────────────────────────────────────────────────────────────
 | SECTION A — SAFE STRING HELPER
 *──────────────────────────────────────────────────────────────*/
static void safe_copy(char *dest, const char *src, int max)
{
    if (!src) {
        dest[0] = '\0';
        return;
    }
    strncpy(dest, src, max - 1);
    dest[max - 1] = '\0';
}

/*──────────────────────────────────────────────────────────────
 | SECTION B — CORE NODES
 *──────────────────────────────────────────────────────────────*/
ASTNode *new_number_node(double value)
{
    ASTNode *node = calloc(1, sizeof(ASTNode));
    if (!node) return NULL;

    node->type = AST_NUMBER;
    node->number_value = value;
    node->chakra_state = CHAKRA_OK;
    return node;
}

ASTNode *new_bool_node(int value)
{
    ASTNode *node = calloc(1, sizeof(ASTNode));
    if (!node) return NULL;

    node->type = AST_BOOL;
    node->bool_value = value;
    node->chakra_state = CHAKRA_OK;

    return node;
}

ASTNode *new_identifier_node(const char *name)
{
    ASTNode *node = calloc(1, sizeof(ASTNode));
    if (!node) return NULL;

    node->type = AST_IDENTIFIER;
    safe_copy(node->name, name, 128);
    node->chakra_state = CHAKRA_OK;
    return node;
}

ASTNode *new_string_node(const char *str)
{
    ASTNode *node = calloc(1, sizeof(ASTNode));
    if (!node) return NULL;

    node->type = AST_STRING;
    safe_copy(node->string_value, str, 256);
    node->chakra_state = CHAKRA_OK;
    return node;
}

ASTNode *new_binary_node(char op, ASTNode *l, ASTNode *r)
{
    ASTNode *node = calloc(1, sizeof(ASTNode));
    if (!node) return NULL;

    node->type = AST_BINARY;
    node->op[0] = op;
    node->left = l;
    node->right = r;
    node->chakra_state = CHAKRA_OK;
    return node;
}

ASTNode *new_unary_node(Token op, ASTNode *expr)
{
    ASTNode *node = calloc(1, sizeof(ASTNode));
    if (!node) return NULL;

    node->type = AST_UNARY;

    /* store operator */
    node->op[0] = op.start[0];
    node->op[1] = '\0';

    node->left = expr;

    node->chakra_state = CHAKRA_OK;
    return node;
}
/*──────────────────────────────────────────────────────────────
 | SECTION B.1 — ASSIGNMENT / UPDATE
 *──────────────────────────────────────────────────────────────*/
ASTNode *new_assignment_node(const char *name, ASTNode *val)
{
    ASTNode *node = calloc(1, sizeof(ASTNode));
    if (!node) return NULL;

    node->type = AST_ASSIGNMENT;
    safe_copy(node->name, name, 128);
    node->value = val;
    node->chakra_state = CHAKRA_OK;
    return node;
}

ASTNode *new_update_node(const char *name, ASTNode *val)
{
    ASTNode *node = calloc(1, sizeof(ASTNode));
    if (!node) return NULL;

    node->type = AST_UPDATE;
    safe_copy(node->name, name, 128);
    node->value = val;
    node->chakra_state = CHAKRA_OK;
    return node;
}

/*──────────────────────────────────────────────────────────────
 | SECTION C — AI NODES
 *──────────────────────────────────────────────────────────────*/
#define AI_NODE(fn, type_id)                      \
ASTNode *fn(ASTNode *msg) {                       \
    ASTNode *node = calloc(1, sizeof(ASTNode));   \
    if (!node) return NULL;                       \
    node->type = type_id;                         \
    node->value = msg;                            \
    node->chakra_state = CHAKRA_OK;               \
    return node;                                  \
}

AI_NODE(new_sakhi_node, AST_SAKHI)
AI_NODE(new_niyu_node,  AST_NIYU)
AI_NODE(new_shiri_node, AST_SHIRI)
AI_NODE(new_mira_node,  AST_MIRA)
AI_NODE(new_kavya_node, AST_KAVYA)

/*──────────────────────────────────────────────────────────────
 | SECTION D — COMMAND + IMPORT
 *──────────────────────────────────────────────────────────────*/
ASTNode *new_command_node(const char *cmd, ASTNode *arg)
{
    ASTNode *node = calloc(1, sizeof(ASTNode));
    if (!node) return NULL;

    node->type = AST_COMMAND;
    safe_copy(node->command_name, cmd, 128);
    node->value = arg;
    node->chakra_state = CHAKRA_OK;
    return node;
}

ASTNode *new_import_node(const char *module_name)
{
    ASTNode *node = calloc(1, sizeof(ASTNode));
    if (!node) return NULL;

    node->type = AST_IMPORT;
    safe_copy(node->name, module_name, 128);
    node->chakra_state = CHAKRA_OK;
    return node;
}

/*──────────────────────────────────────────────────────────────
 | SECTION D.1 — BLOCK / PROGRAM
 *──────────────────────────────────────────────────────────────*/
ASTNode *new_block_node(ASTNode **stmts, int count)
{
    ASTNode *node = calloc(1, sizeof(ASTNode));
    if (!node) return NULL;

    node->type = AST_BLOCK;
    node->statements = stmts;
    node->stmt_count = count;
    node->chakra_state = CHAKRA_OK;
    return node;
}

ASTNode *new_program_node(ASTNode **stmts, int count)
{
    ASTNode *node = calloc(1, sizeof(ASTNode));
    if (!node) return NULL;

    node->type = AST_PROGRAM;
    node->statements = stmts;
    node->stmt_count = count;
    node->chakra_state = CHAKRA_OK;
    return node;
}

/*──────────────────────────────────────────────────────────────
 | SECTION D.2 — CONTROL FLOW
 *──────────────────────────────────────────────────────────────*/
ASTNode *new_if_node(ASTNode *condition, ASTNode *then_block)
{
    ASTNode *node = calloc(1, sizeof(ASTNode));
    if (!node) return NULL;

    node->type = AST_IF;
    node->condition = condition;
    node->left = then_block;     /* then-block */
    node->chakra_state = CHAKRA_OK;
    return node;
}

ASTNode *new_while_node(ASTNode *condition, ASTNode *body)
{
    ASTNode *node = calloc(1, sizeof(ASTNode));
    if (!node) return NULL;

    node->type = AST_WHILE;
    node->condition = condition;
    node->body = body;
    node->chakra_state = CHAKRA_OK;
    return node;
}

/*──────────────────────────────────────────────────────────────
 | SECTION D.2.3 — LOOP CONTROL (rukja / chalu)
 *──────────────────────────────────────────────────────────────*/
ASTNode *new_break_node(void)
{
    ASTNode *node = calloc(1, sizeof(ASTNode));
    if (!node) return NULL;

    node->type = AST_BREAK;
    node->chakra_state = CHAKRA_OK;
    return node;
}

ASTNode *new_continue_node(void)
{
    ASTNode *node = calloc(1, sizeof(ASTNode));
    if (!node) return NULL;

    node->type = AST_CONTINUE;
    node->chakra_state = CHAKRA_OK;
    return node;
}

/*──────────────────────────────────────────────────────────────
 | SECTION D.3 — FUNCTIONS
 *──────────────────────────────────────────────────────────────*/
ASTNode *new_function_node(
    const char *name,
    ASTNode **params,
    int param_count,
    ASTNode *body
) {
    ASTNode *node = calloc(1, sizeof(ASTNode));
    if (!node) return NULL;

    node->type = AST_FUNCTION;
    safe_copy(node->function_name, name, 128);

    node->params = params;
    node->param_count = param_count;
    node->body = body;

    node->chakra_state = CHAKRA_OK;
    return node;
}

ASTNode *new_rachna_node(
    const char *name,
    ASTNode **params,
    int param_count,
    ASTNode *body
)
{
    ASTNode *node = calloc(1, sizeof(ASTNode));
    if (!node) return NULL;

    node->type = AST_RACHNA;

    /* name */
    safe_copy(node->rachna_name, name, 128);

    /* params */
    node->rachna_params = params;
    node->rachna_param_count = param_count;

    /* body */
    node->rachna_body = body;

    node->chakra_state = CHAKRA_OK;
    return node;
}

ASTNode *new_call_node(
    const char *name,
    ASTNode **args,
    int arg_count
) {
    ASTNode *node = calloc(1, sizeof(ASTNode));
    if (!node) return NULL;

    node->type = AST_CALL;
    safe_copy(node->function_name, name, 128);

    node->args = args;
    node->arg_count = arg_count;

    node->chakra_state = CHAKRA_OK;
    return node;
}

ASTNode *new_return_node(ASTNode *expr)
{
    ASTNode *node = calloc(1, sizeof(ASTNode));
    if (!node) return NULL;

    node->type = AST_RETURN;
    node->return_expr = expr;
    node->chakra_state = CHAKRA_OK;
    return node;
}

/*──────────────────────────────────────────────
 | SECTION E — MEMORY CLEANUP
 *──────────────────────────────────────────────*/
void ast_free(ASTNode *node)
{
    if (!node) return;

    ast_free(node->left);
    ast_free(node->right);
    ast_free(node->value);
    ast_free(node->condition);
    ast_free(node->else_block);

    ast_free(node->body);
    ast_free(node->return_expr);

 /* LIST */
    if (node->elements) {
    for (int i = 0; i < node->element_count; i++) {
       ast_free(node->elements[i]);
    }
       free(node->elements);
    }
    if (node->statements) {
        for (int i = 0; i < node->stmt_count; i++)
            ast_free(node->statements[i]);
        free(node->statements);
    }

    if (node->params) {
        for (int i = 0; i < node->param_count; i++)
            ast_free(node->params[i]);
        free(node->params);
    }

    if (node->args) {
        for (int i = 0; i < node->arg_count; i++)
            ast_free(node->args[i]);
        free(node->args);
    }
/* DICT */
    if (node->dict_keys) {
        for (int i = 0; i < node->dict_count; i++) {
            ast_free(node->dict_keys[i]);
        }
        free(node->dict_keys);
    }

    if (node->dict_values) {
        for (int i = 0; i < node->dict_count; i++) {
            ast_free(node->dict_values[i]);
        }
        free(node->dict_values);
    }
    free(node);
}

/*──────────────────────────────────────────────
 | SECTION F — LOGICAL NOT NODE (nahi)
 *──────────────────────────────────────────────*/
ASTNode *new_not_node(ASTNode *expr)
{
    ASTNode *node = calloc(1, sizeof(ASTNode));
    if (!node) return NULL;

    node->type = AST_NOT;
    node->left = expr;
    node->right = NULL;
    node->chakra_state = CHAKRA_OK;
    return node;
}

/*────────────────────────────────────────
  Inc / Dec (++ / --)
────────────────────────────────────────*/

ASTNode *new_inc_node(ASTNode *expr, int is_prefix)
{
    ASTNode *n = calloc(1, sizeof(ASTNode));
    n->type = AST_INC;

    n->left = expr;          // operand (++x or x++)
    n->is_prefix = is_prefix;

    return n;
}

ASTNode *new_dec_node(ASTNode *expr, int is_prefix)
{
    ASTNode *n = calloc(1, sizeof(ASTNode));
    n->type = AST_DEC;

    n->left = expr;          // operand (--x or x--)
    n->is_prefix = is_prefix;

    return n;
}

/*──────────────────────────────────────────────
  LIST / INDEX NODES
──────────────────────────────────────────────*/

ASTNode *new_list_node(ASTNode **elements, int count)
{
    ASTNode *node = (ASTNode *)malloc(sizeof(ASTNode));
    if (!node) return NULL;

    memset(node, 0, sizeof(ASTNode));
    node->type = AST_LIST;

    node->elements = elements;
    node->element_count = count;
    node->chakra_state = CHAKRA_OK;
    return node;
}

ASTNode *new_index_node(ASTNode *target, ASTNode *index_expr)
{
    ASTNode *node = (ASTNode *)malloc(sizeof(ASTNode));
    if (!node) return NULL;

    memset(node, 0, sizeof(ASTNode));
    node->type = AST_INDEX;

    node->index_target = target;
    node->index_expr   = index_expr;

    return node;
}

ASTNode *new_index_update_node(ASTNode *target, ASTNode *index_expr, ASTNode *value)
{
    ASTNode *node = (ASTNode *)calloc(1, sizeof(ASTNode));
    node->type = AST_INDEX_UPDATE;

    node->index_target = target;
    node->index_expr = index_expr;
    node->index_value = value;

    return node;
}
/*──────────────────────────────────────────────
  DICT / MAP NODES
──────────────────────────────────────────────*/
ASTNode *new_dict_node(ASTNode **keys, ASTNode **values, int count)
{
    ASTNode *node = (ASTNode *)malloc(sizeof(ASTNode));
    if (!node) return NULL;

    memset(node, 0, sizeof(ASTNode));
    node->type = AST_DICT;

    node->dict_keys = keys;
    node->dict_values = values;
    node->dict_count = count;
    node->chakra_state = CHAKRA_OK;
    return node;
}
nd{verbatim}
\section*{File: ./src/fix_interactive.c}
egin{verbatim}
#include <stdio.h>
#include <string.h>

#include "../include/fix_interactive.h"
#include "../include/error.h"

/* =========================================================
   MEMORY — LAST INPUT TRACKING
========================================================= */

static char last_input[512] = {0};

/* =========================================================
   SAFE LINE INPUT
========================================================= */

static int read_line(char *buffer, int size)
{
    if (!fgets(buffer, size, stdin))
        return 0;

    buffer[strcspn(buffer, "\n")] = '\0';
    return 1;
}

/* =========================================================
   SAME INPUT GUARD
========================================================= */

static int is_same_input(const char *input)
{
    if (strcmp(last_input, input) == 0)
    {
        printf("\n[SYSTEM] Same input detected. Please modify it.\n");
        return 1;
    }

    strncpy(last_input, input, sizeof(last_input) - 1);
    last_input[sizeof(last_input) - 1] = '\0';

    return 0;
}

/* =========================================================
   FIX MENU CORE (SHARED)
========================================================= */

static int show_fix_menu(const char *input, char *output, int size)
{
    while (1)
    {
        printf("\nFix Options:\n\n");
        printf("M) Manual edit\n");
        printf("E) Quick edit\n");
        printf("I) Ignore\n");

        printf("\nSelect: ");

        char choice = getchar();
        getchar(); /* consume newline */

        /* ========================
           M — Manual Edit
        ======================== */
        if (choice == 'M' || choice == 'm')
        {
            printf("Enter new expression: ");

            if (!read_line(output, size))
                return 0;

            if (is_same_input(output))
                continue;

            return 1;
        }

        /* ========================
           E — Quick Edit
        ======================== */
        if (choice == 'E' || choice == 'e')
        {
            printf("Edit:\n%s\n> ", input);

            if (!read_line(output, size))
                return 0;

            if (is_same_input(output))
                continue;

            return 1;
        }

        /* ========================
           I — Ignore
        ======================== */
        if (choice == 'I' || choice == 'i')
        {
            return 0;
        }

        printf("Invalid choice. Try again.\n");
    }
}

/* =========================================================
   PUBLIC WRAPPERS
========================================================= */

int shriji_interactive_double_operator_fix(
    const char *input,
    char *output,
    int size)
{
    return show_fix_menu(input, output, size);
}

int shriji_interactive_bracket_fix(
    const char *input,
    char *output,
    int size)
{
    return show_fix_menu(input, output, size);
}

int shriji_interactive_operator_fix(
    const char *input,
    char *output,
    int size)
{
    return show_fix_menu(input, output, size);
}
nd{verbatim}
\section*{File: ./src/main.c}
egin{verbatim}
/*
──────────────────────────────────────────────
    SHRIJILANG — MAIN ENTRY (FINAL REPL 🌸)
──────────────────────────────────────────────
*/

#include <stdio.h>
#include <stdlib.h>
#include <string.h>

#ifndef _WIN32
#include <readline/readline.h>
#include <readline/history.h>
#endif

#include "../include/env.h"
#include "../include/smriti.h"
#include "../include/smriti_personal.h"
#include "../include/smriti_session.h"
#include "../include/runtime.h"

#include "../krst/krst_core.h"
#include "../include/krst_router.h"
#include "../include/user_config.h"

#include "../include/normalize_engine.h"
#include "../include/nirman.h"
#include "../include/nirman_flow_engine.h"

#include "../include/parser.h"

/* GLOBAL ENV */
Env *GLOBAL_ENV = NULL;

/*──────────────────────────────────────────────*/
static void shriji_init(void)
{
    config_init();
    GLOBAL_ENV = new_env();

    smriti_init();
    smriti_session_clear();
    smriti_personal_load();
}

/*──────────────────────────────────────────────*/
static void run_repl_mode(void)
{
    printf("──────────────────────────────────────────────\n");
    printf("        🌸 ShrijiLang Shell 🌸\n");
    printf("──────────────────────────────────────────────\n");

    printf("Type 'exit' to quit\n\n");

    ShrijiRuntime rt;
    runtime_init(&rt);
    current_runtime = &rt;

    /* MULTI-LINE BUFFER 🌸 */
    char buffer[8192];
    buffer[0] = '\0';

    /* BRACE TRACKER 🌸 */
    int brace_count = 0;

    while (1)
    {
#ifdef _WIN32
        char line_buffer[1024];

        printf("Shiri> ");
        if (!fgets(line_buffer, sizeof(line_buffer), stdin))
            break;

        line_buffer[strcspn(line_buffer, "\n")] = '\0';
        char *line = line_buffer;
#else
        char *line = readline("Shiri> ");
        if (!line) break;

        if (nirman_is_active())
            normalize_input(line);
#endif

        runtime_reset(&rt);

#ifndef _WIN32
        if (*line) add_history(line);
#endif

        /* EXIT */
        if (strcmp(line, "exit") == 0)
        {
            printf("🌸 Radhe Radhe 🌸\n");
            smriti_session_clear();
#ifndef _WIN32
            free(line);
#endif
            break;
        }

        /* DEV MODE SYSTEM 🌸 */

        if (strcmp(line, "dev") == 0)
        {
            printf("Developer mode: %s\n", DEV_MODE ? "ON" : "OFF");
#ifndef _WIN32
            free(line);
#endif
            continue;
        }

        if (strcmp(line, "dev on") == 0)
        {
            DEV_MODE = 1;
            printf("Developer mode enabled\n");
#ifndef _WIN32
            free(line);
#endif
            continue;
        }

        if (strcmp(line, "dev off") == 0)
        {
            DEV_MODE = 0;
            printf("Developer mode disabled\n");
#ifndef _WIN32
            free(line);
#endif
            continue;
        }

        /* EMPTY INPUT */
        if (line[0] == '\0')
        {
#ifndef _WIN32
            free(line);
#endif
            continue;
        }

        /* NIRMAN COMMAND */
        if (strcmp(line, "nirman") == 0)
        {
            nirman_start();
            printf("Nirman mode activated\n");
            printf("Welcome to Shriji World\n");
#ifndef _WIN32
            free(line);
#endif
            continue;
        }

        if (strcmp(line, "exit nirman") == 0 ||
            strcmp(line, "nirman exit") == 0)
        {
            nirman_stop();
            printf("Nirman mode deactivated\n");
#ifndef _WIN32
            free(line);
#endif
            continue;
        }

        /* NIRMAN FLOW */
        if (nirman_is_active())
        {
            nirman_flow_process(line);
#ifndef _WIN32
            free(line);
#endif
            continue;
        }

        /* APPEND TO BUFFER 🌸 */
        strcat(buffer, line);
        strcat(buffer, "\n");

        /* COUNT BRACES 🌸 */
        for (int i = 0; line[i]; i++)
        {
            if (line[i] == '{') brace_count++;
            if (line[i] == '}') brace_count--;
        }

        /* WAIT UNTIL BLOCK COMPLETE 🌸 */
        if (brace_count != 0)
        {
#ifndef _WIN32
            free(line);
#endif
            continue;
        }

        /* PARSE */
        ASTNode *tree = parse_program(buffer);

        if (!tree)
        {
#ifndef _WIN32
            free(line);
#endif
            buffer[0] = '\0';
            brace_count = 0;
            continue;
        }

        /* EXECUTE */
        smriti_session_set_last_input(buffer);

        int fix_attempt = 0;
        const int MAX_FIX_ATTEMPT = 5;

        while (fix_attempt < MAX_FIX_ATTEMPT)
        {
            KRSTRequest req;

            req.input_text = buffer;
            req.has_correction = 0;
            req.corrected_text[0] = '\0';

            krst_route_request(&req);

            if (!req.has_correction)
                break;

            strcpy(buffer, req.corrected_text);
            fix_attempt++;
        }

        /* RESET BUFFER 🌸 */
        buffer[0] = '\0';
        brace_count = 0;

#ifndef _WIN32
        free(line);
#endif
    }
}

/*──────────────────────────────────────────────*/
static int run_file_mode(const char *path)
{
    FILE *fp = fopen(path, "rb");

    if (!fp)
    {
        printf("File not found: %s\n", path);
        return 1;
    }

    fseek(fp, 0, SEEK_END);
    long size = ftell(fp);
    fseek(fp, 0, SEEK_SET);

    char *buffer = malloc(size + 1);
    if (!buffer)
    {
        fclose(fp);
        return 1;
    }

    size_t read_bytes = fread(buffer, 1, size, fp);
    if (read_bytes != (size_t)size)
    {
        printf("File read error\n");
        free(buffer);
        fclose(fp);
        return 1;
    }

    buffer[size] = '\0';
    fclose(fp);

    smriti_session_set_last_input(buffer);

    ShrijiRuntime rt;
    runtime_init(&rt);
    current_runtime = &rt;

    KRSTRequest req;
    req.input_text = buffer;
    req.has_correction = 0;

    krst_route_request(&req);

    if (req.has_correction)
    {
        req.input_text = req.corrected_text;
        krst_route_request(&req);
    }

    free(buffer);
    return 0;
}

/*──────────────────────────────────────────────*/
int main(int argc, char *argv[])
{
    shriji_init();

    if (argc > 1)
        return run_file_mode(argv[1]);

    run_repl_mode();

    return 0;
}
nd{verbatim}
\section*{File: ./src/command.c}
egin{verbatim}
#include <stdio.h>
#include <string.h>
#include <time.h>

#include "../include/command.h"
#include "../include/event.h"

/*──────────────────────────────────────────────────────────────
 | COMMAND ALIAS TABLE
 | 108 public + remaining locked (sample shown)
 *──────────────────────────────────────────────────────────────*/
static CommandAlias COMMAND_ALIASES[] = {

    /* LIST */
    { "ls",        CMD_LIST,   CMD_PUBLIC },
    { "list",      CMD_LIST,   CMD_PUBLIC },
    { "radha",     CMD_LIST,   CMD_PUBLIC },
    { "r",         CMD_LIST,   CMD_PUBLIC },

    /* CD */
    { "cd",        CMD_CD,     CMD_PUBLIC },
    { "chalo",     CMD_CD,     CMD_PUBLIC },

    /* MKDIR */
    { "mkdir",     CMD_MKDIR,  CMD_PUBLIC },
    { "bana",      CMD_MKDIR,  CMD_PUBLIC },

    /* BOLO (PRINT) */
    { "bolo",      CMD_BOLO,   CMD_PUBLIC },

    /* SAMAY (TIME) ✅ */
    { "samay",     CMD_SAMAY,  CMD_PUBLIC },
    { "time",      CMD_SAMAY,  CMD_PUBLIC },
    { "clock",     CMD_SAMAY,  CMD_PUBLIC },

/* AI Modules */
    { "sakhi",     CMD_SAKHI,  CMD_PUBLIC },
    { "niyu",      CMD_NIYU,   CMD_PUBLIC },
    { "mira",      CMD_MIRA,   CMD_PUBLIC },
    { "kavya",     CMD_KAVYA,  CMD_PUBLIC },
    { "shiri",     CMD_SHIRI,  CMD_PUBLIC },

    /* LOCKED / RESERVED */
    { "shyam",     CMD_DELETE, CMD_LOCKED },
    { "keshav",    CMD_COPY,   CMD_LOCKED },

    { NULL, CMD_UNKNOWN, CMD_LOCKED }
};

/*──────────────────────────────────────────────────────────────
 | Resolve user-facing name → internal command
 *──────────────────────────────────────────────────────────────*/
CommandType resolve_command(const char *name)
{
    if (!name)
        return CMD_UNKNOWN;

    for (int i = 0; COMMAND_ALIASES[i].alias; i++) {

        if (strcmp(COMMAND_ALIASES[i].alias, name) == 0) {

            if (COMMAND_ALIASES[i].access == CMD_LOCKED) {
                event_fire(EVENT_ERROR, "command locked");
                return CMD_UNKNOWN;
            }

            return COMMAND_ALIASES[i].cmd;
        }
    }

    return CMD_UNKNOWN;
}

/*──────────────────────────────────────────────────────────────
 | Internal command name (debug / AI explain)
 *──────────────────────────────────────────────────────────────*/
const char *command_name(CommandType cmd)
{
    switch (cmd) {
        case CMD_LIST:   return "CMD_LIST";
        case CMD_CD:     return "CMD_CD";
        case CMD_MKDIR:  return "CMD_MKDIR";
        case CMD_BOLO:   return "CMD_BOLO";
        case CMD_SAMAY:  return "CMD_SAMAY";

        case CMD_SAKHI:  return "CMD_SAKHI";
        case CMD_NIYU:   return "CMD_NIYU";
        case CMD_MIRA:   return "CMD_MIRA";
        case CMD_KAVYA:  return "CMD_KAVYA";
        case CMD_SHIRI:  return "CMD_SHIRI";

        case CMD_DELETE: return "CMD_DELETE";
        case CMD_COPY:   return "CMD_COPY";
        case CMD_MOVE:   return "CMD_MOVE";

        default:         return "CMD_UNKNOWN";
    }
}

/*──────────────────────────────────────────────────────────────
 | Execute internal command
 | NOTE:
 |  • BOLO printing is handled in interpreter (AST_COMMAND)
 |  • This file only supports commands that need no args
 *──────────────────────────────────────────────────────────────*/
int execute_command(CommandType cmd)
{
    switch (cmd) {

        case CMD_LIST:
            event_fire(EVENT_EXECUTION_SUCCESS, "list executed");
            printf("[LIST] executed (stub)\n");
            return 0;

        case CMD_CD:
            event_fire(EVENT_EXECUTION_SUCCESS, "cd executed");
            printf("[CD] executed (stub)\n");
            return 0;

        case CMD_MKDIR:
            event_fire(EVENT_EXECUTION_SUCCESS, "mkdir executed");
            printf("[MKDIR] executed (stub)\n");
            return 0;

        case CMD_BOLO:
            /* printing is done in interpreter */
            event_fire(EVENT_EXECUTION_SUCCESS, "bolo executed");
            return 0;

        case CMD_SAMAY: {
            /* ✅ print system time in 12-hour format (AM/PM) */
            event_fire(EVENT_EXECUTION_SUCCESS, "samay executed");

            time_t now = time(NULL);
            struct tm *t = localtime(&now);

            if (!t) {
                printf("00:00 AM\n");
                return 0;
            }

            int hour = t->tm_hour;
            int minute = t->tm_min;

            const char *ampm = "AM";
            if (hour >= 12) ampm = "PM";

            hour = hour % 12;
            if (hour == 0) hour = 12;

            printf("%02d:%02d %s\n", hour, minute, ampm);
            return 0;
        }

        case CMD_SAKHI:
        case CMD_NIYU:
        case CMD_MIRA:
        case CMD_KAVYA:
        case CMD_SHIRI:
            /* handled inside interpreter because requires args */
            event_fire(EVENT_ERROR, "AI command requires args");
            return -1;
        default:
            event_fire(EVENT_ERROR, "command not implemented");
            return -1;
    }
}
nd{verbatim}
\section*{File: ./src/atma_lekha.c}
egin{verbatim}
#include "../include/atma_lekha.h"
#include <stdio.h>
#include <time.h>

static const char *ATMA_FILE = ".shriji/atma_lekha.log";

void atma_lekha_init(void)
{
    FILE *fp = fopen(ATMA_FILE, "a");
    if (fp) fclose(fp);
}

void atma_lekha_note(AtmaEventType type,
                     const char *component,
                     const char *detail)
{
    FILE *fp = fopen(ATMA_FILE, "a");
    if (!fp) return;

    time_t now = time(NULL);
    struct tm *t = localtime(&now);

    fprintf(fp,
        "[%04d-%02d-%02d %02d:%02d:%02d] | %d | %s | %s\n",
        t->tm_year + 1900,
        t->tm_mon + 1,
        t->tm_mday,
        t->tm_hour,
        t->tm_min,
        t->tm_sec,
        type,
        component ? component : "unknown",
        detail ? detail : "-"
    );

    fclose(fp);
}
nd{verbatim}
\section*{File: ./src/typo_engine.c}
egin{verbatim}
#include <string.h>
#include "../include/typo_engine.h"

/*
    Simple exit typo detection

    Rule:
    - Length must be 4
    - At least 3 characters match with "exit"
*/

int shriji_check_typo(
    const char *input,
    char *suggestion,
    int suggestion_size
)
{
    if (!input || !suggestion)
        return 0;

    const char *target = "exit";

    if (strlen(input) != 4)
        return 0;

    int match = 0;

    for (int i = 0; i < 4; i++) {
        if (input[i] == target[i])
            match++;
    }

    if (match >= 3) {
        strncpy(suggestion, target, suggestion_size - 1);
        suggestion[suggestion_size - 1] = '\0';
        return 1;
    }

    return 0;
}
nd{verbatim}
\section*{File: ./src/parser.c}
egin{verbatim}
#include <stdio.h>
#include <stdlib.h>
#include <string.h>

#include "../include/parser.h"
#include "../include/token.h"
#include "../include/ast.h"
#include "../include/error.h"

#include "../include/example_builder.h"

#include "lang/grammar.h"

#ifdef SHRIJI_ENABLE_MASTER_TOKENS
#include "lang/token_master.h"
#endif



/*──────────────────────────────────────────────
 | SHRIJILANG PARSER — V1 STABLE (FULL)
 | Supports:
 |   - mavi (declaration)
 |   - update (x = expr)
 |   - if/else (agar/warna)
 |   - while (jabtak)
 |   - break/continue (rukja/chalu)
 |   - return (wapas)
 |   - functions (kaam)
 |   - function call: add(10, 20)
 |   - newline aware
 |   - list literal: [1, 2, 3]
 |   - indexing: a[0]
 *──────────────────────────────────────────────*/




static ASTNode *parse_statement(void);
static ASTNode *parse_block(void);
static ASTNode *parse_if(void);
static ASTNode *parse_while(void);
static ASTNode *parse_return(void);
static ASTNode *parse_function(void);

static ASTNode *parse_rachna_definition(void);

static ASTNode *parse_assignment(void);

/* expression chain */
static ASTNode *parse_comparison(void);
static ASTNode *parse_expression(void);
static ASTNode *parse_term(void);
static ASTNode *parse_unary(void);
static ASTNode *parse_primary(void);

static ASTNode *parse_value(void);

/* list support */
static ASTNode *parse_list_literal(void);
static ASTNode *parse_dict_literal(void);
static ASTNode *parse_postfix(ASTNode *base);

static ASTNode *parse_import(void);

static Token current;

static Token prev_token;

static const char *current_source = NULL;

static void parser_error(
    Token tok,
    ShrijiErrorCode code,
    const char *context,
    const char *message,
    const char *hint
)
{
    shriji_error_at_full(tok, code, context, __func__, message, hint);
}







/*──────────────────────────────────────────────
 | HELPERS
 *──────────────────────────────────────────────*/
static void advance(void)
{
    prev_token = current;

    do {
        current = scan_token();
    } while (current.type == TOKEN_NEWLINE);
}

static int expect(TokenType t, ShrijiErrorCode c,
                  const char *ctx, const char *msg, const char *hint)
{
    if (current.type == t) {
        advance();
        return 1;
    }

   parser_error(current, c, ctx, msg, hint);
    return 0;
}

static void skip_separators(void)
{
    while (current.type == TOKEN_NEWLINE) {
        advance();
    }
}

static void unescape_into(char *out, int outcap, const char *in, int inlen)
{
    int oi = 0;

    for (int i = 0; i < inlen && oi < outcap - 1; i++) {
        char c = in[i];

        if (c == '\\' && i + 1 < inlen) {
            char n = in[i + 1];
            i++;

            switch (n) {
                case 'n':  out[oi++] = '\n'; break;
                case 't':  out[oi++] = '\t'; break;
                case 'r':  out[oi++] = '\r'; break;
                case '"':  out[oi++] = '"';  break;
                case '\\': out[oi++] = '\\'; break;
                default:
                    out[oi++] = n;
                    break;
            }
            continue;
        }

        out[oi++] = c;
    }

    out[oi] = 0;
}






/*──────────────────────────────────────────────
 | PARSER SYNCHRONIZE
 *──────────────────────────────────────────────*/







static void parser_sync(void)
{
    while (current.type != TOKEN_NEWLINE &&
           current.type != TOKEN_EOF)
    {
        advance();
    }
}








/*──────────────────────────────────────────────
 | PROGRAM
 *──────────────────────────────────────────────*/

ASTNode *parse_program(const char *source)
{
   current_source = source;

#ifdef SHRIJI_ENABLE_MASTER_TOKENS
    shriji_register_grammar_core();   // 👈 यही FINAL जगह है
#endif

    init_tokenizer(source);
    advance();

    ASTNode **stmts = NULL;
    int count = 0;

    while (current.type != TOKEN_EOF) {
   error_reported = 0;
        skip_separators();
        if (current.type == TOKEN_EOF)
            break;






ASTNode *s = parse_statement();

if (!s) {
    // 🔥 ERROR RECOVERY
    while (current.type != TOKEN_NEWLINE &&
           current.type != TOKEN_EOF) {
        advance();
    }

    skip_separators();
    continue;
}
        ASTNode **tmp = realloc(stmts, sizeof(ASTNode*) * (count + 1));
        if (!tmp) {
            free(stmts);
            return NULL;
        }
        stmts = tmp;

        stmts[count++] = s;

        skip_separators();
    }

    return new_program_node(stmts, count);
}










/*──────────────────────────────────────────────
  | STATEMENT
  *──────────────────────────────────────────────*/
static ASTNode *parse_statement(void)
{
    ASTNode *node = NULL;

    if (current.type == TOKEN_LEFT_BRACE) return parse_block();

     if (current.type == TOKEN_RACHNA)
    return parse_rachna_definition();

    if (current.type == TOKEN_AGAR) {
        ASTNode *i = parse_if();
        if (!i) printf("DEBUG: IF NULL\n");
        return i;
    }

    if (current.type == TOKEN_JABTAK) {
        ASTNode *w = parse_while();
        if (!w) printf("DEBUG: WHILE NULL\n");
        return w;
    }

    if (current.type == TOKEN_WAPAS)      return parse_return();
    if (current.type == TOKEN_KAAM)       return parse_function();
    if (current.type == TOKEN_RACHNA)    return parse_rachna_definition();
    if (current.type == TOKEN_MAVI)       return parse_assignment();
    if (current.type == TOKEN_IMPORT)     return parse_import();
    /* samay */
    if (current.type == TOKEN_SAMAY) {
        advance();
        skip_separators();
        return new_command_node("samay", NULL);
    }

    /* rukja */
    if (current.type == TOKEN_RUKJA) {
        advance();
        skip_separators();
        return new_break_node();
    }

    /* chalu */
    if (current.type == TOKEN_CHALU) {
        advance();
        skip_separators();
        return new_continue_node();
    }

/* bolo <expr> */
if (current.type == TOKEN_BOLO) {
    advance();
    skip_separators();

    node = parse_expression();

    if (!node) {
        /* expression layer already handled error */
        return NULL;
    }

    skip_separators();
    return new_command_node("bolo", node);
}

/* AI commands */
if (current.type == TOKEN_SAKHI ||
    current.type == TOKEN_NIYU  ||
    current.type == TOKEN_MIRA  ||
    current.type == TOKEN_KAVYA ||
    current.type == TOKEN_SHIRI) {

    char cmd[32];
    int len = current.length;

    if (len >= (int)sizeof(cmd))
        len = (int)sizeof(cmd) - 1;

    strncpy(cmd, current.start, len);
    cmd[len] = '\0';   /* ✅ VERY IMPORTANT */

    advance();
    skip_separators();

    ASTNode *arg = parse_value();
    if (!arg) return NULL;

    skip_separators();
    return new_command_node(cmd, arg);
}









if (current.type == TOKEN_IDENTIFIER) {

    Token name_token = current;

    advance();
    skip_separators();

    /* CASE 1: UPDATE (x = ...) */
    if (current.type == TOKEN_EQUAL) {

        advance();
        skip_separators();

        ASTNode *val = parse_expression();
        if (!val) return NULL;

        char name[128];
        strncpy(name, name_token.start, name_token.length);
        name[name_token.length] = '\0';

        skip_separators();
        return new_update_node(name, val);
    }

    /* CASE 2: FUNCTION CALL (test(...)) */
    if (current.type == TOKEN_LEFT_PAREN) {

        /* rollback */
        current = name_token;

        ASTNode *call = parse_expression();
        if (!call) return NULL;

        skip_separators();
        return call;
    }

    /* CASE 3: VARIABLE */
    current = name_token;

    ASTNode *expr = parse_expression();
    if (!expr) return NULL;

    skip_separators();
    return expr;
}







/* raw expression */
 if (current.type == TOKEN_NUMBER ||
    current.type == TOKEN_STRING ||
    current.type == TOKEN_TRUE   ||
    current.type == TOKEN_FALSE  ||
    current.type == TOKEN_LEFT_PAREN ||
    current.type == TOKEN_NAHI   ||
    current.type == TOKEN_LEFT_BRACKET ||
    current.type == TOKEN_PLUS ||
    current.type == TOKEN_MINUS ||
    current.type == TOKEN_STAR ||
    current.type == TOKEN_SLASH)
{

    node = parse_expression();

    if (!node) {
        if (!error_reported) {
            parser_error(
                current,
                E_PARSE_INVALID_TOKEN,
                "expr",
                "expression samajh nahi aaya",
                "example: 5 + 2"
            );
        }
        return NULL;
    }

    /*  FIXED CHECK */
    if (current.type != TOKEN_NEWLINE &&
        current.type != TOKEN_EOF &&
        current.type != TOKEN_RIGHT_BRACE &&
        current.type != TOKEN_RIGHT_PAREN)
    {
        char example[64];
        shriji_build_example(current_source, example, sizeof(example));

        shriji_error_at(
            current,
            E_PARSE_OPERATOR_CHAIN,
            "expr",
            "Expression ke baad extra tokens bach gaye hain",
            example
        );

        return NULL;
    }

    return node;
}







    /* '=' can never start a statement */
    if (current.type == TOKEN_EQUAL) {

char example[64];
shriji_build_example(current_source, example, sizeof(example));

shriji_error_at(
            current,
            E_PARSE_INVALID_TOKEN,
            "statement",
            "Assignment yahin se start nahi ho sakta. Variable aur '=' ek hi line me likho.",
            example
        );

        return NULL;
    }









/* FINAL FALLBACK — ONLY IF NO ERROR ALREADY REPORTED */
if (!error_reported) {
    parser_error(
        current,
        E_PARSE_INVALID_TOKEN,
        "statement",
        "statement samajh nahi aaya",
        "use: bolo / mavi / agar"
    );
}

return NULL;
}






/*──────────────────────────────────────────────
 | BLOCK { ... }
 *──────────────────────────────────────────────*/
static ASTNode *parse_block(void)
{
    if (!expect(TOKEN_LEFT_BRACE, E_PARSE_BRACKET_MISSING, "{", "missing {", "{ ... }"))
        return NULL;

    ASTNode **stmts = NULL;
    int count = 0;

    while (current.type != TOKEN_RIGHT_BRACE &&
           current.type != TOKEN_EOF)
    {

//  FORCE SKIP INVALID TOKENS (CRITICAL FIX)
if (current.type == TOKEN_ERROR) {
    advance();
    continue;
}

        skip_separators();

        if (current.type == TOKEN_RIGHT_BRACE)
            break;

        ASTNode *s = parse_statement();

        if (!s) {
            printf("DEBUG: statement NULL at block\n");

            while (current.type != TOKEN_NEWLINE &&
                   current.type != TOKEN_EOF &&
                   current.type != TOKEN_RIGHT_BRACE)
            {
                advance();
            }

            skip_separators();
            continue;
        }

        ASTNode **tmp = realloc(stmts, sizeof(ASTNode*) * (count + 1));
        if (!tmp) {
            free(stmts);
            return NULL;
        }

        stmts = tmp;
        stmts[count++] = s;

        skip_separators();
    }

    if (!expect(TOKEN_RIGHT_BRACE, E_PARSE_BRACKET_MISSING, "}", "missing }", "{ ... }")) {
        free(stmts);
        return NULL;
    }

    skip_separators();

    return new_block_node(stmts, count);
}









/*──────────────────────────────────────────────
 | IF / ELSE
 *──────────────────────────────────────────────*/
static ASTNode *parse_if(void)
{
    advance(); /* consume agar */

    ASTNode *cond = parse_expression();
    if (!cond) return NULL;

    skip_separators();

    ASTNode *then_block = parse_block();
    if (!then_block) return NULL;

    skip_separators();

    ASTNode *node = new_if_node(cond, then_block);

    if (current.type == TOKEN_WARNA) {
        advance(); /* consume warna */
        skip_separators();

        node->else_block = parse_block();
        if (!node->else_block) return NULL;

        skip_separators();
    }

    return node;
}








/*──────────────────────────────────────────────
 | WHILE
 *──────────────────────────────────────────────*/
static ASTNode *parse_while(void)
{
    advance(); /* consume jabtak */

    ASTNode *cond = parse_expression();
    if (!cond) return NULL;

    skip_separators();

    ASTNode *body = parse_block();
    if (!body) return NULL;

    return new_while_node(cond, body);
}









/*──────────────────────────────────────────────
 | IMPORT "file.sri"
 *──────────────────────────────────────────────*/
static ASTNode *parse_import(void)
{
    advance(); /* consume import */
    skip_separators();

    if (current.type != TOKEN_STRING) {
char example[64];
shriji_build_example(current_source, example, sizeof(example));

shriji_error_at(
            current,
            E_IMPORT_PATH_INVALID,
            "import",
            "expected string path",
            example
        );

        return NULL;
    }

    char mod[256];
    int len = current.length;

    const char *p = current.start;
    if (len >= 2 && p[0] == '"' && p[len - 1] == '"') {
        p++;
        len -= 2;
    }

    if (len >= (int)sizeof(mod)) len = (int)sizeof(mod) - 1;
    strncpy(mod, p, len);
    mod[len] = 0;

    advance(); /* consume string */

    skip_separators();
    return new_import_node(mod);
}










/*──────────────────────────────────────────────
 | RETURN / WAPAS
 *──────────────────────────────────────────────*/
static ASTNode *parse_return(void)
{
    advance(); /* consume wapas */

    if (current.type == TOKEN_NEWLINE ||
        current.type == TOKEN_RIGHT_BRACE ||
        current.type == TOKEN_EOF) {

        skip_separators();
        return new_return_node(NULL);
    }

       ASTNode *val = parse_value();
    if (!val) return NULL;

    skip_separators();
    return new_return_node(val);
}










/*──────────────────────────────────────────────
 | FUNCTION / KAAM
 *──────────────────────────────────────────────*/
static ASTNode *parse_function(void)
{
    advance(); /* consume kaam */
    skip_separators();

    if (current.type != TOKEN_IDENTIFIER) {
        shriji_error_at(current, E_FUNCTION_NAME_INVALID, "kaam", "expected function name",
                        "kaam add(a, b) { ... }");
        return NULL;
    }

    char fname[128];
    strncpy(fname, current.start, current.length);
    fname[current.length] = 0;
    advance(); /* consume function name */

    skip_separators();

    if (!expect(TOKEN_LEFT_PAREN, E_PARSE_02, "(", "missing (",
                "kaam add(a, b) { ... }"))
        return NULL;

    ASTNode **params = NULL;
    int param_count = 0;

    skip_separators();

    if (current.type != TOKEN_RIGHT_PAREN) {
        while (1) {

            if (current.type != TOKEN_IDENTIFIER) {
                shriji_error_at(current, E_FUNCTION_PARAM_INVALID, "params",
                                "expected parameter name",
                                "kaam add(a, b) { ... }");
                free(params);
                return NULL;
            }

            char p[128];
            strncpy(p, current.start, current.length);
            p[current.length] = 0;
            advance();

            ASTNode **tmp = realloc(params, sizeof(ASTNode*) * (param_count + 1));
            if (!tmp) {
                free(params);
                return NULL;
            }
            params = tmp;

            params[param_count++] = new_identifier_node(p);

            skip_separators();

            if (current.type == TOKEN_COMMA) {
                advance();
                skip_separators();
            }

            if (current.type == TOKEN_RIGHT_PAREN)
                break;
        }
    }

    if (!expect(TOKEN_RIGHT_PAREN, E_PARSE_BRACKET_MISSING, ")", "missing )",
                "kaam add(a, b) { ... }")) {
        free(params);
        return NULL;
    }

    skip_separators();

    ASTNode *body = parse_block();
    if (!body) return NULL;

    skip_separators();
    return new_function_node(fname, params, param_count, body);
}

static ASTNode *parse_rachna_definition(void)
{
    advance(); // consume 'rachna'
    skip_separators();

    /* name */
    if (current.type != TOKEN_IDENTIFIER) {
        parser_error(current, E_PARSE_INVALID_TOKEN,
                     "rachna",
                     "name expected",
                     "example: rachna add() { }");
        return NULL;
    }

    char name[128];
    strncpy(name, current.start, current.length);
    name[current.length] = '\0';

    advance();
    skip_separators();

    /* '(' */
    if (!expect(TOKEN_LEFT_PAREN,
                E_PARSE_INVALID_TOKEN,
                "rachna",
                "'(' expected",
                "example: rachna add() { }"))
        return NULL;

    skip_separators();

    /* params */
    ASTNode **params = NULL;
    int param_count = 0;

    if (current.type != TOKEN_RIGHT_PAREN) {
        while (1) {

            if (current.type != TOKEN_IDENTIFIER) {
                parser_error(current, E_PARSE_INVALID_TOKEN,
                             "rachna",
                             "parameter name expected",
                             "example: rachna add(a,b) { }");
                return NULL;
            }

            char p[128];
            strncpy(p, current.start, current.length);
            p[current.length] = '\0';

            ASTNode **tmp = realloc(params, sizeof(ASTNode*) * (param_count + 1));
            if (!tmp) return NULL;
            params = tmp;

            params[param_count++] = new_identifier_node(p);

            advance();
            skip_separators();

            if (current.type == TOKEN_COMMA) {
                advance();
                skip_separators();
                continue;
            }

            break;
        }
    }

    if (!expect(TOKEN_RIGHT_PAREN,
                E_PARSE_INVALID_TOKEN,
                "rachna",
                "')' expected",
                "example: rachna add(a,b) { }"))
        return NULL;

    skip_separators();

    /* body */
    ASTNode *body = parse_block();
    if (!body) return NULL;

    skip_separators();

    return new_rachna_node(name, params, param_count, body);
}

















/*──────────────────────────────────────────────
 | IDENTIFIER / UPDATE
 *──────────────────────────────────────────────*/

static ASTNode *parse_value(void)
{
    return parse_comparison();
}












/*──────────────────────────────────────────────
  ASSIGNMENT (MAVI)
──────────────────────────────────────────────*/
static ASTNode *parse_assignment(void)
{
    advance(); /* consume 'mavi' */
    skip_separators();

    if (current.type != TOKEN_IDENTIFIER) {
        shriji_error_at(
            current,
            E_ASSIGN_01,
            "mavi",
            "Declaration samajh li, par variable ka naam clear nahi hua.",
            "mavi x  ya  mavi x = 10"
        );
        return NULL;
    }

    char name[128];
    strncpy(name, current.start, current.length);
    name[current.length] = 0;
    advance();

    skip_separators();

    if (current.type == TOKEN_EQUAL) {
        advance();
        skip_separators();

        if (current.type == TOKEN_NEWLINE ||
            current.type == TOKEN_EOF ||
            current.type == TOKEN_RIGHT_BRACE) {

            shriji_error_at(
                current,
                E_ASSIGN_02,
                "assignment",
                "'=' ke baad value missing hai",
                "example: mavi x = 10"
            );
            return NULL;
        }

        ASTNode *val = parse_value();
        if (!val) return NULL;

        skip_separators();
        return new_assignment_node(name, val);
    }

    return new_assignment_node(name, NULL);
}






static ASTNode *parse_logical_and(void)
{
    ASTNode *node = parse_comparison();

    while (current.type == TOKEN_AND)
    {
        char op = '&';
        advance();

        ASTNode *right = parse_comparison();

        node = new_binary_node(op, node, right);
    }

    return node;
}


static ASTNode *parse_logical_or(void)
{
    ASTNode *node = parse_logical_and();

    while (current.type == TOKEN_OR)
    {
        char op = '|';
        advance();

        ASTNode *right = parse_logical_and();
        node = new_binary_node(op, node, right);
    }

    return node;
}

/*──────────────────────────────────────────────
 | EXPRESSIONS
 | precedence:
 |   primary -> unary -> term(* / %) -> expression(+ -) -> comparison
 *──────────────────────────────────────────────*/

static ASTNode *parse_comparison(void)
{
    ASTNode *n = parse_term();
    if (!n) return NULL;

    while (current.type == TOKEN_GT  ||
           current.type == TOKEN_LT  ||
           current.type == TOKEN_GTE ||
           current.type == TOKEN_LTE ||
           current.type == TOKEN_EQEQ ||
           current.type == TOKEN_NEQ)
    {
        char op =
            (current.type == TOKEN_GT)    ? '>' :
            (current.type == TOKEN_LT)    ? '<' :
            (current.type == TOKEN_GTE)   ? 'G' :
            (current.type == TOKEN_LTE)   ? 'L' :
            (current.type == TOKEN_EQEQ)  ? '=' :
            (current.type == TOKEN_NEQ)   ? '!' :
            '?';

        advance();

        ASTNode *right = parse_term();
        if (!right) return NULL;

        n = new_binary_node(op, n, right);
    }

    return n;
}


static ASTNode *parse_expression(void)
{
   ASTNode *n = parse_logical_or();
    if (!n) return NULL;




/*──────────────────────────────────────────────
 | OPERATOR START DETECTION
 *──────────────────────────────────────────────*/
if (prev_token.type == TOKEN_NEWLINE ||
    prev_token.type == TOKEN_EOF)
{
    if (current.type == TOKEN_PLUS ||
        current.type == TOKEN_MINUS ||
        current.type == TOKEN_STAR ||
        current.type == TOKEN_SLASH)
    {

        char example[64];
        shriji_build_example(current_source, example, sizeof(example));

        shriji_error_at(
            current,
            E_PARSE_OPERATOR_START,
            "expr",
            "Expression operator se start nahi ho sakta",
            example
        );

        parser_sync();
        return NULL;
    }
}



while (current.type == TOKEN_PLUS  ||
       current.type == TOKEN_MINUS)
{
    char op;

    if (current.type == TOKEN_PLUS) op = '+';
    else op = '-';

    advance();







/* operator classification */

if (current.type == TOKEN_PLUS  ||
    current.type == TOKEN_MINUS ||
    current.type == TOKEN_STAR  ||
    current.type == TOKEN_SLASH ||
    current.type == TOKEN_AND   ||
    current.type == TOKEN_OR)
{
    ShrijiErrorCode code = E_PARSE_DOUBLE_OPERATOR;

    if (prev_token.type == current.type)
    {
        code = E_PARSE_DOUBLE_OPERATOR;
    }
    else
    {
        code = E_PARSE_OPERATOR_CHAIN;
    }






char example[64];
shriji_build_example(current_source, example, sizeof(example));

shriji_error_at(
    current,
    code,
    "expr",
    "Operator chain detect hui hai.",
    example
);

return NULL;

}












/* operator end detection */

if (current.type == TOKEN_NEWLINE ||
    current.type == TOKEN_EOF ||
    current.type == TOKEN_RIGHT_BRACE)
{
char example[64];
shriji_build_example(current_source, example, sizeof(example));

shriji_error_at(
    current,
    E_PARSE_OPERATOR_END,
    "expr",
    "Operator ke baad value missing hai",
    example
);

    return NULL;
}

        ASTNode *right = parse_term();
        if (!right) return NULL;

        n = new_binary_node(op, n, right);
    }

    return n;
}


static ASTNode *parse_term(void)
{
    ASTNode *n = parse_unary();
    if (!n) return NULL;

    while (current.type == TOKEN_STAR ||
           current.type == TOKEN_SLASH ||
           current.type == TOKEN_MOD)
    {
        char op =
            (current.type == TOKEN_STAR)  ? '*' :
            (current.type == TOKEN_SLASH) ? '/' : '%';

        advance();










/* double operator detection */

if (current.type == TOKEN_STAR ||
    current.type == TOKEN_SLASH ||
    current.type == TOKEN_MOD)
{
    char example[64];
    shriji_build_example(current_source, example, sizeof(example));

    shriji_error_at(
        current,
        E_PARSE_DOUBLE_OPERATOR,
        "expr",
        "Do operators ek saath aa gaye hain.",
        example
    );

    return NULL;
}









/* operator end detection for * / % */

if (current.type == TOKEN_NEWLINE ||
    current.type == TOKEN_EOF ||
    current.type == TOKEN_RIGHT_BRACE)
{

char example[64];
shriji_build_example(current_source, example, sizeof(example));

shriji_error_at(
    current,
    E_PARSE_OPERATOR_END,
    "expr",
    "Expression operator par end ho gaya",
    example
);

skip_separators();
return NULL;

}

        ASTNode *right = parse_unary();
        if (!right) return NULL;

        n = new_binary_node(op, n, right);
    }

    return n;
}









/*──────────────────────────────────────────────
 | UNARY
 |   - prefix ++x / --x
 |   - unary + / -
 |   - nahi (logical not)
 *──────────────────────────────────────────────*/
static ASTNode *parse_unary(void)
{
    /* prefix ++ */
    if (current.type == TOKEN_PLUSPLUS) {
        advance();

        ASTNode *expr = parse_unary();
        if (!expr) return NULL;

        Token fake = {0};
        fake.start = "+";
        fake.length = 1;

        return new_unary_node(fake, expr);
    }

    /* prefix -- */
    if (current.type == TOKEN_MINUSMINUS) {
        advance();

        ASTNode *expr = parse_unary();
        if (!expr) return NULL;

        Token fake = {0};
        fake.start = "-";
        fake.length = 1;

        return new_unary_node(fake, expr);
    }

    /* unary + / - */
    if (current.type == TOKEN_PLUS || current.type == TOKEN_MINUS) {

        Token op = current;
        advance();

        /* prevent invalid operator chain */
        if (current.type == TOKEN_PLUS ||
            current.type == TOKEN_MINUS ||
            current.type == TOKEN_STAR ||
            current.type == TOKEN_SLASH)
        {
            shriji_error_at(
                current,
                E_PARSE_OPERATOR_CHAIN,
                "expr",
                "operator chain galat hai",
                "example: +5 ya -3"
            );
            return NULL;
        }

        ASTNode *expr = parse_primary();
        if (!expr) return NULL;

        return new_unary_node(op, expr);
    }

    /* nahi (logical not) */
    if (current.type == TOKEN_NAHI) {
        advance();

        ASTNode *expr = parse_unary();
        if (!expr) return NULL;

        return new_not_node(expr);
    }

    /* default primary */
    ASTNode *node = parse_primary();
    if (!node) return NULL;

    /* detect missing operator: 6(5+9) */
    if (current.type == TOKEN_LEFT_PAREN)
    {
        char example[64];
        shriji_build_example(current_source, example, sizeof(example));

        shriji_error_at(
            current,
            E_PARSE_MISSING_OPERATOR,
            "expr",
            "Do values ke beech operator missing hai",
            example
        );

        return NULL;
    }

    return node;
}







/*──────────────────────────────────────────────
 | LIST LITERAL: [a, b, c]
 *──────────────────────────────────────────────*/
static ASTNode *parse_list_literal(void)
{
    if (!expect(TOKEN_LEFT_BRACKET, E_LIST_SYNTAX_ERROR, "[", "missing [", "[1,2,3]"))
        return NULL;

    skip_separators();

    ASTNode **elements = NULL;
    int count = 0;

    /* empty list: [] */
    if (current.type != TOKEN_RIGHT_BRACKET) {

        while (1) {

            ASTNode *elem = parse_expression();
            if (!elem) {
                free(elements);
                return NULL;
            }

            ASTNode **tmp = realloc(elements, sizeof(ASTNode*) * (count + 1));
            if (!tmp) {
                free(elements);
                return NULL;
            }
            elements = tmp;
            elements[count++] = elem;

            skip_separators();

            if (current.type == TOKEN_RIGHT_BRACKET)
                break;

            if (!expect(TOKEN_COMMA, E_LIST_SYNTAX_ERROR, ",", "expected ','", "[1, 2, 3]")) {
                free(elements);
                return NULL;
            }

            skip_separators();
        }
    }

    if (!expect(TOKEN_RIGHT_BRACKET, E_PARSE_BRACKET_MISSING, "]", "missing ]", "[1,2,3]")) {
        free(elements);
        return NULL;
    }

    skip_separators();
    return new_list_node(elements, count);
}












/*──────────────────────────────────────────────
 | POSTFIX:
 |   - indexing: a[0]
 |   - postfix inc/dec: x++ , x--
 *──────────────────────────────────────────────*/
static ASTNode *parse_postfix(ASTNode *base)
{
    ASTNode *node = base;

    while (1) {

        /* ───── indexing: a[expr] ───── */
        if (current.type == TOKEN_LEFT_BRACKET) {

            advance(); /* consume '[' */
            skip_separators();

            ASTNode *idx = parse_expression();
            if (!idx) return NULL;

            skip_separators();

            if (!expect(TOKEN_RIGHT_BRACKET, E_PARSE_BRACKET_MISSING,
                        "]", "missing ]", "a[0]"))
                return NULL;

            skip_separators();

            node = new_index_node(node, idx);
            continue;
        }

        break;
    }

    return node;
}










/*──────────────────────────────────────────────
 | PRIMARY
 *──────────────────────────────────────────────*/
static ASTNode *parse_primary(void)
{
if (current.type == TOKEN_EOF ||
        current.type == TOKEN_NEWLINE)
    {
        char example[64];
        shriji_build_example(current_source, example, sizeof(example));

        shriji_error_at(
            current,
            E_PARSE_OPERATOR_END,
            "expr",
            "Expression incomplete hai",
            example
        );

        return NULL;
    }


    if (current.type == TOKEN_TRUE) {
         advance();
           return new_bool_node(1);
        }

   if (current.type == TOKEN_FALSE) {
        advance();
           return new_bool_node(0);
        }








if (current.type == TOKEN_NUMBER) {

    char temp[64];
    int len = current.length;

    if (len >= 63)
        len = 63;

    strncpy(temp, current.start, len);
    temp[len] = '\0';

    double v = atof(temp);

    advance();

    return new_number_node(v);
}









if (current.type == TOKEN_STRING) {
    char s[256];
    int len = current.length;

    const char *p = current.start;
    if (len >= 2 && p[0] == '"' && p[len - 1] == '"') {
        p++;
        len -= 2;
    }

    unescape_into(s, sizeof(s), p, len);

    advance();
    return new_string_node(s);
}


    /* list literal: [ ... ] */
    if (current.type == TOKEN_LEFT_BRACKET) {
        ASTNode *lst = parse_list_literal();
        if (!lst) return NULL;
        return parse_postfix(lst);
    }





/* dict literal: { ... } */
if (current.type == TOKEN_LEFT_BRACE) {
    ASTNode *d = parse_dict_literal();
    if (!d) return NULL;
    ASTNode *node = parse_postfix(d);
if (!node) return NULL;
return node;
}

    if (current.type == TOKEN_IDENTIFIER) {

        char name[128];
        strncpy(name, current.start, current.length);
        name[current.length] = 0;
        advance();






/* function call: name(...) */
if (current.type == TOKEN_LEFT_PAREN) {

    advance(); /* consume '(' */
    skip_separators();

    ASTNode **args = NULL;
    int arg_count = 0;

    if (current.type != TOKEN_RIGHT_PAREN) {
        while (1) {

            ASTNode *arg = parse_expression();
            if (!arg) {
                free(args);
                return NULL;
            }

            ASTNode **tmp = realloc(args, sizeof(ASTNode*) * (arg_count + 1));
            if (!tmp) {
                free(args);
                return NULL;
            }
            args = tmp;
            args[arg_count++] = arg;

            skip_separators();

            if (current.type == TOKEN_RIGHT_PAREN)
                break;

            if (current.type != TOKEN_COMMA) {
                shriji_error_at(
                    current,
                    E_FUNCTION_PARAM_INVALID,
                    "call",
                    "expected ',' between arguments",
                    "example: add(5, 6)"
                );
                free(args);
                return NULL;
            }

            advance();
            skip_separators();
        }
    }

    if (!expect(TOKEN_RIGHT_PAREN, E_PARSE_01, ")", "missing ')'",
                "example: add(10, 20)")) {
        free(args);
        return NULL;
    }

    skip_separators();

    ASTNode *node = new_call_node(name, args, arg_count);
    node = parse_postfix(node);
    if (!node) return NULL;

    return node;
}


        ASTNode *node = new_identifier_node(name);
node = parse_postfix(node);
if (!node) return NULL;
return node;
    }

    if (current.type == TOKEN_LEFT_PAREN) {
        advance();
        ASTNode *n = parse_comparison();
        if (!n) return NULL;

        if (!expect(TOKEN_RIGHT_PAREN, E_PARSE_UNMATCHED_PAREN, ")", "missing )", "(expr)"))
            return NULL;

        ASTNode *node = parse_postfix(n);
if (!node) return NULL;
return node;
    }

    shriji_error_at(current, E_PARSE_INVALID_TOKEN, "expr", "bad token", "check syntax");
    return NULL;
}










/*──────────────────────────────────────────────
 | DICT / MAP LITERAL: { "a": 1, "b": 2 }
 *──────────────────────────────────────────────*/
static ASTNode *parse_dict_literal(void)
{
    if (!expect(TOKEN_LEFT_BRACE, E_PARSE_BRACKET_MISSING, "{", "missing {", "{\"a\":1}"))
        return NULL;

    skip_separators();

    ASTNode **keys = NULL;
    ASTNode **vals = NULL;
    int count = 0;

    /* empty dict: {} */
    if (current.type != TOKEN_RIGHT_BRACE) {

        while (1) {

            /* key must be string for v1 */
            if (current.type != TOKEN_STRING) {
                shriji_error_at(
                    current,
                    E_DICT_KEY_INVALID,
                    "dict",
                    "dict key must be string",
                    "example: {\"a\": 1, \"b\": 2}"
                );
                free(keys);
                free(vals);
                return NULL;
            }

            /* parse key as string node */
            ASTNode *k = parse_primary();
            if (!k) {
                free(keys);
                free(vals);
                return NULL;
            }

            skip_separators();

            if (!expect(TOKEN_COLON, E_DICT_SYNTAX_ERROR, ":", "missing ':'",
                        "example: {\"a\": 1}")) {
                free(keys);
                free(vals);
                return NULL;
            }

            skip_separators();

            /* parse value expression */
            ASTNode *v = parse_expression();
            if (!v) {
                free(keys);
                free(vals);
                return NULL;
            }








/* push pair */
ASTNode **tmpk = realloc(keys, sizeof(ASTNode*) * (count + 1));
if (!tmpk) {
    free(keys);
    free(vals);
    return NULL;
}
keys = tmpk;

ASTNode **tmpv = realloc(vals, sizeof(ASTNode*) * (count + 1));
if (!tmpv) {
    free(vals);
    return NULL;
}
vals = tmpv;

keys[count] = k;
vals[count] = v;
count++;

skip_separators();






/* end? */
if (current.type == TOKEN_RIGHT_BRACE)
    break;

if (!expect(TOKEN_COMMA, E_DICT_SYNTAX_ERROR, ",", "expected ','",
            "example: {\"a\": 1, \"b\": 2}")) {
    free(keys);
    free(vals);
    return NULL;
}

        }
    }

    if (!expect(TOKEN_RIGHT_BRACE, E_PARSE_BRACKET_MISSING, "}", "missing '}'",
                "example: {\"a\": 1, \"b\": 2}")) {
        free(keys);
        free(vals);
        return NULL;
    }

    skip_separators();
    return new_dict_node(keys, vals, count);
}

nd{verbatim}
\section*{File: ./src/fix_engine.c}
egin{verbatim}
/* SHREE_RADHIKA_RANI */

#include <stdio.h>
#include <string.h>
#include <ctype.h>

#include "../include/parser.h"
#include "../include/error.h"
#include "../include/fix_engine.h"
#include "../include/fix_rules.h"

#define MAX_WORKING_BUFFER 512


/*──────────────────────────────────────────────
   PARSE ONLY (STRICT MODE)
──────────────────────────────────────────────*/

ASTNode *parse_once(const char *input)
{
    error_reported = 0;
    shriji_set_error_mode(ERROR_MODE_IMMEDIATE);

    return parse_program(input);
}


/*──────────────────────────────────────────────
   MAIN EXECUTION (SAFE VERSION)
──────────────────────────────────────────────*/

ASTNode *language_execute_with_fix(
    const char *input,
    int *was_fixed,
    int *penalty_out
)
{
    if (!input)
        return NULL;

    if (was_fixed)
        *was_fixed = 0;

    if (penalty_out)
        *penalty_out = 0;

    error_reported = 0;
    shriji_set_error_mode(ERROR_MODE_IMMEDIATE);

    /* FIRST PARSE */
    ASTNode *root = parse_program(input);

    /* ERROR CASE */
    if (!root || error_reported)
    {
        char fixed[MAX_WORKING_BUFFER];
        memset(fixed, 0, sizeof(fixed));

        FixType type = fix_apply(input, fixed);

        /* SAFE FIX */
        if (type == FIX_SAFE)
        {
            if (was_fixed)
                *was_fixed = 1;

            error_reported = 0;

            ASTNode *new_root = parse_program(fixed);

            if (!new_root || error_reported)
                return NULL;

            return new_root;
        }

        /* no safe fix */
        return NULL;
    }

    return root;
}


/* ================================
   HELPER: invalid token
================================ */

static int has_invalid_char(const char *input)
{
    for (int i = 0; input[i]; i++)
    {
        char c = input[i];

        if (isdigit(c) || strchr("+-*/(). ", c))
            continue;

        return 1;
    }
    return 0;
}


/* ================================
   HELPER: operator chain length (FIXED)
================================ */

static int operator_chain_len(const char *input, int pos)
{
    int len = 1;

    while (input[pos + len] != '\0' &&
           strchr("+-*/", input[pos + len]))
    {
        len++;
    }

    return len;
}


/* ================================
   MAIN FIX ENGINE
================================ */

FixType fix_apply(const char *input, char *output)
{
    if (!input || !*input)
        return FIX_NONE;

    if (has_invalid_char(input))
        return FIX_REJECT;

    int i, j = 0;
    int changed = 0;

    for (i = 0; input[i]; i++)
    {
        /* implicit multiplication */
        if (isdigit(input[i]) && input[i+1] == '(')
        {
            output[j++] = input[i];
            output[j++] = '*';
            changed = 1;
            continue;
        }

        if (input[i] == ')' && isdigit(input[i+1]))
        {
            output[j++] = input[i];
            output[j++] = '*';
            changed = 1;
            continue;
        }

        /* operator chain */
        if (strchr("+-*/", input[i]))
        {
            int len = operator_chain_len(input, i);

            if (len == 2 && input[i] == input[i+1])
            {
                output[j++] = input[i];
                i++;
                changed = 1;
                continue;
            }

            if (len > 2 || (len == 2 && input[i] != input[i+1]))
            {
                return FIX_DOUBTFUL;
            }
        }

        output[j++] = input[i];
    }

    output[j] = '\0';

    if (changed)
        return FIX_SAFE;

    return FIX_NONE;
}
nd{verbatim}
\section*{File: ./src/nirman.c}
egin{verbatim}
#include "../include/nirman.h"
#include <string.h>

/* ===============================
   GLOBAL CONTEXT
   =============================== */
NirmanContext NIRMAN_CTX = {0};

/* ===============================
   LIFECYCLE
   =============================== */

void nirman_start(void)
{
    NIRMAN_CTX.active = 1;

    /* reset flow */
    NIRMAN_CTX.flow_stage = 0;
    NIRMAN_CTX.goal[0] = '\0';
}

void nirman_stop(void)
{
    NIRMAN_CTX.active = 0;

    /* reset flow */
    NIRMAN_CTX.flow_stage = 0;
    NIRMAN_CTX.goal[0] = '\0';
}

int nirman_is_active(void)
{
    return NIRMAN_CTX.active;
}

/* ===============================
   INTERNAL HELPERS
   =============================== */

static void to_lower(char *s)
{
    for (int i = 0; s[i]; i++)
    {
        if (s[i] >= 'A' && s[i] <= 'Z')
            s[i] += 32;
    }
}

/* ===============================
   INTENT ENGINE
   =============================== */

int nirman_detect_intent(const char *input)
{
    if (!input)
        return 0;

    char temp[256];
    strncpy(temp, input, sizeof(temp) - 1);
    temp[255] = '\0';

    to_lower(temp);

    /* 1 = TALK */
    if (strstr(temp, "kaise") ||
        strstr(temp, "hello") ||
        strstr(temp, "hi") ||
        strstr(temp, "hy") ||
        strstr(temp, "radhe"))
    {
        return 1;
    }

    /* 2 = BUILD */
    if (strstr(temp, "bana") ||
        strstr(temp, "create") ||
        strstr(temp, "app") ||
        strstr(temp, "software"))
    {
        return 2;
    }

    /* 3 = ERROR / HELP */
    if (strstr(temp, "error") ||
        strstr(temp, "problem") ||
        strstr(temp, "issue"))
    {
        return 3;
    }

    return 0;
}
nd{verbatim}
\section*{File: ./src/fix_rules.c}
egin{verbatim}
#include "../include/fix_rules.h"
#include "../include/fix_interactive.h"

#include <string.h>
#include <stdio.h>
#include <ctype.h>

/*──────────────────────────────────────────────
    FIX FUNCTIONS
──────────────────────────────────────────────*/

/* FIX: double operator (6++6 → 6+6) */
static int fix_double_operator(
    char *buffer,
    size_t buffer_size,
    const ShrijiErrorInfo *err)
{
    (void)buffer_size;
    (void)err;

    int changed = 0;

    for (int i = 0; buffer[i]; i++)
    {
        if ((buffer[i] == '+' || buffer[i] == '-' ||
             buffer[i] == '*' || buffer[i] == '/')
            && buffer[i] == buffer[i+1])
        {
            memmove(&buffer[i], &buffer[i+1], strlen(&buffer[i+1]) + 1);
            changed = 1;
            i--;
        }
    }

    return changed;
}

/* FIX: missing operator (6(7+7) → 6*(7+7)) */
static int fix_missing_operator(
    char *buffer,
    size_t buffer_size,
    const ShrijiErrorInfo *err)
{
    (void)err;

    for (int i = 0; buffer[i]; i++)
    {
        if (isdigit(buffer[i]) && buffer[i+1] == '(')
        {
            if (strlen(buffer) + 2 >= buffer_size)
                return 0;

            memmove(&buffer[i+2], &buffer[i+1], strlen(&buffer[i+1]) + 1);
            buffer[i+1] = '*';

            return 1;
        }
    }

    return 0;
}

/* FIX: missing value after '=' */
static int fix_missing_value_after_equal(
    char *buffer,
    size_t buffer_size,
    const ShrijiErrorInfo *err)
{
    (void)err;

    size_t len = strlen(buffer);

    if (len + 3 >= buffer_size)
        return 0;

    while (len > 0 && isspace((unsigned char)buffer[len - 1]))
        len--;

    if (len > 0 && buffer[len - 1] == '=')
    {
        strcat(buffer, " 0");
        return 1;
    }

    return 0;
}

/*──────────────────────────────────────────────
    FIX RULE TABLE
──────────────────────────────────────────────*/

static FixRule fix_rules[] = {

    /* deterministic: always correct */
    {
        .error_code = E_PARSE_MISSING_OPERATOR,
        .category = FIXCAT_MISSING_OPERATOR,
        .safety = FIX_SAFE_DETERMINISTIC,
        .confidence_penalty = 5,
        .max_attempt = 1,
        .reparse_required = 1,
        .apply_fix = fix_missing_operator
    },

    /* deterministic: assignment must have value */
    {
        .error_code = E_ASSIGN_02,
        .category = FIXCAT_MISSING_VALUE,
        .safety = FIX_SAFE_DETERMINISTIC,
        .confidence_penalty = 6,
        .max_attempt = 1,
        .reparse_required = 1,
        .apply_fix = fix_missing_value_after_equal
    },

    /* structural: operator cleanup (not always guaranteed) */
    {
        .error_code = E_PARSE_OPERATOR_CHAIN,
        .category = FIXCAT_EXTRA_OPERATOR,
        .safety = FIX_SAFE_STRUCTURAL,
        .confidence_penalty = 8,
        .max_attempt = 1,
        .reparse_required = 1,
        .apply_fix = fix_double_operator
    }

};

/*──────────────────────────────────────────────
    RULE LOOKUP
──────────────────────────────────────────────*/

static int fix_rule_count =
    sizeof(fix_rules) / sizeof(FixRule);

FixRule *shriji_get_rule_for_error(ShrijiErrorCode code)
{
    for (int i = 0; i < fix_rule_count; i++)
    {
        if (fix_rules[i].error_code == code)
            return &fix_rules[i];
    }

    return NULL;
}
nd{verbatim}
\section*{File: ./src/smriti.c}
egin{verbatim}
#include "../include/smriti.h"
#include <string.h>

/* ------------------------------
   INTERNAL STATE (SILENT)
   ------------------------------ */

/* continuation (existing) */
static int    has_cont = 0;
static double cont_val = 0.0;
static char   cont_op  = 0;

/* STEP-11 context */
static unsigned long turn_counter = 0;
static char last_input_buf[512];
static int  last_intent_val = 0;

/* ------------------------------
   LIFECYCLE
   ------------------------------ */
void smriti_init(void)
{
    has_cont = 0;
    cont_val = 0.0;
    cont_op  = 0;

    turn_counter = 0;
    last_input_buf[0] = '\0';
    last_intent_val = 0;
}

/* ------------------------------
   CONTINUATION (AS-IS)
   ------------------------------ */
int smriti_has_continuation(void)
{
    return has_cont;
}

void smriti_set_continuation(double value, char op)
{
    has_cont = 1;
    cont_val = value;
    cont_op  = op;
}

void smriti_clear_continuation(void)
{
    has_cont = 0;
    cont_val = 0.0;
    cont_op  = 0;
}

double smriti_get_continuation_value(void)
{
    return cont_val;
}

char smriti_get_continuation_op(void)
{
    return cont_op;
}

/* ------------------------------
   STEP-11 CONTEXT (PASSIVE)
   ------------------------------ */
void smriti_next_turn(void)
{
    turn_counter++;
}

unsigned long smriti_get_turn(void)
{
    return turn_counter;
}

void smriti_set_last_input(const char *text)
{
    if (!text)
    {
        last_input_buf[0] = '\0';
        return;
    }
    /* safe snapshot */
    strncpy(last_input_buf, text, sizeof(last_input_buf) - 1);
    last_input_buf[sizeof(last_input_buf) - 1] = '\0';
}

const char* smriti_get_last_input(void)
{
    return last_input_buf;
}

void smriti_set_last_intent(int intent)
{
    last_intent_val = intent;
}

int smriti_get_last_intent(void)
{
    return last_intent_val;
}
nd{verbatim}
\section*{File: ./src/translator.c}
egin{verbatim}
#include <string.h>
#include <stdlib.h>
#include <ctype.h>
#include "../include/translator.h"

// --------------------------------------------------------
// INTERNAL STRUCT FOR KEYWORD MAPPING
// --------------------------------------------------------
typedef struct {
    const char *world;
    const char *shriji;
} TranslateEntry;

// --------------------------------------------------------
// WORLD → SHRIJILANG KEYWORD MAP (SMART ROUTING V3)
// --------------------------------------------------------
static TranslateEntry keyword_map[] = {

    // PRINT / OUTPUT
    {"print", "aaradhya"},
    {"echo", "aaradhya"},
    {"say", "sakhi"},
    {"speak", "sakhi"},
    {"tell", "sakhi"},

    // MIRA — Mathematical Solver
    {"solve", "mira"},
    {"calculate", "mira"},
    {"compute", "mira"},

    // SHIRI — Soft Guidance
    {"guide", "shiri"},
    {"help", "shiri"},
    {"direction", "shiri"},

    // NIYU — Logic / Why AI
    {"why", "niyu"},
    {"explain", "niyu"},
    {"reason", "niyu"},

    // KAVYA — Creative Engine
    {"create", "kavya"},
    {"write", "kavya"},
    {"story", "kavya"},

    // DEVOTIONAL → DIRECT KAVYA
    {"radhe", "kavya"},
    {"hare", "kavya"},
    {"krishna", "kavya"},
    {"shyam", "kavya"},

    // VARIABLE SYSTEM
    {"var", "mavi"},
    {"let", "mavi"},
    {"int", "mavi"},
    {"float", "mavi"}
};

// --------------------------------------------------------
// GET FIRST WORD
// --------------------------------------------------------
char* get_first_word(const char *source) {
    while (*source == ' ') source++;

    const char *start = source;
    while (*source && !isspace(*source))
        source++;

    int len = source - start;
    char *word = malloc(len + 1);
    strncpy(word, start, len);
    word[len] = '\0';

    return word;
}

// --------------------------------------------------------
// MAIN KEYWORD TRANSLATION
// --------------------------------------------------------
const char* translate_keyword(const char *word) {

for (size_t i = 0; i < sizeof(keyword_map)/sizeof(keyword_map[0]); i++) {

        if (strcmp(word, keyword_map[i].world) == 0)
            return keyword_map[i].shriji;
    }

    return NULL;
}

// --------------------------------------------------------
// SMART AUTOCORRECT FOR NATURAL INPUT
// --------------------------------------------------------
const char* mira_autocorrect(const char *word) {

    // Logical → Mira
    if (strcmp(word, "prit") == 0) return "mira";

    // Devotional → Kavya
    if (strcmp(word, "radhe") == 0) return "kavya";
    if (strcmp(word, "hare") == 0) return "kavya";
    if (strcmp(word, "krishna") == 0) return "kavya";

    // Emotional pain → Sakhi
    if (strcmp(word, "mera") == 0) return "sakhi";
    if (strcmp(word, "mann") == 0) return "sakhi";
    if (strcmp(word, "toot") == 0) return "sakhi";
    if (strcmp(word, "dukhi") == 0) return "sakhi";
    if (strcmp(word, "hurt") == 0) return "sakhi";
    if (strcmp(word, "sad") == 0) return "sakhi";

    // Creative autocorrect
    if (strcmp(word, "creat") == 0) return "kavya";

    // Shiri autocorrect
    if (strcmp(word, "gide") == 0) return "shiri";

    return NULL;
}
nd{verbatim}
\section*{File: ./src/ai_router.c}
egin{verbatim}
#include <stdio.h>
#include <string.h>

#include "../include/ai_router.h"
#include "../include/ai_intent.h"

#include "../pragya/contracts/l3_request.h"
#include "../pragya/contracts/l3_response.h"
#include "../pragya/contracts/l3_intent.h"

/*
   Legacy AI Router

   NOTE:
   ShrijiLang now uses KRST + Pragya Router
   for all intelligence routing.

   This module is preserved only for
   future conversational AI layer.
*/

L3Response ai_router_dispatch(const ShrijiBridgePacket *pkt)
{
    L3Response resp;

    resp.text = NULL;
    resp.success = 0;

    if (!pkt || !pkt->text)
        return resp;

    /* For now just acknowledge input */

    resp.text = "AI router currently inactive.";
    resp.success = 1;

    return resp;
}
nd{verbatim}
\section*{File: ./src/state.c}
egin{verbatim}
#include <stdio.h>
#include "../include/state.h"

/*──────────────────────────────────────────────────────────────
 | SHRIJILANG — EXECUTION STATE ENGINE
 |
 | Purpose:
 |   - Runtime safety monitoring
 |   - Confidence / risk state management
 |   - Execution limits (anti-freeze)
 |   - Human confirmation escalation policy
 |
 | NOTE:
 |   This layer must remain deterministic.
 |   No external dependencies, no printing from core logic.
 *──────────────────────────────────────────────────────────────*/

/*──────────────────────────────────────────────
 | INTERNAL CONSTANTS (Tuning)
 *──────────────────────────────────────────────*/

/* Global execution budget:
 * Prevents accidental freeze when user writes an infinite loop.
 * Keep large enough so normal programs can run.
 */
#define SHRIJI_DEFAULT_MAX_STEPS   100000000LL  /* 10 crore */

/* Error escalation thresholds */
#define SHRIJI_WARN_THRESHOLD      3
#define SHRIJI_CRITICAL_THRESHOLD  5

/*──────────────────────────────────────────────
 | INIT
 *──────────────────────────────────────────────*/
void state_init(ExecutionState *state)
{
    if (!state) return;

    /* runtime safety status */
    state->safety = STATE_SAFE;

    /* confidence model */
    state->confidence = CONFIDENCE_HIGH;

    /* pressure model */
    state->error_pressure = 0;

    /* human control mode */
    state->human = HUMAN_AUTO_OK;

    /* execution budget */
    state->max_steps  = SHRIJI_DEFAULT_MAX_STEPS;
    state->steps_used = 0;
}

/*──────────────────────────────────────────────
 | RESET
 *──────────────────────────────────────────────*/
void state_reset(ExecutionState *state)
{
    state_init(state);
}

/*──────────────────────────────────────────────
 | ERROR HOOK
 *──────────────────────────────────────────────*/
void state_on_error(ExecutionState *state)
{
    if (!state) return;

    /* increase error pressure */
    if (state->error_pressure < 1000000)
        state->error_pressure++;

    /* stage-1: warning / risk */
    if (state->error_pressure >= SHRIJI_WARN_THRESHOLD) {
        state->confidence = CONFIDENCE_LOW;
        state->safety = STATE_RISK;
        state->human = HUMAN_CONFIRM_REQUIRED;
    }

    /* stage-2: critical stop */
    if (state->error_pressure >= SHRIJI_CRITICAL_THRESHOLD) {
        state->safety = STATE_CRITICAL;
        state->human = HUMAN_MANDATORY;
    }
}

/*──────────────────────────────────────────────
 | SUCCESS HOOK
 *──────────────────────────────────────────────*/
void state_on_success(ExecutionState *state)
{
    if (!state) return;

    /* decay pressure gradually */
    if (state->error_pressure > 0)
        state->error_pressure--;

    /* fully recover when stable again */
    if (state->error_pressure == 0) {
        state->confidence = CONFIDENCE_HIGH;
        state->safety = STATE_SAFE;
        state->human = HUMAN_AUTO_OK;
    }
}
nd{verbatim}
\section*{File: ./src/token.c}
egin{verbatim}
/*=============================================================
  SHRIJILANG — TOKENIZER (PRODUCTION STABLE VERSION)
=============================================================*/

#include <stdio.h>
#include <string.h>
#include <ctype.h>

#include "../include/token.h"

/*──────────────────────────────────────────────
  INTERNAL STATE
──────────────────────────────────────────────*/

static const char *source = NULL;
static int current = 0;
static int line = 1;
static int col  = 1;

/*──────────────────────────────────────────────
  HELPERS
──────────────────────────────────────────────*/

static char peek(void)
{
    return source[current];
}

static char peek_next(void)
{
    if (source[current] == '\0') return '\0';
    return source[current + 1];
}

static char advance(void)
{
    char c = source[current++];

    if (c == '\n') {
        line++;
        col = 1;
    } else {
        col++;
    }

    return c;
}

static int match(char expected)
{
    if (peek() != expected) return 0;
    advance();
    return 1;
}

static Token make_token(TokenType type, const char *start, int length)
{
    Token tok;
    tok.type   = type;
    tok.start  = start;
    tok.length = length;
    tok.line   = line;
    tok.col    = col - length;
    return tok;
}

/*──────────────────────────────────────────────
  SKIP WHITESPACE + COMMENTS
──────────────────────────────────────────────*/

static void skip_whitespace(void)
{
    for (;;) {
        char c = peek();

        switch (c) {

            case ' ':
            case '\r':
            case '\t':
                advance();
                break;

            case '\n':
                advance();
                return;

            case '#':
                while (peek() != '\n' && peek() != '\0') {
                    advance();
                }
                if (peek() == '\n') advance();
                break;

            default:
                return;
        }
    }
}

/*──────────────────────────────────────────────
  NUMBER
──────────────────────────────────────────────*/

static Token number(void)
{
    const char *start = &source[current - 1];

    while (isdigit(peek()))
        advance();

    if (peek() == '.' && isdigit(peek_next())) {
        advance();
        while (isdigit(peek()))
            advance();
    }

    int length = &source[current] - start;
    return make_token(TOKEN_NUMBER, start, length);
}

/*──────────────────────────────────────────────
  IDENTIFIER / KEYWORD
──────────────────────────────────────────────*/

static TokenType check_keyword(const char *start, int length)
{
    if (length == 4 && strncmp(start, "mavi", 4) == 0) return TOKEN_MAVI;
    if (length == 4 && strncmp(start, "agar", 4) == 0) return TOKEN_AGAR;
    if (length == 5 && strncmp(start, "warna", 5) == 0) return TOKEN_WARNA;
    if (length == 4 && strncmp(start, "kaam", 4) == 0) return TOKEN_KAAM;
    if (length == 6 && strncmp(start, "jabtak", 6) == 0) return TOKEN_JABTAK;
    if (length == 5 && strncmp(start, "rukja", 5) == 0) return TOKEN_RUKJA;
    if (length == 5 && strncmp(start, "chalu", 5) == 0) return TOKEN_CHALU;
    if (length == 4 && strncmp(start, "nahi", 4) == 0) return TOKEN_NAHI;
    if (length == 4 && strncmp(start, "sahi", 4) == 0) return TOKEN_TRUE;
    if (length == 5 && strncmp(start, "galat", 5) == 0) return TOKEN_FALSE;


    if (length == 4 && strncmp(start, "bolo", 4) == 0) return TOKEN_BOLO;
    if (length == 6 && strncmp(start, "rachna", 6) == 0) return TOKEN_RACHNA;
    if (length == 5 && strncmp(start, "wapas", 5) == 0) return TOKEN_WAPAS;


    if (length == 5 && strncmp(start, "sakhi", 5) == 0) return TOKEN_SAKHI;
    if (length == 4 && strncmp(start, "niyu", 4) == 0) return TOKEN_NIYU;
    if (length == 4 && strncmp(start, "mira", 4) == 0) return TOKEN_MIRA;
    if (length == 5 && strncmp(start, "kavya", 5) == 0) return TOKEN_KAVYA;
    if (length == 5 && strncmp(start, "shiri", 5) == 0) return TOKEN_SHIRI;

    if (length == 6 && strncmp(start, "import", 6) == 0) return TOKEN_IMPORT;

    return TOKEN_IDENTIFIER;
}

static Token identifier(void)
{
    const char *start = &source[current - 1];

    while (isalnum(peek()) || peek() == '_')
        advance();

    int length = &source[current] - start;
    TokenType type = check_keyword(start, length);

    return make_token(type, start, length);
}

/*──────────────────────────────────────────────
  STRING
──────────────────────────────────────────────*/

static Token string(void)
{
    const char *start = &source[current - 1];

    while (peek() != '"' && peek() != '\0') {
        advance();
    }

    if (peek() == '\0') {
        return make_token(TOKEN_ERROR, start, 0);
    }

    advance();

    int length = &source[current] - start;
    return make_token(TOKEN_STRING, start, length);
}

/*──────────────────────────────────────────────
  INIT
──────────────────────────────────────────────*/

void init_tokenizer(const char *src)
{
    source  = src;
    current = 0;
    line    = 1;
    col     = 1;
}

/*──────────────────────────────────────────────
  MAIN SCANNER
──────────────────────────────────────────────*/

Token scan_token(void)
{
    skip_whitespace();

    const char *start = &source[current];

    if (peek() == '\0')
        return make_token(TOKEN_EOF, start, 0);

    char c = advance();

    if (isdigit(c))
        return number();

    if (isalpha(c) || c == '_')
        return identifier();

    switch (c) {

        case '+':
         if (match('+')) return make_token(TOKEN_PLUSPLUS, start, 2);
                   return make_token(TOKEN_PLUS, start, 1);

        case '-':
         if (match('-')) return make_token(TOKEN_MINUSMINUS, start, 2);
                   return make_token(TOKEN_MINUS, start, 1);

        case '*': return make_token(TOKEN_STAR, start, 1);
        case '/': return make_token(TOKEN_SLASH, start, 1);
        case '&':
    if (match('&')) return make_token(TOKEN_AND, start, 2);
             break;

        case '|':
    if (match('|')) return make_token(TOKEN_OR, start, 2);
             break;

        case '%': return make_token(TOKEN_MOD, start, 1);

        case '(': return make_token(TOKEN_LEFT_PAREN, start, 1);
        case ')': return make_token(TOKEN_RIGHT_PAREN, start, 1);
        case '{': return make_token(TOKEN_LEFT_BRACE, start, 1);
        case '}': return make_token(TOKEN_RIGHT_BRACE, start, 1);
        case '[': return make_token(TOKEN_LEFT_BRACKET, start, 1);
        case ']': return make_token(TOKEN_RIGHT_BRACKET, start, 1);

        case ',': return make_token(TOKEN_COMMA, start, 1);
        case ':': return make_token(TOKEN_COLON, start, 1);

        case '=':
            return make_token(match('=') ? TOKEN_EQEQ : TOKEN_EQUAL, start, 1);

        case '!':
            return make_token(match('=') ? TOKEN_NEQ : TOKEN_NOT, start, 1);

        case '>':
            return make_token(match('=') ? TOKEN_GTE : TOKEN_GT, start, 1);

        case '<':
            return make_token(match('=') ? TOKEN_LTE : TOKEN_LT, start, 1);

        case '\n':
            return make_token(TOKEN_NEWLINE, start, 1);

        case '"':
            return string();
    }

    /* Unknown / Unicode characters skip */
    return scan_token();
}
nd{verbatim}
\section*{File: ./src/user_config.c}
egin{verbatim}
#include "../include/user_config.h"

/* ===== GLOBAL CONFIG ===== */

/* 🌸 DEV MODE */
int DEV_MODE = 0;  // 0 = USER | 1 = DEV

/* 🌸 LANGUAGE MODE */
int CURRENT_MODE = MODE_INDIA;

/* 🌸 PERSONALITY */
int CURRENT_PERSONALITY = P_SAKHI;

/* ===== INIT ===== */
void config_init()
{
    /* future config load */
}

/* ===== SETTERS ===== */
void config_set_personality(int p)
{
    CURRENT_PERSONALITY = p;
}

void config_set_mode(int mode)
{
    CURRENT_MODE = mode;
}
nd{verbatim}
\section*{File: ./fix_rules.h}
egin{verbatim}
#ifndef SHRIJI_FIX_RULES_H
#define SHRIJI_FIX_RULES_H

#include "../include/error.h"

/*
    SHRIJILANG — FIX CATEGORY SYSTEM
    ---------------------------------
    Ye layer error code ko fix category me convert karegi.

    IMPORTANT:
    Fix engine directly error code check nahi karega.
    Pehle category milegi.
    Fir category ke basis pe rule apply hoga.
*/

/* ─────────────────────────────────────
   FIX CATEGORY ENUM
───────────────────────────────────── */

typedef enum {

    FIX_NONE = 0,              // No fix possible
    FIX_MISSING_VALUE,         // '=' ke baad value missing
    FIX_BAD_TOKEN,             // Unexpected / invalid token
    FIX_STRUCTURAL_SAFE        // Safe structural auto-correction

} FixCategory;


/* ─────────────────────────────────────
   CLASSIFICATION FUNCTION
───────────────────────────────────── */

/*
    Error code → Fix category mapping
*/
FixCategory shriji_classify_error(ShrijiErrorCode code);


#endif
nd{verbatim}
\section*{File: ./gyaan/data/gyaan_rules.c}
egin{verbatim}
/*=============================================================
  SHRIJI GYAAN ENGINE — FINAL COMPLETE RULE SET
=============================================================*/

#include <stddef.h>
#include "../../include/error.h"

/*──────────────────────────────────────────────
  STRUCT
──────────────────────────────────────────────*/
typedef struct {
    int code;

    const char *module;
    const char *file;
    const char *function;

    const char *rule;
    const char *explain;
    const char *hint;

} GyaanRule;

/*=============================================================
  🧠 RULE DATABASE START
=============================================================*/
static GyaanRule GYAAN_RULES[] = {

/*──────────── PARSER CORE ────────────*/

{E_PARSE_OPERATOR_START,"parser","parser.c","parse_expression",
"expression cannot start with operator",
"operator se expression start nahi hota",
"example: 10 + 2"},

{E_PARSE_OPERATOR_END,"parser","parser.c","parse_expression",
"operator must have value after it",
"operator ke baad value missing hai",
"example: 10 + 2"},

{E_PARSE_DOUBLE_OPERATOR,"parser","parser.c","parse_expression",
"two operators cannot come together",
"do operator ek saath nahi aate",
"example: 10 + 2"},

{E_PARSE_OPERATOR_CHAIN,"parser","parser.c","parse_expression",
"operator chain is invalid",
"operator chain galat hai",
"example: 10 + 2"},

{E_PARSE_INVALID_TOKEN,"parser","parser.c","parse_statement",
"statement must start with valid token",
"line ki shuruaat galat token se hui hai",
"use: bolo / mavi / agar"},

{E_PARSE_BRACKET_MISSING,"parser","parser.c","parse_primary",
"closing bracket missing",
"bracket close nahi hua",
"example: (10 + 2)"},

/*──────────── PARSER ADVANCED ────────────*/

{E_PARSE_EXTRA_OPERATOR,"parser","parser.c","parse_expression",
"extra operator found",
"extra operator detect hua",
"example: 10 + 2"},

{E_PARSE_MISSING_OPERATOR,"parser","parser.c","parse_expression",
"operator missing between values",
"values ke beech operator missing hai",
"example: 10 + 2"},

{E_PARSE_UNMATCHED_PAREN,"parser","parser.c","parse_primary",
"parentheses mismatch",
"bracket match nahi hua",
"example: (10 + 2)"},

{E_PARSE_BRACKET_EXTRA,"parser","parser.c","parse_primary",
"extra closing bracket",
"extra bracket mila",
"bracket remove karein"},

{E_PARSE_MISSING_OPERAND,"parser","parser.c","parse_expression",
"operand missing",
"value missing hai operator ke liye",
"example: 10 + 2"},

/*──────────── ASSIGNMENT ────────────*/

{E_ASSIGN_01,"runtime","interpreter.c","AST_IDENTIFIER",
"variable must be declared before use",
"variable declare nahi hua",
"use: mavi x = 10"},

{E_ASSIGN_02,"parser","parser.c","parse_assignment",
"assignment must have value",
"assignment me value missing hai",
"example: mavi x = 10"},

/*──────────── FUNCTION ────────────*/

{E_PARSE_01,"parser","parser.c","general",
"syntax error",
"syntax samajh nahi aaya",
"statement check karein"},

{E_PARSE_02,"system","multiple","general",
"invalid usage",
"usage galat hai",
"syntax check karein"},

/*──────────── CONTROL FLOW ────────────*/

{E_IF_01,"parser","parser.c","parse_if",
"invalid if condition",
"if condition galat hai",
"example: agar x > 0 { ... }"},

{E_RUNTIME_LOOP_LIMIT,"runtime","interpreter.c","AST_WHILE",
"loop exceeded safe limit",
"loop limit cross ho gaya",
"condition fix karein"},

/*──────────── DATA STRUCTURES ────────────*/

{E_RUNTIME_INDEX_ERROR,"runtime","interpreter.c","AST_INDEX",
"index out of bounds",
"index range ke bahar hai",
"example: a[0]"},

/*──────────── RUNTIME ────────────*/

{E_RUNTIME_DIV_ZERO,"runtime","interpreter.c","safe_div",
"division by zero not allowed",
"0 se divide nahi kar sakte",
"example: 10 / 2"},

{E_RUNTIME_TYPE_MISMATCH,"runtime","interpreter.c","AST_BINARY",
"type mismatch in expression",
"different type values use hue hain",
"same type use karein"},

{E_RUNTIME_01,"runtime","interpreter.c","general",
"runtime error occurred",
"runtime me issue aaya",
"logic check karein"}

};

/*=============================================================
  🧠 RULE DATABASE END
=============================================================*/

/*──────────────────────────────────────────────
  COUNT
──────────────────────────────────────────────*/
static int GYAAN_COUNT =
sizeof(GYAAN_RULES) / sizeof(GYAAN_RULES[0]);

/*──────────────────────────────────────────────
  LOOKUP
──────────────────────────────────────────────*/
const GyaanRule* gyaan_lookup(int code)
{
    for (int i = 0; i < GYAAN_COUNT; i++) {
        if (GYAAN_RULES[i].code == code) {
            return &GYAAN_RULES[i];
        }
    }
    return NULL;
}
nd{verbatim}
\section*{File: ./gyaan/core/gyaan_engine.h}
egin{verbatim}
#ifndef GYAAN_ENGINE_H
#define GYAAN_ENGINE_H

#include "../../include/error.h"

typedef struct {

    ShrijiErrorCode error;

    const char *module;
    const char *file;
    const char *function;

    const char *rule;
    const char *explain;
    const char *hint;

} GyaanEntry;

const GyaanEntry* gyaan_lookup(ShrijiErrorCode code);

void gyaan_print(const ShrijiErrorInfo *err);

#endif
nd{verbatim}
\section*{File: ./gyaan/core/gyaan_engine.c}
egin{verbatim}
/*=============================================================
  SHRIJI GYAAN ENGINE — CORE (HINGLISH + CLEAR UX VERSION)
=============================================================*/

#include <stdio.h>
#include <string.h>
#include "../../include/error.h"

typedef struct {
    int code;

    const char *module;
    const char *file;
    const char *function;

    const char *rule;
    const char *explain;
    const char *hint;

} GyaanRule;

extern const GyaanRule* gyaan_lookup(int code);

static const char* safe(const char *s)
{
    return (s && *s) ? s : "N/A";
}

void gyaan_print(const ShrijiErrorInfo *err)
{
    if (!err) return;

    const GyaanRule *r = gyaan_lookup(err->code);

    /* SOFT FUNCTION MATCH */
    if (r && err->function && r->function) {
        if (strcmp(err->function, r->function) != 0) {
            /* allow fallback */
        }
    }

    /* FALLBACK (NO RULE FOUND) */
    if (!r) {
        printf("\nSHRIJI INSIGHT\n");
        printf("──────────────────────────────\n");

        printf("System ko is error ka exact explanation nahi mila\n");

        printf("Samjho   : input ka structure (syntax) sahi format me nahi hai\n");

        printf("Hint     : statement ko proper format me likho\n");
        printf("Example  : bolo 5 + 2\n");

        if (err->has_location) {
            printf("\nLocation : line %d, col %d\n", err->line, err->col);
        }

        printf("──────────────────────────────\n");
        return;
    }

    /* NORMAL RULE PRINT */
    printf("\nSHRIJI INSIGHT\n");
    printf("──────────────────────────────\n");

    printf("Module   : %s\n", safe(r->module));
    printf("File     : %s\n", safe(r->file));
    printf("Function : %s\n\n", safe(r->function));

    printf("Rule     : %s\n", safe(r->rule));
    printf("Samjho   : %s\n", safe(r->explain));
    printf("Hint     : %s\n", safe(r->hint));

    switch (err->code)
    {
        case E_PARSE_MISSING_OPERATOR:
            printf("\nNote     : yahan values ke beech operator missing tha\n");
            break;

        case E_PARSE_DOUBLE_OPERATOR:
            printf("\nNote     : do operator ek saath aa gaye the\n");
            break;

        case E_PARSE_OPERATOR_START:
            printf("\nNote     : expression operator se start nahi hota\n");
            break;

        case E_PARSE_OPERATOR_END:
            printf("\nNote     : expression operator par khatam ho gaya\n");
            break;

        default:
            break;
    }

    if (err->has_location) {
        printf("\nLocation : line %d, col %d\n", err->line, err->col);
    }

    printf("──────────────────────────────\n");
}
nd{verbatim}
\section*{File: ./krst/krst_router.c}
egin{verbatim}
#include <stdio.h>
#include <string.h>

#include "../include/krst_router.h"
#include "../krst/krst_core.h"
#include "../include/pragya_avastha.h"

#include "../include/error.h"
#include "../include/error_intelligence.h"

/* ===============================
   INPUT NORMALIZATION (SAFE)
   =============================== */
static void normalize_input(const char *input, char *output)
{
    if (!input || !output) return;

    int i = 0;
    for (; input[i]; i++) {
        if (input[i] >= 'A' && input[i] <= 'Z')
            output[i] = input[i] + 32;
        else
            output[i] = input[i];
    }
    output[i] = '\0';
}

/* ===============================
   COMMAND DETECTION
   =============================== */
static char detect_command(const char *input)
{
    if (!input) return 'G';

    if (input[0] == 'g' && input[1] == ' ')
        return 'G';

    if (input[0] == 'e' && input[1] == ' ')
        return 'E';

    return 'G';
}

/* ===============================
   CONTEXT REWRITE (C2 COMPLETE)
   =============================== */
static int rewrite_context_input(KRSTRequest *req)
{
    if (!req || !req->input_text) return 0;

    int num = 0;
    static char buffer[128];

    const char *input = req->input_text;

    /* ===== ADD ===== */
    if ((strstr(input, "add") || strstr(input, "jodo")) &&
        sscanf(input, "usme %d", &num) == 1) {

        snprintf(buffer, sizeof(buffer), "bolo last + %d", num);
        req->input_text = buffer;
        return 1;
    }

    /* ===== MINUS ===== */
    if ((strstr(input, "minus") || strstr(input, "ghatao") || strstr(input, "subtract")) &&
        (sscanf(input, "usme %d", &num) == 1 ||
         sscanf(input, "usme se %d", &num) == 1)) {

        snprintf(buffer, sizeof(buffer), "bolo last - %d", num);
        req->input_text = buffer;
        return 1;
    }

    /* ===== MULTIPLY ===== */
    if ((strstr(input, "multiply") || strstr(input, "guna") || strstr(input, "into")) &&
        sscanf(input, "usme %d", &num) == 1) {

        snprintf(buffer, sizeof(buffer), "bolo last * %d", num);
        req->input_text = buffer;
        return 1;
    }

    /* ===== DIVIDE ===== */
    if ((strstr(input, "divide") || strstr(input, "bhaag") || strstr(input, "divided")) &&
        (sscanf(input, "usme %d", &num) == 1 ||
         sscanf(input, "usme se %d", &num) == 1)) {

        snprintf(buffer, sizeof(buffer), "bolo last / %d", num);
        req->input_text = buffer;
        return 1;
    }

    return 0;
}

/* ===============================
   MAIN ROUTER
   =============================== */
void krst_route_request(KRSTRequest *req)
{
    if (!req || !req->input_text)
        return;

    /* ===== SAFE NORMALIZATION ===== */
    static char normalized[256];

    normalize_input(req->input_text, normalized);
    req->input_text = normalized;

    /* ===== COMMAND DETECTION ===== */
    char cmd = detect_command(req->input_text);

    /* ===== COMMAND ROUTING ===== */

     if (cmd == 'E') {
    req->input_text += 2;

    static char explain_buf[256];
   snprintf(explain_buf, sizeof(explain_buf), "bolo %.240s", req->input_text);
    req->input_text = explain_buf;

    printf("[EXPLAIN MODE]\n");
}

    else if (cmd == 'G') {
        if (req->input_text[0] == 'g' && req->input_text[1] == ' ')
            req->input_text += 2;
    }

    /* ===== CONTEXT REWRITE ===== */
    rewrite_context_input(req);

    /* ===== CREATE AVASTHA ===== */
    PragyaAvastha avastha;
    memset(&avastha, 0, sizeof(PragyaAvastha));

    avastha.raw_text = req->input_text;

      avastha.raw_text = req->input_text;

/*  NEW: pass input mode */
    avastha.is_program_mode = req->is_program_mode;

/* ===== KRST CORE ===== */
krst_process(&avastha);

/* ===== ERROR INTELLIGENCE FALLBACK ===== */
if (!avastha.has_correction && error_reported)
{
    shriji_error_intelligence(&avastha, &LAST_ERROR, NULL);
}

/* ===== RETURN CORRECTION ===== */
if (avastha.has_correction)
{
    req->has_correction = 1;

    strncpy(req->corrected_text,
            avastha.corrected_text,
            sizeof(req->corrected_text) - 1);

    req->corrected_text[
        sizeof(req->corrected_text) - 1] = '\0';
}
}
nd{verbatim}
\section*{File: ./krst/krst_core.c}
egin{verbatim}
/* RADHE_RADHE_SHREEJI_RADHE_RANI */

#include <stdio.h>
#include <string.h>
#include <ctype.h>

#include "../include/interpreter.h"
#include "../include/error.h"
#include "../include/fix_engine.h"
#include "../include/nirman.h"
#include "../include/parser.h"
#include "../include/value.h"
#include "../include/env.h"
#include "../include/runtime.h"
#include "../include/pragya_avastha.h"
#include "../include/decision_engine.h"

#include "krst_core.h"
#include "krst_types.h"

#include "../pragya/router/l3_router.h"
#include "../include/error_intelligence.h"

#include "../include/user_config.h"

extern Env *GLOBAL_ENV;

/* ================= SESSION ================= */

static int session_confidence = 100;
static int session_risk = 4;
static int success_streak = 0;

/* ================= LABEL ================= */

static const char* teach_label(int level)
{
    switch(level)
    {
        case KRST_TEACH_SILENT:  return "SILENT";
        case KRST_TEACH_HINT:    return "HINT";
        case KRST_TEACH_EXPLAIN: return "EXPLAIN";
        case KRST_TEACH_DEEP:    return "DEEP";
        case KRST_TEACH_TRAIN:   return "TRAIN";
        default: return "SILENT";
    }
}

static const char* tone_label(int tone)
{
    switch(tone)
    {
        case KRST_TONE_NEUTRAL: return "NEUTRAL";
        case KRST_TONE_STRICT:  return "STRICT";
        case KRST_TONE_ALERT:   return "ALERT";
        default: return "NEUTRAL";
    }
}

/* ================= THRESHOLD ================= */

static void apply_thresholds(KRSTDecision *d)
{
    int c = d->confidence_score;

    if (c <= 20)      d->teaching_level = KRST_TEACH_TRAIN;
    else if (c <= 40) d->teaching_level = KRST_TEACH_DEEP;
    else if (c <= 60) d->teaching_level = KRST_TEACH_EXPLAIN;
    else if (c <= 80) d->teaching_level = KRST_TEACH_HINT;
    else              d->teaching_level = KRST_TEACH_SILENT;

    if (session_risk >= 76)
    {
        d->tone = KRST_TONE_ALERT;
        d->escalate = 1;
    }
    else if (session_risk >= 30)
    {
        d->tone = KRST_TONE_STRICT;
        d->escalate = 0;
    }
    else
    {
        d->tone = KRST_TONE_NEUTRAL;
        d->escalate = 0;
    }
}

/* ================= MAIN ================= */

int krst_process(PragyaAvastha *avastha)
{
    if (!avastha || !avastha->raw_text)
        return 1;

const char *input = avastha->raw_text;

/*  NIRMAN START */
if (strcmp(input, "nirman") == 0) {
    nirman_start();
    printf("Nirman mode activated\n");
    printf("Welcome to Shriji World\n");
    return 0;
}

/*  NIRMAN STOP */
if (strcmp(input, "exit nirman") == 0) {
    nirman_stop();
    printf("Nirman mode deactivated\n");
    return 0;
}

/*  NIRMAN MODE ROUTING */
if (nirman_is_active()) {

    int intent = nirman_detect_intent(input);

    printf("SHRIJI (intent=%d)\n", intent);

    return 0;
}


    if (DEV_MODE)
        printf("[KRST] Input: %s\n", input);

    error_reported = 0;
    avastha->stop_execution = 0;

    avastha->confidence = session_confidence;
    avastha->risk = session_risk;

    /* ================= PARSE ================= */

    int was_fixed = 0;
    int fix_penalty = 0;

    ASTNode *root = NULL;

    if (avastha->is_program_mode)
        root = parse_program(input);
    else
        root = language_execute_with_fix(input, &was_fixed, &fix_penalty);

    /* ================= ERROR ================= */

    if (!root || error_reported)
    {
        avastha->ast = NULL;
        avastha->stop_execution = 1;
        success_streak = 0;

        const ShrijiErrorInfo *err_ptr = shriji_last_error();

        if (!err_ptr)
            return 0;

        ShrijiErrorInfo err_safe = *err_ptr;

        DecisionType d_type = shriji_take_decision(
            session_confidence,
            session_risk,
            err_safe.code
        );

        char fixed[256];
        FixType ftype = fix_apply(input, fixed);

        /* AUTO FIX ONLY */
        if (ftype == FIX_SAFE && d_type == DECISION_AUTO_FIX)
        {
            printf("\n[Shriji] auto fix → %s\n\n", fixed);

            avastha->has_correction = 1;

            snprintf(avastha->corrected_text,
                     sizeof(avastha->corrected_text),
                     "%s", fixed);

            return 0;
        }

        /*  NO EI CALL HERE */
        return 0;
    }

    /* ================= SUCCESS ================= */
   if (!root || error_reported || avastha->stop_execution)
{
    return 0;
}
    avastha->ast = root;

    ShrijiRuntime runtime;
    runtime_init(&runtime);

    Value result = run_program(root, GLOBAL_ENV, &runtime);

    if (!error_reported && !runtime.error_flag)
{
        success_streak++;

        int recovery = 5 + (success_streak * 2);
        if (recovery > 25) recovery = 25;

        session_confidence += recovery;
        session_risk -= recovery;

        if (was_fixed)
        {
            session_confidence -= fix_penalty;
            session_risk += fix_penalty;
        }

        if (result.type == VAL_NUMBER)
            printf("%g\n", result.number);
        else if (result.type == VAL_STRING && result.string)
            printf("%s\n", result.string);
    }

    value_free(&result);

    if (session_confidence < 0) session_confidence = 0;
    if (session_confidence > 100) session_confidence = 100;

    if (session_risk < 4) session_risk = 4;
    if (session_risk > 100) session_risk = 100;

    KRSTDecision decision = {0};
    decision.confidence_score = session_confidence;

    apply_thresholds(&decision);

    if (DEV_MODE)
    {
        printf("[KRST] Confidence: %d | Risk: %d | Teach: %s | Tone: %s | Esc: %d\n",
               session_confidence,
               session_risk,
               teach_label(decision.teaching_level),
               tone_label(decision.tone),
               decision.escalate);
    }

    avastha->confidence = session_confidence;
    avastha->risk = session_risk;
    avastha->teach_level = decision.teaching_level;
    avastha->tone = decision.tone;

    pragya_route(avastha);

    return 0;
}
nd{verbatim}
\section*{File: ./krst/krst_types.h}
egin{verbatim}
#ifndef KRST_TYPES_H
#define KRST_TYPES_H

#include "../include/env.h"

/* ==============================
   KRST ACTION ENUM
============================== */

typedef enum {
    KRST_ACTION_ALLOW = 0,
    KRST_ACTION_BLOCK,
    KRST_ACTION_TEACH,
    KRST_ACTION_ESCALATE
} KRSTAction;

/* ==============================
   KRST TEACHING LEVEL ENUM
============================== */

typedef enum {
    KRST_TEACH_SILENT = 0,
    KRST_TEACH_HINT,
    KRST_TEACH_EXPLAIN,
    KRST_TEACH_DEEP,
    KRST_TEACH_TRAIN
} KRSTTeachingLevel;

/* ==============================
   KRST TONE ENUM
============================== */

typedef enum {
    KRST_TONE_NEUTRAL = 0,
    KRST_TONE_SUPPORTIVE,
    KRST_TONE_GUIDING,
    KRST_TONE_STRICT,
    KRST_TONE_ALERT
} KRSTTone;

/* ==============================
   KRST REQUEST STRUCT
============================== */
typedef struct
{
    const char *input_text;

     int is_program_mode;
    /* correction support */

    char corrected_text[512];
    int has_correction;

} KRSTRequest;
/* ==============================
   KRST DECISION STRUCT
============================== */

typedef struct {
    KRSTAction action;

    int confidence_score;   /* 0–100 */
    int risk_score;         /* 0–100 */

    int allow_execution;
    int reason_code;

    KRSTTeachingLevel teaching_level;
    KRSTTone tone;

    int escalate;   /* 1 = escalation required */

} KRSTDecision;

#endif
nd{verbatim}
\section*{File: ./krst/krst_core.h}
egin{verbatim}
#ifndef KRST_CORE_H
#define KRST_CORE_H

#include "../include/pragya_avastha.h"

/*
    KRST Core Interface

    Role:
    - Main processing engine
    - Takes PragyaAvastha as input
    - Updates it (error, correction, teaching, etc.)
*/

/* Main processing function */
int krst_process(PragyaAvastha *avastha);

#endif
nd{verbatim}
\section*{File: ./krst/niyu_engine.c}
egin{verbatim}
#include <string.h>
#include "niyu_engine.h"

int niyu_analyze(const char *input)
{
    if (!input)
        return 0;

    if (strlen(input) == 0)
        return 0;

    return 0;
}
nd{verbatim}
\section*{File: ./Makefile}
egin{verbatim}
# ============================================================
# SHRIJILANG — UNIFIED BUILD SYSTEM (KRST + PRAGYA)
# ============================================================

CC = gcc

CFLAGS = -Wall -Wextra -std=c11 -O2 \
	-Iinclude \
	-Isrc \
	-Ienv \
	-Ipragya \
	-Ikrst

# ============================================================
# CORE SOURCE FILES
# ============================================================

SRC = \
	src/token.c \
	src/ast.c \
	src/parser.c \
	gyaan/core/gyaan_engine.c \
	gyaan/data/gyaan_rules.c \
	src/interpreter.c \
	src/runtime.c \
	src/error.c \
	src/state.c \
	src/event.c \
	src/command.c \
	src/translator.c \
	src/example_builder.c \
	src/ai_router.c \
	src/ai_intent.c \
	src/atma_lekha.c \
	src/smriti.c \
        src/smriti_session.c \
	src/smriti_personal.c \
	src/main.c \
	env/env.c \
	src/lang/grammar/grammar_registry.c \
	src/lang/grammar/grammar_core.c \
	pragya/router/l3_router.c \
	pragya/niyu/niyu_core.c \
	pragya/sakhi/sakhi_core.c \
	pragya/mira/mira_core.c \
	krst/krst_core.c \
	krst/krst_router.c \
	krst/niyu_engine.c \
	src/fix_engine.c \
	src/fix_rules.c \
	src/fix_interactive.c \
	src/typo_engine.c \
	src/error_intelligence.c \
	src/decision_engine.c \
src/user_config.c \
src/nirman.c \
src/normalize_engine.c \
src/nirman_flow_engine.c \

# ============================================================
# OUTPUT
# ============================================================

OUT = shrijilang

# ============================================================
# BUILD RULES
# ============================================================

$(OUT): $(SRC)
	$(CC) $(CFLAGS) $(SRC) -o $(OUT) -lreadline

clean:
	rm -f $(OUT)
	rm -f *.o
	rm -f src/*.o
	rm -f env/*.o
	rm -f src/lang/grammar/*.o
	rm -f pragya/router/*.o
	rm -f pragya/niyu/*.o
	rm -f pragya/sakhi/*.o
	rm -f pragya/mira/*.o
	rm -f krst/*.o

run: $(OUT)
	./$(OUT)

rebuild: clean $(OUT)
nd{verbatim}
\section*{File: ./pragya/sakhi/sakhi_core.c}
egin{verbatim}
#include "sakhi_core.h"
#include <stdio.h>

L3Response sakhi_speak(const L3Request *request,
                       const NiyuResult *result)
{
    (void)result;

    L3Response res;

    if (!request || !request->raw_input)
    {
        res.text = "Input samajh nahi aaya.";
        res.success = 0;
        return res;
    }

    printf("[PRAGYA] Main sun rahi hoon. Thoda clear likhoge to aur achha ho jayega 🙂\n");

    res.text = request->raw_input;
    res.success = 1;

    return res;
}
nd{verbatim}
\section*{File: ./pragya/sakhi/sakhi_core.h}
egin{verbatim}
#ifndef SAKHI_CORE_H
#define SAKHI_CORE_H

#include "sakhi_contract.h"

/*
 * SAKHI CORE
 *
 * Implementation file will:
 * - handle tone / language / clarity
 * - never inspect raw logic
 *
 * This header intentionally exposes no helpers.
 */

#endif /* SAKHI_CORE_H */
nd{verbatim}
\section*{File: ./pragya/sakhi/sakhi_contract.h}
egin{verbatim}
#ifndef SAKHI_CONTRACT_H
#define SAKHI_CONTRACT_H

#include "../contracts/l3_request.h"
#include "../contracts/l3_response.h"
#include "../niyu/niyu_contract.h"

/*
 * SAKHI CONTRACT
 *
 * Sakhi is the voice of the system.
 * - Receives decided intent + thinking result (opaque)
 * - Produces final human-readable response
 *
 * Sakhi never:
 * - thinks or calculates
 * - decides intent
 * - teaches (that is Mira)
 */

/* Forward declaration (opaque thinking result) */
/*
 * sakhi_speak
 *
 * Final response generator.
 * - Called only by Router
 * - Uses intent already decided
 * - Formats language and tone
 */
L3Response sakhi_speak(const L3Request *request,
                       const NiyuResult *result);

#endif /* SAKHI_CONTRACT_H */
nd{verbatim}
\section*{File: ./pragya/contracts/l3_intent.h}
egin{verbatim}
#ifndef L3_INTENT_H
#define L3_INTENT_H

/* 
 * L3 Intent:
 * Router decides intent.
 * Niyu never decides intent.
 * Sakhi/Mira never guess intent.
 */

typedef enum {
    L3_INTENT_UNKNOWN = 0,

    /* Explanation / reasoning */
    L3_INTENT_EXPLAIN,
    L3_INTENT_WHY,

    /* Pure calculation / logic */
    L3_INTENT_CALC,

    /* Teaching / learning */
    L3_INTENT_TEACH,

    /* Emotional / tone-based */
    L3_INTENT_EMOTION,

    /* Explicit command */
    L3_INTENT_COMMAND

} L3Intent;

#endif /* L3_INTENT_H */
nd{verbatim}
\section*{File: ./pragya/contracts/l3_rules.h}
egin{verbatim}
#ifndef L3_RULES_H
#define L3_RULES_H

/*
 * L3 NON-NEGOTIABLE LAWS
 *
 * 1. Router decides intent.
 * 2. Niyu THINKS, never SPEAKS.
 * 3. Sakhi SPEAKS, never THINKS.
 * 4. Mira TEACHES only when intent == L3_INTENT_TEACH.
 *    No external permission, no memory switch.
 * 5. No direct Niyu <-> Sakhi calls.
 * 6. All communication flows through Router.
 * 7. Contracts are stable forever.
 * 8. Engine may change, behavior must not.
 */

#define L3_RULE_ROUTER_IS_AUTHORITY 1
#define L3_RULE_NIYU_SILENT         1
#define L3_RULE_SAKHI_NO_LOGIC      1
#define L3_RULE_MIRA_EXPLICIT_ONLY  1

#endif /* L3_RULES_H */
nd{verbatim}
\section*{File: ./pragya/contracts/l3_request.h}
egin{verbatim}
#ifndef L3_REQUEST_H
#define L3_REQUEST_H

#include "l3_intent.h"

/*
 * L3Request:
 * Immutable input structure.
 * Router decides intent.
 * No teaching_mode flag here.
 * Behavior intent se hi control hoga.
 */

typedef struct {
    const char *raw_input;     /* original user text */
    L3Intent    intent;        /* decided by router */
} L3Request;

#endif /* L3_REQUEST_H */
nd{verbatim}
\section*{File: ./pragya/contracts/l3_response.h}
egin{verbatim}
#ifndef L3_RESPONSE_H
#define L3_RESPONSE_H

/*
 * L3Response:
 * Only Sakhi or Mira can create final response.
 * Niyu can never write directly here.
 */

typedef struct {
    const char *text;      /* final output */
    int         success;   /* logical success/failure */
} L3Response;

#endif /* L3_RESPONSE_H */
nd{verbatim}
\section*{File: ./pragya/mira/mira_core.c}
egin{verbatim}
// JAI SHREE KRISHNA
#include "mira_contract.h"
#include "../niyu/niyu_contract.h"

#include "../../include/error.h"
#include "../../include/pragya_avastha.h"

#include <stdio.h>
#include <string.h>


/* =========================================
   POINTER PRINTER
========================================= */

static void print_pointer(const char *text, int pos)
{
    if (!text || pos < 0)
        return;

    printf("%s\n", text);

    for (int i = 0; i < pos; i++)
        printf(" ");

    printf("^\n\n");
}


/* =========================================
   MIRA TEACHING ENGINE
========================================= */

void mira_teach_avastha(PragyaAvastha *a, NiyuResult *res)
{
    if (!a || !res)
        return;

    const char *text = a->raw_text;


    /* valid input → silent */

    if (res->error_type == NIYU_ERR_VALID)
        return;


    switch (a->error_code)
    {

    /* =========================================
       DOUBLE OPERATOR
    ========================================= */

    case E_PARSE_DOUBLE_OPERATOR:
    case E_PARSE_OPERATOR_CHAIN:

        printf("[SHRI_JI] Lagta hai yahaan extra operator aa gaya hai \n\n");

        print_pointer(text, res->error_pos);

        printf("Do operators ek saath allowed nahi hote.\n\n");

        printf("Example:\n");
        printf("6 + 5\n");

        break;


    /* =========================================
       OPERATOR AT START
    ========================================= */

    case E_PARSE_OPERATOR_START:

        printf("[SHRI_JI] Expression operator se start nahi hota \n\n");

        print_pointer(text, res->error_pos);

        printf("Pehle ek value honi chahiye.\n\n");

        printf("Example:\n");
        printf("6 + 5\n");

        break;


    /* =========================================
       OPERATOR AT END
    ========================================= */

    case E_PARSE_OPERATOR_END:

        printf("[SHRI_JI] Expression incomplete lag raha hai \n\n");

        print_pointer(text, res->error_pos);

        printf("Operator ke baad ek value expected hoti hai.\n\n");

        printf("Example:\n");
        printf("6 + 5\n");

        break;


    /* =========================================
       MISSING OPERATOR
    ========================================= */

    case E_PARSE_MISSING_OPERATOR:

        printf("[SHRI_JI] Lagta hai operator missing hai \n\n");

        print_pointer(text, res->error_pos);

        printf("Do values ke beech operator hona chahiye.\n\n");

        printf("Example:\n");
        printf("5 + 6\n");

        break;


    /* =========================================
       ASSIGNMENT ERROR
    ========================================= */

    case E_ASSIGN_02:

        printf("[SHRI_JI] Assignment incomplete lag raha hai.\n\n");

        print_pointer(text, res->error_pos);

        printf("'=' ke baad value expected hai.\n\n");

        printf("Example:\n");
        printf("mavi x = 10\n");

        break;


    /* =========================================
       GENERIC SYNTAX ERROR
    ========================================= */

    case E_PARSE_02:

        printf("[SHRI_JI] Syntax error detect hua hai.\n\n");

        print_pointer(text, res->error_pos);

        printf("Yahan invalid token ya unsupported symbol mila hai.\n\n");

        printf("Example valid syntax:\n");
        printf("6 + 5\n");

        break;


    /* =========================================
       DEFAULT
    ========================================= */

default:
    return;
    }
}
nd{verbatim}
\section*{File: ./pragya/mira/mira_contract.h}
egin{verbatim}
#ifndef MIRA_CONTRACT_H
#define MIRA_CONTRACT_H

#include "../contracts/l3_request.h"
#include "../contracts/l3_response.h"
#include "../niyu/niyu_contract.h"
#include "../../include/pragya_avastha.h"

/*
 * MIRA CONTRACT
 *
 * Mira is the teaching layer.
 * - Activated only when teaching_mode == 1
 * - Provides structured explanation / lesson
 *
 * Mira never:
 * - handles normal conversation
 * - decides intent
 * - performs deep logic (that is Niyu)
 */

/*
 * Legacy interface (old L3 system)
 */
L3Response mira_teach(const L3Request *request);

/*
 * New ShrijiLang teaching interface
 * Used by Pragya Router
 */

void mira_teach_avastha(
    PragyaAvastha *avastha,
    NiyuResult *result
);

#endif /* MIRA_CONTRACT_H */
nd{verbatim}
\section*{File: ./pragya/mira/mira_core.h}
egin{verbatim}
#ifndef MIRA_CORE_H
#define MIRA_CORE_H

#include "mira_contract.h"

/*
 * MIRA CORE
 *
 * Implementation file will:
 * - manage lesson flow
 * - provide explanations
 *
 * No other module may call Mira directly.
 */

#endif /* MIRA_CORE_H */
nd{verbatim}
\section*{File: ./pragya/README_L3.txt}
egin{verbatim}
L3 (KR̥ṢṬ) — FOUNDATION DOCUMENT
===============================

This document is the final, non-negotiable foundation of L3.
It defines roles, laws, flow, and future guarantees.
Implementation may evolve. Behavior must not.

----------------------------------------------------------------
1. PURPOSE
----------------------------------------------------------------
L3 is the intelligence orchestration layer of ShrijiLang.

It separates:
- Thinking (logic, calculation)
- Speaking (human-facing output)
- Teaching (explicit learning mode)

L3 is designed to work today with lightweight logic
and tomorrow with machine-backed, high-compute intelligence
without breaking language behavior.

----------------------------------------------------------------
2. CORE ROLES (FROZEN)
----------------------------------------------------------------

Router
------
- Central authority of L3
- Decides intent
- Controls full request lifecycle
- Chooses which module participates

Router never:
- Thinks deeply
- Formats language
- Teaches directly

Niyu (THINK)
------------
- Silent thinking engine
- Performs logic, reasoning, calculation
- Produces internal result only

Niyu never:
- Speaks to user
- Decides intent
- Formats language

Sakhi (SPEAK)
-------------
- Final voice of the system
- Produces human-readable output
- Applies tone and clarity

Sakhi never:
- Thinks or calculates
- Decides intent
- Teaches

Mira (TEACH)
------------
- Teaching and explanation layer
- Activated only when explicitly allowed
- Provides structured learning responses

Mira never:
- Handles normal conversation
- Decides intent
- Performs deep logic

----------------------------------------------------------------
3. FLOW (CANONICAL)
----------------------------------------------------------------

User Input
   |
   v
Router
   |
   |-- decide intent
   |-- check teaching_mode
   |
   |--> Niyu (if logic / calculation needed)
   |        |
   |        v
   |   NiyuResult (internal)
   |
   |--> Sakhi (normal response)
   |        OR
   |--> Mira  (teaching response)
   |
   v
Final L3Response to User

No module may bypass the Router.

----------------------------------------------------------------
4. CONTRACTS (IMMUTABLE)
----------------------------------------------------------------
- l3_request.h  : defines input shape
- l3_intent.h   : defines intent space
- l3_response.h : defines output shape
- l3_rules.h    : defines laws

Contracts are stable forever.
Engines may change. Contracts must not.

----------------------------------------------------------------
5. NON-NEGOTIABLE LAWS
----------------------------------------------------------------
1. Router is the sole authority.
2. Niyu thinks, never speaks.
3. Sakhi speaks, never thinks.
4. Mira teaches only when explicitly enabled.
5. No direct Niyu <-> Sakhi communication.
6. All flows pass through Router.
7. Behavior consistency is mandatory.
8. Future machine intelligence must fit this structure.

----------------------------------------------------------------
6. FUTURE GUARANTEE
----------------------------------------------------------------
L3 is intentionally designed as a replaceable intelligence layer.

Today:
- Local logic
- Lightweight reasoning

Tomorrow:
- Dedicated machines
- Heavy computation
- Real-world modeling

The language behavior, contracts, and flow
must remain unchanged across this evolution.

----------------------------------------------------------------
7. FINAL NOTE
----------------------------------------------------------------
This document is the NĒEV (foundation) of L3.

Any change that violates this document
is considered a breaking change and is not allowed.

Implementation follows this document.
This document never follows implementation.
nd{verbatim}
\section*{File: ./pragya/router/l3_router.c}
egin{verbatim}
#include "l3_router.h"

#include "../niyu/niyu_contract.h"
#include "../sakhi/sakhi_contract.h"
#include "../mira/mira_contract.h"

#include "../../include/pragya_avastha.h"

#include <stdio.h>

/*
=========================================================
   PRAGYA ROUTER (L3)

   ROLE:
   - KRST se aayi hui avastha ko process karna
   - Niyu → logic analysis
   - Mira → teaching layer
   - Sakhi → user response

   NOTE:
   - stop_execution flag ka respect MUST hai
=========================================================
*/

void pragya_route(PragyaAvastha *avastha)
{
    /* ========================= */
    /* BASIC SAFETY CHECK        */
    /* ========================= */

    if (!avastha)
        return;

    /* ========================= */
    /* 🔥 STOP EXECUTION GUARD   */
    /* ========================= */

    if (avastha->stop_execution)
        return;

    /* ========================= */
    /* SILENT MODE               */
    /* ========================= */

    if (avastha->teach_level == 0)
        return;

    /* ========================= */
    /* NIYU ANALYSIS             */
    /* ========================= */

    NiyuResult *result = niyu_think_avastha(avastha);

    if (!result)
        return;

    /* ========================= */
    /* MIRA TEACHING             */
    /* ========================= */

    if (avastha->teach_level > 1)
    {
        mira_teach_avastha(avastha, result);
    }

    /* ========================= */
    /* SAKHI RESPONSE            */
    /* ========================= */

    if (avastha->error_code != 0)
    {
        L3Request request;

        request.raw_input = avastha->raw_text;
        request.intent = 0;

        sakhi_speak(&request, result);
    }

    /* ========================= */
    /* CLEANUP                   */
    /* ========================= */

    niyu_free_result(result);
}
nd{verbatim}
\section*{File: ./pragya/router/l3_router.h}
egin{verbatim}
#ifndef PRAGYA_ROUTER_H
#define PRAGYA_ROUTER_H

#include "../../include/pragya_avastha.h"

void pragya_route(PragyaAvastha *avastha);

#endif
nd{verbatim}
\section*{File: ./pragya/niyu/niyu_contract.h}
egin{verbatim}
#ifndef NIYU_CONTRACT_H
#define NIYU_CONTRACT_H

#include "../../include/pragya_avastha.h"

/* =============================== */
/* NIYU ERROR TYPES                */
/* =============================== */

typedef enum {

    NIYU_ERR_VALID = 0,
    NIYU_ERR_SYNTAX,
    NIYU_ERR_MISSING_OPERAND,
    NIYU_ERR_EXTRA_OPERATOR,
    NIYU_ERR_INVALID_IDENTIFIER,
    NIYU_ERR_UNDEFINED_VARIABLE,
    NIYU_ERR_TYPE_MISMATCH,
    NIYU_ERR_DIV_ZERO,
    NIYU_ERR_UNEXPECTED_TOKEN,
    NIYU_ERR_EMPTY_INPUT,
    NIYU_ERR_INVALID_COMMAND,
    NIYU_ERR_RUNTIME

} NiyuErrorType;


/* =============================== */
/* NIYU INTENT TYPES               */
/* =============================== */

typedef enum {

    NIYU_INTENT_CALCULATION = 0,
    NIYU_INTENT_DECLARATION,
    NIYU_INTENT_PRINT,
    NIYU_INTENT_UNKNOWN

} NiyuIntent;


/* =============================== */
/* NIYU RESULT STRUCT              */
/* =============================== */

typedef struct {

    NiyuIntent intent;
    NiyuErrorType error_type;

    int severity;

    int error_pos;

} NiyuResult;


/* =============================== */
/* NIYU CORE API                   */
/* =============================== */

NiyuResult* niyu_think_avastha(PragyaAvastha *avastha);

void niyu_free_result(NiyuResult *res);

#endif
nd{verbatim}
\section*{File: ./pragya/niyu/niyu_core.c}
egin{verbatim}
#include "niyu_contract.h"
#include "../../include/pragya_avastha.h"

#include <stdlib.h>
#include <string.h>
#include <ctype.h>

/* ============================= */
/* Helper functions              */
/* ============================= */

static int is_operator(char c)
{
    return (c=='+' || c=='-' || c=='*' || c=='/');
}

static int is_digit_char(char c)
{
    return isdigit((unsigned char)c);
}

static int starts_with(const char *text, const char *word)
{
    return strncmp(text, word, strlen(word)) == 0;
}

/* ============================= */
/* Skip spaces                   */
/* ============================= */

static int next_non_space(const char *t, int i)
{
    while (t[i] && isspace((unsigned char)t[i]))
        i++;

    return i;
}


/* ============================= */
/* NIYU THINK ENGINE             */
/* ============================= */

NiyuResult* niyu_think_avastha(PragyaAvastha *a)
{
    if (!a || !a->raw_text)
        return NULL;

    const char *t = a->raw_text;

    NiyuResult *res = malloc(sizeof(NiyuResult));

    if (!res)
        return NULL;

    res->intent = NIYU_INTENT_UNKNOWN;
    res->error_type = NIYU_ERR_VALID;
    res->severity = 0;
    res->error_pos = -1;

    int len = strlen(t);

    if (len == 0)
    {
        res->error_type = NIYU_ERR_EMPTY_INPUT;
        res->severity = 1;
        return res;
    }

    /* ============================= */
    /* INTENT DETECTION              */
    /* ============================= */

    if (starts_with(t, "mavi"))
    {
        res->intent = NIYU_INTENT_DECLARATION;
        return res;
    }

    if (starts_with(t, "bolo"))
    {
        res->intent = NIYU_INTENT_PRINT;
        return res;
    }

    res->intent = NIYU_INTENT_CALCULATION;

    /* ============================= */
    /* ERROR PATTERN DETECTION       */
    /* ============================= */

    for (int i = 0; i < len; i++)
    {
        int j = next_non_space(t, i+1);

        if (!t[j])
            break;

        char c1 = t[i];
        char c2 = t[j];

        /* operator operator */

        if (is_operator(c1) && is_operator(c2))
        {
            res->error_type = NIYU_ERR_EXTRA_OPERATOR;
            res->severity = 2;
            res->error_pos = j;
            return res;
        }

        /* digit digit (missing operator) */

        if (is_digit_char(c1) && is_digit_char(c2))
        {
            if (j > i+1)
            {
                res->error_type = NIYU_ERR_SYNTAX;
                res->severity = 2;
                res->error_pos = j;
                return res;
            }
        }
    }

    /* ============================= */
    /* Missing operand               */
    /* ============================= */

    int last = len-1;

    while (last >= 0 && isspace((unsigned char)t[last]))
        last--;

    if (last >= 0 && is_operator(t[last]))
    {
        res->error_type = NIYU_ERR_MISSING_OPERAND;
        res->severity = 2;
        res->error_pos = last;
        return res;
    }

    /* ============================= */
    /* Undefined variable heuristic  */
    /* ============================= */

    if (isalpha((unsigned char)t[0]))
    {
        res->error_type = NIYU_ERR_UNDEFINED_VARIABLE;
        res->severity = 3;
        res->error_pos = 0;
        return res;
    }

    return res;
}


/* ============================= */
/* FREE RESULT                   */
/* ============================= */

void niyu_free_result(NiyuResult *res)
{
    if (res)
        free(res);
}
nd{verbatim}
\section*{File: ./pragya/niyu/niyu_core.h}
egin{verbatim}
#ifndef NIYU_CORE_H
#define NIYU_CORE_H

#include "niyu_contract.h"

/*
 * NIYU CORE
 *
 * Implementation file will:
 * - define NiyuResult structure
 * - implement internal logic
 *
 * This header exposes NOTHING extra.
 * Future machine-backed engine can replace niyu_core.c
 * without touching any other file.
 */

#endif /* NIYU_CORE_H */
nd{verbatim}
\section*{File: ./include/user_config.h}
egin{verbatim}
#ifndef USER_CONFIG_H
#define USER_CONFIG_H

// Personality mapping
#define P_SAKHI 1
#define P_NIYU  2
#define P_MIRA  3
#define P_KAVYA 4
#define P_SHIRI 5
#define P_ALL   6   // Unified AI

// Mode mapping
#define MODE_INDIA  1   // Hindi-English mix
#define MODE_GLOBAL 2   // Full English

// Developer mode flag (ON by default)
extern int DEV_MODE;

// Current personality (default = Sakhi)
extern int CURRENT_PERSONALITY;

// Current mode (default = INDIA)
extern int CURRENT_MODE;

void config_init();
void config_set_personality(int p);
void config_set_mode(int mode);

#endif
nd{verbatim}
\section*{File: ./include/niyu_engine.h}
egin{verbatim}
#ifndef NIYU_ENGINE_H
#define NIYU_ENGINE_H

int niyu_analyze(const char *input);

#endif
nd{verbatim}
\section*{File: ./include/event.h}
egin{verbatim}
#ifndef EVENT_H
#define EVENT_H

#include "runtime.h"
/*──────────────────────────────────────────────────────────────
 |  SHRIJILANG — EVENT SYSTEM
 |
 |  Purpose:
 |   • Represent what just happened in the system
 |   • Bridge Interpreter → State → AI
 |
 |  Introduced in STEP-6
 |  Extended in STEP-7.3 (COMMAND events)
 *────────────────────────────────────────────────────────────*/

/*──────────────────────────────────────────────────────────────
 |  EVENT TYPES
 *────────────────────────────────────────────────────────────*/
typedef enum {

    EVENT_NONE = 0,

    /* Execution lifecycle */
    EVENT_EXECUTION_START,
    EVENT_EXECUTION_BLOCKED,
    EVENT_EXECUTION_SUCCESS,

    /* Variable & logic */
    EVENT_ASSIGNMENT,
    EVENT_SHUNYA_ACCESS,

    /* Command system (STEP-7.3) */
    EVENT_COMMAND,          /* ls / cd / mkdir / aliases */

    /* Errors */
    EVENT_ERROR,
    EVENT_FATAL_ERROR

} EventType;


/*──────────────────────────────────────────────────────────────
 |  EVENT STRUCTURE
 *────────────────────────────────────────────────────────────*/
typedef struct {

    EventType type;

    /* Optional short context (safe capped) */
    char context[64];

} Event;


/*──────────────────────────────────────────────────────────────
 |  EVENT API
 *────────────────────────────────────────────────────────────*/
void event_fire(EventType type, const char *context);
void event_bind_runtime(ShrijiRuntime *rt);

#endif /* EVENT_H */
nd{verbatim}
\section*{File: ./include/fix_engine.h}
egin{verbatim}
/* RADHE_RADHE_SHREEJI */

/* =========================================================
   SHRIJILANG — EXECUTION ENGINE (CORE CONTRACT)
   ========================================================= */

#ifndef SHRIJI_ENGINE_H
#define SHRIJI_ENGINE_H

#include "ast.h"
#include "value.h"
#include "env.h"
#include "runtime.h"
#include "error.h"

/* =========================================================
   ENGINE CONFIG
   ========================================================= */

/* Maximum number of fix attempts allowed */
#define ENGINE_MAX_FIX_ATTEMPTS 5

/* =========================================================
   ENGINE STATUS
   ========================================================= */

typedef enum
{
    ENGINE_OK = 0,
    ENGINE_PARSE_ERROR,
    ENGINE_RUNTIME_ERROR,
    ENGINE_FIX_FAILED
} EngineStatus;

/* =========================================================
   ENGINE RESULT STRUCTURE
   ========================================================= */

typedef struct
{
    /* ---------- Language Layer ---------- */

    ASTNode *ast;              /* Final AST (after fix if any) */

    /* ---------- Runtime Layer ---------- */

    Value result;              /* Final evaluated result */

    /* ---------- Fix Engine Meta ---------- */

    int was_fixed;             /* 1 if any fix applied */
    int fix_attempts;          /* number of attempts used */
    int confidence_penalty;    /* total penalty from fixes */

    /* ---------- Status ---------- */

    EngineStatus status;       /* final execution status */

} EngineResult;

/* =========================================================
   ENGINE ENTRY POINT
   ========================================================= */

/*
   Main execution pipeline:

   Input
     → Parse
     → Fix (if needed)
     → Re-parse
     → Runtime
     → Result

   NOTE:
   - This function is PURE (no printing ideally in future)
   - KRST should control how it is used
*/
EngineResult shriji_engine_execute(
    const char *input,
    Env *env
);

/* =========================================================
   LOW-LEVEL CONTROL (INTERNAL USE)
   ========================================================= */

/* Single parse attempt (no fix) */
ASTNode *engine_parse_once(const char *input);

/* Apply fix rules */
int engine_apply_fix(
    char *buffer,
    size_t size,
    const ShrijiErrorInfo *err
);

/* =========================================================
   SAFETY UTILITIES
   ========================================================= */

/* Reset result safely */
void engine_result_init(EngineResult *res);

/* Free result safely */
void engine_result_free(EngineResult *res);

/* =========================================================
   LEGACY FIX ENGINE (KRST COMPAT)
   ========================================================= */

ASTNode *language_execute_with_fix(
    const char *input,
    int *was_fixed,
    int *penalty_out
);
typedef enum {
    FIX_NONE = 0,
    FIX_SAFE,
    FIX_DOUBTFUL,
    FIX_REJECT
} FixType;

FixType fix_apply(const char *input, char *output);

#endif /* SHRIJI_ENGINE_H */
nd{verbatim}
\section*{File: ./include/typo_engine.h}
egin{verbatim}
#ifndef SHRIJI_TYPO_ENGINE_H
#define SHRIJI_TYPO_ENGINE_H

/*
    Phase-1 Typo Suggestion Engine
    Only checks for "exit" typo
*/

int shriji_check_typo(
    const char *input,
    char *suggestion,
    int suggestion_size
);

#endif
nd{verbatim}
\section*{File: ./include/smriti_personal.h}
egin{verbatim}
#ifndef SMRITI_PERSONAL_H
#define SMRITI_PERSONAL_H

/* ===============================
   PERSONAL SMRITI API
   =============================== */

/* load from file (startup) */
void smriti_personal_load(void);

/* save one key */
void smriti_personal_set(const char *key, const char *value);

/* get value */
const char* smriti_personal_get(const char *key);

/* clear memory (RAM only) */
void smriti_personal_clear(void);

/* force save to file */
void smriti_personal_save(void);

#endif
nd{verbatim}
\section*{File: ./include/ast.h}
egin{verbatim}
#ifndef AST_H
#define AST_H
#include "../include/token.h"
#ifdef SHRIJI_ENABLE_MASTER_TOKENS
#include "lang/token_master.h"
#endif

/*──────────────────────────────────────────────────────────────
 |  SHRIJILANG — ABSTRACT SYNTAX TREE (AST)
 |
 |  Used by:
 |   • Parser
 |   • Interpreter
 |   • AI Engines (Sakhi / Niyu / Mira / Kavya / Shiri)
 |   • Universal Command Engine (sri_*)
 |   • Future Grammar (block, if, loop, function)
 *────────────────────────────────────────────────────────────*/

/*──────────────────────────────────────────────────────────────
 |  SECTION A — CHAKRA DIAGNOSTIC ENGINE
 *────────────────────────────────────────────────────────────*/
typedef enum {
    CHAKRA_OK = 0,
    CHAKRA_MOOLADHARA,
    CHAKRA_SWADHISTHAN,
    CHAKRA_MANIPUR,
    CHAKRA_ANAHATA,
    CHAKRA_VISHUDDHI,
    CHAKRA_AJNA,
    CHAKRA_SAHASRARA
} ChakraState;


/*──────────────────────────────────────────────────────────────
 |  SECTION B — AST NODE TYPES
 *────────────────────────────────────────────────────────────*/
typedef enum {

    /* Core literals */
    AST_NUMBER,
    AST_IDENTIFIER,
    AST_STRING,
    AST_BOOL,        // true / false
    /* Expressions */
    AST_BINARY,
    AST_UNARY,
    AST_ASSIGNMENT,
    AST_UPDATE,        // x = expression
    AST_LIST,          // [a, b, c]
    AST_INDEX,         // a[0]
    AST_DICT,          // {"a":1, "b":2}
    AST_INDEX_UPDATE,
/* Control / Grammar */
    AST_IF,
    AST_BLOCK,
    AST_PROGRAM,     // script / file root
    AST_IMPORT,      // module import
    AST_LOOP,        // future
    AST_WHILE,       // jabtak loop
    AST_BREAK,       // rukja (break)
    AST_CONTINUE,    // chalu (continue)
    AST_RETURN,      // wapas (return)
    AST_NOT,    //logical NOT (Nahi)
    AST_INC,         // ++
    AST_DEC,         // --

    /* AI Nodes */
    AST_SAKHI,
    AST_NIYU,
    AST_SHIRI,
    AST_MIRA,
    AST_KAVYA,

/* Functions */
    AST_FUNCTION,      // function definition
    AST_CALL,          // function call
    AST_RACHNA,
    /* Universal command (sri_*) */
    AST_COMMAND

} ASTNodeType;


/*─────────────────────────────────────────────────────────────
 | SECTION C — AST NODE STRUCTURE
 *─────────────────────────────────────────────────────────────*/
typedef struct ASTNode {

    ASTNodeType type;

    /*────────────────────────
      Primitive values
    ────────────────────────*/
    double  number_value;
    int  bool_value;     // 1 = true, 0 = false
    char string_value[256];
    char name[128];

    /*────────────────────────
      Binary expression
    ────────────────────────*/
    struct ASTNode *left;
    struct ASTNode *right;
    char op[4];                 // + - * / > < >= <= == !=

    int is_prefix;   /* for ++x / x++ / --x / x-- */

    /*────────────────────────
      Assignment / Command
    ────────────────────────*/
    struct ASTNode *value;
    char command_name[128];

    /*────────────────────────
      IF / ELSE
    ────────────────────────*/
    struct ASTNode *condition;   // agar condition
    struct ASTNode *else_block;  // warna block

    /*────────────────────────
      Block
    ────────────────────────*/
    struct ASTNode **statements;
    int stmt_count;

/*────────────────────────
      Function
    ────────────────────────*/
    char function_name[128];

    struct ASTNode **params;
    int param_count;

    struct ASTNode *body;

/*────────────────────────
  Rachna
────────────────────────*/
char rachna_name[128];
struct ASTNode **rachna_params;
int rachna_param_count;
struct ASTNode *rachna_body;

/*────────────────────────
      List / Index
    ────────────────────────*/
    struct ASTNode **elements;
    int element_count;

    struct ASTNode *index_target;   // a[...]
    struct ASTNode *index_expr;     // ... inside []

    struct ASTNode *index_value;    // a[expr] = value
/*────────────────────────
      Dict / Map
    ────────────────────────*/
    struct ASTNode **dict_keys;
    struct ASTNode **dict_values;
    int dict_count;

    /*────────────────────────
      Call
    ────────────────────────*/
    struct ASTNode **args;
    int arg_count;

    /*────────────────────────
      Return
    ────────────────────────*/
    struct ASTNode *return_expr;

    /*────────────────────────
      Diagnostics
    ────────────────────────*/
    ChakraState chakra_state;

} ASTNode;


/*──────────────────────────────────────────────────────────────
 |  SECTION D — CONSTRUCTORS
 *────────────────────────────────────────────────────────────*/

/* Core */
ASTNode *new_number_node(double value);
ASTNode *new_bool_node(int value);
ASTNode *new_identifier_node(const char *name);
ASTNode *new_string_node(const char *str);

/* Expressions */
ASTNode *new_binary_node(char op, ASTNode *left, ASTNode *right);
ASTNode *new_unary_node(Token op, ASTNode *expr);
ASTNode *new_assignment_node(const char *name, ASTNode *value);
ASTNode *new_update_node(const char *name, ASTNode *value);

ASTNode *new_list_node(ASTNode **elements, int count);
ASTNode *new_index_node(ASTNode *target, ASTNode *index_expr);

ASTNode *new_dict_node(ASTNode **keys, ASTNode **values, int count);

ASTNode *new_index_update_node(ASTNode *target, ASTNode *index_expr, ASTNode *value);

/* AI */
ASTNode *new_sakhi_node(ASTNode *msg);
ASTNode *new_niyu_node(ASTNode *msg);
ASTNode *new_shiri_node(ASTNode *msg);
ASTNode *new_mira_node(ASTNode *msg);
ASTNode *new_kavya_node(ASTNode *msg);

/* Universal command */
ASTNode *new_command_node(const char *cmd_name, ASTNode *arg);
ASTNode *new_import_node(const char *module_name);

/* Block */
ASTNode *new_block_node(ASTNode **stmts, int count);
ASTNode *new_program_node(ASTNode **stmts, int count);

/* Grammar */
ASTNode *new_if_node(ASTNode *condition, ASTNode *then_block);
ASTNode *new_while_node(ASTNode *condition, ASTNode *body);
ASTNode *new_break_node(void);
ASTNode *new_continue_node(void);
ASTNode *new_return_node(ASTNode *expr);

/* Functions */
ASTNode *new_function_node(
    const char *name,
    ASTNode **params,
    int param_count,
    ASTNode *body
);

ASTNode *new_rachna_node(
    const char *name,
    ASTNode **params,
    int param_count,
    ASTNode *body
);
ASTNode *new_call_node(
    const char *name,
    ASTNode **args,
    int arg_count
);

ASTNode *new_return_node(ASTNode *expr);

/* Cleanup */
void ast_free(ASTNode *node);

/* Logical */
ASTNode *new_not_node(ASTNode *expr);

/* Inc / Dec */
ASTNode *new_inc_node(ASTNode *expr, int is_prefix);
ASTNode *new_dec_node(ASTNode *expr, int is_prefix);

#endif /* AST_H */
nd{verbatim}
\section*{File: ./include/ai_intent.h}
egin{verbatim}
#ifndef AI_INTENT_H
#define AI_INTENT_H

/*──────────────────────────────────────────────
  AI INTENT — L3 (KR̥ṢṬ)
  Purpose:
  - User ke words ko canonical intent me normalize karna
  - Routing decision ko simplify karna
  - Interpreter / execution ko bilkul affect NAHI karta
──────────────────────────────────────────────*/

typedef enum {
    AI_INTENT_UNKNOWN = 0,
    AI_INTENT_EXPLAIN,   /* samjhao, explain, batao */
    AI_INTENT_WHY,       /* why, kyun, kyu */
    AI_INTENT_CALC,      /* calculate, solve, math */
    AI_INTENT_EMOTION    /* sad, dukhi, lonely */
} AIIntent;

/* raw text -> detected intent */
AIIntent ai_detect_intent(const char *message);

#endif /* AI_INTENT_H */
nd{verbatim}
\section*{File: ./include/decision_engine.h}
egin{verbatim}
#ifndef DECISION_ENGINE_H
#define DECISION_ENGINE_H

#include "error.h"

typedef enum {
    DECISION_AUTO_FIX,
    DECISION_EDIT,
    DECISION_REJECT
} DecisionType;

/* classification */
int is_safe_error(ShrijiErrorCode code);
int is_ambiguous_error(ShrijiErrorCode code);

/* main decision */
DecisionType shriji_take_decision(
    int confidence,
    int risk,
    ShrijiErrorCode code
);

#endif
nd{verbatim}
\section*{File: ./include/lang/lang_config.h}
egin{verbatim}
#ifndef SHRIJI_LANG_CONFIG_H
#define SHRIJI_LANG_CONFIG_H

#define SHRIJI_ENABLE_MASTER_TOKENS

#endif
nd{verbatim}
\section*{File: ./include/lang/token_master.h}
egin{verbatim}
#ifndef SHRIJI_TOKEN_MASTER_H
#define SHRIJI_TOKEN_MASTER_H

/*
  SHRIJILANG — MASTER TOKEN REGISTRY
  ---------------------------------
  This file is a DESIGN + FUTURE registry.

  RULES:
  1. token.h is the runtime truth
  2. Only tokens present in token.h are ACTIVE
  3. Others are FUTURE / PARKED (no runtime meaning yet)
*/

/* =====================================================
   CORE / ACTIVE TOKENS (MIRROR OF token.h)
   ===================================================== */

typedef enum {

    /* ----- Special ----- */
    MT_TOKEN_EOF = 0,
    MT_TOKEN_ERROR,
    MT_TOKEN_NEWLINE,

    /* ----- Literals ----- */
    MT_TOKEN_NUMBER,
    MT_TOKEN_IDENTIFIER,
    MT_TOKEN_STRING,
    MT_TOKEN_TRUE,
    MT_TOKEN_FALSE,

    /* ----- Core Keywords ----- */
    MT_TOKEN_MAVI,
    MT_TOKEN_AGAR,
    MT_TOKEN_WARNA,
    MT_TOKEN_KAAM,
    MT_TOKEN_JABTAK,
    MT_TOKEN_RUKJA,
    MT_TOKEN_CHALU,
    MT_TOKEN_NAHI,
    MT_TOKEN_BOLO,
    MT_TOKEN_WAPAS,
    MT_TOKEN_IMPORT,

    /* ----- AI / System Commands ----- */
    MT_TOKEN_SAKHI,
    MT_TOKEN_NIYU,
    MT_TOKEN_MIRA,
    MT_TOKEN_KAVYA,
    MT_TOKEN_SHIRI,
    MT_TOKEN_SAMAY,

    /* ----- Operators ----- */
    MT_TOKEN_PLUS,
    MT_TOKEN_MINUS,
    MT_TOKEN_STAR,
    MT_TOKEN_SLASH,
    MT_TOKEN_MOD,

    MT_TOKEN_EQUAL,
    MT_TOKEN_EQEQ,
    MT_TOKEN_NEQ,
    MT_TOKEN_GT,
    MT_TOKEN_LT,
    MT_TOKEN_GTE,
    MT_TOKEN_LTE,

    /* ----- Increment / Decrement ----- */
    MT_TOKEN_PLUSPLUS,
    MT_TOKEN_MINUSMINUS,

    /* ----- Logical ----- */
    MT_TOKEN_AND,
    MT_TOKEN_OR,
    MT_TOKEN_NOT,

    /* ----- Punctuation ----- */
    MT_TOKEN_LEFT_PAREN,
    MT_TOKEN_RIGHT_PAREN,
    MT_TOKEN_LEFT_BRACE,
    MT_TOKEN_RIGHT_BRACE,
    MT_TOKEN_LEFT_BRACKET,
    MT_TOKEN_RIGHT_BRACKET,
    MT_TOKEN_COMMA,
    MT_TOKEN_COLON,

/* =====================================================
   FUTURE / PARKED TOKENS (NOT IN token.h YET)
   ===================================================== */

    /* ----- Declarations ----- */
    MT_TOKEN_CONST,
    MT_TOKEN_LET,

    /* ----- Types ----- */
    MT_TOKEN_TYPE_INT,
    MT_TOKEN_TYPE_FLOAT,
    MT_TOKEN_TYPE_STRING,
    MT_TOKEN_TYPE_BOOL,
    MT_TOKEN_TYPE_LIST,
    MT_TOKEN_TYPE_DICT,
    MT_TOKEN_TYPE_ANY,

    /* ----- Control Flow (future) ----- */
    MT_TOKEN_FOR,
    MT_TOKEN_DO,
    MT_TOKEN_SWITCH,
    MT_TOKEN_CASE,
    MT_TOKEN_DEFAULT,

    /* ----- Assignment Variants ----- */
    MT_TOKEN_PLUSEQ,
    MT_TOKEN_MINUSEQ,
    MT_TOKEN_STAREQ,
    MT_TOKEN_SLASHEQ,
    MT_TOKEN_MODEQ,

    /* ----- Strict / Advanced Compare ----- */
    MT_TOKEN_STRICT_EQ,
    MT_TOKEN_STRICT_NEQ,

    /* ----- Bitwise ----- */
    MT_TOKEN_BIT_AND,
    MT_TOKEN_BIT_OR,
    MT_TOKEN_BIT_XOR,
    MT_TOKEN_BIT_NOT,
    MT_TOKEN_SHIFT_LEFT,
    MT_TOKEN_SHIFT_RIGHT,

    /* ----- Ternary ----- */
    MT_TOKEN_QUESTION,

    /* ----- Structure (future sugar) ----- */
    MT_TOKEN_DOT,
    MT_TOKEN_ARROW,

    /* ----- Data Structures ----- */
    MT_TOKEN_LIST,
    MT_TOKEN_DICT,
    MT_TOKEN_SET,
    MT_TOKEN_TUPLE,

    /* ----- Sentinel ----- */
    MT_TOKEN__COUNT

} ShrijiMasterToken;

#endif /* SHRIJI_TOKEN_MASTER_H */
nd{verbatim}
\section*{File: ./include/lang/grammar_core.h}
egin{verbatim}
#ifndef SHRIJI_LANG_GRAMMAR_CORE_H
#define SHRIJI_LANG_GRAMMAR_CORE_H

#include "../token.h"

/*
 * Phase-1 Grammar Core
 * Only classification — no parsing logic
 */

typedef enum {
    GRAMMAR_STATEMENT,
    GRAMMAR_EXPRESSION,
    GRAMMAR_OPERATOR,
    GRAMMAR_LITERAL,

    GRAMMAR_INCDEC,     /* ++ -- */

    GRAMMAR_RESERVED,
    GRAMMAR_UNKNOWN
} GrammarKind;

/* Query API */
GrammarKind shriji_grammar_kind(TokenType t);
const char *shriji_grammar_kind_name(GrammarKind k);

#endif /* SHRIJI_LANG_GRAMMAR_CORE_H */
nd{verbatim}
\section*{File: ./include/lang/grammar.h}
egin{verbatim}
#ifndef SHRIJI_LANG_GRAMMAR_H
#define SHRIJI_LANG_GRAMMAR_H

/* Grammar extension registry
 * Phase-1: only hooks, no grammar rules
 */

/* Called once during parser init */
void shriji_register_grammar_core(void);

/* Optional future grammars */
void shriji_register_grammar_logic(void);
void shriji_register_grammar_incdec(void);
void shriji_register_grammar_loop(void);
void shriji_register_grammar_func(void);

#endif /* SHRIJI_LANG_GRAMMAR_H */
nd{verbatim}
\section*{File: ./include/value.h}
egin{verbatim}
#ifndef VALUE_H
#define VALUE_H

#include <string.h>
#include <stdlib.h>

#include "ast.h"

/*──────────────────────────────────────────────
  VALUE TYPES
──────────────────────────────────────────────*/
typedef enum {
    VAL_NULL = 0,
    VAL_NUMBER,
    VAL_STRING,
    VAL_BOOL,
    VAL_LIST,
    VAL_FUNCTION,
    VAL_DICT
} ValueType;

typedef struct Value {
    ValueType type;

    double number;
    char *string;
    int boolean;
    /* ✅ list value */
    struct Value *items;
    int count;

    /* ✅ function value */
    ASTNode *function;

    /* ✅ dict value */
    struct Value *dict_keys;
    struct Value *dict_values;
    int dict_count;

} Value;

/*──────────────────────────────────────────────
  CONSTRUCTORS
──────────────────────────────────────────────*/
static inline Value value_null(void)
{
    Value v;
    v.type = VAL_NULL;

    v.number = 0;
    v.string = NULL;

    v.items = NULL;
    v.count = 0;

    v.function = NULL;

    v.dict_keys = NULL;
    v.dict_values = NULL;
    v.dict_count = 0;

    return v;
}

static inline Value value_number(double n)
{
    Value v = value_null();
    v.type = VAL_NUMBER;
    v.number = n;
    return v;
}

static inline Value value_string(const char *s)
{
    Value v = value_null();
    v.type = VAL_STRING;

    if (!s) {
        v.string = NULL;
        return v;
    }

    v.string = (char *)malloc(strlen(s) + 1);
    if (!v.string) return value_null();

    strcpy(v.string, s);
    return v;
}

/* BOOL SUPPORT */
static inline Value value_bool(int b)
{
    Value v = value_null();
    v.type = VAL_BOOL;
    v.boolean = b ? 1 : 0;
    return v;
}

static inline int value_is_truthy(Value v)
{
    switch (v.type) {
        case VAL_BOOL:   return v.boolean;
        case VAL_NUMBER: return v.number != 0;
        case VAL_STRING: return v.string && v.string[0] != '\0';
        case VAL_NULL:   return 0;
        default:         return 1;  // list, dict, function → truthy
    }
}


/* ✅ list constructor */
static inline Value value_list(Value *items, int count)
{
    Value v = value_null();
    v.type = VAL_LIST;
    v.items = items;
    v.count = count;
    return v;
}

/* ✅ function constructor */
static inline Value value_function(ASTNode *fn)
{
    Value v = value_null();
    v.type = VAL_FUNCTION;
    v.function = fn;
    return v;
}

/* ✅ dict constructor */
static inline Value value_dict(Value *keys, Value *values, int count)
{
    Value v = value_null();
    v.type = VAL_DICT;

    v.dict_keys = keys;
    v.dict_values = values;
    v.dict_count = count;

    return v;
}

/*──────────────────────────────────────────────
  COPY
──────────────────────────────────────────────*/
static inline Value value_copy(Value src)
{
    if (src.type == VAL_STRING) {
        return value_string(src.string ? src.string : "");
    }

    if (src.type == VAL_NUMBER) {
        return value_number(src.number);
    }

    if (src.type == VAL_BOOL) {
         return value_bool(src.boolean);
    }


    if (src.type == VAL_LIST) {

        if (src.count <= 0 || !src.items) {
            return value_list(NULL, 0);
        }

        Value *items = (Value *)malloc(sizeof(Value) * src.count);
        if (!items) return value_null();

        for (int i = 0; i < src.count; i++) {
            items[i] = value_copy(src.items[i]);
        }

        return value_list(items, src.count);
    }

    if (src.type == VAL_FUNCTION) {
        /* function node pointer copy (AST is owned elsewhere) */
        return value_function(src.function);
    }

    if (src.type == VAL_DICT) {

        if (src.dict_count <= 0 || !src.dict_keys || !src.dict_values) {
            return value_dict(NULL, NULL, 0);
        }

        Value *keys = (Value *)malloc(sizeof(Value) * src.dict_count);
        Value *values = (Value *)malloc(sizeof(Value) * src.dict_count);

        if (!keys || !values) {
            if (keys) free(keys);
            if (values) free(values);
            return value_null();
        }

        for (int i = 0; i < src.dict_count; i++) {
            keys[i] = value_copy(src.dict_keys[i]);
            values[i] = value_copy(src.dict_values[i]);
        }

        return value_dict(keys, values, src.dict_count);
    }

    return value_null();
}

/*──────────────────────────────────────────────
  FREE
──────────────────────────────────────────────*/
static inline void value_free(Value *v)
{
    if (!v) return;

    if (v->type == VAL_STRING && v->string) {
        free(v->string);
        v->string = NULL;
    }

    if (v->type == VAL_LIST && v->items) {
        for (int i = 0; i < v->count; i++) {
            value_free(&v->items[i]);
        }
        free(v->items);
        v->items = NULL;
        v->count = 0;
    }

    if (v->type == VAL_DICT) {
        if (v->dict_keys) {
            for (int i = 0; i < v->dict_count; i++) {
                value_free(&v->dict_keys[i]);
            }
            free(v->dict_keys);
            v->dict_keys = NULL;
        }

        if (v->dict_values) {
            for (int i = 0; i < v->dict_count; i++) {
                value_free(&v->dict_values[i]);
            }
            free(v->dict_values);
            v->dict_values = NULL;
        }

        v->dict_count = 0;
    }

    /* function is not freed here */
    v->function = NULL;

    /* reset */
    v->type = VAL_NULL;
    v->number = 0.0;
    v->string = NULL;

    v->items = NULL;
    v->count = 0;

    v->dict_keys = NULL;
    v->dict_values = NULL;
    v->dict_count = 0;
}


#endif
nd{verbatim}
\section*{File: ./include/log_mode.h}
egin{verbatim}
#ifndef LOG_MODE_H
#define LOG_MODE_H

#include <stdio.h>
#include "user_config.h"

/* ---------------------------------------------------------
 * LOGGING MACROS
 * DEV_MODE=1  => verbose prints
 * DEV_MODE=0  => silent
 * --------------------------------------------------------- */

#define NIYU_LOG(...)  do { if (DEV_MODE) printf(__VA_ARGS__); } while (0)
#define SAKHI_LOG(...) do { if (DEV_MODE) printf(__VA_ARGS__); } while (0)
#define MIRA_LOG(...)  do { if (DEV_MODE) printf(__VA_ARGS__); } while (0)
#define KAVYA_LOG(...) do { if (DEV_MODE) printf(__VA_ARGS__); } while (0)

#endif
nd{verbatim}
\section*{File: ./include/ai_router.h}
egin{verbatim}
#ifndef SHRIJILANG_AI_ROUTER_H
#define SHRIJILANG_AI_ROUTER_H

#include "../pragya/contracts/l3_response.h"

/*──────────────────────────────────────────────
  SHRIJILANG AI ROUTER — PUBLIC INTERFACE (L3)
──────────────────────────────────────────────*/

typedef enum {
    SHRIJI_LANG_AUTO = 0,
    SHRIJI_LANG_HINDI_MIX,
    SHRIJI_LANG_ENGLISH
} ShrijiLangHint;

typedef enum {
    SHRIJI_SRC_FILE = 0,
    SHRIJI_SRC_REPL
} ShrijiSource;

typedef struct {
    const char      *text;
    ShrijiLangHint   lang;
    ShrijiSource     source;
} ShrijiBridgePacket;

/* NEW: L3 dispatcher */
L3Response ai_router_dispatch(const ShrijiBridgePacket *packet);

/* old compatibility */
void ai_router_process(const char *message);

#endif /* SHRIJILANG_AI_ROUTER_H */
nd{verbatim}
\section*{File: ./include/atma_lekha.h}
egin{verbatim}
#ifndef SHRIJILANG_ATMA_LEKHA_H
#define SHRIJILANG_ATMA_LEKHA_H

/*──────────────────────────────────────────────────────────────
  ATMA-LEKHA — SYSTEM SELF AUDIT (L3-C3)

  Purpose:
  - Record system-side failures & misrouting
  - NO user data
  - Silent by default
  - No behavior change

  This is NOT Smriti.
──────────────────────────────────────────────────────────────*/

typedef enum {
    ATMA_ROUTER_FALLBACK = 1,
    ATMA_INTENT_UNKNOWN,
    ATMA_COMMAND_UNKNOWN,
    ATMA_INTERNAL_MISROUTE
} AtmaEventType;

/* initialize system self-audit */
void atma_lekha_init(void);

/* record an event (silent) */
void atma_lekha_note(AtmaEventType type,
                     const char *component,
                     const char *detail);

#endif
nd{verbatim}
\section*{File: ./include/error.h}
egin{verbatim}
#ifndef ERROR_H
#define ERROR_H

#include "../include/token.h"

typedef enum {

    E_UNKNOWN = 0,

    E_ASSIGN_01,
    E_ASSIGN_02,

    E_PARSE_01,
    E_PARSE_02,

    E_PARSE_EXTRA_OPERATOR,
    E_PARSE_MISSING_OPERATOR,
    E_PARSE_DOUBLE_OPERATOR,
    E_PARSE_OPERATOR_START,
    E_PARSE_OPERATOR_END,
    E_PARSE_OPERATOR_CHAIN,
    E_PARSE_MISSING_OPERAND,

    E_PARSE_BRACKET_MISSING,
    E_PARSE_BRACKET_EXTRA,
    E_PARSE_UNMATCHED_PAREN,
    E_PARSE_UNMATCHED_BRACKET,
    E_PARSE_UNMATCHED_BRACE,

    E_PARSE_INVALID_TOKEN,
    E_PARSE_EMPTY_EXPRESSION,

    E_FUNCTION_NAME_INVALID,
    E_FUNCTION_PARAM_INVALID,

    E_IMPORT_PATH_INVALID,

    E_LIST_SYNTAX_ERROR,
    E_DICT_KEY_INVALID,
    E_DICT_SYNTAX_ERROR,

    E_IF_01,

    E_RUNTIME_01,
    E_RUNTIME_DIV_ZERO,
    E_RUNTIME_TYPE_MISMATCH,
    E_RUNTIME_LOOP_LIMIT,
    E_RUNTIME_INDEX_ERROR

} ShrijiErrorCode;


/*──────────────────────────────────────────────*/

typedef struct {

    ShrijiErrorCode code;

    const char *context;
    const char *message;
    const char *hint;

    const char *function;   // NEW (optional)

    int has_location;
    int line;
    int col;

} ShrijiErrorInfo;


/*──────────────────────────────────────────────*/

/* OLD STYLE (5 ARGUMENT) */
void shriji_error_at(
    Token tok,
    ShrijiErrorCode code,
    const char *context,
    const char *message,
    const char *hint
);

/* NEW STYLE (6 ARGUMENT) */
void shriji_error_at_full(
    Token tok,
    ShrijiErrorCode code,
    const char *context,
    const char *function,
    const char *message,
    const char *hint
);

const ShrijiErrorInfo *shriji_last_error(void);

extern int error_reported;
extern ShrijiErrorInfo LAST_ERROR;

/*──────────────────────────────────────────────*/

typedef enum {
    ERROR_MODE_IMMEDIATE,
    ERROR_MODE_SILENT,
    ERROR_MODE_COLLECT
} ShrijiErrorMode;

void shriji_set_error_mode(ShrijiErrorMode mode);

void shriji_print_all_errors(void);

void shriji_error(
    ShrijiErrorCode code,
    const char *context,
    const char *message,
    const char *hint
);

#endif
nd{verbatim}
\section*{File: ./include/token_list.h}
egin{verbatim}
// SHRIJI_RADHR RADHR

#ifndef SHRIJI_TOKEN_LIST_H
#define SHRIJI_TOKEN_LIST_H

/*
──────────────────────────────────────────────
 ShrijiLang Token System V2
 Dynamic Token Storage (Compiler Grade)
──────────────────────────────────────────────
*/

#include "token.h"

/*──────────────────────────────────────────────
  TOKEN LIST STRUCTURE
──────────────────────────────────────────────*/

typedef struct
{
    Token *data;      // token array
    int count;        // number of tokens used
    int capacity;     // allocated memory

} TokenList;


/*──────────────────────────────────────────────
  TOKEN LIST API
──────────────────────────────────────────────*/

/* create new list */
TokenList token_list_create(void);

/* add token */
void token_list_push(TokenList *list, Token tok);

/* get token */
Token token_list_get(TokenList *list, int index);

/* number of tokens */
int token_list_size(TokenList *list);

/* free memory */
void token_list_free(TokenList *list);


#endif
nd{verbatim}
\section*{File: ./include/translator.h}
egin{verbatim}
#ifndef TRANSLATOR_H
#define TRANSLATOR_H

// Translate English/world command → ShrijiLang keyword
const char* translate_keyword(const char *word);

// Autocorrect system for mistyped user commands
const char* mira_autocorrect(const char *word);

// Extract first word from input (helper for parser)
char* get_first_word(const char *source);

#endif
nd{verbatim}
\section*{File: ./include/nirman_flow_engine.h}
egin{verbatim}
#ifndef NIRMAN_FLOW_ENGINE_H
#define NIRMAN_FLOW_ENGINE_H

void nirman_flow_process(char *line);

#endif
nd{verbatim}
\section*{File: ./include/nirman.h}
egin{verbatim}
#ifndef NIRMAN_H
#define NIRMAN_H

/* ===============================
   CONTEXT STRUCT
   =============================== */
typedef struct {
    int active;           /* nirman mode ON/OFF */
    int flow_stage;       /* state machine stage */
    char goal[256];       /* user goal */
} NirmanContext;

/* ===============================
   GLOBAL ACCESS
   =============================== */
extern NirmanContext NIRMAN_CTX;

/* ===============================
   LIFECYCLE
   =============================== */
void nirman_start(void);
void nirman_stop(void);
int nirman_is_active(void);

/* ===============================
   INTENT ENGINE
   =============================== */
int nirman_detect_intent(const char *input);

#endif
nd{verbatim}
\section*{File: ./include/pragya_avastha.h}
egin{verbatim}
#ifndef PRAGYA_AVASTHA_H
#define PRAGYA_AVASTHA_H

#include "ast.h"
#include "error.h"
#include "value.h"

/*
   PragyaAvastha

   Shriji system ki vartamaan sthiti.
   Yeh packet KRST se Pragya tak jata hai
   aur fir Niyu analysis ke liye use karti hai.
*/

typedef struct {

/* User input */
const char *raw_text;

/*  NEW: input mode */
int is_program_mode;

/* NEW: correction system */
int has_correction;
char corrected_text[256];
int stop_execution;

    /* Parsed structure */
    ASTNode *ast;

    /* Error state */
    ShrijiErrorCode error_code;

    int error_pos;
    int error_len;

    /* KRST metrics */
    int confidence;
    int risk;

    int teach_level;
    int tone;

    /* Runtime result */
    Value runtime_result;

} PragyaAvastha;

#endif
nd{verbatim}
\section*{File: ./include/example_builder.h}
egin{verbatim}
#ifndef SHRIJI_EXAMPLE_BUILDER_H
#define SHRIJI_EXAMPLE_BUILDER_H

void shriji_build_example(
    const char *input,
    char *out,
    int size
);

#endif
nd{verbatim}
\section*{File: ./include/normalize_engine.h}
egin{verbatim}
#ifndef NORMALIZE_ENGINE_H
#define NORMALIZE_ENGINE_H

void normalize_input(char *input);

#endif
nd{verbatim}
\section*{File: ./include/state.h}
egin{verbatim}
#ifndef STATE_H
#define STATE_H

/*──────────────────────────────────────────────────────────────
 |  SHRIJILANG — EXECUTION STATE SYSTEM
 |
 |  Purpose:
 |   • Track current execution condition
 |   • Provide memory of risk, confidence, and errors
 |
 |  Used by:
 |   • Interpreter
 |   • Decision Engine (future)
 *────────────────────────────────────────────────────────────*/

/*──────────────────────────────────────────────────────────────
 |  SAFETY LEVEL — how risky the current situation is
 *────────────────────────────────────────────────────────────*/
typedef enum {
    STATE_SAFE = 0,
    STATE_ALERT,
    STATE_RISK,
    STATE_CRITICAL
} SafetyLevel;


/*──────────────────────────────────────────────────────────────
 |  CONFIDENCE LEVEL — how sure the system is
 *────────────────────────────────────────────────────────────*/
typedef enum {
    CONFIDENCE_HIGH = 0,
    CONFIDENCE_MEDIUM,
    CONFIDENCE_LOW
} ConfidenceLevel;


/*──────────────────────────────────────────────────────────────
 |  HUMAN DEPENDENCY — when human input is required
 *────────────────────────────────────────────────────────────*/
typedef enum {
    HUMAN_AUTO_OK = 0,
    HUMAN_CONFIRM_REQUIRED,
    HUMAN_MANDATORY
} HumanDependency;


/*──────────────────────────────────────────────────────────────
 |  STATE STRUCTURE — complete execution snapshot
 *────────────────────────────────────────────────────────────*/
typedef struct {
    SafetyLevel      safety;
    ConfidenceLevel  confidence;
    int              error_pressure;
    HumanDependency  human;
long long max_steps;
long long steps_used;

} ExecutionState;


/*──────────────────────────────────────────────────────────────
 |  STATE API (implemented in state.c)
 *────────────────────────────────────────────────────────────*/
void state_init(ExecutionState *state);
void state_reset(ExecutionState *state);
void state_on_error(ExecutionState *state);
void state_on_success(ExecutionState *state);

#endif /* STATE_H */
nd{verbatim}
\section*{File: ./include/fix_rules.h}
egin{verbatim}
#ifndef SHRIJI_FIX_RULES_H
#define SHRIJI_FIX_RULES_H

#include <stddef.h>
#include "error.h"

/* ─────────────────────────────────────────────
   FIX CATEGORY
───────────────────────────────────────────── */

typedef enum {
    FIXCAT_NONE = 0,
    FIXCAT_MISSING_VALUE,
    FIXCAT_BAD_TOKEN,
    FIXCAT_MISSING_OPERATOR,
    FIXCAT_MISSING_IDENTIFIER,
    FIXCAT_EXTRA_OPERATOR
} FixCategory;

/* ─────────────────────────────────────────────
   FIX SAFETY LEVEL
───────────────────────────────────────────── */

typedef enum {
    FIX_SAFE_DETERMINISTIC = 0,   /* 100% sure → auto apply */
    FIX_SAFE_STRUCTURAL,          /* safe but not guaranteed */
    FIX_RISKY_CONTEXTUAL          /* user decision required */
} FixSafetyLevel;

/* ─────────────────────────────────────────────
   FIX FUNCTION SIGNATURE
───────────────────────────────────────────── */

typedef int (*FixApplyFunction)(
    char *buffer,
    size_t buffer_size,
    const ShrijiErrorInfo *err
);

/* ─────────────────────────────────────────────
   FIX RULE STRUCTURE
───────────────────────────────────────────── */

typedef struct {

    ShrijiErrorCode error_code;

    FixCategory category;
    FixSafetyLevel safety;

    int confidence_penalty;
    int max_attempt;
    int reparse_required;

    FixApplyFunction apply_fix;

} FixRule;

/* ─────────────────────────────────────────────
   RULE LOOKUP
───────────────────────────────────────────── */

FixRule *shriji_get_rule_for_error(ShrijiErrorCode code);

#endif
nd{verbatim}
\section*{File: ./include/sakhi_engine.h}
egin{verbatim}
#ifndef SAKHI_ENGINE_H
#define SAKHI_ENGINE_H

void sakhi_speak(const char *msg);

#endif
nd{verbatim}
\section*{File: ./include/env.h}
egin{verbatim}
#ifndef ENV_H
#define ENV_H

#include "value.h"

/*──────────────────────────────────────────────
 | SHRIJILANG — ENVIRONMENT (VARIABLE STORE)
 |
 | Env = Scope stack
 |  - Global scope always exists
 |  - Blocks { } push/pop scopes
 |
 | Rules:
 |  - env_set()     : create/update in CURRENT scope only
 |  - env_update()  : update in nearest existing scope (top → bottom)
 |  - env_get()     : read from nearest scope (top → bottom)
 |  - env_exists()  : check existence (top → bottom)
 *──────────────────────────────────────────────*/

#define ENV_MAX_VARS     1024
#define ENV_MAX_SCOPES   64

/* one variable */
typedef struct {
    char name[64];
    Value value;
} EnvVar;

/* one scope frame */
typedef struct {
    EnvVar vars[ENV_MAX_VARS];
    int count;
} EnvFrame;

/* environment = stack of frames */
typedef struct {
    EnvFrame frames[ENV_MAX_SCOPES];
    int depth; /* >= 1 always (global scope) */
} Env;

/* create */
Env *new_env(void);

/* scope control */
void env_push_scope(Env *env);
void env_pop_scope(Env *env);

/* variables */
void env_set(Env *env, const char *name, Value v);       /* current scope only */
void env_update(Env *env, const char *name, Value v);    /* nearest scope update */
Value env_get(Env *env, const char *name);               /* nearest scope read */
int env_exists(Env *env, const char *name);

#endif /* ENV_H */
nd{verbatim}
\section*{File: ./include/error_intelligence.h}
egin{verbatim}
#ifndef SHRIJI_ERROR_INTELLIGENCE_H
#define SHRIJI_ERROR_INTELLIGENCE_H

#include "error.h"
#include "../pragya/niyu/niyu_contract.h"
#include "../include/pragya_avastha.h"

/* error clarity levels */

#define ERR_CONFUSED 0
#define ERR_PARTIAL  1
#define ERR_CLEAR    2


/* understanding engine */

int shriji_understand_error(
    ShrijiErrorCode code,
    const NiyuResult *niyu
);


/* pointer extractor */

void shriji_extract_pointer(
    const ShrijiErrorInfo *err,
    int *pos,
    int *len
);


/* main intelligence entry */

void shriji_error_intelligence(
    PragyaAvastha *avastha,
    const ShrijiErrorInfo *err,
    const NiyuResult *niyu
);

int shriji_get_example(const char *input, int pos, char *out);
#endif
nd{verbatim}
\section*{File: ./include/token.h}
egin{verbatim}
#ifndef TOKEN_H
#define TOKEN_H

/*──────────────────────────────────────────────
  SHRIJILANG — TOKEN SYSTEM (LOCKED)
  Now supports: line/col for professional errors
──────────────────────────────────────────────*/

typedef enum {

    TOKEN_EOF = 0,

    /* Literals */
    TOKEN_NUMBER,
    TOKEN_IDENTIFIER,
    TOKEN_STRING,
    TOKEN_TRUE,
    TOKEN_FALSE,

    /* Keywords */
    TOKEN_MAVI,
    TOKEN_AGAR,
    TOKEN_WARNA,
    TOKEN_KAAM,
    TOKEN_JABTAK,
    TOKEN_RUKJA,
    TOKEN_CHALU,
    TOKEN_NAHI,
    TOKEN_BOLO,

    TOKEN_SAKHI,
    TOKEN_NIYU,
    TOKEN_MIRA,
    TOKEN_KAVYA,
    TOKEN_SHIRI,

    TOKEN_SAMAY,
    TOKEN_WAPAS,

    TOKEN_IMPORT,
    TOKEN_RACHNA,

    /* Operators */
    TOKEN_PLUS,
    TOKEN_MINUS,
    TOKEN_STAR,
    TOKEN_SLASH,
    TOKEN_MOD,

    TOKEN_EQUAL,
    TOKEN_EQEQ,
    TOKEN_NEQ,
    TOKEN_GT,
    TOKEN_LT,
    TOKEN_GTE,
    TOKEN_LTE,

/* Increment / Decrement */
TOKEN_PLUSPLUS,     // ++
TOKEN_MINUSMINUS,   // --

/* Logical Operators */
TOKEN_AND,          // &&
TOKEN_OR,           // ||
TOKEN_NOT,          // !

    /* Punctuation */
    TOKEN_LEFT_PAREN,
    TOKEN_RIGHT_PAREN,
    TOKEN_LEFT_BRACKET,
    TOKEN_RIGHT_BRACKET,
    TOKEN_LEFT_BRACE,
    TOKEN_RIGHT_BRACE,
    TOKEN_COMMA,
    TOKEN_COLON,
    TOKEN_NEWLINE,

    /* Special */
    TOKEN_ERROR

} TokenType;

typedef struct {
    TokenType type;
    const char *start;
    int length;
    int line;
    int col;
} Token;

/* Scanner API */
void init_tokenizer(const char *source);
Token scan_token(void);

//* #ifdef SHRIJI_ENABLE_MASTER_TOKENS
//* #include "lang/token_master.h"
//* #endif

#endif // TOKEN_H
nd{verbatim}
\section*{File: ./include/fix_interactive.h}
egin{verbatim}
#ifndef SHRIJI_FIX_INTERACTIVE_H
#define SHRIJI_FIX_INTERACTIVE_H

#include <stddef.h>
#include "error.h"


/* =========================================
   INTERACTIVE OPERATOR FIX
   Example: 5 5 → 5 + 5
   ========================================= */

int shriji_interactive_operator_fix(
    const char *input,
    char *output,
    int size
);


/* =========================================
   INTERACTIVE BRACKET FIX
   Example: (5+6 → (5+6)
            (5+6 → 5+6
   ========================================= */

int shriji_interactive_bracket_fix(
    const char *input,
    char *output,
    int size
);

int shriji_interactive_double_operator_fix(
    const char *input,
    char *output,
    int size
);

int fix_missing_value_after_equal_interactive(
    char *buffer,
    size_t buffer_size,
    const ShrijiErrorInfo *err
);
#endif
nd{verbatim}
\section*{File: ./include/mira_engine.h}
egin{verbatim}
#ifndef MIRA_ENGINE_H
#define MIRA_ENGINE_H

void mira_teach(const char *msg);

#endif
nd{verbatim}
\section*{File: ./include/interpreter.h}
egin{verbatim}
#ifndef INTERPRETER_H
#define INTERPRETER_H

#include "env.h"
#include "ast.h"
#include "value.h"
#include "runtime.h"   // 🔥 NEW — runtime struct access

//──────────────────────────────────────────────────────────────
// SHRIJI TRUTH DIARY (SINGLE SOURCE OF TRUTH)
//──────────────────────────────────────────────────────────────
typedef struct {
    char last_var[64];
    long long last_value;
    int has_assignment;
    int had_runtime_error;
} ShrijiTruthDiary;

// defined in src/interpreter.c
extern ShrijiTruthDiary SHRIJI_DIARY;


//──────────────────────────────────────────────────────────────
// MAIN EVALUATOR
//──────────────────────────────────────────────────────────────
//
// 🔥 IMPORTANT CHANGE:
// eval now receives runtime pointer
//
Value eval(ASTNode *node, Env *env, ShrijiRuntime *rt);


//──────────────────────────────────────────────────────────────
// PROGRAM ENTRY
//──────────────────────────────────────────────────────────────
//
// 🔥 IMPORTANT CHANGE:
// run_program now receives runtime pointer
//
Value run_program(ASTNode *program, Env *env, ShrijiRuntime *rt);


#endif /* INTERPRETER_H */
nd{verbatim}
\section*{File: ./include/krst_router.h}
egin{verbatim}
#ifndef KRST_ROUTER_H
#define KRST_ROUTER_H

#include "../krst/krst_types.h"

void krst_route_request(KRSTRequest *req);

#endif
nd{verbatim}
\section*{File: ./include/parser.h}
egin{verbatim}
#ifndef PARSER_H
#define PARSER_H

//──────────────────────────────────────────────────────────────
//  SHRIJILANG PARSER — HEADER (UNIVERSE-CLASS VERSION)
//
//  This module converts raw source code strings into AST Nodes.
//  
//  Core features supported in V1:
//    ✔ mavi assignment       → mavi x = 10 + 20
//    ✔ arithmetic expressions
//    ✔ sakhi "msg"
//    ✔ niyu "msg"
//    ✔ shiri "msg"
//    ✔ mira "msg"
//    ✔ kavya "msg"
//    ✔ sri_<command> "msg"
//    ✔ GAURI Super-List Engine (V1.5)
//
//  Upcoming in V2:
//    ✔ multi-statement blocks { }
//    ✔ loops
//    ✔ conditions
//──────────────────────────────────────────────────────────────

#include "ast.h"

//──────────────────────────────────────────────────────────────
//  PARSE ENTRY POINT
//──────────────────────────────────────────────────────────────
//
//  parse_program()
//      Converts a full input string into a root AST node.
//      Returns NULL if syntax error.
//
ASTNode *parse_program(const char *source);


//──────────────────────────────────────────────────────────────
//  GAURI DATA-STRUCTURE PARSER DECLARATIONS (NEW)
//──────────────────────────────────────────────────────────────
//
//  Gauri commands allow creation + mutation of "living lists"
//  Example:
//      gauri_list x = "[1,2,3]"
//      gauri_push x 50
//      gauri_pop x
//      gauri_get x 2
//      gauri_set x 2 99
//      gauri_sum x
//      gauri_sort x
//      gauri_reverse x
//
ASTNode* parse_gauri_list();
ASTNode* parse_gauri_push();
ASTNode* parse_gauri_pop();
ASTNode* parse_gauri_get();
ASTNode* parse_gauri_set();
ASTNode* parse_gauri_sum();
ASTNode* parse_gauri_sort();
ASTNode* parse_gauri_reverse();

#endif
nd{verbatim}
\section*{File: ./include/runtime.h}
egin{verbatim}
#ifndef RUNTIME_H
#define RUNTIME_H

#include "state.h"
#include "value.h"

/*──────────────────────────────────────────────
 | SHRIJI RUNTIME — ISOLATED EXECUTION BRAIN
 *──────────────────────────────────────────────*/

typedef struct {

    ExecutionState state;

    int break_flag;
    int continue_flag;
    int loop_depth;

    int return_flag;
    Value return_value;

    int call_depth;

    int error_flag;   // NEW

} ShrijiRuntime;

/* GLOBAL POINTER */
extern ShrijiRuntime *current_runtime;

/* Init once */
void runtime_init(ShrijiRuntime *rt);

/* Reset before each run */
void runtime_reset(ShrijiRuntime *rt);

#endif
nd{verbatim}
\section*{File: ./include/smriti.h}
egin{verbatim}
#ifndef SMRITI_H
#define SMRITI_H

/* lifecycle */
void smriti_init(void);

/* math continuation (existing behavior) */
int  smriti_has_continuation(void);
void smriti_set_continuation(double value, char op);
void smriti_clear_continuation(void);

double smriti_get_continuation_value(void);
char   smriti_get_continuation_op(void);

/* ===== STEP-11: context (NO logic, NO decision) ===== */

/* turn tracking */
void smriti_next_turn(void);
unsigned long smriti_get_turn(void);

/* last input snapshot */
void smriti_set_last_input(const char *text);
const char* smriti_get_last_input(void);

/* last intent snapshot (L3 intent enum value passed as int) */
void smriti_set_last_intent(int intent);
int  smriti_get_last_intent(void);

#endif
nd{verbatim}
\section*{File: ./include/smriti_session.h}
egin{verbatim}
#ifndef SMRITI_SESSION_H
#define SMRITI_SESSION_H

/* ===============================
   SESSION MEMORY (SHORT TERM)
   =============================== */

#include "value.h"

/* last input */
void smriti_session_set_last_input(const char *text);
const char* smriti_session_get_last_input(void);

/* last output */
void smriti_session_set_last_output(const char *text);
const char* smriti_session_get_last_output(void);

/* last evaluated value (UPGRADED) */
void smriti_session_set_last(Value v);
Value smriti_session_get_last(void);

/* clear session */
void smriti_session_clear(void);

#endif
nd{verbatim}
\section*{File: ./include/command.h}
egin{verbatim}
#ifndef COMMAND_H
#define COMMAND_H

/*──────────────────────────────────────────────────────────────
 | INTERNAL COMMAND IDENTIFIERS (TRUTH)
 *──────────────────────────────────────────────────────────────*/
typedef enum {

    CMD_UNKNOWN = 0,

    /* Filesystem */
    CMD_LIST,        /* ls */
    CMD_CD,          /* cd */
    CMD_MKDIR,       /* mkdir */

    /* Output / I/O */
    CMD_BOLO,        /* bolo (print) */
    CMD_SAMAY,

/* AI Modules */
    CMD_SAKHI,
    CMD_NIYU,
    CMD_MIRA,
    CMD_KAVYA,
    CMD_SHIRI,

    /* Future expansion */
    CMD_DELETE,
    CMD_COPY,
    CMD_MOVE



} CommandType;


/*──────────────────────────────────────────────────────────────
 | COMMAND ACCESS POLICY
 *──────────────────────────────────────────────────────────────*/
typedef enum {
    CMD_PUBLIC = 1,     /* visible to everyone */
    CMD_LOCKED = 0      /* reserved / internal */
} CommandAccess;


/*──────────────────────────────────────────────────────────────
 | COMMAND ALIAS (1008 NAMES LIVE HERE)
 *──────────────────────────────────────────────────────────────*/
typedef struct {
    const char     *alias;     /* user-facing name */
    CommandType     cmd;       /* internal truth */
    CommandAccess   access;    /* public / locked */
} CommandAlias;


/*──────────────────────────────────────────────────────────────
 | API
 *──────────────────────────────────────────────────────────────*/
CommandType resolve_command(const char *name);
const char *command_name(CommandType cmd);
int execute_command(CommandType cmd);

#endif /* COMMAND_H */
nd{verbatim}
\section*{File: ./release/common/docs/README.txt}
egin{verbatim}
ShrijiLang v1.0 (Core Stable)

──────────────────────────────────────────────
Introduction
──────────────────────────────────────────────
ShrijiLang ek custom programming language hai
jo Hinglish based syntax par kaam karti hai.

ShrijiLang is a custom programming language
built with Hinglish-style syntax for intuitive learning.

Ye version "Core Stable" hai.
This version is the "Core Stable" release.

──────────────────────────────────────────────
Features
──────────────────────────────────────────────
✔ rachna (functions)
✔ multi-argument support
✔ nested function calls
✔ recursion support
✔ KRST intelligence system
✔ error insight system

✔ Function support (rachna)
✔ Multiple arguments
✔ Nested execution
✔ Recursion support
✔ KRST intelligent execution layer
✔ Error insight system

──────────────────────────────────────────────
How to Run (Linux)
──────────────────────────────────────────────
1. Terminal open karo / Open terminal
2. Navigate to bin folder:

   cd release/linux/bin

3. Run:

   ./shrijilang

──────────────────────────────────────────────
Example
──────────────────────────────────────────────
rachna add(a, b) {
    wapas a + b
}

bolo add(5, 3)

Output:
8

──────────────────────────────────────────────
Limitations (Current Version)
──────────────────────────────────────────────
- Named arguments supported nahi hain
  (example: add(a=5, b=10))

- Default parameters supported nahi hain
  (example: rachna add(a, b=5))

- Named arguments are not supported yet
- Default parameters are not supported yet

──────────────────────────────────────────────
Future Updates
──────────────────────────────────────────────
v1.1 → Named arguments
v1.2 → Default parameters
v2.0 → Advanced execution engine

──────────────────────────────────────────────
Author
──────────────────────────────────────────────
MR._MR Manikpuri

Mission:
Thee Angel 😇

Vision:
To build a new generation intelligent programming ecosystem.
nd{verbatim}
nd{document}
